% 纲要标题：# 第三卷　黎曼几何与几何分析
% 纲要位置：几何学计划纲要.md，第 2118 行。
% 正文状态：未撰写；本文件目前是框架。

% 纲要内容开始
% 
% > 由局部度量和曲率进入比较，再建立足够的函数空间、椭圆和抛物工具。调和映射、极小对象、几何测度、几何流和典范度量形成可分别进入的深度路线。分析工具在本卷承担证明，不转移给尚未核验的其他系列。
% 
% 纲要内容结束

% !TeX TXS-program:compile = cmd /c build-Book3.bat -NoPause
% !TeX TXS-program:view = txs:///view-pdf-internal --embedded "Book3.pdf"

\documentclass[UTF8,zihao=-4,openany,fontset=fandol]{ctexbook}
\newcommand{\BookTitle}{黎曼几何与几何分析}
\newcommand{\BookVolumeName}{第三卷}
\newcommand{\BookNumber}{3}
\newcommand{\BookDraftDate}{2026年10月}
% 《几何学》五卷共用的基础排版与语义宏。
% 各卷主文件须先定义 \BookTitle、\BookVolumeName 和 \BookDraftDate。
\newcommand{\BookAuthor}{R. EverStarine}
% 稿件状态由各卷元数据指定；尚未开写的卷沿用框架状态。
\providecommand{\BookStatus}{框架预览·正文未撰写}

\usepackage[a4paper,inner=28mm,outer=24mm,top=27mm,bottom=28mm,
  headheight=16pt,headsep=8mm,footskip=12mm]{geometry}

\usepackage{amsmath}
\usepackage{mathtools}
\usepackage{amsthm}
\usepackage{aliascnt}
\usepackage{unicode-math}
% 数学字母使用默认数学字体内的对应字形；禁止 mathrm 等默认继承正文族。
% 运算符（sin、log、operatorname 等）随 mathrm 使用数学正体。
\unimathsetup{mathrm=sym,mathit=sym,mathbf=sym,mathsf=sym,mathtt=sym}
\setmainfont{texgyrepagella-regular.otf}[
  BoldFont=texgyrepagella-bold.otf,
  ItalicFont=texgyrepagella-italic.otf,
  BoldItalicFont=texgyrepagella-bolditalic.otf
]
\setsansfont{texgyreheros-regular.otf}[
  BoldFont=texgyreheros-bold.otf,
  ItalicFont=texgyreheros-italic.otf,
  BoldItalicFont=texgyreheros-bolditalic.otf
]
\setmonofont{lmmonolt10-regular.otf}[
  BoldFont=lmmonolt10-bold.otf,
  ItalicFont=lmmonolt10-oblique.otf,
  BoldItalicFont=lmmonolt10-boldoblique.otf
]
% 人物注释中的古希腊文、斯拉夫文字、日文异体字与希伯来文原名使用有相应字形的字体。
\newfontfamily\GeometryOriginalNameFont{Times New Roman}
\NewDocumentCommand{\OriginalName}{m}{{\GeometryOriginalNameFont #1}}
\newCJKfontfamily\GeometryOriginalHanFont{SimSun}
\NewDocumentCommand{\OriginalHanName}{m}{{\GeometryOriginalHanFont #1}}
\newfontfamily\GeometryHebrewFont{Times New Roman}[Script=Hebrew]
\TeXXeTstate=1
\NewDocumentCommand{\OriginalHebrewName}{m}{{\GeometryHebrewFont\beginR#1\endR}}
% 行内、行间及标题内公式均使用 unicode-math 默认 Latin Modern Math。
% 公式之外的正文西文仍由 setmainfont 使用 Pagella；黑板粗体单独采用 AMS。
\let\mathbb\relax
\DeclareMathAlphabet{\mathbb}{U}{msb}{m}{n}
% unicode-math 的 Latin Modern Math 将 \varnothing 映为瘦长椭圆；空集单独
% 采用传统 AMS MSBM 的圆形斜杠字形，不改变其余数学符号的字体。
\DeclareSymbolFont{GeometryAMSb}{U}{msb}{m}{n}
\DeclareMathSymbol{\GeometryRoundEmptySet}{\mathord}{GeometryAMSb}{"3F}
% 当前 Latin Modern Math 缺少横截符号 U+22D4；仅此关系符使用 AMSa 的现有字形。
\DeclareSymbolFont{GeometryAMSa}{U}{msa}{m}{n}
\DeclareMathSymbol{\GeometryTransverse}{\mathrel}{GeometryAMSa}{"74}
\AtBeginDocument{%
  \let\varnothing\GeometryRoundEmptySet
  \let\emptyset\GeometryRoundEmptySet
  \let\pitchfork\GeometryTransverse
}

\usepackage{microtype}
\usepackage{xcolor}
\usepackage{graphicx}
\usepackage{booktabs}
\usepackage{longtable}
\usepackage{array}
\usepackage{tabularx}
\usepackage{algorithm}
\usepackage{algpseudocode}
\usepackage{enumitem}
\usepackage{fancyhdr}
% 双面排版补齐奇数页时，空白页不残留上一章的页眉或页码。
\usepackage{emptypage}
\usepackage[most]{tcolorbox}
\usepackage{tikz-cd}
\usepackage{csquotes}
\usepackage[
  backend=biber,
  style=numeric-comp,
  sorting=nyt,
  giveninits=true,
  maxbibnames=99,
  doi=true,
  url=true,
  isbn=true
]{biblatex}
\addbibresource{../Shared/References.bib}
% 骨架没有引文时保留书目入口；写入真实条目后由完整构建自动排入书目。
\newif\ifGeometryHasBibliography
\AtDataInput{\global\GeometryHasBibliographytrue}
\NewDocumentCommand{\PrintBookBibliography}{}{%
  \ifGeometryHasBibliography
    \printbibliography[heading=bibintoc,title={参考文献}]%
  \else
    \chapter*{参考文献}%
    \addcontentsline{toc}{chapter}{参考文献}%
    \noindent 待撰写\par
  \fi
}
% 访问日期保留在书目库供内部核验，成品参考文献不显示 visited on。
\renewbibmacro*{url+urldate}{\usebibmacro{url}}
% biblatex 在书目内重绑定格式宏；在其作用域中将强调统一为正体，
% 同时避免书名中的中文及标点仍继承旧的斜体样式。
\AtBeginBibliography{\let\emph\textup\let\textit\textup}
% 中文条目的责任者连接符：不可套用西文的 “and”，按 GB/T 7714—2015 的著录
% 惯例改用逗号分隔。中文条目在 Shared/References.bib 中以 langid = {chinese}
% 标注（含译本，因其著录信息为中文形式），西文条目保持 “and” 不变。
% 按条目切换须在每条书目项处理时判定；\DeclareDelimFormat 的可选参数不是
% 语言选择器（实测无效），故用 \AtEveryBibitem 配合 \iffieldequalstr。
\AtEveryBibitem{%
  \iffieldequalstr{langid}{chinese}%
    {\DeclareDelimFormat{finalnamedelim}{\addcomma\space}%
     \DeclareDelimFormat{multinamedelim}{\addcomma\space}}%
    {\DeclareDelimFormat{finalnamedelim}{\addspace\bibstring{and}\space}%
     \DeclareDelimFormat{multinamedelim}{\addcomma\space}}%
}
\usepackage{etoolbox}
\usepackage{xparse}
\usepackage[noautomatic]{imakeidx}

\definecolor{GeometryBlue}{HTML}{244A64}
\definecolor{GeometryBlueLight}{HTML}{EEF4F7}
\definecolor{GeometryGray}{HTML}{4F5962}
\definecolor{GeometryGrayLight}{HTML}{F4F5F6}

\ctexset{
  part = {
    % 编题页分两行：第一行“第X编”，与封面“第一卷”同款（大号无衬线灰色）；
    % 第二行编名，与封面“几何学”同款（大号无衬线蓝色粗体），字号保持原样。
    % 第一行与行距由 number／aftername 给出，format 只负责编名的字体，
    % 以免字号被重置或被重复输出。
    format = {\centering\fontsize{30}{38}\selectfont\sffamily\bfseries\color{GeometryBlue}},
    name = {,},
    number = {\centering\normalfont\normalsize{\fontsize{24}{30}\selectfont\sffamily\color{GeometryGray}第\chinese{part}编}},
    aftername = {\par\vspace{0.7cm}},
    beforeskip = 0pt,
    afterskip = 2.5\baselineskip,
    pagestyle = plain
  },
  chapter = {
    format = \raggedright\Huge\bfseries\color{GeometryBlue},
    name = {第,章},
    number = \chinese{chapter},
    aftername = \quad,
    beforeskip = 0pt,
    afterskip = 30pt,
    pagestyle = plain
  },
  section = {
    format = \Large\bfseries\color{GeometryBlue},
    beforeskip = 2.4ex plus .6ex minus .2ex,
    afterskip = 1.1ex plus .2ex
  },
  subsection = {
    format = \large\bfseries\color{GeometryGray},
    beforeskip = 2ex plus .4ex minus .2ex,
    afterskip = .8ex plus .2ex
  }
}

\setlength{\parindent}{2em}
\setlength{\parskip}{0pt}
\linespread{1.18}
% 仅用于两个连续的普通正文段落；列表、公式、定理、标题等结构边界不用。
% 段间额外留白统一由这一长度控制，避免改动全局 \parskip 后扰乱各类环境。
\newlength{\GeometryProseParagraphGap}
\AtBeginDocument{%
  \setlength{\GeometryProseParagraphGap}{0.40\baselineskip}%
}
\NewDocumentCommand{\ProseParagraphBreak}{}{%
  \par\addvspace{\GeometryProseParagraphGap}%
}
% 与代数学保持一致：小节编号并进入目录，便于按具体主题定位。
\setcounter{secnumdepth}{3}
\setcounter{tocdepth}{2}
\allowdisplaybreaks[2]
\emergencystretch=2em
\clubpenalty=10000
\widowpenalty=10000
\displaywidowpenalty=10000

% 固定列表内部段距，避免 flushbottom 将默认的弹性 \parsep 拉大到超过正文段距。
\setlist[itemize]{leftmargin=2.2em,itemsep=.35ex,parsep=0pt,topsep=.6ex}
\setlist[enumerate]{leftmargin=2.5em,itemsep=.35ex,parsep=0pt,topsep=.6ex,
  font=\normalfont\upshape}
\setlist[enumerate,1]{label=\arabic*.,ref=\arabic*}
\setlist[enumerate,2]{label=\arabic*),ref=\arabic*}
\setlist[enumerate,3]{label=\alph*),ref=\alph*}
\setlist[description]{leftmargin=2.5em,labelindent=0pt,itemsep=.35ex,parsep=0pt,topsep=.6ex}

% “下列命题等价”专用列表；普通罗列仍使用 enumerate。
\newlist{equivalences}{enumerate}{2}
\setlist[equivalences,1]{leftmargin=2.8em,itemsep=.35ex,parsep=0pt,topsep=.6ex,
  label=\Roman*.,ref=\Roman*,font=\normalfont\upshape}
\setlist[equivalences,2]{leftmargin=2.8em,itemsep=.35ex,parsep=0pt,topsep=.6ex,
  label=\roman*),ref=\roman*,font=\normalfont\upshape}

\pagestyle{fancy}
\fancyhf{}
\fancyhead[LE]{\small\color{GeometryGray}\thepage\quad\leftmark}
% 奇数页采用页前最后有效的节标记，避免页中开始新节时提前更换页眉。
% 新节或习题的标记从后续页生效；章首页使用 plain，不显示上一章的标记。
\fancyhead[RO]{\small\color{GeometryGray}\TopMark{2e-right}\quad\thepage}
\renewcommand{\headrulewidth}{0.35pt}
% 保留数学标题的原始大小写，避免 book 类的 \MakeUppercase
% 与 unicode-math 命令在页眉标记中发生冲突。
% 偶数页页眉显示“章序＋章标题”（如“第一章 代数对象、映射与基本构造”）；
% 章序与章标题之间以 \quad 分隔，间距与奇数页节眉“节号＋节标题”之间一致；
% 附录章的章序由 \CTEXthechapter 按本卷附录字母（如 A）给出。
\renewcommand{\chaptermark}[1]{\markboth{\CTEXthechapter\quad #1}{}}
\renewcommand{\sectionmark}[1]{\markright{\thesection\quad #1}}
% 章末习题不编号，但进入习题页后应覆盖上一节留下的奇数页页眉；
% 同时写入目录项。PDF 书签由 hyperref 对 \section*{习题} 自动生成，
% 不再另行 \pdfbookmark，避免“习题”下再嵌套一个“习题”。
\newcommand{\BookExercises}{%
  \section*{习题}%
  \markright{习题}%
  \addcontentsline{toc}{section}{习题}%
}
\fancypagestyle{plain}{
  \fancyhf{}
  \fancyfoot[C]{\small\color{GeometryGray}\thepage}
  \renewcommand{\headrulewidth}{0pt}
}

\usepackage[unicode,colorlinks=true,linkcolor=GeometryBlue,
  citecolor=GeometryBlue,urlcolor=GeometryBlue,pdfborder={0 0 0}]{hyperref}
\usepackage{bookmark}
% PDF 书签沿用目录条目中的编号：正文章为“第几章”，附录章为“附录 A”，
% 节与小节分别保留 x.x、x.x.x；无编号的导读、习题和卷末单元仍不加号。
\bookmarksetup{numbered}

% unicode-math 会调整原始 Unicode 希腊字母的类别码。若把这些字母直接写进
% \texorpdfstring 的书签分支，辅助文件回读时类别码不同，LaTeX 会在每一轮都
% 误报标签发生变化。书签分支统一写成下列可展开宏，既保留希腊字母显示，
% 又使写入和回读的标签文本稳定。
\newcommand{\GeometryPdfAlpha}{α}
\newcommand{\GeometryPdfGamma}{Γ}
\newcommand{\GeometryPdfLambda}{λ}
\newcommand{\GeometryPdfMu}{μ}
\newcommand{\GeometryPdfPi}{π}
\newcommand{\GeometryPdfSigma}{σ}
\newcommand{\GeometryPdfTheta}{θ}
\newcommand{\GeometryPdfXi}{ξ}
\newcommand{\GeometryPdfZeta}{ζ}
\newcommand{\GeometryPdfCapitalLambda}{Λ}
\newcommand{\GeometryPdfOmega}{Ω}
\newcommand{\GeometryPdfCapitalSigma}{Σ}
\newcommand{\GeometryPdfDelta}{δ}
\newcommand{\GeometryPdfEpsilon}{ε}
\newcommand{\GeometryPdfIota}{ι}
\newcommand{\GeometryPdfSmallOmega}{ω}
\newcommand{\GeometryPdfEll}{ℓ}

\usepackage[nameinlink,noabbrev]{cleveref}
\newcommand{\crefpairconjunction}{和}
\newcommand{\crefmiddleconjunction}{、}
\newcommand{\creflastconjunction}{和}

% CTeX 2.5.10 的标题前缀宏带一个参数，而 TeX Live 2025 中 cleveref
% 仍按普通前缀宏展开它们，会把节号写成含 \endcsname 的损坏元数据。
% 改为恒等包装后，CTeX 传入的标题编号保持不变，\ref 与 \cref 均可正常工作。
\makeatletter
\def\p@part#1{#1}
\def\p@chapter#1{#1}
\def\p@section#1{#1}
\def\p@subsection#1{#1}
\def\p@subsubsection#1{#1}
\def\p@paragraph#1{#1}
\def\p@subparagraph#1{#1}
\makeatother
% 非数学西文字母统一正体；不改变括号类别、标点或数学字体。
% 2026-09-11：非数学模式的西文字母统一使用 Pagella 正体，保留粗细区别。
% 保留字母级边界，不重分类数字、括号、标点及中文。
% 此设置也覆盖列表序号字母和术语原文；数学模式保持原样。
\newfontfamily\GeometryLetterFont{texgyrepagella-regular.otf}[
  BoldFont=texgyrepagella-bold.otf,
  ItalicFont=texgyrepagella-italic.otf,
  BoldItalicFont=texgyrepagella-bolditalic.otf
]
\ExplSyntaxOn
\bool_new:N \l_geometry_text_letters_bool
\bool_set_true:N \l_geometry_text_letters_bool
\cs_new_protected:Npn \GeometryPreserveMathText
  { \bool_set_false:N \l_geometry_text_letters_bool }
\cs_new_protected:Npn \geometry_letter_group_begin:
  {
    \group_begin:
    \bool_if:NT \l_geometry_text_letters_bool
      { \GeometryLetterFont \upshape }
  }
% 继承 Default 的全部 xeCJK 转移，保留中西文间距及标点处理；
% 不能用只有开关分组的转移覆盖 Boundary 或 CJK 的原处理。
\xeCJK_new_class:n { GeometryLatinLetter }
\seq_map_inline:Nn \g__xeCJK_class_seq
  {
    \str_if_eq:nnF {#1} { GeometryLatinLetter }
      {
        \xeCJK_copy_inter_class_toks:nnnn
          {#1} {GeometryLatinLetter} {#1} {Default}
        \xeCJK_copy_inter_class_toks:nnnn
          {GeometryLatinLetter} {#1} {Default} {#1}
        \xeCJK_app_inter_class_toks:nnn {#1} {GeometryLatinLetter}
          { \geometry_letter_group_begin: }
        \xeCJK_pre_inter_class_toks:nnn {GeometryLatinLetter} {#1}
          { \group_end: }
      }
  }
% ASCII 字母及常用拉丁扩展字母；明确跳过乘号、除号，不重分类标点。
\xeCJKDeclareCharClass { GeometryLatinLetter }
  { "0041->"005A, "0061->"007A, "00C0->"00D6, "00D8->"00F6,
    "00F8->"024F, "1E00->"1EFF }
\ExplSyntaxOff
% 数学组内局部关闭字母切换，连公式中 \text{...} 的字体也保持原样。
% 只追加，不覆盖 amsmath、unicode-math、xeCJK 已有的数学初始化。
\AtBeginDocument{%
  \everymath\expandafter{\the\everymath\GeometryPreserveMathText}%
  \everydisplay\expandafter{\the\everydisplay\GeometryPreserveMathText}%
}


% 定理类环境（定理、定义、命题、引理、推论、断言、猜想、开放问题、注记、
% 习题）的正文：中文用楷体，西文保持正体，数学字体不受影响。
% 注意 \normalfont 只重置西文字族，中文字族仍随默认（宋体），故须显式 \kaishu；
% 不能退回 \itshape——那会把西文一并变斜。
\newtheoremstyle{geometryplain}
  {1.2ex}{1.2ex}{\normalfont\upshape\kaishu}{}
  {\color{GeometryBlue}\bfseries}{.}{0.6em}{}
\newtheoremstyle{geometrydefinition}
  {1.2ex}{1.2ex}{\normalfont\kaishu}{}
  {\color{GeometryBlue}\bfseries}{.}{0.6em}{}
\newtheoremstyle{geometryremark}
  {1.0ex}{1.0ex}{\normalfont\kaishu}{}
  {\color{GeometryGray}\bfseries}{.}{0.6em}{}

\theoremstyle{geometryplain}
\newtheorem{theorem}{定理}[section]
% 正文默认按“章.节.流水号”编号；未分节的附录省略中间的 0。
\renewcommand{\thetheorem}{%
  \ifnum\value{section}=0
    \thechapter.\arabic{theorem}%
  \else
    \thesection.\arabic{theorem}%
  \fi
}
\newaliascnt{lemma}{theorem}
\newtheorem{lemma}[lemma]{引理}
\aliascntresetthe{lemma}
\newaliascnt{proposition}{theorem}
\newtheorem{proposition}[proposition]{命题}
\aliascntresetthe{proposition}
\newaliascnt{corollary}{theorem}
\newtheorem{corollary}[corollary]{推论}
\aliascntresetthe{corollary}
\newaliascnt{claim}{theorem}
\newtheorem{claim}[claim]{断言}
\aliascntresetthe{claim}
\newaliascnt{conjecture}{theorem}
\newtheorem{conjecture}[conjecture]{猜想}
\aliascntresetthe{conjecture}
\newaliascnt{openproblem}{theorem}
\newtheorem{openproblem}[openproblem]{开放问题}
\aliascntresetthe{openproblem}

\theoremstyle{geometrydefinition}
\newaliascnt{definition}{theorem}
\newtheorem{definition}[definition]{定义}
\aliascntresetthe{definition}
\newaliascnt{example}{theorem}
\newtheorem{example}[example]{例}
\aliascntresetthe{example}
% 习题按章单独编号，不占用定理类环境的流水号。
\newtheorem{exercise}{习题}[chapter]

\theoremstyle{geometryremark}
\newaliascnt{remark}{theorem}
\newtheorem{remark}[remark]{注记}
\aliascntresetthe{remark}
\renewcommand{\proofname}{证明}
% 仅在证明紧接定理类陈述时收紧间距，并按陈述结尾是否为行间公式调整。
% 证明标题（证明／解／证明概要）的中文用楷体，
% 西文保持正体；正文不施加楷体，仍为默认宋体。不能退回 \itshape。
\newif\ifGeometryDirectProof
\newif\ifGeometryCoupledProof
\newif\ifGeometryStatementEndsInDisplay
\newskip\GeometryProofTrailingSkip
% 行间公式结束后，TeX 留下空的水平列表；在定理环境结束前辨认这一状态。
% 证明跟在公式之后时不再减去定理尾距，否则公式下方会比上方更紧。
\newcommand{\GeometryRecordStatementTail}{%
  \ifhmode
    \ifnum\lastnodetype=-1\relax
      \global\GeometryStatementEndsInDisplaytrue
    \else
      \global\GeometryStatementEndsInDisplayfalse
    \fi
  \else
    \global\GeometryStatementEndsInDisplayfalse
  \fi
}
\AddToHook{para/begin}{\global\GeometryDirectProoffalse\global\GeometryStatementEndsInDisplayfalse}
\AddToHook{env/theorem/end}{\GeometryRecordStatementTail}
\AddToHook{env/lemma/end}{\GeometryRecordStatementTail}
\AddToHook{env/proposition/end}{\GeometryRecordStatementTail}
\AddToHook{env/corollary/end}{\GeometryRecordStatementTail}
\AddToHook{env/claim/end}{\GeometryRecordStatementTail}
\AddToHook{env/exercise/end}{\GeometryRecordStatementTail}
\AddToHook{env/theorem/after}{\global\GeometryDirectProoftrue}
\AddToHook{env/lemma/after}{\global\GeometryDirectProoftrue}
\AddToHook{env/proposition/after}{\global\GeometryDirectProoftrue}
\AddToHook{env/corollary/after}{\global\GeometryDirectProoftrue}
\AddToHook{env/claim/after}{\global\GeometryDirectProoftrue}
\AddToHook{env/exercise/after}{\global\GeometryDirectProoftrue}
\AddToHook{env/figure/before}{\global\GeometryDirectProoffalse}
\AddToHook{env/table/before}{\global\GeometryDirectProoffalse}
\makeatletter
\renewenvironment{proof}[1][\proofname]{\par
  \pushQED{\qed}%
  \normalfont \topsep6\p@\@plus6\p@\relax
  \GeometryCoupledProoffalse
  \ifGeometryDirectProof
    \GeometryCoupledProoftrue
    \GeometryProofTrailingSkip=\topsep
    \advance\GeometryProofTrailingSkip by \partopsep
    \ifGeometryStatementEndsInDisplay\else
      \kern-5\p@% 只在陈述以普通文字结尾时抵消部分尾距。
    \fi
    \topsep2\p@\partopsep\z@\relax
  \fi
  \global\GeometryDirectProoffalse
  \global\GeometryStatementEndsInDisplayfalse
  \trivlist
  \item[\hskip\labelsep\upshape\kaishu #1\@addpunct{.}]\ignorespaces
}{%
  \popQED
  \ifGeometryCoupledProof\@topsepadd\GeometryProofTrailingSkip\fi
  \endtrivlist\@endpefalse
}
\makeatother
\newenvironment{solution}
  {\begin{proof}[解]}
  {\end{proof}}
\newenvironment{proofsketch}
  {\begin{proof}[证明概要]}
  {\end{proof}}

% 正文选读内容的星号属于编号而不属于标题；第二参数统一接收标签。
\newcommand{\GeometryOptionalMark}{\texorpdfstring{\textsuperscript{*}}{*}}
% 选做习题沿用同一编号星号。可选参数仍只表示习题自身的简短题名，
% 不得再写“选读”或“选做”；星号由本环境自动附在题号右侧。
\NewDocumentEnvironment{optionalexercise}{o}{%
  \begingroup
  \let\GeometrySavedTheExercise\theexercise
  \renewcommand{\theexercise}{\GeometrySavedTheExercise\GeometryOptionalMark}%
  \IfNoValueTF{#1}{\begin{exercise}}{\begin{exercise}[#1]}%
}{%
  \end{exercise}%
  \endgroup
}
% 可选参数只供目录与 PDF 书签使用；正文题名可借此按语义手动换行，
% 又不把换行命令写入目录或书签。
\NewDocumentCommand{\OptionalChapter}{O{}mm}{%
  \begingroup
  \let\GeometrySavedTheChapter\thechapter
  \renewcommand{\thechapter}{\GeometrySavedTheChapter\GeometryOptionalMark}%
  \ctexset{chapter={name={第,章\GeometryOptionalMark},
    number=\chinese{chapter},aftername=\quad}}%
  \ifstrempty{#1}{\chapter{#2}}{\chapter[#1]{#2}}\label{#3}%
  \endgroup
}
\NewDocumentCommand{\OptionalSection}{mm}{%
  \begingroup
  \let\GeometrySavedTheSection\thesection
  \renewcommand{\thesection}{\GeometrySavedTheSection\GeometryOptionalMark}%
  \ctexset{section={number=\thesection}}%
  \section{#1}\label{#2}%
  \endgroup
}
\NewDocumentCommand{\OptionalSubsection}{mm}{%
  \begingroup
  \let\GeometrySavedTheSubsection\thesubsection
  \renewcommand{\thesubsection}{\GeometrySavedTheSubsection\GeometryOptionalMark}%
  \ctexset{subsection={number=\thesubsection}}%
  \subsection{#1}\label{#2}%
  \endgroup
}

% 编题页必须从奇数页开始，题页仅在页脚显示页码，其后保留一张空白页。
\makeatletter
\NewDocumentCommand{\BookPart}{mm}{%
  \cleardoublepage
  \begingroup
  \def\GeometryPartLabel{#2}%
  % \part 会在返回前翻页，故须在 \@endpart 开始处记录题页标签。
  \pretocmd{\@endpart}{\label{\GeometryPartLabel}}{}%
    {\PackageError{geometry}{无法为编题页写入标签}{请检查当前文档类的 \string\@endpart 定义。}}%
  \part{#1}%
  \endgroup
  \thispagestyle{empty}%
  \null
  \clearpage
}
% 各编的“本编导读”是无编号章，在总目录和 PDF 书签中与所属编的各章并列。
% addcontentsline 同时建立目录项和书签，不再另加同名 pdfbookmark。
\NewDocumentCommand{\BookPartGuide}{O{本编导读}}{%
  \chapter*{#1}%
  \addcontentsline{toc}{chapter}{#1}%
}

% “附录”单独排无编号编题页，使用与正文编名相同的字号、字体和颜色。
% 题页从奇数页开始，其后保留一张空白页，再进入第一个附录章。
% startatroot 使根级书签与书后顶层节点分离；phantomsection 须先于
% addcontentsline，目录中的“附录”行才能获得独立锚点而非沿用前一节的
% section*.N。根级“附录”书签连同其子章由 addcontentsline 的 part 行自动
% 建立，不得再手写 \pdfbookmark（实测会产生重复的“附录”节点）。
\NewDocumentCommand{\BookAppendixPart}{m}{%
  \cleardoublepage
  \begingroup
  \def\GeometryPartLabel{#1}%
  % 在 part* 的题页尚未翻走时建立唯一的目录与书签目标。
  \pretocmd{\@endpart}{%
    \bookmarksetup{startatroot}%
    \phantomsection
    \addcontentsline{toc}{part}{附录}%
    \label{\GeometryPartLabel}%
  }{}{\PackageError{geometry}{无法为附录编题页写入目录与标签}{请检查当前文档类的 \string\@endpart 定义。}}%
  \part*{附录}%
  \endgroup
  \thispagestyle{empty}%
  \null
  \clearpage
}
\makeatother

% 列表项的括号与编号整体承担链接，避免只让中间的编号可点击。
\NewDocumentCommand{\itemref}{m}{\hyperref[#1]{(\ref*{#1})}}

% 章、节、小节的引用：整个“第 X 章／第 X 节／第 X 小节”都是跳转链接。
% 链接必须走 hyperref 的标签解析机制（目标名 chapter.N／section.N.M）；
% 早先用 \hyperlink{裸标签名} 指向不存在的目标，链接全部失效。
%
% 章号用汉字（“第五章”）且必须取“被引用那一章”的编号；\thechapter 是当前
% 计数器，在编导读、前言或跨章引用时会给出错误章号（表现为“第零章”）。
% hyperref 下 \ref 受保护、\getrefnumber 在 \edef 中又失真，故改由 zref 记录：
% 各章在 \chapter 之后调用 \BookChapterNumber，把汉字号存为 zref 属性；
% \zref@extract 可展开，能在 \edef 中安全取出。
\usepackage[user,hyperref]{zref}
\makeatletter
\zref@newprop{geometrychapternum}{}
% 在 \chapter 之后调用，参数即该章标签；把本章汉字号存为 zref 属性
\NewDocumentCommand{\BookChapterNumber}{m}{%
  \zref@setcurrent{geometrychapternum}{\chinese{chapter}}%
  \zref@labelbyprops{zch@#1}{geometrychapternum}%
}
% 取出某章标签对应的汉字号（\zref@extract 可展开，能在 \edef 中安全使用）
\NewDocumentCommand{\GeometryChineseChapterOf}{m}{%
  \zref@extract{zch@#1}{geometrychapternum}}
\makeatother
% \ref* 的分组会遮住编号末端的西文字符类别；编号前的间距可由 xeCJK
% 自动产生，编号后的同类间距须在完整链接内部显式补回。
\NewDocumentCommand{\secref}{m}{%
  \hyperref[#1]{第\ref*{#1}\CJKecglue 节}}
\NewDocumentCommand{\subsecref}{m}{%
  \hyperref[#1]{第\ref*{#1}\CJKecglue 小节}}
\NewDocumentCommand{\chapref}{m}{%
  \begingroup
  \xeCJKsetup{CJKecglue={}}%
  \edef\GeometryChapterDisp{第\GeometryChineseChapterOf{#1}章}%
  \hyperref[#1]{\GeometryChapterDisp}%
  \endgroup}

% 算法采用带编号的步骤，并明确标出输入与输出。
\floatname{algorithm}{算法}
\renewcommand{\listalgorithmname}{算法目录}
\algrenewcommand\algorithmicrequire{\textbf{输入：}}
\algrenewcommand\algorithmicensure{\textbf{输出：}}
\numberwithin{algorithm}{chapter}
% 完整中文编号进入题注、各自目录和引用；附录自动使用字母章号。
% 内部 \nobreakspace 与 cleveref 默认的定理引用一致；引用整串置于同一链接，
% 链接后的中文间距由下方 \cref 包装补齐。
\renewcommand{\thealgorithm}{算法\nobreakspace\thechapter-\arabic{algorithm}}
\renewcommand{\thefigure}{图\nobreakspace\thechapter-\arabic{figure}}
\renewcommand{\thetable}{表\nobreakspace\thechapter-\arabic{table}}
\renewcommand{\figurename}{}
\renewcommand{\tablename}{}
\makeatletter
\renewcommand{\fnum@figure}{\thefigure}
\renewcommand{\fnum@table}{\thetable}
\renewcommand{\fnum@algorithm}{\thealgorithm}
\makeatother
\NewDocumentEnvironment{algorithmsteps}{mm}
  {\begin{algorithmic}[1]\Require #1\Ensure #2}
  {\end{algorithmic}}

\newtcolorbox{learninggoals}{
  enhanced,breakable,colback=GeometryBlueLight,colframe=GeometryBlue,
  boxrule=.45pt,arc=1mm,left=1.2mm,right=1.2mm,top=1mm,bottom=1mm,
  title=本章目标,fonttitle=\bfseries
}
\newtcolorbox{partguide}{
  enhanced,breakable,colback=GeometryGrayLight,colframe=GeometryGray,
  boxrule=.35pt,arc=1mm,left=1.2mm,right=1.2mm,top=1mm,bottom=1mm,
  title=本编导读,fonttitle=\bfseries
}
\newtcolorbox{chaptersummary}{
  enhanced,breakable,colback=GeometryBlueLight,colframe=GeometryBlue,
  boxrule=.45pt,arc=1mm,left=1.2mm,right=1.2mm,top=1mm,bottom=1mm,
  title=本章小结,fonttitle=\bfseries,
  % 与阅读提示框一致：仅本框内容为多段散文，须恢复首行缩进。
  before upper={\parindent=2em\parskip=0pt}
}
% 仅用于尚待高阶数学复核的核心断言；成稿前必须处理或移除。
\newtcolorbox{coreverification}{
  enhanced,breakable,colback=yellow!10,colframe=orange!85!black,
  boxrule=.55pt,arc=1mm,left=1.2mm,right=1.2mm,top=1mm,bottom=1mm,
  title=待数学复核,fonttitle=\bfseries
}
% 兼容旧稿的占位命令：只输出普通文字，不另加方框。
\NewDocumentCommand{\GeometryPlaceholder}{m}{待撰写}
\NewDocumentEnvironment{chapterabstract}{}
  {\small\list{}{\rightmargin\leftmargin
    \listparindent=2\ccwd\itemindent=\listparindent}%
    \item\relax\color{GeometryGray}}
  {\endlist}
% 引诗与短引文统一居中，中文使用定理类环境的楷体；章摘要已在上方独立定义，
% 不继承这里的字形与对齐方式。
\RenewDocumentEnvironment{quote}{}
  {\list{}{\rightmargin\leftmargin}\item\relax\centering\kaishu}
  {\endlist}
% 引诗环境保留为语义清楚的兼容入口，版式与 quote 完全一致。
\NewDocumentEnvironment{poemquote}{}
  {\begin{quote}}
  {\end{quote}}
% 引诗末尾的作者或出处在引诗区域内右对齐。
\NewDocumentCommand{\quoteattribution}{m}{%
  \par\smallskip
  {\raggedleft #1\par}%
}

\makeindex[name=chinese,title=中文索引,columns=2,intoc,options=-q]
\makeindex[name=foreign,title=外文索引,columns=2,intoc,options=-q]
\makeindex[name=symbols,title=符号索引,columns=2,intoc,options=-q]
% 索引采用稍小字号与自然栏底，字母组间仅留少量空白；末页平衡不得越过正文高度。
\indexsetup{othercode={\small\raggedcolumns\maxbalancingoverflow=0pt
  \renewcommand{\indexspace}{\par\addvspace{0.2\baselineskip}}}}

% 有条目时读入 makeindex 生成的真实索引；尚无条目时仍保留书后结构与目录入口。
% 一键脚本会在 .idx 为空时删除旧 .ind，故这里不会误排上一次构建的陈旧索引。
\NewDocumentCommand{\PrintBookIndex}{m m}{%
  \IfFileExists{#1.ind}{%
    % 双栏输出可能丢失页眉标记；在索引范围内直接绑定标题并完成分页。
    \cleardoublepage
    \begingroup
      \fancyhead[LE]{\small\color{GeometryGray}\thepage\quad #2}%
      \fancyhead[RO]{\small\color{GeometryGray}#2\quad\thepage}%
      \printindex[#1]%
      \clearpage
    \endgroup
  }{%
    \chapter*{#2}%
    \addcontentsline{toc}{chapter}{#2}%
    \noindent 待撰写\par
  }%
}

% 所有术语在每卷首次正式定义或引入时都显示完整原文，并登记中外文索引。
% 星号形式与缺失原文均视为错误，防止旧的隐藏原文规则被重新带回源码。
% 定理类正文处于楷体；Fandol 楷体没有粗体字形，直接 \textbf 会退回普通楷体。
% 术语主称因此局部恢复正文的宋体字族，再使用与普通正文一致的粗体。
\NewDocumentCommand{\term}{s O{} m o}{%
  \IfBooleanT{#1}{%
    \PackageError{geometry}{\string\term* 已取消}{请改用 \string\term，并保留末尾的原文注解。}%
  }%
  {\songti\textbf{\textup{#3}}}%
  \IfNoValueTF{#4}{%
    \PackageError{geometry}{术语缺少原文注解}{每个术语首次正式定义或引入时都必须写成 \string\term[排序键]{中文主称}[原文]。}%
  }{%
    \textnormal{(\textup{#4})}%
  }%
  \ifstrempty{#2}{\index[chinese]{#3}}{\index[chinese]{#2@#3}}%
  \IfNoValueF{#4}{\index[foreign]{#4}}%
}
\NewDocumentCommand{\termidx}{O{}m}{%
  \ifstrempty{#1}{\index[chinese]{#2}}{\index[chinese]{#1@#2}}%
}
% 别称与主称一样用 \termidx 登记实际页码，不生成“见某某”的转引条目。
\NewDocumentCommand{\foreignidx}{O{}m}{%
  \ifstrempty{#1}{\index[foreign]{#2}}{\index[foreign]{#1@#2}}%
}
\newcounter{geometrySymbolIndexOrder}
\NewDocumentCommand{\symidx}{m}{%
  \stepcounter{geometrySymbolIndexOrder}%
  \edef\GeometrySymbolSortKey{%
    \number\numexpr10000+\value{geometrySymbolIndexOrder}\relax}%
  \index[symbols]{\GeometrySymbolSortKey@#1}%
}

% 一般强调与术语粗体分离。
\NewDocumentCommand{\keypoint}{m}{\CJKunderdot{#1}}
\NewDocumentCommand{\foreignkeypoint}{m}{\underline{\textup{#1}}}


\crefname{chapter}{章}{章}
\crefname{part}{编}{编}
\crefname{section}{节}{节}
\crefname{subsection}{小节}{小节}
\crefname{theorem}{定理}{定理}
\crefname{lemma}{引理}{引理}
\crefname{proposition}{命题}{命题}
\crefname{corollary}{推论}{推论}
\crefname{claim}{断言}{断言}
\crefname{conjecture}{猜想}{猜想}
\crefname{openproblem}{开放问题}{开放问题}
\crefname{definition}{定义}{定义}
\crefname{example}{例}{例}
\crefname{remark}{注记}{注记}
\crefname{exercise}{习题}{习题}
\crefname{equation}{式}{式}
\crefname{figure}{}{}
\crefname{table}{}{}
\crefname{algorithm}{}{}
\crefformat{figure}{#2#1#3}
\Crefformat{figure}{#2#1#3}
\crefformat{table}{#2#1#3}
\Crefformat{table}{#2#1#3}
\crefformat{algorithm}{#2#1#3}
\Crefformat{algorithm}{#2#1#3}

% cleveref 会在自己的分组中排出“定义 3.1.5”一类完整链接，类别名与编号之间
% 使用 \nobreakspace。分组结束后，xeCJK 无法再从末端识别出西文数字，
% 因而不会在紧随其后的汉字前插入 \CJKecglue。这里只在下一个源码记号确为
% CJK 正文字时补回同一种间距；空格、控制序列以及中英文标点均不处理，
% 从而不在“定义3.1.5，”或“定义3.1.5.”前制造额外空白。
\NewCommandCopy{\GeometryOriginalCref}{\cref}
\NewCommandCopy{\GeometryOriginalCrefCapital}{\Cref}
\ExplSyntaxOn
\cs_new_protected:Npn \geometry_cref_maybe_cjk_glue:
  { \peek_after:Nw \geometry_cref_maybe_cjk_glue_aux: }
\cs_new_protected:Npn \geometry_cref_maybe_cjk_glue_aux:
  {
    \token_if_space:NF \l_peek_token
      {
        \token_if_cs:NF \l_peek_token
          {
            \xeCJK_if_same_class:NNT \l_peek_token 中
              { \CJKecglue }
          }
      }
  }
\RenewDocumentCommand{\cref}{s m}
  {
    \IfBooleanTF{#1}
      { \GeometryOriginalCref*{#2} }
      { \GeometryOriginalCref{#2} }
    \geometry_cref_maybe_cjk_glue:
  }
\RenewDocumentCommand{\Cref}{s m}
  {
    \IfBooleanTF{#1}
      { \GeometryOriginalCrefCapital*{#2} }
      { \GeometryOriginalCrefCapital{#2} }
    \geometry_cref_maybe_cjk_glue:
  }
% 附录标签沿用字母编号；名称与编号整体进入同一链接，并与 \cref
% 采用相同的链接后中西文间距规则。
\NewDocumentCommand{\appref}{m}
  {
    \hyperref[#1]{附录\CJKecglue\ref*{#1}}
    \geometry_cref_maybe_cjk_glue:
  }
\ExplSyntaxOff

\newcommand{\MakeBookTitle}[1][封面]{%
  % 书名页在 PDF 侧栏中作为顶级入口；卷级文件可传入更具辨识性的书签名。
  \bookmarksetup{startatroot}%
  \pdfbookmark[0]{#1}{GeometryCover}%
  \begin{titlepage}
    \centering
    \vspace*{2.2cm}
    {\fontsize{36}{46}\selectfont\sffamily 几何学\par}
    \vspace{0.7cm}
    {\Large\sffamily\color{GeometryGray}\BookVolumeName\par}
    \vspace{1.0cm}
    {\fontsize{24}{32}\selectfont\bfseries\color{GeometryBlue}\BookTitle\par}
    \vfill
    {\large \BookAuthor\ 编写\par}
    \vspace{0.5cm}
    {\normalsize\BookDraftDate\par}
    \vspace{0.7cm}
    {\large\sffamily\bfseries\color{GeometryGray}\BookStatus\par}
    \vspace*{1.8cm}
  \end{titlepage}
}

% 目录也应有独立的顶级 PDF 书签。使用公共命令避免五卷分别遗漏。
\newcommand{\MakeBookContents}{%
  \cleardoublepage
  \bookmarksetup{startatroot}%
  \pdfbookmark[0]{目录}{GeometryContents}%
  \tableofcontents
}

% 图、表、算法目录不进入总目录，也不建立 PDF 书签；各自页面使用原有章级标题。
% 只另起一页而不补单双页。
\makeatletter
\newcommand{\GeometryListStart}{%
  \clearpage
}
% 各列表命令自带 \chapter*，局部消除其中的再次换页，保留原有标题与列表排版。
\newcommand{\GeometryListOnCurrentPage}[1]{%
  \begingroup
    \let\clearpage\relax
    \let\cleardoublepage\relax
    #1%
  \endgroup
}
\newcommand{\MakeBookListOfFigures}{%
  \GeometryListStart
  \GeometryListOnCurrentPage{\listoffigures}%
}
\newcommand{\MakeBookListOfTables}{%
  \begingroup
    \renewcommand{\listtablename}{表目录}%
    \GeometryListStart
    \GeometryListOnCurrentPage{\listoftables}%
  \endgroup
}
\newcommand{\MakeBookListOfAlgorithms}{%
  \begingroup
    \GeometryListStart
    % float.sty 在 listof 内重设编号栏；在标题回调中再加宽它。
    \let\GeometryOriginalFloatListHead\float@listhead
    \renewcommand{\float@listhead}[1]{%
      \renewcommand*\l@algorithm{\@dottedtocline{1}{1.5em}{5em}}%
      \GeometryOriginalFloatListHead{##1}%
    }%
    \GeometryListOnCurrentPage{\listofalgorithms}%
  \endgroup
}
\makeatother

% 无编号章（序言、目录、编导读、后记、参考文献、索引等）同样写入页眉标记，
% 使后续偶数页页眉显示该单元标题。置于 hyperref 之后，保证保留 hyperref 的章锚点。
\makeatletter
\let\Geometry@orig@schapter\@schapter
\def\@schapter#1{%
  \markboth{#1}{#1}%
  \Geometry@orig@schapter{#1}%
}
% 目录按层级着色：编为黑色、小节为深灰；章与节保持默认链接蓝。
\pretocmd{\l@part}{\hypersetup{linkcolor=black}}{}{%
  \PackageWarning{geometry}{无法为目录编条目着色}}
\pretocmd{\l@chapter}{\hypersetup{linkcolor=GeometryBlue}}{}{%
  \PackageWarning{geometry}{无法为目录章条目着色}}
\pretocmd{\l@section}{\hypersetup{linkcolor=GeometryBlue}}{}{%
  \PackageWarning{geometry}{无法为目录节条目着色}}
\pretocmd{\l@subsection}{\hypersetup{linkcolor=GeometryGray}}{}{%
  \PackageWarning{geometry}{无法为目录小节条目着色}}
% 编的目录条目与 PDF 书签须用纯文本“第X编”。ctex 的 \CTEX@part@tocline
% 取 \CTEXthepart，而该宏带上编题页的颜色与字体设置，会把格式代码写进目录，
% 并使书签出现 “=GeometryGray第一编” 之类乱码。此覆盖须置于 \ctexset 之后，
% 否则会被 ctex 重新定义；\CTEXifname 是 ctex 的内部开关，此处沿用。
\def\CTEX@part@tocline#1#2{%
  \CTEXifname
    {第\chinese{part}编\hspace{1em}}%
    {}%
  #2}
\makeatother

\hypersetup{pdftitle={几何学·\BookVolumeName：\BookTitle},pdfauthor={\BookAuthor},pdfsubject={几何学书系框架预览}}

\begin{document}
\frontmatter
% !TeX root = ../Book3.tex
\MakeBookTitle[\BookTitle]
\clearpage
\thispagestyle{empty}
\null
\clearpage
% !TeX root = ../Book3.tex
\chapter*{序言}
\addcontentsline{toc}{chapter}{序言}
本卷为《几何学》\BookVolumeName《\BookTitle》的框架预览。目录依据《几何学计划纲要》建立，选读标记与原纲要一致。

待撰写

\MakeBookContents
% 有实际图、表或算法后再启用相应目录。
% \MakeBookListOfFigures
% \MakeBookListOfTables
% \MakeBookListOfAlgorithms


\mainmatter
% !TeX root = ../Book3.tex
% 纲要标题：# 第一编　度量、测地线与曲率结构
% 纲要位置：几何学计划纲要.md，第 2019 行。
% 正文状态：未撰写；本文件目前是框架。

% 纲要内容开始
% 
% 纲要内容结束

\BookPart{度量、测地线与曲率结构}{g3:part:01}
\BookPartGuide
待撰写

% !TeX root = ../../Book3.tex
% 纲要标题：## 第1章　黎曼度量与 Levi-Civita 联络
% 纲要位置：几何学计划纲要.md，第 2124 行。
% 正文状态：初稿；尚未正式审核。

% 纲要内容开始
% 
% > 本章目标与证明范围：从长度和度量相容性导出唯一无挠联络。公式必须能在坐标标架、正交标架和具体度量中互相核验；不把法坐标提前到指数映射之前。
% 
% 纲要内容结束

\chapter{黎曼度量与 Levi-Civita 联络}
\label{g3:ch:01}
\BookChapterNumber{g3:ch:01}

\begin{chapterabstract}
	光滑结构允许求导，黎曼度量进一步规定切向量的长度和夹角。本章由局部正定矩阵构造整体度量，建立梯度、体积密度、散度及分部积分公式；再从度量相容和无挠两个要求导出 Koszul 公式，证明 Levi-Civita 联络的存在唯一性。坐标标架与正交标架给出同一联络的两种表达，第一结构方程把无挠条件写成微分形式等式。诱导度量、乘积、共形变换和翘曲积展示这些公式的计算方式，二维模型与常数缩放则确定本卷的曲率符号和尺度约定。
\end{chapterabstract}

\begin{learninggoals}
	\begin{itemize}
		\item 从局部度量矩阵计算长度、梯度、体积密度与散度，并说明换坐标时这些几何量为何保持不变。
		\item 从度量相容和无挠条件推导 Koszul 公式，逐项验证它定义联络所需的线性及乘积法则。
		\item 在坐标标架和正交标架之间转换联络系数，用第一结构方程核验计算。
		\item 计算诱导、共形和翘曲积度量的联络，比较常数缩放对长度、体积和曲率的影响。
	\end{itemize}
\end{learninggoals}

在平面中，同一条曲线可以用直角坐标表示，也可以用极坐标表示。若曲线写成 \((r(t),\theta(t))\)，其速度的长度平方是
\[
	\dot r(t)^2+r(t)^2\dot\theta(t)^2.
\]
角度变化相同，在较大的圆周上走过的弧长就较长。这里的系数 \(r^2\) 记录坐标方向的实际尺度；它并不表示平面本身发生了弯曲。推广到流形时，我们在每个切空间上指定内积，并要求其系数随位置光滑变化。这使长度的定义脱离某一次坐标选择，也使没有预先嵌入空间的流形能够讨论体积和梯度。

\begin{figure}[htbp]
	\centering
	% !TeX root = ../Book3.tex
\begin{tikzpicture}[scale=1.15,>=stealth]
	\draw[->,gray] (-0.25,0)--(3.65,0) node[right]{\(x\)};
	\draw[->,gray] (0,-0.25)--(0,3.85) node[above]{\(y\)};
	\draw[gray!55] (0:1.3) arc[start angle=0,end angle=87,radius=1.3];
	\draw[gray!55] (0:2.6) arc[start angle=0,end angle=87,radius=2.6];
	\draw[dashed,gray] (0,0)--(30:3.6);
	\draw (0.53,0) arc[start angle=0,end angle=30,radius=0.53];
	\node at (15:0.77){\(\theta\)};
	\coordinate (P) at (30:1.3);
	\coordinate (Q) at (30:2.6);
	\fill (P) circle (1.5pt) node[below right]{\(p\)};
	\fill (Q) circle (1.5pt) node[below right]{\(q\)};
	\draw[->,thick,blue!65!black] (P)--++(30:0.8) node[midway,below right]{\(\partial_r\)};
	\draw[->,thick,blue!65!black] (Q)--++(30:0.8) node[above right]{\(\partial_r\)};
	\draw[->,thick,red!65!black] (P)--++(120:1.04) node[above left]{\(\partial_\theta|_p\)};
	\draw[->,thick,red!65!black] (Q)--++(120:2.08) node[above right]{\(\partial_\theta|_q\)};
	\node[below left] at (0,0){\(O\)};
\end{tikzpicture}

	\caption[极坐标方向的长度]{极坐标方向的长度。图中 \(q\) 的半径是 \(p\) 的两倍；用相同绘图比例表示向量时，两个径向坐标向量长度相同，角向坐标向量的长度则相差两倍。}
	\label{g3:fig:polar-coordinate-vectors}
\end{figure}

\cref{g3:fig:polar-coordinate-vectors}显示了度量如何记录坐标方向的尺度。有了长度，还需要比较邻近点的向量。在平面的极坐标标架中，径向和角向单位向量都随角度转动，所以只对向量的分量求导会漏掉标架本身的变化。联络把这部分变化计入导数。我们将要求这种导数满足内积的乘积法则，并使两个方向求导的反对称部分与向量场的括号一致。两个要求共同确定一个联络，其坐标表达和标架表达可以互相检验。

\ProseParagraphBreak

本章使用第一卷第11章的坐标图与切空间、第12章的张量和光滑单位分解，以及第14章的微分形式、密度积分和 Stokes 公式；每次实际调用时会恢复所需条件。\footnote{George Gabriel Stokes（1819--1903），生于爱尔兰的英国数学家、物理学家，研究流体力学与光学。}第一卷第15章学过的一般联络可帮助理解这些构造，本章仍给出切丛联络的定义和度量决定联络的证明。阅读本章无需先学第二卷的同调与同伦理论。度量给出曲线长度，长度的下确界如何成为距离、测地线如何产生法坐标，将在\chapref{g3:ch:02}展开。

\ProseParagraphBreak

以下光滑指 \(C^\infty\)，流形取有限维、Hausdorff、第二可数；除讨论边界积分时另作说明外，取实维数 \(n\ge1\) 且无边界。度量不要求可定向、紧或完备。坐标写为 \(x^1,\ldots,x^n\)，\(\partial_i\) 表示坐标向量场；同一单项式中重复的一个上指标和一个下指标表示从 \(1\) 到 \(n\) 求和。需要对照其他表述时，可参看 Petersen 第3版第1.1--1.2节和第1.4.2--1.4.4节的度量与标架、第2.2节和第2.4节的联络，以及第4.2.1节、第4.2.3节和第4.3.1节的球面与翘曲积模型\cite[第1.1--1.2节; 第1.4.2--1.4.4节; 第2.2节; 第2.4节; 第4.2.1节; 第4.2.3节; 第4.3.1节]{Petersen2016}。

% !TeX root = ../../Book3.tex
% 纲要标题：### 1.1 度量与几何量
% 纲要位置：几何学计划纲要.md，第 2128 行。
% 正文状态：初稿；尚未正式审核。

% 纲要内容开始
% 
% * 度量的存在及局部矩阵
% * 音乐同构、梯度和对偶度量
% * 体积密度、散度及分部积分
% * 坐标变化下的自然性
% 
% 纲要内容结束

\section{度量与几何量}
\label{g3:sec:0101}

一张坐标图能把曲线写成若干实函数，却没有规定哪个方向长、哪个方向短。即使在同一个平面上，\(\mathrm{d}x^2+\mathrm{d}y^2\) 与 \(4\,\mathrm{d}x^2+\mathrm{d}y^2\) 也给同一水平线段不同的长度。度量把这种测量规则放在每个切空间中，并要求它随点光滑变化。本节先说明度量怎样产生长度、梯度和体积，再由体积的变化定义散度；这些构造只需切丛、微分和积分，尚不需要比较不同点的向量。

\subsection{度量的存在及局部矩阵}
\label{g3:subsec:0101:01}

\begin{definition}
	\label{g3:def:riemannian-metric}
	设 \(M\) 是有限维、Hausdorff、第二可数的光滑流形。若光滑对称二阶协变张量场 \(g\) 在每个 \(p\in M\) 上给出正定双线性型 \(g_p\)，则称 \(g\) 为 \(M\) 上的\term[limanduliang]{黎曼度量}[Riemannian metric]，称 \((M,g)\) 为\term[limanliuxing]{黎曼流形}[Riemannian manifold]。向量的长度与非零向量之间的角度由
	\[
		|v|_g=\sqrt{g_p(v,v)},\qquad
		\cos\angle_g(v,w)=\frac{g_p(v,w)}{|v|_g|w|_g}
	\]
	确定。
\end{definition}

这个名称纪念 Riemann。\footnote{Riemann，德文全名 Georg Friedrich Bernhard Riemann（1826--1866），德国数学家，在复分析、几何及分析学中作出基础性贡献。}本章讨论 \(n\geq1\) 维光滑流形，通常无边界；涉及边界时另行说明。取局部坐标 \((x^1,\ldots,x^n)\)，写 \(\partial_i=\partial/\partial x^i\)。度量具有表达
\begin{equation}
	\label{g3:eq:metric-local-matrix}
	g=g_{ij}\,\mathrm{d}x^i\otimes\mathrm{d}x^j,\qquad
	g_{ij}=g(\partial_i,\partial_j).
\end{equation}
这里及后文对重复的一上一下指标求和；完整求和有时仍写出以便辨认。矩阵 \(G=(g_{ij})\) 对称正定，各分量光滑。反之，一组这样的局部矩阵，只要在交叠处遵守张量变换规律，就确定一个度量。\(\mathrm{d}x^i\mathrm{d}x^j\) 等常见的线元简写在本章均按相应对称双线性型理解；特别地，\(2g_{12}\,\mathrm{d}x^1\mathrm{d}x^2\) 包含两个对称交叉项。

\begin{proposition}
	\label{g3:prop:metric-existence}
	每个有限维、Hausdorff、第二可数的光滑流形都存在黎曼度量，允许流形有边界。若 \(g_0,g_1\) 是两个这样的度量，则 \((1-t)g_0+tg_1\) 在 \(0\leq t\leq1\) 时仍是黎曼度量。
\end{proposition}
\begin{proof}
	第一卷第12.3节的光滑单位分解定理（定理 12.3.5）说明：对任意坐标覆盖 \(\{U_\alpha\}\)，可取非负光滑函数 \(\rho_\alpha\)，使支集族局部有限、\(\operatorname{supp}\rho_\alpha\subset U_\alpha\)，且 \(\sum_\alpha\rho_\alpha=1\)。在 \(U_\alpha\) 上把坐标标架规定为正交单位标架，得到局部度量 \(h_\alpha\)。令
	\[
		g=\sum_\alpha\rho_\alpha h_\alpha,
	\]
	其中每项在 \(U_\alpha\) 外补零。支集包含关系使补零光滑，局部有限性使和在每点附近只有有限项。对非零 \(v\in\mathrm{T}_pM\)，各非零权重对应的 \(h_\alpha(v,v)\) 都为正，又至少有一个权重为正，故 \(g_p(v,v)>0\)。这正是第一卷丛度量构造在切丛上的应用。最后，两个正数的非负凸组合仍正，直接给出度量族的正定性。
\end{proof}

存在性依靠正定型的凸性，并没有给流形指定一个首选度量。度量一旦选定，曲线的长度便有确定含义。

\begin{definition}
	\label{g3:def:riemannian-length}
	设 \((M,g)\) 为黎曼流形，\(\gamma:[a,b]\to M\) 是分片 \(C^1\) 曲线。其\term[quxianchangdu]{长度}[length]定义为
	\[
		L_g(\gamma)=\int_a^b|\dot\gamma(t)|_g\,\mathrm{d}t
		=\int_a^b\sqrt{g_{ij}(\gamma(t))\dot x^i(t)\dot x^j(t)}\,\mathrm{d}t,
	\]
	右侧坐标表达在曲线的各个坐标片上分别计算并相加。
\end{definition}

被积量是切向量的内在长度，因而换坐标不改变结果。若 \(t=\phi(s)\) 是区间之间的 \(C^1\) 微分同胚，则链式法则给出 \(|(\gamma\circ\phi)'(s)|_g=|\phi'(s)|\,|\dot\gamma(\phi(s))|_g\)，换元后得到长度不变；递增与递减参数均可。正定性保证非零速度的长度为正。用所有连接两点的曲线长度取下确界会得到内在距离，其正定性与距离拓扑将在\chapref{g3:ch:02}证明。

\begin{proposition}
	\label{g3:prop:metric-coordinate-bounds}
	设 \(g\) 为光滑流形上的黎曼度量，\(K\) 是包含于一张坐标图 \(U\) 的非空紧集。存在常数 \(0<c\leq C<\infty\)，使对 \(p\in K\) 和 \(v=v^i\partial_i\in\mathrm{T}_pM\)，有
	\[
		c\sum_i(v^i)^2\leq g_p(v,v)\leq C\sum_i(v^i)^2.
	\]
\end{proposition}
\begin{proof}
	函数 \((p,a)\mapsto g_{ij}(p)a^ia^j\) 在紧集 \(K\times S^{n-1}\) 上连续且处处为正，故取得正的最小值 \(c\) 和有限最大值 \(C\)。对非零坐标向量先除以其 Euclid 长度，再利用二次齐次性；零向量时两式均成立。\footnote{Euclid（\OriginalName{Εὐκλείδης}，拉丁转写 Eukleidēs，约公元前300年活动，确切生卒年不详），古希腊数学家，活动于 Alexandria，以《几何原本》对几何与数论的系统论述著称。}
\end{proof}

这个估计把几何长度与坐标长度联系起来：曲线留在 \(K\) 时，两种长度分别相差至多因子 \(\sqrt c\) 与 \(\sqrt C\)。紧性使常数能够统一选择；在整个非紧坐标域上，未必存在同一对常数。例如 \(\mathbb R\) 上的 \(g=e^{2x}\mathrm{d}x^2\) 处处正定，但相对于 \(\mathrm{d}x^2\) 既没有全局正下界，也没有全局有限上界。

\subsection{音乐同构、梯度和对偶度量}
\label{g3:subsec:0101:02}

微分 \(\mathrm{d}f_p\) 接受方向并返回函数的变化率，是协变向量；描述“增长最快的方向”却需要一个向量。度量通过内积配对把这两类对象联系起来。

\begin{definition}
	\label{g3:def:musical-isomorphisms}
	黎曼流形 \((M,g)\) 的\term[yinyuetonggou]{音乐同构}[musical isomorphisms]是互逆的光滑丛同构
	\[
		\flat_g:\mathrm{T}M\longrightarrow\mathrm{T}^*M,
		\quad v^{\flat_g}=g(v,\mathord\cdot),
		\qquad
		\sharp_g:\mathrm{T}^*M\longrightarrow\mathrm{T}M.
	\]
	\(\sharp_g\) 由 \(g(\alpha^{\sharp_g},v)=\alpha(v)\) 刻画。度量明确时省略符号中的下标 \(g\)。
\end{definition}

正定性使 \(v\mapsto g(v,\mathord\cdot)\) 单射，有限维且两边维数相同，所以它可逆。记 \((g^{ij})=G^{-1}\)，则
\begin{equation}
	\label{g3:eq:musical-coordinates}
	(v^\flat)_i=g_{ij}v^j,\qquad
	(\alpha^\sharp)^i=g^{ij}\alpha_j.
\end{equation}
逆矩阵可由余子式除以 \(\det G>0\) 表示，所以逆映射也光滑。名称中的两个符号分别来自降号与升号，恰好对应把分量的指标降低和升高；它们依赖 \(g\)，并非切空间与余切空间原有的识别。

\begin{definition}
	\label{g3:def:dual-metric-gradient}
	黎曼流形 \((M,g)\) 的\term[duiouduliang]{对偶度量}[dual metric]是余切丛上的内积
	\[
		g^{-1}(\alpha,\beta)=g(\alpha^\sharp,\beta^\sharp)
		=g^{ij}\alpha_i\beta_j.
	\]
	对光滑实函数 \(f\)，其\term[tidu]{梯度}[gradient]为 \(\operatorname{grad}_g f=(\mathrm{d}f)^\sharp\)，即满足
	\begin{equation}
		\label{g3:eq:gradient-definition}
		g(\operatorname{grad}_g f,X)=Xf,\qquad
		\operatorname{grad}_g f=g^{ij}(\partial_jf)\partial_i
	\end{equation}
	的唯一光滑向量场。
\end{definition}

\symidx{\(\flat_g,\sharp_g\)}
\symidx{\(\operatorname{grad}_g f\)}

在 \(\mathrm{d}f_p\ne0\) 处，对所有单位向量 \(v\) 有
\[
	\mathrm{d}f_p(v)=g_p(\operatorname{grad}_g f,v)
	\leq|\operatorname{grad}_g f|_g,
\]
等号恰在 \(v=\operatorname{grad}_g f/|\operatorname{grad}_g f|_g\) 时成立。这来自内积不等式及其等号条件，也说明“最快”是以 \(g\) 的单位长度为尺度。若 \(M\) 无边界，且 \(\mathrm{d}f\) 在非空水平集 \(S=f^{-1}(c)\) 上处处非零，则第一卷第11.4节的正则水平集定理（定理 11.4.4）给出 \(\mathrm{T}_pS=\ker\mathrm{d}f_p\)，故梯度与水平集正交；在临界点，梯度为零，不再给出一个单位法向。

\begin{example}
	\label{g3:exm:gradient-metric-dependence}
	在 \(\mathbb R^2\) 上取 \(g=4\,\mathrm{d}x^2+\mathrm{d}y^2\) 与 \(f=x+y\)。则
	\[
		\mathrm{d}f=\mathrm{d}x+\mathrm{d}y,\qquad
		\operatorname{grad}_g f=\tfrac14\partial_x+\partial_y,
		\qquad |\mathrm{d}f|_{g^{-1}}^2=\tfrac54.
	\]
	水平线 \(x+y=c\) 的切向量 \(\partial_x-\partial_y\) 与这个梯度的 \(g\) 配对为 \(4\cdot\tfrac14-1=0\)。若换成 Euclid 度量，梯度变为 \(\partial_x+\partial_y\)，而 \(\mathrm{d}f\) 本身不变。逆矩阵出现在梯度公式中，正是为了补偿各方向长度尺度的不同。
\end{example}

同一个度量还在各阶张量丛上给出内积。例如对二阶协变张量 \(A\)，有 \(|A|_g^2=g^{ik}g^{j\ell}A_{ij}A_{k\ell}\)；对微分 \(k\) 形式，若写成 \(\alpha=\frac1{k!}\alpha_{i_1\cdots i_k}\mathrm{d}x^{i_1}\wedge\cdots\wedge\mathrm{d}x^{i_k}\)，则
\[
	|\alpha|_g^2=\frac1{k!}
	g^{i_1j_1}\cdots g^{i_kj_k}
	\alpha_{i_1\cdots i_k}\alpha_{j_1\cdots j_k}.
\]
因子 \(1/k!\) 消去同一有序指标组的排列重复，使正交余标架中递增指标的楔积成为正交单位基。

\subsection{体积密度、散度及分部积分}
\label{g3:subsec:0101:03}

切空间中的内积还确定平行多面体的体积。对切向量组 \((v_1,\ldots,v_n)\)，其体积为内积矩阵 \((g(v_i,v_j))\) 行列式的非负平方根。换基矩阵若为 \(A\)，内积矩阵变为 \(A^{\mathsf T}GA\)，体积就乘以 \(|\det A|\)。这提示应使用绝对行列式变换的密度，而不必先选择定向。

\begin{definition}
	\label{g3:def:metric-volume-density}
	\(n\) 维黎曼流形 \((M,g)\) 的\term[tijimidu]{体积密度}[volume density]记为 \(\mu_g\)，在坐标中写为
	\begin{equation}
		\label{g3:eq:volume-density}
		\mu_g=\sqrt{\det G}\,|\mathrm{d}x^1\cdots\mathrm{d}x^n|.
	\end{equation}
	若另外给定定向，则在正向坐标中
	\[
		\operatorname{vol}_g=\sqrt{\det G}\,\mathrm{d}x^1\wedge\cdots\wedge\mathrm{d}x^n
	\]
	定义相应的\term[tijixingshi]{体积形式}[volume form]。
\end{definition}

\symidx{\(\mu_g\)}

密度符号中的竖线表示换坐标时取绝对行列式，不是对微分符号求绝对值。第一卷第14.3节的积分良定义定理（定理 14.3.4）已证明：将紧支光滑密度用单位分解拆到坐标片后，按通常多元积分求和，结果与坐标和分解无关。因此 \(\int_M f\mu_g\) 对紧支光滑 \(f\) 有定义，不要求 \(M\) 可定向。一般非负可测函数的体积积分可由这一光滑正密度所给的局部测度定义；本章的分部积分只用紧支光滑情形。定向反转会改变 \(\operatorname{vol}_g\) 的符号，却不改变 \(\mu_g\)。

\ProseParagraphBreak

散度测量向量场的流对体积的一阶改变。设 \(X\) 是无边界流形上的光滑向量场。第一卷第13.1节的最大局部流定理（定理 13.1.7）给出随 \((t,p)\) 光滑变化的局部微分同胚 \(\Phi_t\)，满足 \(\Phi_0=\operatorname{id}\) 与 \(\partial_t\Phi_t=X\circ\Phi_t\)。密度的\term[liedaoshu]{Lie 导数}[Lie derivative]定义为 \(\mathcal L_X\mu_g=\left.\frac{\mathrm{d}}{\mathrm{d}t}\right|_{t=0}\Phi_t^*\mu_g\)。\footnote{Lie，Marius Sophus Lie（1842--1899），挪威数学家，创立连续变换群理论，并研究其与微分方程的联系。}这里只需在每一点附近对足够小的 \(t\) 定义流，无需假设流对所有时间存在。由于 \(\mu_g\) 处处正，存在唯一光滑函数满足
\begin{equation}
	\label{g3:eq:divergence-definition}
	\mathcal L_X\mu_g=(\operatorname{div}_gX)\mu_g.
\end{equation}
这个函数称为 \(X\) 的\term[sandu]{散度}[divergence]。

\begin{proposition}
	\label{g3:prop:divergence-coordinate}
	在无边界黎曼流形 \((M,g)\) 的任一坐标图中，对光滑向量场 \(X=X^i\partial_i\)，有
	\begin{equation}
		\label{g3:eq:divergence-coordinate}
		\operatorname{div}_gX
		=\frac1{\sqrt{\det G}}\partial_i\bigl(\sqrt{\det G}\,X^i\bigr).
	\end{equation}
	因此，对光滑实函数 \(f\)，有
	\begin{equation}
		\label{g3:eq:divergence-product}
		\operatorname{div}_g(fX)=\mathrm{d}f(X)+f\operatorname{div}_gX.
	\end{equation}
	公式也唯一地延伸到带边界流形上光滑至边界的向量场。
\end{proposition}
\begin{proof}
	记 \(m=\sqrt{\det G}\)。局部流拉回密度的系数为 \(m(\Phi_t(x))|\det\mathrm{D}\Phi_t(x)|\)。在 \(t=0\) 时，\(\Phi_0=\operatorname{id}\)，且 \(\left.\partial_t\Phi_t\right|_0=X\)。行列式在单位矩阵处的一阶项是迹，故
	\[
		\left.\frac{\mathrm{d}}{\mathrm{d}t}\right|_0
		\bigl(m(\Phi_t(x))|\det\mathrm{D}\Phi_t(x)|\bigr)
		=X^i\partial_i m+m\partial_iX^i=\partial_i(mX^i).
	\]
	除以 \(m>0\) 即得坐标公式；对 \(fX\) 展开偏导数得到乘积公式。若流形有边界，坐标系数在半空间边界处仍光滑；不同边界图中的表达在稠密的内部相同，连续性使它们在边界也相同。这样定义边界处的散度，无须要求 \(X\) 的双侧流保留流形。
\end{proof}

在定向坐标片中，直接将 \(\iota_X\operatorname{vol}_g\) 写成缺少一个坐标微分的各项可得
\begin{equation}
	\label{g3:eq:divergence-exterior}
	\mathrm{d}(\iota_X\operatorname{vol}_g)
	=(\operatorname{div}_gX)\operatorname{vol}_g.
\end{equation}
这里 \(\iota_X\) 将 \(X\) 插入形式的第一个位置；各项的排列符号恰好使外微分的系数成为 \(\partial_i(mX^i)\)。由此可以把散度积分转成边界通量。

\ProseParagraphBreak

若 \(M\) 有光滑边界，限制 \(g\) 到 \(\mathrm{T}\partial M\) 就给出边界度量，记它的体积密度为 \(\mu_{\partial g}\)。在边界坐标中取定义函数 \(r\)，使 \(M\) 对应 \(r\geq0\)、边界对应 \(r=0\)，且边界上 \(\mathrm{d}r\ne0\)。于是 \(\nu=-\operatorname{grad}_g r/|\operatorname{grad}_g r|_g\) 是外单位法向。它由正交、单位长度和向外这三个条件唯一确定，所以不同边界图中的选择相容。

\begin{theorem}[散度定理]
	\label{g3:thm:divergence-theorem}
	设 \((M^n,g)\) 为有限维、Hausdorff、第二可数的光滑黎曼流形，\(n\geq1\)，允许有光滑边界，无须可定向。若 \(X\) 是光滑至边界的紧支向量场，\(\nu\) 是边界的外单位法向，\(\mu_{\partial g}\) 是边界诱导度量的体积密度，则
	\begin{equation}
		\label{g3:eq:divergence-theorem}
		\int_M\operatorname{div}_gX\,\mu_g
		=\int_{\partial M}g(X,\nu)\,\mu_{\partial g}.
	\end{equation}
	无边界时右侧为零；\(n=1\) 时边界积分按各端点的计数密度求和。
\end{theorem}
\begin{proof}
	先假设 \(X\) 的支集位于一张坐标片 \(U\) 内并紧含于该片。给 \(U\) 任选定向，对紧支形式 \(\iota_X\operatorname{vol}_g\) 使用第一卷第14.4节的 Stokes 定理（定理 14.4.3）：外微分的积分等于在外法向优先定向下的边界积分。由\eqref{g3:eq:divergence-exterior}，左边就是散度乘以正体积密度的积分。

	当 \(n\geq2\) 时，在边界点取正交单位基 \((\nu,e_1,\ldots,e_{n-1})\)，使它符合 \(U\) 的定向。把 \(X\) 分成法向部分 \(g(X,\nu)\nu\) 和切向部分，切向部分与其余 \(n-1\) 个切向量线性相关，因而在顶次形式中贡献为零。故
	\[
		(\iota_X\operatorname{vol}_g)(e_1,\ldots,e_{n-1})=g(X,\nu).
	\]
	所以边界项正是 \(g(X,\nu)\mu_{\partial g}\) 的积分。\(n=1\) 时，在正向单位切向量为 \(e\) 的记号下，若 \(\nu=\varepsilon e\)，边界点的诱导定向符号就是 \(\varepsilon\)；其贡献为 \(\varepsilon\operatorname{vol}_g(X)=g(X,\nu)\)，结论相同。反转所选定向会同时反转体积形式与边界定向，两个密度积分都不变，所得局部公式不依赖这次选择。

	一般情形取从属于坐标覆盖的光滑单位分解 \(\{\rho_\alpha\}\)。紧支性及局部有限性保证只有有限个 \(\rho_\alpha X\) 非零，且各项支集紧含于相应坐标片。逐项应用局部公式并相加。乘积公式给出
	\[
		\sum_\alpha\operatorname{div}_g(\rho_\alpha X)
		=\Bigl(\sum_\alpha\rho_\alpha\Bigr)\operatorname{div}_gX
		+X\Bigl(\sum_\alpha\rho_\alpha\Bigr)
		=\operatorname{div}_gX,
	\]
	边界通量的和也等于 \(g(X,\nu)\)。这证明全局公式，包括不可定向情形。
\end{proof}

把定理应用于 \(fX\)，得到分部积分的基本形式。若 \(f\) 光滑、\(X\) 光滑，且其中至少一个紧支，则
\begin{equation}
	\label{g3:eq:integration-by-parts}
	\int_M f\operatorname{div}_gX\,\mu_g
	=-\int_M g(\operatorname{grad}_g f,X)\,\mu_g
	+\int_{\partial M}f g(X,\nu)\,\mu_{\partial g}.
\end{equation}
在带边界情形，所有数据都要求光滑至边界。无边界、\(f\) 在边界为零，或 \(X\) 沿边界切向时，边界项消失。非紧流形上的支集条件则控制无穷远的通量，不能从公式中直接删除。

\begin{definition}
	\label{g3:def:laplace-beltrami}
	对黎曼流形上的光滑实函数 \(f\)，采用\term[laplacebeltramisuanzi]{Laplace--Beltrami 算子}[Laplace--Beltrami operator]的非负号约定
	\begin{equation}
		\label{g3:eq:laplace-beltrami}
		\Delta_g f\coloneqq-\operatorname{div}_g(\operatorname{grad}_g f)
		=-\frac1{\sqrt{\det G}}
		\partial_i\bigl(\sqrt{\det G}\,g^{ij}\partial_jf\bigr).
	\end{equation}
\end{definition}

\symidx{\(\operatorname{div}_g X\)}
\symidx{\(\Delta_g\)}

这个名称沿用 Laplace 和 Beltrami 的姓氏。\footnote{Laplace，法文全名 Pierre-Simon Laplace（1749--1827），法国数学家和天文学家，研究概率、势函数与天体力学。Beltrami，意大利文名 Eugenio Beltrami（1835--1900），意大利数学家，研究微分几何、非 Euclid 几何模型及微分算子。}具体说，无边界时，对 \(u,v\in C_c^\infty(M)\)，分部积分给出
\begin{equation}
	\label{g3:eq:laplace-energy}
	\int_M u\Delta_gv\,\mu_g
	=\int_M g(\operatorname{grad}_g u,\operatorname{grad}_g v)\,\mu_g
	=\int_M v\Delta_gu\,\mu_g.
\end{equation}
令 \(u=v\)，便有 \(\int_Mu\Delta_gu\,\mu_g\geq0\)。这说明本书的非负号约定指积分二次型非负，不是说每个非负函数的 \(\Delta_g\) 仍逐点非负。带边界时，对光滑至边界的紧支函数 \(u,v\)，相应公式为
\[
	\int_Mu\Delta_gv\,\mu_g
	=\int_M g(\operatorname{grad}_g u,\operatorname{grad}_g v)\,\mu_g
	-\int_{\partial M}u\partial_\nu v\,\mu_{\partial g},
\]
其中 \(\partial_\nu v=\mathrm{d}v(\nu)\)。例如两函数都在边界为零时，仍得到\eqref{g3:eq:laplace-energy}的对称关系。

\begin{example}
	\label{g3:exm:polar-metric-operators}
	在平面去掉原点后的一个角度坐标片上，Euclid 度量写为
	\[
		g=\mathrm{d}r^2+r^2\mathrm{d}\theta^2,\qquad
		\mu_g=r|\mathrm{d}r\,\mathrm{d}\theta|.
	\]
	因此
	\[
		\begin{aligned}
			\operatorname{grad}_g f
			&=f_r\partial_r+r^{-2}f_\theta\partial_\theta,\\
			\operatorname{div}_g X
			&=r^{-1}\partial_r(rX^r)+\partial_\theta X^\theta,\\
			\Delta_g f
			&=-f_{rr}-r^{-1}f_r-r^{-2}f_{\theta\theta}.
		\end{aligned}
	\]
	这里 \(X^\theta\) 是相对于 \(\partial_\theta\) 的系数；若改用单位向量 \(r^{-1}\partial_\theta\) 的系数，散度的角向项也必须相应改写。例如 \(X=\partial_r\) 满足 \(\operatorname{div}_gX=1/r\)：向外移动厚度相同的小扇环，其周长随半径增长。原点不属于这张坐标图，公式中的 \(1/r\) 不表示原平面的度量在原点有奇性。
\end{example}

\subsection{坐标变化下的自然性}
\label{g3:subsec:0101:04}

前面的公式都含矩阵逆、行列式和偏导数。它们能否给出同一个几何对象，取决于这些变化是否彼此抵消。先直接核对换坐标，再把结论表述为局部等距下的相容性。

\ProseParagraphBreak

在交叠坐标 \(x=x(y)\) 中，记 \(J^i{}_a=\partial x^i/\partial y^a\)。由 \(\partial/\partial y^a=J^i{}_a\partial_i\)，新矩阵满足
\begin{equation}
	\label{g3:eq:metric-coordinate-transformation}
	G_y=J^{\mathsf T}G_xJ,\qquad
	G_y^{-1}=J^{-1}G_x^{-1}(J^{\mathsf T})^{-1},\qquad
	\sqrt{\det G_y}=|\det J|\sqrt{\det G_x}.
\end{equation}
第一式是双线性型的换基，第二式保证升指标后得到向量的变换规律，第三式恰是密度的变换规律。因此 \(g^{-1}\)、梯度和 \(\mu_g\) 的局部表达在交叠处相容；散度由局部流定义，也与坐标选择无关。

\begin{definition}
	\label{g3:def:local-isometry}
	设 \((M,g)\)、\((N,h)\) 是黎曼流形。若光滑映射 \(F:M\to N\) 是局部微分同胚且 \(F^*h=g\)，即
	\[
		h_{F(p)}(\mathrm{d}F_pv,\mathrm{d}F_pw)=g_p(v,w),
	\]
	则称 \(F\) 为\term[jubudengju]{局部等距}[local isometry]。若它还是微分同胚，则称为\term[dengju]{等距}[isometry]。
\end{definition}

\begin{proposition}
	\label{g3:prop:metric-naturality}
	设 \(F:(M,g)\to(N,h)\) 是无边界黎曼流形间的局部等距，\(f\in C^\infty(N)\)。则
	\[
		F^*\mu_h=\mu_g,\qquad
		\mathrm{d}F\bigl(\operatorname{grad}_g(f\circ F)\bigr)
		=(\operatorname{grad}_h f)\circ F,
	\]
	且对任意分片 \(C^1\) 曲线 \(\gamma\)，有 \(L_h(F\circ\gamma)=L_g(\gamma)\)。若向量场 \(X\) 与 \(Y\) 满足 \(\mathrm{d}F(X)=Y\circ F\)，则
	\[
		\operatorname{div}_gX=(\operatorname{div}_hY)\circ F,
		\qquad \Delta_g(f\circ F)=(\Delta_h f)\circ F.
	\]
\end{proposition}
\begin{proof}
	密度等式由\eqref{g3:eq:metric-coordinate-transformation}的行列式关系得出，长度等式由 \(\mathrm{d}F\) 逐点保持速度的内积得出。对任意 \(v\in\mathrm{T}_pM\)，有
	\[
		\begin{aligned}
			h\bigl(\mathrm{d}F\operatorname{grad}_g(f\circ F),\mathrm{d}Fv\bigr)
			&=\mathrm{d}(f\circ F)(v)\\
			&=\mathrm{d}f(\mathrm{d}Fv)
			=h\bigl(\operatorname{grad}_h f,\mathrm{d}Fv\bigr).
		\end{aligned}
	\]
	因 \(\mathrm{d}F_p\) 满射且 \(h\) 非退化，得到梯度关系。

	最后把 \(F\) 限制在可逆的邻域中，在目标选择坐标 \(y\)，源上取 \(x=y\circ F\)。相关向量场在这两套对应坐标中的分量相同，\(G(x)=H(y)\) 也对应相同。因此\eqref{g3:eq:divergence-coordinate}在对应点给出同一实数，证明散度关系。将它用于刚得到的梯度关系，再带入 \(\Delta=-\operatorname{div}\operatorname{grad}\)，便得到最后一个等式。
\end{proof}

局部等距可以重复覆盖同一个区域。例如 \(\mathbb R\to S^1\)、\(t\mapsto(\cos t,\sin t)\)，对 \(\mathrm{d}t^2\) 与单位圆的诱导度量是局部等距，它保持每条参数曲线的长度，却不保持整个空间的总体积。积分换元还需整体条件：若 \(F\) 是等距、\(f\in C_c^\infty(N)\)，则第一卷定理 14.3.6 的密度换元给出 \(\int_M(f\circ F)\mu_g=\int_Nf\mu_h\)；若仅有局部等距，重叠覆盖的次数可能改变积分。矩阵变换、局部几何量和整体积分因而各自有清楚的适用范围。

% !TeX root = ../../Book3.tex
% 纲要标题：### 1.2 Levi-Civita 定理
% 纲要位置：几何学计划纲要.md，第 2135 行。
% 正文状态：未撰写；本文件目前是框架。

% 纲要内容开始
% 
% * 度量相容与无挠的独立要求
% * Koszul 公式的推导
% * 由公式构造联络并验证性质
% * 存在唯一性及 Christoffel 表达
% 
% 纲要内容结束

\section{Levi-Civita 定理}
\label{g3:sec:0102}

待撰写

\subsection{度量相容与无挠的独立要求}
\label{g3:subsec:0102:01}
待撰写

\subsection{Koszul 公式的推导}
\label{g3:subsec:0102:02}
待撰写

\subsection{由公式构造联络并验证性质}
\label{g3:subsec:0102:03}
待撰写

\subsection{存在唯一性及 Christoffel 表达}
\label{g3:subsec:0102:04}
待撰写
% 纲要条目排为编号小节标题，并进入目录与书签；说明段仍保留于纲要注释中。

% !TeX root = ../../Book3.tex
% 纲要标题：### 1.3 标架与结构方程
% 纲要位置：几何学计划纲要.md，第 2142 行。
% 正文状态：初稿；尚未正式审核。

% 纲要内容开始
% 
% * 正交标架和联络一形式
% * Cartan 第一结构方程
% * 曲率形式的预告与符号约定
% 
% 纲要内容结束

\section{标架与结构方程}
\label{g3:sec:0103}

坐标标架的优点是向量场两两交换，代价是度量矩阵往往随点变化。若在每点把基改成正交单位基，度量矩阵就恒为单位矩阵，但新的基向量场通常不再交换。联络一形式和结构方程把这两种变化一起记录下来，使度量的计算可以在适合几何形状的标架中进行。

\subsection{正交标架和联络一形式}
\label{g3:subsec:0103:01}

给定局部光滑标架 \((v_1,\ldots,v_n)\)，依次作正交化：
\[
	e_1=\frac{v_1}{|v_1|_g},\qquad
	w_j=v_j-\sum_{a<j}g(v_j,e_a)e_a,\qquad
	e_j=\frac{w_j}{|w_j|_g}.
\]
各 \(w_j\) 处处非零，因为原标架线性无关；平方范数为正的光滑函数，其平方根和倒数光滑。因此这给出光滑的\term[zhengjiaodanweibiaojia]{正交单位标架}[orthonormal frame]。局部存在这样的标架不要求流形可定向。

\begin{definition}\label{g3:def:connection-one-forms}
	设 \((M,g)\) 为光滑黎曼流形，\(\nabla\) 为其 Levi-Civita 联络。在开集 \(U\) 上取光滑行标架 \(e=(e_1,\ldots,e_n)\)，以
	\[
		\nabla_Xe_j=e_i\omega^i{}_j(X)
	\]
	定义\term[lianluoyixingshi]{联络一形式}[connection one-forms] \(\omega^i{}_j\)，并记矩阵 \(\omega=(\omega^i{}_j)\)。若 \(u\) 为分量列，则
	\begin{equation}\label{g3:eq:frame-connection}
		\nabla_X(eu)=e\bigl(X(u)+\omega(X)u\bigr),\qquad
		\nabla e=e\omega.
	\end{equation}
	对偶余标架写为列 \(\theta=(\theta^1,\ldots,\theta^n)^{\mathsf T}\)，满足 \(\theta^i(e_j)=\delta^i_j\)，故 \(X=e\theta(X)\)。
\end{definition}

上下指标的位置记录输入与输出：\(j\) 指被微分的标架向量，\(i\) 指结果的分量；求导方向留给一形式接受。在坐标标架中，\(\omega^k{}_j=\Gamma^k_{ij}\,\mathrm{d}x^i\)。这与第一卷第15章的行标架、列分量约定一致，该卷将一般丛的联络矩阵记作 \(A\)。

\begin{proposition}\label{g3:prop:frame-metric-and-change}
	设 \((M,g)\) 为光滑黎曼流形，\(\nabla\) 为与 \(g\) 相容的光滑联络。在局部标架 \(e\) 中记 \(G=(g(e_i,e_j))\)，联络矩阵为 \(\omega\)。则
	\[
		\mathrm{d}G=\omega^{\mathsf T}G+G\omega.
	\]
	特别地，在正交单位标架中，\(\omega^i{}_j=-\omega^j{}_i\)。若 \(e'=eQ\)，其中 \(Q:U\to\operatorname{O}(n)\) 为光滑正交矩阵函数，则
	\begin{equation}\label{g3:eq:orthonormal-frame-change}
		\omega'=Q^{-1}\omega Q+Q^{-1}\mathrm{d}Q,\qquad
		\theta'=Q^{-1}\theta.
	\end{equation}
\end{proposition}
\begin{proof}
	度量相容性给出
	\[
		\mathrm{d}\bigl(g(e_i,e_j)\bigr)
		=\omega^k{}_i g(e_k,e_j)+g(e_i,e_k)\omega^k{}_j,
	\]
	即所述矩阵等式。若 \(G=I\)，左侧为零，便得到斜对称性。换标架时，对 \(eQ\) 使用乘积规则：
	\[
		\nabla e'=e(\omega Q+\mathrm{d}Q)
		=e'Q^{-1}(\omega Q+\mathrm{d}Q).
	\]
	与 \(\nabla e'=e'\omega'\) 比较即得第一式。由 \(X=e\theta(X)=eQ\theta'(X)\) 得第二式。
\end{proof}

正交标架的度量系数虽为常数，联络系数仍可非零。若 \([e_i,e_j]=c^k{}_{ij}e_k\)，将三个标架向量代入 Koszul 公式，得到
\begin{equation}\label{g3:eq:orthonormal-koszul}
	2\omega^k{}_j(e_i)=c^k{}_{ij}-c^i{}_{jk}+c^j{}_{ki}.
\end{equation}
度量的变化现在由标架向量的括号表达。特别地，若正交单位标架的所有括号都为零，则此标架中的联络矩阵为零。

\subsection{Cartan 第一结构方程}
\label{g3:subsec:0103:02}

余标架 \(\theta\) 把实际切向量转为标架分量。它的外微分能检测标架的不交换性，因为
\[
	\mathrm{d}\theta^k(e_i,e_j)
	=e_i\bigl(\theta^k(e_j)\bigr)-e_j\bigl(\theta^k(e_i)\bigr)
	-\theta^k([e_i,e_j])=-c^k{}_{ij}.
\]
将这部分变化与联络一形式合并，就得到 Cartan 第一结构方程。\footnote{Élie Joseph Cartan（1869--1951），法国数学家，在 Lie 群与 Lie 代数、微分方程和微分几何之间建立了深刻联系，发展了移动标架与外微分方法。}

\begin{proposition}[Cartan 第一结构方程]\label{g3:prop:first-structure-equation}
	设 \(M\) 为光滑 \(n\) 维无边界流形，\(\nabla\) 为光滑切丛联络，\(e\) 为一个局部光滑标架，\(\theta\) 为其对偶余标架，\(\omega\) 为相应联络矩阵。定义挠率分量二形式 \(\Theta^i(X,Y)=\theta^i(T^\nabla(X,Y))\)，则
	\begin{equation}\label{g3:eq:first-structure-equation}
		\Theta^i=\mathrm{d}\theta^i+\omega^i{}_j\wedge\theta^j.
	\end{equation}
	特别地，黎曼度量的 Levi-Civita 联络在正交单位标架中满足
	\begin{equation}\label{g3:eq:levi-civita-structure-equations}
		\mathrm{d}\theta^i+\omega^i{}_j\wedge\theta^j=0,
		\qquad \omega^i{}_j+\omega^j{}_i=0.
	\end{equation}
	给定正交单位余标架 \(\theta\) 后，这两组等式唯一确定联络一形式 \(\omega\)。
\end{proposition}
\begin{proof}
	由式\eqref{g3:eq:frame-connection}，\(\nabla_XY\) 的第 \(i\) 个分量为
	\[
		X(\theta^i(Y))+\omega^i{}_j(X)\theta^j(Y).
	\]
	交换 \(X,Y\) 再减去 \([X,Y]\) 的分量，得到
	\begin{align*}
		\Theta^i(X,Y)
		={}&X(\theta^i(Y))-Y(\theta^i(X))-\theta^i([X,Y])\\
		&+\omega^i{}_j(X)\theta^j(Y)-\omega^i{}_j(Y)\theta^j(X).
	\end{align*}
	前三项是 \(\mathrm{d}\theta^i(X,Y)\)，后两项是 \((\omega^i{}_j\wedge\theta^j)(X,Y)\)，故第一式成立。无挠性与度量相容性给出式\eqref{g3:eq:levi-civita-structure-equations}。

	反过来，一组光滑一形式 \(\omega^i{}_j\) 通过式\eqref{g3:eq:frame-connection}定义局部联络：一形式对方向的线性和普通导数的乘积规则保证式\eqref{g3:eq:connection-rules}。若它满足所写的两组等式，第一式保证无挠，第二式保证保持 \(g=\sum_i\theta^i\otimes\theta^i\)。由\cref{g3:thm:levi-civita}的局部唯一性，它就是该区域上的 Levi-Civita 联络。
\end{proof}

\begin{example}\label{g3:exm:polar-moving-frame}
	继续\cref{g3:exm:polar-christoffel}的极坐标平面。在 \(r>0\) 的局部区域取
	\[
		e_1=\partial_r,\qquad e_2=r^{-1}\partial_\varphi,
		\qquad \theta^1=\mathrm{d}r,\quad\theta^2=r\,\mathrm{d}\varphi.
	\]
	这是正交单位标架，且 \([e_1,e_2]=-r^{-1}e_2\)。斜对称性使联络矩阵只有一个独立一形式：记 \(\omega^2{}_1=\alpha\)、\(\omega^1{}_2=-\alpha\)。第一结构方程为
	\[
		\mathrm{d}\theta^1-\alpha\wedge\theta^2=0,
		\qquad \mathrm{d}\theta^2+\alpha\wedge\theta^1=0.
	\]
	设 \(\alpha=a\,\mathrm{d}r+b\,\mathrm{d}\varphi\)。第一式给出 \(a=0\)，第二式给出 \(b=1\)，故
	\[
		\omega=
		\begin{pmatrix}0&-\mathrm{d}\varphi\\ \mathrm{d}\varphi&0\end{pmatrix}.
	\]
	于是 \(\nabla_{e_1}e_2=0\)、\(\nabla_{e_2}e_1=r^{-1}e_2\)、\(\nabla_{e_2}e_2=-r^{-1}e_1\)，与坐标系数换算所得相同。特别地，\(\mathrm{d}\theta^2=\mathrm{d}r\wedge\mathrm{d}\varphi\ne0\)，但它恰被 \(\omega^2{}_1\wedge\theta^1\) 抵消；无挠不要求各余标架分量都闭。
\end{example}

这个例子还可以直接通过换架得到。笛卡儿标架的联络矩阵为零，极坐标正交标架是将其旋转角度 \(\varphi\) 得到的，因此式\eqref{g3:eq:orthonormal-frame-change}给出 \(\omega=Q^{-1}\mathrm{d}Q\)。一般地，在定向黎曼曲面上，若正向正交单位标架旋转了光滑角度 \(\psi\)，同样计算可得
\begin{equation}\label{g3:eq:surface-connection-rotation}
	\alpha'=\alpha+\mathrm{d}\psi,
	\qquad \alpha=\omega^2{}_1.
\end{equation}
联络一形式随标架转动增加了角度的微分，而 \(\mathrm{d}\alpha\) 保持不变。这个二形式将记录曲面的曲率。

\subsection{曲率形式的预告与符号约定}
\label{g3:subsec:0103:03}

对向量场依次沿两个方向求协变导数，次序改变以后可以产生新的向量。方向场的括号本身也会贡献一次求导，因此定义
\begin{equation}\label{g3:eq:curvature-definition}
	R(X,Y)Z\coloneqq\nabla_X\nabla_YZ-\nabla_Y\nabla_XZ-\nabla_{[X,Y]}Z.
\end{equation}
这沿用第一卷第15.3节的丛联络曲率约定；第三章将系统讨论黎曼曲率的对称性与收缩。这里先得到计算模型所需的矩阵公式，并固定截面曲率的正负号。

\begin{proposition}\label{g3:prop:curvature-form}
	设 \(M\) 为光滑无边界流形，\(\nabla\) 为光滑切丛联络。式\eqref{g3:eq:curvature-definition}对 \(X,Y,Z\) 均为 \(C^\infty(M)\) 线性。在局部标架 \(e\) 中，若 \(\nabla e=e\omega\)，则
	\begin{equation}\label{g3:eq:curvature-matrix}
		\Omega=\mathrm{d}\omega+\omega\wedge\omega,
		\qquad
		R(X,Y)(eu)=e\Omega(X,Y)u,
	\end{equation}
	其中 \((\omega\wedge\omega)^i{}_j=\sum_k\omega^i{}_k\wedge\omega^k{}_j\)。换标架 \(e'=eQ\) 后，\(\Omega'=Q^{-1}\Omega Q\)。对黎曼度量的 Levi-Civita 联络，在正交单位标架中还有 \(\Omega^i{}_j=-\Omega^j{}_i\)。
\end{proposition}
\begin{proof}
	用式\eqref{g3:eq:frame-connection}计算两次导数，并将结果的分量写出。在 \(R(X,Y)(eu)\) 中，对 \(u\) 的二阶导数合成 \((XY-YX-[X,Y])u=0\)，含 \(X(u)\) 或 \(Y(u)\) 的交叉项也相消，剩下
	\[
		\bigl(X(\omega(Y))-Y(\omega(X))-\omega([X,Y])
		+\omega(X)\omega(Y)-\omega(Y)\omega(X)\bigr)u.
	\]
	前三项是 \(\mathrm{d}\omega(X,Y)\)，后两项是 \((\omega\wedge\omega)(X,Y)\)。于是得到式\eqref{g3:eq:curvature-matrix}，且右侧对三个输入都为函数线性，证明张量性。

	同一向量的分量满足 \(u=Qu'\)，由 \(e\Omega u=e'\Omega'u'\) 得 \(\Omega'=Q^{-1}\Omega Q\)。若 \(\omega\) 斜对称，则
	\[
		\Omega^j{}_i
		=-\mathrm{d}\omega^i{}_j
		+\sum_k\omega^k{}_j\wedge\omega^i{}_k
		=-\mathrm{d}\omega^i{}_j-\sum_k\omega^i{}_k\wedge\omega^k{}_j
		=-\Omega^i{}_j.
	\]
	第二个等号使用了一形式楔积的反交换性。
\end{proof}

矩阵乘法参与楔积，故 \(\omega\wedge\omega\) 一般不为零。另一方面，曲率没有换架公式中的附加微分项，因而它是在不同标架中表达同一个张量。极坐标平面的 \(\omega\) 非零，但 \(\mathrm{d}\omega=0\)、\(\omega\wedge\omega=0\)，所以 \(R=0\)。

\begin{definition}\label{g3:def:sectional-curvature}
	设 \((M,g)\) 为维数至少为二的光滑黎曼流形，\(R\) 为其 Levi-Civita 联络的曲率。对线性无关的 \(v,w\in\mathrm{T}_pM\)，定义
	\begin{equation}\label{g3:eq:sectional-curvature}
		K(v,w)=\frac{g(R(v,w)w,v)}{g(v,v)g(w,w)-g(v,w)^2}.
	\end{equation}
	它称为二维平面 \(\operatorname{span}\{v,w\}\) 的\term[jiemianqulv]{截面曲率}[sectional curvature]。
\end{definition}

分母是 \(v,w\) 张成平行四边形的面积平方。这个值确实只依赖二维平面：若 \(v'=av+bw\)、\(w'=cv+dw\)，则 \(R(v',w')=(ad-bc)R(v,w)\)。又因 \(R(v,w)\) 关于 \(g\) 反对称，分子和分母在换基后都乘以 \((ad-bc)^2\)，比值不变。

\ProseParagraphBreak

在定向曲面上取正向正交单位标架，记 \(\alpha=\omega^2{}_1\)。由于二维斜对称矩阵只有一个独立分量，\((\omega\wedge\omega)^2{}_1=0\)。又有
\[
	K=g(R(e_1,e_2)e_2,e_1)=\Omega^1{}_2(e_1,e_2),
\]
从而
\begin{equation}\label{g3:eq:surface-curvature-form}
	\mathrm{d}\alpha=\Omega^2{}_1=-K\,\theta^1\wedge\theta^2.
\end{equation}
这与第一卷第9章的 \(\mathrm{d}\omega=-K\,\mathrm{d}A\) 相同，其中该卷的单个一形式 \(\omega\) 对应这里的 \(\alpha\)。下一节将以旋转型度量计算正、负曲率模型，检验这个符号如何进入实际公式。

% !TeX root = ../../Book3.tex
% 纲要标题：### 1.4 度量的实际计算
% 纲要位置：几何学计划纲要.md，第 2045 行。
% 正文状态：未撰写；本文件目前是框架。

% 纲要内容开始
% 
% * 诱导度量和乘积度量
% * 共形变换的联络
% * 翘曲积与旋转模型
% * 不同曲率归一化的比较
% 
% ---
% 
% 纲要内容结束

\section{度量的实际计算}
\label{g3:sec:0104}

待撰写

\subsection{诱导度量和乘积度量}
\label{g3:subsec:0104:01}
待撰写

\subsection{共形变换的联络}
\label{g3:subsec:0104:02}
待撰写

\subsection{翘曲积与旋转模型}
\label{g3:subsec:0104:03}
待撰写

\subsection{不同曲率归一化的比较}
\label{g3:subsec:0104:04}
待撰写
% 纲要条目排为编号小节标题，并进入目录与书签；说明段仍保留于纲要注释中。


\begin{chaptersummary}
	黎曼度量逐点给出正定内积，单位分解把局部度量粘合为整体度量。音乐同构将一形式变为向量，因而把函数的微分变为梯度；行列式给出按绝对 Jacobian 变换的体积密度。散度描述密度沿流的变化，积分公式把它与边界通量相连，并确定非负 Laplace 算子 \(\Delta_g=-\operatorname{div}_g\operatorname{grad}_g\) 的符号。

	\ProseParagraphBreak

	度量相容保证微分与内积相容，无挠把联络的反对称部分固定为向量场括号。Koszul 公式同时落实这两个要求，既给出唯一性，也构造出所需联络。坐标中的 Christoffel 系数和正交标架中的反对称联络矩阵描述同一个对象；标架变化时出现的导数项，解释了为何联络系数本身不是张量。

	\ProseParagraphBreak

	诱导度量可用环境导数的切向投影计算，乘积度量把两个因子的联络分开。共形因子的微分以及翘曲函数的对数导数，则分别控制相应的联络修正。二维翘曲模型给出 \(K=-f''/f\)，而常数缩放 \(g\mapsto c^2g\) 保持联络并使截面曲率乘以 \(c^{-2}\)。这些计算将在测地线和曲率理论中继续作为可直接核验的模型。
\end{chaptersummary}

\BookExercises

\begin{exercise}\label{g3:ex:metric-convex-comparison}
	设 \(g,h\) 为同一 \(n\) 维光滑流形 \(M\) 上的黎曼度量。
	\begin{enumerate}
		\item 若 \(M\) 紧，证明存在常数 \(0<c\le C<\infty\)，使 \(c\,g(v,v)\le h(v,v)\le C\,g(v,v)\) 对一切切向量成立，并比较两种度量下的曲线长度、对偶度量和体积密度。
		\item 在 \(\mathbb R^2\) 上构造体积密度相同、却没有上述整体比较常数的两个度量。
	\end{enumerate}
\end{exercise}
\begin{solution}
	1. 每一点附近可取一个闭包紧含于坐标图的小邻域。在这个闭包与坐标单位球面的乘积上，连续正函数 \(h(v,v)/g(v,v)\) 有正的最小值和有限最大值。紧性使有限个这样的小邻域覆盖 \(M\)；取各下界的最小值、各上界的最大值，即得 \(c,C\)。对分片光滑曲线逐点开平方再积分，得到
	\[
		\sqrt c\,L_g(\gamma)\le L_h(\gamma)\le\sqrt C\,L_g(\gamma).
	\]
	逐点选 \(g\) 正交单位基，并将 \(h\) 的正定矩阵正交对角化，其特征值 \(\lambda_i\) 都落在 \([c,C]\) 内。对偶度量的特征值是 \(\lambda_i^{-1}\)，密度之比是 \(\sqrt{\lambda_1\cdots\lambda_n}\)，所以
	\[
		C^{-1}g^{-1}\le h^{-1}\le c^{-1}g^{-1},\qquad
		c^{n/2}\mu_g\le\mu_h\le C^{n/2}\mu_g.
	\]
	前一组不等式按二次型理解，后一组按正密度的逐点大小理解。

	2. 取 \(g=\mathrm dx^2+\mathrm dy^2\)、\(h=e^{2x}\mathrm dx^2+e^{-2x}\mathrm dy^2\)。二者的矩阵行列式均为一，体积密度相同；但在向量 \(\partial_x\) 上的比值为 \(e^{2x}\)，既无全局正下界，也无有限上界。
\end{solution}

\begin{exercise}\label{g3:ex:non-diagonal-gradient}
	在 \(\mathbb R^2\) 上令
	\[
		g=(\mathrm dx+x\,\mathrm dy)^2+\mathrm dy^2.
	\]
	证明它是黎曼度量，求其体积密度、\(\operatorname{grad}_g y\) 和 \(\Delta_g y\)。解释为何后一个结果不与紧流形上 Laplace 算子的积分性质矛盾。
\end{exercise}
\begin{solution}
	若 \(v=a\partial_x+b\partial_y\)，则 \(g(v,v)=(a+xb)^2+b^2\)，仅在 \(a=b=0\) 时为零。矩阵及逆矩阵为
	\[
		G=\begin{pmatrix}1&x\\x&1+x^2\end{pmatrix},\qquad
		G^{-1}=\begin{pmatrix}1+x^2&-x\\-x&1\end{pmatrix},\qquad
		\det G=1.
	\]
	所以 \(\mu_g=|\mathrm dx\,\mathrm dy|\)，且
	\[
		\operatorname{grad}_g y=-x\partial_x+\partial_y,\qquad
		\Delta_g y=-\bigl(\partial_x(-x)+\partial_y1\bigr)=1.
	\]
	这里底空间非紧，函数及其梯度都不紧支；无边界紧流形上的积分结论不能用于这组数据。体积密度与普通面积相同，也不意味着梯度相同。
\end{solution}

\begin{exercise}\label{g3:ex:conformal-laplacian}
	设 \((M^n,g)\) 为无边界黎曼流形，\(u,f\in C^\infty(M)\)，\(\widetilde g=e^{2u}g\)。证明
	\[
		\operatorname{grad}_{\widetilde g}f=e^{-2u}\operatorname{grad}_g f,
		\qquad \operatorname{div}_{\widetilde g}X=\operatorname{div}_gX+nX(u),
	\]
	以及
	\[
		\Delta_{\widetilde g}f
		=e^{-2u}\bigl(\Delta_g f-(n-2)g(\operatorname{grad}_g u,\operatorname{grad}_g f)\bigr).
	\]
	推得在二维时，方程 \(\Delta_g f=0\) 及紧支函数的能量 \(\int_M|\operatorname{grad}_g f|_g^2\mu_g\) 均保持共形不变。
\end{exercise}
\begin{solution}
	由梯度的定义，\(\widetilde g(e^{-2u}\operatorname{grad}_g f,Y)=Y(f)\)，得到第一式。行列式给出 \(\mu_{\widetilde g}=e^{nu}\mu_g\)，故散度的局部公式给出
	\[
		\operatorname{div}_{\widetilde g}X
		=\frac{1}{e^{nu}\sqrt{\det G}}\partial_i
		\bigl(e^{nu}\sqrt{\det G}\,X^i\bigr)
		=\operatorname{div}_gX+nX(u).
	\]
	再用 \(\operatorname{div}_g(aX)=a\operatorname{div}_gX+X(a)\)，得到
	\[
		-\operatorname{div}_{\widetilde g}(e^{-2u}\operatorname{grad}_g f)
		=e^{-2u}\Delta_g f+(2-n)e^{-2u}g(\operatorname{grad}_g u,\operatorname{grad}_g f).
	\]
	当 \(n=2\)，额外的一阶项消失，正因子 \(e^{-2u}\) 不改变方程的零解集。一般维数下，能量的被积密度变为
	\[
		|\operatorname{grad}_{\widetilde g}f|_{\widetilde g}^2\mu_{\widetilde g}
		=e^{(n-2)u}|\operatorname{grad}_g f|_g^2\mu_g;
	\]
	二维时因子等于一，紧支性保证积分有限。
\end{solution}

\begin{exercise}\label{g3:ex:compact-harmonic-functions}
	设 \((M,g)\) 是连通、紧且无边界的黎曼流形。证明每个满足 \(\Delta_g f=0\) 的光滑实函数都是常数；并证明方程 \(\Delta_g f=h\) 有光滑解的必要条件是 \(\int_M h\mu_g=0\)。给出非紧或有边界情形下第一条结论失效的例子。
\end{exercise}
\begin{solution}
	分部积分给出
	\[
		0=\int_M f\Delta_g f\,\mu_g
		=\int_M|\operatorname{grad}_g f|_g^2\mu_g.
	\]
	被积函数非负且连续，故处处为零。于是 \(\mathrm df=0\)，函数局部常值，再由连通性成为常数。将 \(\Delta_g f=h\) 积分，散度定理给出 \(\int_Mh\mu_g=0\)。这一步只得到必要条件。普通实线上 \(f(x)=x\) 是非常数调和函数；在带 Euclid 度量的紧区间 \([0,1]\) 上，同一函数仍满足内部的 \(-f''=0\)，其边界值没有被要求相同。
\end{solution}

\begin{exercise}\label{g3:ex:hessian-torsion}
	设 \(\nabla\) 为无边界光滑流形上的联络，\(T\) 为其挠率。对光滑实函数 \(f\) 定义双线性表达式
	\[
		H_f(X,Y)=X(Yf)-(\nabla_XY)f.
	\]
	证明它是光滑协变二阶张量，并求 \(H_f(X,Y)-H_f(Y,X)\)。推得：\(\nabla\) 无挠当且仅当每个局部光滑函数的 \(H_f\) 都对称。写出其坐标表达式。
\end{exercise}
\begin{solution}
	第一变量乘以光滑函数 \(a\) 时，两项均乘以 \(a\)。第二变量变为 \(aY\) 时，\(X(a)Yf\) 与联络乘积法则给出的同一项抵消，也只留下 \(aH_f(X,Y)\)。因此两变量均为 \(C^\infty\) 线性。相减得到
	\[
		H_f(X,Y)-H_f(Y,X)
		=[X,Y]f-(\nabla_XY-\nabla_YX)f
		=-\mathrm df\bigl(T(X,Y)\bigr).
	\]
	无挠即推出对称。反过来，每一点处的余向量都能写成某个局部线性坐标函数的微分；若右侧对所有这样的函数为零，则 \(T(X,Y)=0\)。坐标中
	\[
		(H_f)_{ij}=\partial_i\partial_jf-\Gamma^k{}_{ij}\partial_kf.
	\]
	当 \(\nabla\) 为度量的 Levi-Civita 联络时，由内积乘积法则还得 \(H_f(X,Y)=g(\nabla_X\operatorname{grad}_g f,Y)\)。
\end{solution}

\begin{exercise}\label{g3:ex:metric-connection-three-form}
	设 \((M,g)\) 为无边界黎曼流形，\(\nabla\) 为其 Levi-Civita 联络，\(H\) 为光滑三形式。用
	\[
		g(A_XY,Z)=\tfrac12H(X,Y,Z)
	\]
	定义 \(A\)，并令 \(D_XY=\nabla_XY+A_XY\)。证明 \(D\) 是度量相容联络，且其挠率满足 \(g(T^D(X,Y),Z)=H(X,Y,Z)\)。证明它是唯一具有这两项性质的联络。
\end{exercise}
\begin{solution}
	由度量的非退化性，\(A_XY\) 唯一确定且光滑；三形式的张量性使 \(A\) 对 \(X,Y\) 均为 \(C^\infty\) 线性，故 \(D\) 满足联络法则。由 \(H(X,Y,Z)=-H(X,Z,Y)\)，有
	\[
		g(A_XY,Z)+g(Y,A_XZ)=0,
	\]
	所以给 \(\nabla\) 加上 \(A\) 仍保持度量相容。又因 \(\nabla\) 无挠，
	\[
		g(T^D(X,Y),Z)
		=g(A_XY-A_YX,Z)=H(X,Y,Z).
	\]
	若 \(D'\) 也有这些性质，记 \(B_XY=D'_XY-D_XY\) 和 \(b(X,Y,Z)=g(B_XY,Z)\)。相同挠率给出前两变量对称，度量相容给出后两变量反对称。依次交换变量，得到
	\[
		\begin{aligned}
			b(X,Y,Z)&=b(Y,X,Z)=-b(Y,Z,X)\\
			&=-b(Z,Y,X)=b(Z,X,Y)=b(X,Z,Y)\\
			&=-b(X,Y,Z).
		\end{aligned}
	\]
	所以 \(b=0\)，由非退化性 \(B=0\)，即 \(D'=D\)。
\end{solution}

\begin{exercise}\label{g3:ex:same-connection-two-metrics}
	设 \(g,h\) 是同一无边界光滑流形上的黎曼度量，\(\nabla\) 为 \(g\) 的 Levi-Civita 联络。逐点定义 \(S\) 使 \(h(X,Y)=g(SX,Y)\)。证明 \(S\) 是光滑、正定且关于 \(g\) 自伴的切丛自同态，并证明两种度量有相同的 Levi-Civita 联络当且仅当
	\[
		(\nabla_XS)Y\coloneqq\nabla_X(SY)-S(\nabla_XY)=0
	\]
	对所有 \(X,Y\) 成立。举出两种并非常数倍、却具有相同联络的度量。
\end{exercise}
\begin{solution}
	矩阵表达为 \(S=G^{-1}H\)，因此光滑。由 \(h\) 对称得到 \(g(SX,Y)=g(X,SY)\)，由 \(h(X,X)>0\) 得到正定性。使用 \(\nabla g=0\) 对定义式求导，得
	\[
		Xh(Y,Z)-h(\nabla_XY,Z)-h(Y,\nabla_XZ)
		=g((\nabla_XS)Y,Z).
	\]
	右侧恒零等价于 \(\nabla h=0\)。\(\nabla\) 已无挠，所以由存在唯一性定理，这又等价于 \(\nabla\) 是 \(h\) 的 Levi-Civita 联络。在 \(\mathbb R^2\) 上取 \(g=\mathrm dx^2+\mathrm dy^2\)、\(h=2\mathrm dx^2+3\mathrm dy^2\)，两者系数均为常数，联络系数都为零，但并非常数倍。
\end{solution}

\begin{exercise}\label{g3:ex:shear-coordinates}
	对微分同胚 \(F:\mathbb R^2\to\mathbb R^2\)、\(F(u,v)=(u,v+u^2)\)，令 \(g=F^*(\mathrm dx^2+\mathrm dy^2)\)。计算 \(g\) 的矩阵及全部非零 Christoffel 系数，再在标架
	\[
		e_1=\partial_u-2u\partial_v,\qquad e_2=\partial_v
	\]
	中计算联络形式和曲率形式。
\end{exercise}
\begin{solution}
	\(F^{-1}(x,y)=(x,y-x^2)\)，且 \(\mathrm dx=\mathrm du\)、\(\mathrm dy=\mathrm dv+2u\,\mathrm du\)，所以
	\[
		G=\begin{pmatrix}1+4u^2&2u\\2u&1\end{pmatrix},\qquad \det G=1.
	\]
	用 Christoffel 公式得到唯一非零系数 \(\Gamma^v{}_{uu}=2\)。也可利用 \(\mathrm dF(\partial_u)=\partial_x+2u\partial_y\)、\(\mathrm dF(\partial_v)=\partial_y\)：对前者再沿 \(\partial_u\) 求环境导数，得到 \(2\partial_y\)，其余方向导数为零。

	\(\mathrm dF(e_1)=\partial_x\)、\(\mathrm dF(e_2)=\partial_y\)，故 \((e_1,e_2)\) 正交归一，且每个 \(e_i\) 在所有方向上都平行。因此 \(\omega=0\)、\(\Omega=0\)。直接核对时，\(\nabla_{\partial_u}(\partial_u-2u\partial_v)=2\partial_v-2\partial_v=0\)。这个计算展示了非零坐标联络系数与零曲率可以同时出现。
\end{solution}

\begin{exercise}\label{g3:ex:rotating-frame}
	在带普通度量的平面开集 \(U\) 上取任意光滑函数 \(\varphi\)，令
	\[
		e_1=\cos\varphi\,\partial_x+\sin\varphi\,\partial_y,\qquad
		e_2=-\sin\varphi\,\partial_x+\cos\varphi\,\partial_y.
	\]
	求对偶余标架、联络矩阵和曲率矩阵。判断这组标架何时能在每一点附近成为某组坐标的坐标标架。
\end{exercise}
\begin{solution}
	对偶余标架为 \(\theta^1=\cos\varphi\,\mathrm dx+\sin\varphi\,\mathrm dy\)、\(\theta^2=-\sin\varphi\,\mathrm dx+\cos\varphi\,\mathrm dy\)。逐项求导得 \(\nabla e_1=e_2\mathrm d\varphi\)、\(\nabla e_2=-e_1\mathrm d\varphi\)，故
	\[
		\omega=\begin{pmatrix}0&-\mathrm d\varphi\\\mathrm d\varphi&0\end{pmatrix},\qquad
		\Omega=\mathrm d\omega+\omega\wedge\omega=0.
	\]
	无挠性给出
	\[
		[e_1,e_2]=-e_1(\varphi)e_1-e_2(\varphi)e_2.
	\]
	若为坐标标架，其括号必须为零，故 \(\mathrm d\varphi=0\)。反之，\(\varphi\) 局部常值时，固定旋转后的直角坐标便给出所需坐标标架。因此条件是 \(\varphi\) 在每个连通分支上为常数；任意位置依赖的旋转仍给出正交标架，却一般不是坐标标架。
\end{solution}

\begin{exercise}\label{g3:ex:warped-radial-laplacian}
	设 \((F^m,h)\) 是无边界黎曼流形，\(I\subset\mathbb R\) 是开区间，\(f:I\to(0,\infty)\) 光滑。在 \(I\times F\) 上取 \(g=\mathrm dr^2+f(r)^2h\)。对只依赖 \(r\) 的函数 \(v\)，求 \(\Delta_gv\)，并求出所有满足 \(\Delta_gv=0\) 的这类函数。将结果用于 \(m=n-1\)、\(f(r)=r\)、\(r>0\)，区分 \(n=2\) 与 \(n>2\)。
\end{exercise}
\begin{solution}
	在 \(F\) 的任一坐标中，\(\mu_g=f^m|\mathrm dr|\mu_h\)，而 \(\operatorname{grad}_g v=v'\partial_r\)，所以
	\[
		\Delta_gv=-f^{-m}(f^mv')'
		=-v''-m\frac{f'}f v'.
	\]
	调和方程等价于 \(f^mv'=B\)，故全部解为
	\[
		v(r)=A+B\int_{r_0}^r f(t)^{-m}\,\mathrm dt.
	\]
	在穿孔 Euclid 空间的径向模型中，\(n=2\) 得到 \(A+B\log r\)；\(n>2\) 得到 \(A+Br^{2-n}\)，其中常数已重新命名。非恒定解在原点发散，因此不能直接延拓为整个空间上的光滑调和函数。
\end{solution}

\begin{exercise}\label{g3:ex:torus-revolution}
	设 \(a>b>0\)，用两个模 \(2\pi\) 的角参数表示旋转环面
	\[
		X(\varphi,\theta)=\bigl((a+b\cos\varphi)\cos\theta,
		(a+b\cos\varphi)\sin\theta,b\sin\varphi\bigr).
	\]
	求诱导度量、总面积及曲率 \(K\)，指出正曲率、负曲率和零曲率所在的位置。
\end{exercise}
\begin{solution}
	两个偏导互相垂直，长度分别为 \(b\) 和 \(a+b\cos\varphi>0\)，故
	\[
		g=b^2\mathrm d\varphi^2+(a+b\cos\varphi)^2\mathrm d\theta^2,\qquad
		\mu_g=b(a+b\cos\varphi)|\mathrm d\varphi\,\mathrm d\theta|.
	\]
	在参数基本矩形上积分，得到 \(\operatorname{Area}(M)=4\pi^2ab\)。局部取 \(r=b\varphi\)，则 \(f(r)=a+b\cos(r/b)\)，由 \(K=-f''/f\) 得
	\[
		K=\frac{\cos\varphi}{b(a+b\cos\varphi)}.
	\]
	因此外侧 \(\cos\varphi>0\) 时为正，内侧 \(\cos\varphi<0\) 时为负，顶部和底部的两个平行圆 \(\cos\varphi=0\) 上为零。弧长变量 \(r\) 只需局部选取，公式对参数周期相容，故给出整个环面上的光滑函数。
\end{solution}

\begin{exercise}\label{g3:ex:round-sphere-height}
	令 \(S_a^n=\{x\in\mathbb R^{n+1}:|x|=a\}\)，\(a>0\)，取诱导度量，令 \(q(x)=x_{n+1}\)。用切向投影计算 \(\operatorname{grad}q\)，并证明 \(\Delta q=(n/a^2)q\)。解释该公式在常数缩放下的变化。
\end{exercise}
\begin{solution}
	记最后一个常坐标向量为 \(E\)。球面在 \(x\) 处的单位法向为 \(x/a\)，因此
	\[
		\operatorname{grad}q=E-\frac{q}{a^2}x.
	\]
	对切向量 \(V\) 求环境导数，得到 \(-V(q)x/a^2-qV/a^2\)。第一项法向，投影到切空间后为
	\[
		\nabla_V\operatorname{grad}q=-\frac{q}{a^2}V.
	\]
	在 \(n\) 维切空间上取迹，\(\operatorname{div}\operatorname{grad}q=-nq/a^2\)，依本书的符号约定即得所求。若把度量乘以 \(c^2\)，则梯度乘以 \(c^{-2}\)，散度在常数体积缩放下不变，故 Laplace 算子乘以 \(c^{-2}\)。半径从 \(a\) 变为 \(ca\) 时系数变为 \(n/(c^2a^2)\)，与公式一致。
\end{solution}

% !TeX root = ../../Book3.tex
% 纲要标题：## 第2章　测地线、距离、完备性与割迹
% 纲要位置：几何学计划纲要.md，第 2157 行。
% 正文状态：初稿；尚未正式审核。

% 纲要内容开始
% 
% > 本章目标与证明范围：把局部常微分方程、法邻域极小性和整体完备性连续建立。割点与共轭点的精细关系在 Jacobi 理论后再深化；本章不提前使用指标定理。
% 
% 纲要内容结束

\chapter{测地线、距离、完备性与割迹}
\label{g3:ch:02}
\BookChapterNumber{g3:ch:02}

\begin{chapterabstract}
	联络给出测地线方程，长度则给出端点间的内在距离。本章先用常微分方程构造指数映射，由 Gauss 引理证明短径向测地线最短，并建立法坐标与强凸邻域；随后完整证明 Hopf--Rinow 定理，将度量完备、测地线延拓和闭有界集紧性联系起来。最后沿每个初方向定义割时与割迹，说明最短线何时失去唯一性，证明紧集上的注入半径有正下界，并在球面、平坦环面和实射影空间中计算这些量。
\end{chapterabstract}

\begin{learninggoals}
	\begin{itemize}
		\item 从测地线初值问题构造指数映射，说明其定义域、光滑性、零处微分和仿射缩放性质。
		\item 证明 Gauss 引理、径向极小性和强凸邻域存在性，并区分法邻域、短线唯一性与全局最短线。
		\item 逐步完成 Hopf--Rinow 定理的等价闭环，辨认有限维、连通和无边界条件在证明中的用途。
		\item 用割时记录最大极小区间，计算标准模型的割迹与注入半径，比较局部坐标尺度和整体完备性。
	\end{itemize}
\end{learninggoals}

站在球面上一点，指定初始方向后沿大圆匀速前进，运动规律始终不变；到达对径点后，继续沿原方向走却会比回头走更长。这个例子把三个问题分开：初值是否确定一条曲线，这条曲线能否一直延长，以及它的某一段是否在所有连接端点的曲线中最短。第一个问题由微分方程回答，后两个则要结合长度与整体空间来研究。

\ProseParagraphBreak

本章从上一章的 Levi-Civita 联络出发，将“加速度为零”写成切丛上的一阶方程。第一卷第13.1节的最大局部流定理负责存在唯一性与初值依赖，第11.3节的逆函数定理把指数映射变成局部坐标。随后所有关于长度极小性、距离拓扑、凸邻域和完备性的结论都在本章证明。商空间模型还使用第一卷第5.3节的覆盖道路提升；需要时会说明覆盖和局部等距条件。

\ProseParagraphBreak

以下流形均为有限维、Hausdorff、第二可数的光滑流形，度量光滑且正定。除专门的边界反例外，流形无边界；涉及有限值内在距离时取连通流形。局部指数映射的构造不要求完备，只有到讨论整体延拓与割迹时才增加相应假设。可对照 Petersen 第3版第5.2--5.3节、第5.5.1--5.5.3节，以及第5.7.1节、第5.7.3--5.7.4节阅读测地线、距离、指数坐标、完备性与注入半径\cite[第5.2--5.3节; 第5.5.1--5.5.3节; 第5.7.1节; 第5.7.3--5.7.4节]{Petersen2016}。

% !TeX root = ../../Book3.tex
% 纲要标题：### 2.1 指数映射与法坐标
% 纲要位置：几何学计划纲要.md，第 2161 行。
% 正文状态：初稿；尚未正式审核。

% 纲要内容开始
% 
% * 测地线方程及初值依赖
% * 指数映射的微分
% * Gauss 引理的变分证明
% * 法坐标中的度量与联络
% * 局部径向极小性
% 
% 纲要内容结束

\section{指数映射与法坐标}
\label{g3:sec:0201}

在 Euclid 空间中，从点 \(p\) 沿初速度 \(v\) 直行，时间 \(t\) 后到达 \(p+tv\)。这既确定一条曲线，也把原点附近的向量变成 \(p\) 附近的点。在黎曼流形上，“速度不改变”应写成速度沿自身的协变导数为零。先解这一初值问题，再把时间一的终点作为向量的像，就得到指数映射。它在零向量附近可逆，由此产生的坐标使从中心出发的测地线成为直线；这些曲线的长度极小性还需要度量的进一步分析。

\subsection{测地线方程及初值依赖}
\label{g3:subsec:0201:01}

设 \(\gamma:I\to M\) 为光滑曲线，\(V(t)\in\mathrm{T}_{\gamma(t)}M\) 为沿曲线的光滑向量场。在坐标中写 \(\gamma(t)=(x^i(t))\)、\(V(t)=V^k(t)\partial_k|_{\gamma(t)}\)，沿曲线的协变导数为
\begin{equation}\label{g3:eq:covariant-derivative-along-curve}
	\mathrm{D}_tV
	=\left(\dot V^k+\Gamma^k_{ij}(x(t))\dot x^iV^j\right)\partial_k.
\end{equation}
这是把式\eqref{g3:eq:coordinate-covariant-derivative}限制到曲线后的表达；向量场沿曲线的分量可以独立指定，不要求先在曲线像的邻域上延拓。上一章的 Christoffel 换坐标公式保证这个表达在交叠处一致。度量相容性给出
\[
	\frac{\mathrm{d}}{\mathrm{d}t}g(V,W)
	=g(\mathrm{D}_tV,W)+g(V,\mathrm{D}_tW).
\]

\begin{definition}\label{g3:def:geodesic}
	设 \((M,g)\) 为无边界光滑黎曼流形，\(\nabla\) 为其 Levi-Civita 联络。若光滑曲线 \(\gamma:I\to M\) 满足 \(\mathrm{D}_t\dot\gamma=0\)，则称 \(\gamma\) 为\term[cedixian]{测地线}[geodesic]。常值曲线也包括在内。坐标中的测地线方程是
	\begin{equation}\label{g3:eq:geodesic-equation}
		\ddot x^k+\Gamma^k_{ij}(x)\dot x^i\dot x^j=0,
		\qquad k=1,\ldots,n.
	\end{equation}
\end{definition}

这是一组二阶方程，所以初值同时包括位置 \(p\) 与速度 \(v\in\mathrm{T}_pM\)。把速度也作为未知量，令 \(v^i=\dot x^i\)，就得到切丛坐标 \((x^i,v^i)\) 中的一阶系统
\[
	\dot x^k=v^k,\qquad
	\dot v^k=-\Gamma^k_{ij}(x)v^iv^j.
\]
其右端定义切丛上的光滑向量场
\begin{equation}\label{g3:eq:geodesic-spray}
	S=v^i\frac{\partial}{\partial x^i}
	-\Gamma^k_{ij}(x)v^iv^j\frac{\partial}{\partial v^k}.
\end{equation}
称其为\term[cedixianpenwu]{测地线喷雾}[geodesic spray]。换坐标时，速度分量按 \(\widetilde v^a=(\partial y^a/\partial x^i)v^i\) 变换；对此求时间导数产生的二阶坐标导数，恰与 Christoffel 变换中的二阶项抵消。因此各图中的 \(S\) 表示同一个切丛向量场。它的积分曲线正是 \((\gamma(t),\dot\gamma(t))\)。

\begin{theorem}\label{g3:thm:geodesic-ivp}
	设 \((M,g)\) 为有限维、Hausdorff、第二可数、无边界的光滑黎曼流形。对每个 \(p\in M\)、\(v\in\mathrm{T}_pM\)，存在唯一的最大测地线
	\[
		\gamma_{p,v}:I_{p,v}\longrightarrow M,\qquad
		\gamma_{p,v}(0)=p,\quad \dot\gamma_{p,v}(0)=v,
	\]
	其中 \(I_{p,v}\) 为含零的开区间，任意同初值测地线都是它的限制。集合
	\[
		\mathcal D=\{(t,p,v):v\in\mathrm{T}_pM,\ t\in I_{p,v}\}
		\subset\mathbb R\times\mathrm{T}M
	\]
	开，映射 \((t,p,v)\mapsto(\gamma_{p,v}(t),\dot\gamma_{p,v}(t))\) 在其上光滑。每条测地线速度的长度恒为 \(|v|_g\)。当 \(v=0\) 时，测地线恒为 \(p\)，且 \(I_{p,0}=\mathbb R\)；当 \(a\ne0\) 时，
	\begin{equation}\label{g3:eq:geodesic-scaling}
		I_{p,av}=a^{-1}I_{p,v},\qquad
		\gamma_{p,av}(t)=\gamma_{p,v}(at).
	\end{equation}
	任意仿射重参数 \(t=as+b\) 仍给出测地线。
\end{theorem}
\begin{proof}
	第一卷第13.1节“最大局部流”定理说明：无边界、有限维、Hausdorff、第二可数光滑流形上的光滑向量场，有唯一的最大积分曲线；它们的时间与初值组成开定义域，解映射在该定义域上联合光滑。切丛 \(\mathrm{T}M\) 是满足这些条件的 \(2n\) 维流形，而式\eqref{g3:eq:geodesic-spray}给出其上的光滑向量场。因此该定理适用于 \(S\)。其积分曲线的第一组方程保证第二分量是第一分量的速度，第二组方程正是测地线方程，故得到存在、唯一、最大区间及光滑依赖。
	% 跨卷依赖：g1:thm:maximal-local-flow，已核读其陈述与证明。

	由度量相容性，
	\[
		\frac{\mathrm{d}}{\mathrm{d}t}|\dot\gamma|_g^2
		=2g(\mathrm{D}_t\dot\gamma,\dot\gamma)=0,
	\]
	所以速度长度为常数。零初值速度给出常曲线，唯一性使它就是最大解。

	若 \(a\ne0\)，曲线 \(t\mapsto\gamma_{p,v}(at)\) 的协变加速度是原曲线协变加速度的 \(a^2\) 倍，且初速度为 \(av\)，故它是相应初值解。最大性先给出 \(a^{-1}I_{p,v}\subset I_{p,av}\)；再用缩放因子 \(a^{-1}\) 得反向包含，证明式\eqref{g3:eq:geodesic-scaling}。同一求导计算允许时间平移，因而仿射重参数保持测地线方程。
\end{proof}

对于非恒定测地线，参数中的“匀速”是方程的一部分。若以光滑函数 \(t=h(s)\) 重参数，则
\[
	\mathrm{D}_s\frac{\mathrm{d}}{\mathrm{d}s}\gamma(h(s))
	=h''(s)\dot\gamma(h(s)).
\]
原速度处处非零，因此新曲线仍满足测地线方程当且仅当 \(h''=0\)。曲线的像可以不变，所用参数却未必仍是测地线参数。

\subsection{指数映射的微分}
\label{g3:subsec:0201:02}

初速度的大小同时记录了单位时间内走多远。由缩放公式，沿速度 \(v\) 运动时间 \(t\)，与沿速度 \(tv\) 运动时间一到达同一点。于是可以把时间固定，将所有允许的初速度组织成一个映射。

\begin{definition}\label{g3:def:exponential-map}
	设 \((M,g)\) 为光滑无边界黎曼流形。定义
	\[
		\mathcal E=\{(p,v)\in\mathrm{T}M:1\in I_{p,v}\},
		\qquad
		\mathcal E_p=\mathcal E\cap\mathrm{T}_pM.
	\]
	映射
	\[
		\operatorname{Exp}:\mathcal E\longrightarrow M,\qquad
		\operatorname{Exp}(p,v)=\gamma_{p,v}(1)
	\]
	称为\term[zhishuyingshe]{指数映射}[exponential map]，其在单个切空间上的限制记为 \(\exp_p:\mathcal E_p\to M\)。
\end{definition}

\begin{proposition}\label{g3:prop:exponential-local}
	设 \((M,g)\) 为有限维、Hausdorff、第二可数、无边界光滑黎曼流形。指数映射的定义域 \(\mathcal E\) 是包含零截面的开集，\(\operatorname{Exp}\) 光滑；每个 \(\mathcal E_p\) 关于零向量星形，即 \(v\in\mathcal E_p\)、\(0\le t\le1\) 时 \(tv\in\mathcal E_p\)。在自然识别 \(\mathrm{T}_0(\mathrm{T}_pM)=\mathrm{T}_pM\) 下，
	\begin{equation}\label{g3:eq:exponential-derivative-zero}
		(\mathrm{d}\exp_p)_0=\operatorname{id}_{\mathrm{T}_pM}.
	\end{equation}
	因此存在 \(R>0\)，使 \(\exp_p\) 在切空间开球 \(B_p(0,R)=\{v:|v|_g<R\}\) 上为到其像的微分同胚。此外，端点映射
	\begin{equation}\label{g3:eq:geodesic-endpoint-map}
		F:\mathcal E\longrightarrow M\times M,\qquad
		F(p,v)=(p,\exp_pv)
	\end{equation}
	在每个 \((p,0)\) 附近为局部微分同胚。
\end{proposition}
\begin{proof}
	\(\mathcal E\) 是开集 \(\mathcal D\) 在时间一的切片，故开；指数映射是测地流在该切片上的位置分量，故光滑。零速度解对所有时间存在，所以 \(\mathcal E\) 包含零截面。若 \(v\in\mathcal E_p\)，开区间 \(I_{p,v}\) 包含 \([0,1]\)。对 \(0<t\le1\)，式\eqref{g3:eq:geodesic-scaling}说明 \(1\in I_{p,tv}\)，而 \(t=0\) 已经处理，故纤维星形。

	任取 \(w\in\mathrm{T}_pM\)。当 \(s\) 充分接近零时，缩放公式给出
	\[
		\exp_p(sw)=\gamma_{p,w}(s).
	\]
	在 \(s=0\) 求导，右侧导数是初速度 \(w\)，便得到式\eqref{g3:eq:exponential-derivative-zero}。第一卷定理11.3.3的逆函数定理适用于同维无边界流形之间微分可逆的光滑映射，因此 \(\exp_p\) 在零向量的某个邻域上为微分同胚。取包含于该邻域的小球，即得所述 \(R\)。
	% 跨卷依赖：g1:thm:inverse-function。

	最后看起点也变化的情形。在零截面处有自然分裂
	\[
		\mathrm{T}_{(p,0)}(\mathrm{T}M)
		\simeq\mathrm{T}_pM\oplus\mathrm{T}_pM:
	\]
	第一项由零截面上移动起点得到，第二项由固定起点、改变纤维向量得到。由于 \(\exp_x0=x\)，第一项对终点的导数是恒等；式\eqref{g3:eq:exponential-derivative-zero}说明第二项也是恒等。故
	\[
		\mathrm{d}F_{(p,0)}(\xi,\eta)=(\xi,\xi+\eta).
	\]
	这个线性映射可逆，逆为 \((\xi,\zeta)\mapsto(\xi,\zeta-\xi)\)。再次应用逆函数定理，得到 \((p,0)\) 和 \((p,p)\) 的开邻域之间的光滑局部逆。
\end{proof}

\begin{definition}\label{g3:def:normal-neighborhood}
	若 \(U\subset\mathcal E_p\) 是包含零、关于零星形的开集，且 \(\exp_p|_U\) 为微分同胚，则称 \(\exp_p(U)\) 为 \(p\) 的\term[falinyu]{法邻域}[normal neighborhood]。当 \(U=B_p(0,R)\) 为球时，称其像为 \(p\) 的\term[faqiu]{法球}[normal ball]。
\end{definition}

这里的球先在切空间中按 \(g_p\) 定义，其像与内在距离球的关系将在长度极小性之后建立。\cref{g3:prop:exponential-local}的端点映射还说明：起点与终点同时作小变化时，连接它们的一个小初速度可以光滑地求出。第二节将据此统一邻近起点可用的半径。

\begin{example}\label{g3:exm:interval-exponential-domain}
	在开区间 \(M=(0,1)\) 上取 \(g=\mathrm{d}x^2\)。测地线是 \(\gamma_{p,v}(t)=p+tv\)，只允许取使其仍落在 \((0,1)\) 的最大时间区间。因此
	\[
		\mathcal E_p=(-p,1-p),\qquad \exp_p(v)=p+v.
	\]
	定义域包含零且星形，却通常不是整个切空间。指数映射局部存在无需整体延拓到所有时间。
\end{example}

\begin{example}\label{g3:exm:round-sphere-exponential}
	半径 \(a>0\) 的圆球面 \(S_a^n\subset\mathbb R^{n+1}\) 取诱导度量。对 \(p\in S_a^n\) 及非零 \(v\in\mathrm{T}_pS_a^n\)，令
	\[
		\gamma(t)=\cos(t|v|/a)p
		+a\sin(t|v|/a)\frac{v}{|v|}.
	\]
	因为 \(p\perp v\)、\(|p|=a\)，这条曲线始终在球面上，初速度为 \(v\)，且 \(\gamma''(t)=-(|v|^2/a^2)\gamma(t)\) 沿法向。由\cref{g3:prop:induced-connection-projection}，其协变加速度为零，所以唯一性给出
	\begin{equation}\label{g3:eq:round-sphere-exponential}
		\exp_p(v)=\cos(|v|/a)p
		+a\frac{\sin(|v|/a)}{|v|}v,
		\qquad \exp_p(0)=p.
	\end{equation}
	右端在零处光滑，因为 \(\cos(|v|/a)\) 和 \(\sin(|v|/a)/|v|\) 都是 \(|v|^2\) 的光滑函数。

	这个映射把 \(B_p(0,\pi a)\) 微分同胚地送到 \(S_a^n\setminus\{-p\}\)。事实上，对既非 \(p\) 又非 \(-p\) 的 \(q\)，令
	\[
		\vartheta=\arccos\frac{\langle p,q\rangle}{a^2}\in(0,\pi),
		\qquad
		u=\frac{q-\cos\vartheta\,p}{a\sin\vartheta}.
	\]
	则 \(u\) 是唯一的单位切向量，使 \(q=\exp_p(a\vartheta u)\)。此逆在远离 \(p\) 处光滑，在 \(p\) 附近与逆函数定理给出的光滑逆一致。半径等于 \(\pi a\) 的所有向量则映到同一个对径点 \(-p\)，显示局部可逆的范围会受到整体几何限制。
\end{example}

\subsection{Gauss 引理的变分证明}
\label{g3:subsec:0201:03}

指数映射把每条径向直线变成测地线。若同时改变初速度的方向，这些测地线便形成一族曲线。比较径向运动与这种横向变化，就可以判断指数映射如何处理正交关系；它也是随后证明长度极小性的依据。这个结论称为 Gauss 引理。\footnote{Carl Friedrich Gauss（1777--1855），德国数学家，在数论、微分几何和测地学等领域作出重要贡献。}

\begin{lemma}[Gauss 引理]\label{g3:lem:gauss}
	设 \((M,g)\) 为光滑无边界黎曼流形，\(p\in M\)，\(v\in\mathcal E_p\)，并用向量空间的平移识别 \(\mathrm{T}_v(\mathrm{T}_pM)=\mathrm{T}_pM\)。对任意 \(w\in\mathrm{T}_pM\)，有
	\begin{equation}\label{g3:eq:gauss-lemma}
		g_{\exp_pv}\bigl((\mathrm{d}\exp_p)_vv,(\mathrm{d}\exp_p)_vw\bigr)
		=g_p(v,w).
	\end{equation}
	特别地，指数映射保持径向向量的长度，并把与 \(v\) 正交的向量送到与径向速度正交的向量。
\end{lemma}
\begin{proof}
	\(v=0\) 时等式两边都为零。以下固定 \(v\ne0\)。由于 \(\mathcal E_p\) 开，可取小参数区间，使 \(v+sw\in\mathcal E_p\)。考虑光滑变分
	\[
		H(s,t)=\exp_p\bigl(t(v+sw)\bigr)=\gamma_{p,v+sw}(t),
		\qquad 0\le t\le1.
	\]
	记沿这个变分的向量场 \(T=\partial_tH\)、\(V=\partial_sH\)。每条 \(t\)-曲线都是测地线，故 \(\mathrm{D}_tT=0\)；又因 \(H(s,0)=p\)，有 \(V(s,0)=0\)。

	先补明混合协变导数为何可以交换。在任一覆盖变分像的坐标图内，
	\[
		(\mathrm{D}_tV)^k
		=\partial_t\partial_sH^k+\Gamma^k_{ij}(H)\partial_tH^i\partial_sH^j,
	\]
	而 \((\mathrm{D}_sT)^k\) 由交换 \(s,t\) 得到。普通混合偏导相等，且无挠性给出 \(\Gamma^k_{ij}=\Gamma^k_{ji}\)，所以
	\begin{equation}\label{g3:eq:mixed-variation-derivatives}
		\mathrm{D}_tV=\mathrm{D}_sT.
	\end{equation}
	这个计算只用光滑映射 \(H\) 的坐标表达，不要求 \(H\) 是浸入，因而也适用于 \(t=0\)。

	由度量相容性和式\eqref{g3:eq:mixed-variation-derivatives}，
	\[
		\partial_t g(T,V)
		=g(T,\mathrm{D}_tV)
		=g(T,\mathrm{D}_sT)
		=\frac12\partial_s g(T,T).
	\]
	每条测地线常速，故 \(g(T,T)=g_p(v+sw,v+sw)\)，与 \(t\) 无关。在 \(s=0\) 处，上式右端就是 \(g_p(v,w)\)。从 \(t=0\) 积分到 \(t=1\)，利用初始内积为零，得到
	\[
		g(T(0,1),V(0,1))=g_p(v,w).
	\]
	链式法则给出 \(T(0,1)=(\mathrm{d}\exp_p)_vv\)、\(V(0,1)=(\mathrm{d}\exp_p)_vw\)，即为所求。取 \(w=v\) 得径向长度保持；取 \(w\perp v\) 得正交关系。
\end{proof}

在法球去掉中心的部分，写 \(v=ru\)，其中 \(|u|_g=1\)、\(r>0\)。定义径向单位向量
\[
	\partial_r\big|_{\exp_p(ru)}=(\mathrm{d}\exp_p)_{ru}u.
\]
Gauss 引理说明其长度为一，并与固定半径的球面像正交。因此在这些极坐标中，度量可以写成
\[
	g=\mathrm{d}r^2+g_r,
\]
其中 \(g_r\) 是单位方向球面上的一族正定度量。这里 \(g_r\) 随 \(r\) 及方向变化，中心 \(r=0\) 则应使用下面的光滑法坐标。

\subsection{法坐标中的度量与联络}
\label{g3:subsec:0201:04}

在 \(\mathrm{T}_pM\) 选一组 \(g_p\)-正交单位基 \((E_1,\ldots,E_n)\)。把法球中的点写成
\[
	q=\exp_p(x^iE_i),
\]
便得到以 \(p\) 为中心的\term[fazuobiao]{法坐标}[normal coordinates] \(x^1,\ldots,x^n\)。这些坐标先在切空间中作线性坐标，再经指数映射送到流形；它们的特殊性来自径向直线恰好成为测地线。

\begin{proposition}\label{g3:prop:normal-coordinate-first-order}
	设 \((M,g)\) 为光滑无边界黎曼流形，在 \(p\) 的法邻域中采用由 \(g_p\)-正交单位基产生的法坐标。则在中心 \(p\)，
	\begin{equation}\label{g3:eq:normal-coordinate-first-order}
		g_{ij}(p)=\delta_{ij},\qquad
		\Gamma^k_{ij}(p)=0,\qquad
		\partial_k g_{ij}(p)=0.
	\end{equation}
\end{proposition}
\begin{proof}
	由 \((\mathrm{d}\exp_p)_0=\operatorname{id}\)，有 \(\partial_i|_p=E_i\)，故 \(g_{ij}(p)=\delta_{ij}\)。对任意 \(a=(a^1,\ldots,a^n)\)，曲线 \(t\mapsto\exp_p(ta^iE_i)\) 在足够小时间内是测地线，其法坐标为 \(x^i(t)=ta^i\)。在 \(t=0\) 将它代入测地线方程，得到
	\[
		\Gamma^k_{ij}(p)a^ia^j=0\qquad\text{对一切 }a.
	\]
	先取 \(a\) 为单个坐标基向量，再取两个坐标基向量的和，并用 \(\Gamma^k_{ij}=\Gamma^k_{ji}\)，便得到所有系数都为零。最后，度量相容性写成
	\[
		\partial_k g_{ij}
		=\Gamma^\ell_{ki}g_{\ell j}+\Gamma^\ell_{kj}g_{i\ell};
	\]
	在 \(p\) 求值就得到度量的一阶偏导全部为零。
\end{proof}

所以在中心附近，Taylor 公式给出 \(g_{ij}(x)=\delta_{ij}+O(|x|^2)\)。法坐标消去了中心处的一阶度量变化；二阶及更高阶变化仍记录实际几何。第三章将把二阶项与曲率联系起来。极坐标中的 \(\mathrm{d}r^2+g_r\) 和这里的 \(g_{ij}(x)\) 描述同一度量：前者突出径向正交性，后者在中心仍保持光滑。

\subsection{局部径向极小性}
\label{g3:subsec:0201:05}

考虑从法球中心 \(p\) 走到 \(q=\exp_p(v)\)。径向测地线长度为 \(|v|_g\)。沿其他曲线前进时，速度可以分成径向分量和球面方向分量；Gauss 引理使两部分正交，因此球面方向的运动会增加总长度。还须处理一个整体问题：竞争曲线可能先离开法球，再从外面绕回 \(q\)。缩小一层法球后，首次离开所付出的长度足以排除这种捷径。

\begin{proposition}\label{g3:prop:normal-radial-minimality}
	设 \((M,g)\) 为有限维、Hausdorff、第二可数、无边界光滑黎曼流形，\(p\in M\)、\(R>0\)，并设 \(\exp_p:B_p(0,R)\to U_R\) 为微分同胚。则有以下结论。
	\begin{enumerate}
		\item 对 \(v\in B_p(0,R)\)，从 \(p\) 到 \(q=\exp_pv\) 的任意连续分片 \(C^1\) 曲线 \(c:[a,b]\to M\) 都满足
		\[
			L_g(c)\ge |v|_g.
		\]
		径向测地线 \(t\mapsto\exp_p(tv)\)、\(0\le t\le1\)，达到这个下界。
		\item 若 \(v\ne0\)，等号成立当且仅当
		\[
			c(t)=\exp_p(\lambda(t)v),
		\]
		其中 \(\lambda:[a,b]\to[0,1]\) 为非减的连续分片 \(C^1\) 函数，\(\lambda(a)=0\)、\(\lambda(b)=1\)。因此允许停顿时，最短曲线的像仍是唯一的径向线段；固定参数区间及常速后，曲线也唯一。
		\item 对任意 \(0<r<R\)，曲线从 \(p\) 首次到达 \(U_r=\exp_p(B_p(0,r))\) 的边界时，已经走过的长度至少为 \(r\)。集合 \(\overline{U_r}=\exp_p(\overline{B_p(0,r)})\) 紧。
	\end{enumerate}
\end{proposition}
\begin{proof}
	先在 \(U_R\) 内比较长度。令 \(x(t)=\exp_p^{-1}(c(t))\)、\(\rho(t)=|x(t)|_g\)。在 \(x(t)\ne0\) 处，Gauss 引理与内积不等式给出
	\[
		|\dot\rho(t)|
		=\left|g\left((\mathrm{d}\exp_p)_{x(t)}
			\frac{x(t)}{|x(t)|_g},\dot c(t)\right)\right|
		\le|\dot c(t)|_g.
	\]
	曲线经过 \(p\) 时，\(\rho\) 未必可微。为直接处理这些时刻，取
	\[
		\rho_\varepsilon(t)=\sqrt{|x(t)|_g^2+\varepsilon^2},
		\qquad\varepsilon>0.
	\]
	它在每个 \(C^1\) 片段上都是 \(C^1\)。在 \(x\ne0\) 处，其导数为 \(\rho/\sqrt{\rho^2+\varepsilon^2}\) 乘以 \(\dot\rho\)；在 \(x=0\) 处，直接求导得零。因此处处有 \(|\dot\rho_\varepsilon|\le|\dot c|_g\)。逐片积分并令 \(\varepsilon\to0\)，对任意子区间 \([s,t]\) 得到
	\begin{equation}\label{g3:eq:radial-length-bound}
		L_g(c|_{[s,t]})\ge|\rho(t)-\rho(s)|.
	\end{equation}
	这个论证也包括速度为零和多次经过中心的曲线。

	再排除离开法球的竞争曲线。固定 \(|v|_g<r<R\)。闭切球 \(\overline{B_p(0,r)}\) 紧，其像 \(K_r\) 紧，因 \(M\) 为 Hausdorff 空间而闭；又因 \(\exp_p\) 在更大的开球上为微分同胚，\(K_r\) 的内部是 \(U_r\)，边界是半径 \(r\) 的切球面之像。因此 \(K_r=\overline{U_r}\subset U_R\)。若 \(c\) 离开 \(U_r\)，连续性给出首次到达边界的时刻 \(\tau\)，且 \(c([a,\tau])\subset K_r\)。这段曲线可在 \(U_R\) 中提升，式\eqref{g3:eq:radial-length-bound}给出
	\[
		L_g(c)\ge L_g(c|_{[a,\tau]})\ge r>|v|_g.
	\]
	若 \(c\) 始终在 \(U_r\) 中，直接用式\eqref{g3:eq:radial-length-bound}得到 \(L_g(c)\ge|v|_g\)。径向测地线的速度长度恒为 \(|v|_g\)，故在时间区间 \([0,1]\) 上恰有此长度。这也证明了第三项。

	最后分析等号。设 \(v\ne0\) 且 \(L_g(c)=|v|_g\)。刚才的严格不等式保证曲线始终在 \(U_r\) 中。对任意 \(a\le s\le t\le b\)，分别对三段使用式\eqref{g3:eq:radial-length-bound}，并考虑有向径向变化，得到
	\[
		\begin{aligned}
		L_g(c)
		&=L_g(c|_{[a,s]})+L_g(c|_{[s,t]})+L_g(c|_{[t,b]})\\
		&\ge \rho(s)+\bigl(\rho(t)-\rho(s)\bigr)
			+\bigl(\rho(b)-\rho(t)\bigr)=\rho(b).
		\end{aligned}
	\]
	总和已经取等，三项分别减去对应有向径向变化均非负，因此每项也取等，特别地
	\[
		L_g(c|_{[s,t]})=\rho(t)-\rho(s)\ge0.
	\]
	所以 \(\rho\) 非减，集合 \(\{t:\rho(t)>0\}\) 是离开中心以后的一个区间。在这个区间内写 \(x=\rho u\)、\(|u|_g=1\)。Gauss 引理使径向与角向速度正交，故在每个光滑片段上
	\[
		|\dot c|_g^2
		=\dot\rho^{\,2}
		+\left|(\mathrm{d}\exp_p)_x(\rho\dot u)\right|_g^2.
	\]
	任意这类子区间上的长度等号说明，连续非负函数 \(|\dot c|_g-\dot\rho\) 的积分为零，因而它恒为零。上式与 \(\mathrm{d}\exp_p\) 可逆性给出 \(\dot u=0\)。有限个分段端点处的连续性使这些常方向一致，最终必须为 \(u=v/|v|_g\)。在其余时间曲线恒为 \(p\)，所以 \(c(t)=\exp_p(\lambda(t)v)\)，其中 \(\lambda=\rho/|v|_g\) 非减。由于此时 \(\rho=\langle x,u\rangle_{g_p}\)，\(\lambda\) 在经过中心的分段上也为 \(C^1\)。

	反过来，这种重参数曲线的速度长度是 \(\lambda'(t)|v|_g\)，逐片积分恰得 \(|v|_g\)。若要求常速，\(\lambda'(t)\) 必为 \(1/(b-a)\)，从而 \(\lambda(t)=(t-a)/(b-a)\)。当 \(v=0\) 时，下界为零，只有常曲线取得等号。
\end{proof}

特别地，每条测地线的足够短子段都是端点之间的最短曲线：在子段起点取法球，再用初值唯一性把该段识别为径向测地线。在这里，“足够短”由起点附近指数映射的可逆范围决定。球面例子说明，延长同一条测地线以后可能到达多个径向方向汇合的点；局部的最短性并不自动延续到任意长的区间。

% !TeX root = ../../Book3.tex
% 纲要标题：### 2.2 内在距离和凸邻域
% 纲要位置：几何学计划纲要.md，第 2169 行。
% 正文状态：初稿；尚未正式审核。

% 纲要内容开始
% 
% * 分片光滑长度和距离拓扑
% * 短测地线与强凸邻域
% * 局部唯一性和整体多条最短线
% 
% 纲要内容结束

\section{内在距离和凸邻域}
\label{g3:sec:0202}

从一点出发，足够短的径向测地线已经给出了最短路。现在需要同时移动两个端点：哪些邻近点可以由唯一的短测地线连接，这条线是否仍留在原来的邻域内？回答这些问题之前，先把所有曲线长度的下确界变成一个真正的距离，并核对它与原来的流形拓扑一致。本节的流形均无边界；除讨论连通分支的说明外，均假设连通。

\subsection{分片光滑长度和距离拓扑}
\label{g3:subsec:0202:01}

本章的分片光滑曲线具有有限分割，在每个闭子区间上光滑到端点，允许在分割点有角点或出现常值段。长度沿用\cref{g3:def:riemannian-length}的定义。把曲线分成两段，长度就是两段长度之和；反向参数化不改变长度。这两项性质使“在所有连接曲线中取下确界”能够满足距离的对称性与三角不等式。

\begin{lemma}
	\label{g3:lem:piecewise-smooth-connected}
	连通光滑流形中的任意两点都可以用有限段光滑曲线连接。
\end{lemma}
\begin{proof}
	固定一点 \(p\)，令 \(A\) 为可以从 \(p\) 用分片光滑曲线到达的点集。若 \(q\in A\)，在 \(q\) 附近选坐标球；球内的每一点都可由坐标直线段连接到 \(q\)，接在已有曲线之后便仍可从 \(p\) 到达。因此 \(A\) 开。若 \(q\notin A\)，以 \(q\) 为中心的足够小坐标球也不含 \(A\) 中的点，否则沿反向坐标线段就能到达 \(q\)。故补集也开。由 \(p\in A\) 和连通性，\(A=M\)。
\end{proof}

\begin{definition}
	\label{g3:def:riemannian-distance}
	在连通黎曼流形 \((M,g)\) 上，定义\term[neizaijuli]{内在距离}[intrinsic distance]
	\begin{equation}
		\label{g3:eq:riemannian-distance}
		d_g(p,q)\coloneqq\inf_{\gamma}L_g(\gamma),
	\end{equation}
	其中 \(\gamma:[0,1]\to M\) 遍历所有从 \(p\) 到 \(q\) 的分片光滑曲线。度量明确时写作 \(d\)。满足 \(L_g(\gamma)=d(p,q)\) 的曲线称为两点之间的\term[zuiduanquxian]{最短曲线}[minimizing curve]。
\end{definition}

若流形不连通，在每个连通分支上分别使用这个定义；不同分支间没有连接曲线，规定下确界为 \(+\infty\)。以下关于通常有限值距离的结果都在一个连通分支上使用。

\begin{theorem}
	\label{g3:thm:riemannian-distance}
	设 \((M,g)\) 为连通、有限维、Hausdorff、第二可数的无边界光滑黎曼流形。\(d_g\) 是 \(M\) 上的有限值距离，且它诱导的拓扑恰为原来的流形拓扑。
\end{theorem}
\begin{proof}
	\cref{g3:lem:piecewise-smooth-connected}保证下确界所取的集合非空；每条分片光滑曲线的参数区间紧，长度有限，故 \(d(p,q)<\infty\)。长度非负，常值曲线给出 \(d(p,p)=0\)，反向曲线给出 \(d(p,q)=d(q,p)\)。对 \(\varepsilon>0\)，选择从 \(p\) 到 \(q\)、从 \(q\) 到 \(r\) 的曲线，使长度分别小于 \(d(p,q)+\varepsilon\)、\(d(q,r)+\varepsilon\)。拼接并线性重参数化后，得到
	\[
		d(p,r)\leq d(p,q)+d(q,r)+2\varepsilon.
	\]
	令 \(\varepsilon\downarrow0\) 得三角不等式。这里没有假设下确界已经取到。

	关键是证明不同点的距离为正，同时排除绕出坐标图后任意缩短的可能。固定 \(p\)，选坐标 \(z\)，使 \(z(p)=0\)，且某个闭坐标球 \(\{|z|\leq a\}\) 紧含于图域，其中 \(a>0\)。由\cref{g3:prop:metric-coordinate-bounds}，这个闭球上存在 \(0<c\leq C\)，使
	\[
		c|\xi|^2\leq g_{ij}(z)\xi^i\xi^j\leq C|\xi|^2.
	\]
	记开坐标球为 \(U_a\)。任何从 \(p\) 出发并离开 \(U_a\) 的曲线，第一次到达边界以前的部分都留在闭坐标球中；这部分的坐标长度至少为 \(a\)，所以其几何长度至少为 \(\sqrt c\,a\)。因此
	\begin{equation}
		\label{g3:eq:coordinate-exit-cost}
		d(p,q)\geq\sqrt c\,a\qquad(q\notin U_a).
	\end{equation}
	若 \(q\in U_a\)，留在闭球内的连接曲线长度至少为 \(\sqrt c\,|z(q)|\)；离开的曲线则至少长 \(\sqrt c\,a\)，也不小于这个下界。另一方面，坐标直线段 \(t\mapsto tz(q)\) 的长度至多为 \(\sqrt C\,|z(q)|\)。于是
	\begin{equation}
		\label{g3:eq:distance-coordinate-bounds}
		\sqrt c\,|z(q)|\leq d(p,q)\leq\sqrt C\,|z(q)|
		\qquad(q\in U_a).
	\end{equation}
	由两式可见，只要 \(q\ne p\)，距离就严格为正。

	最后，\eqref{g3:eq:distance-coordinate-bounds}的上界说明足够小的坐标球包含于任意给定的 \(d\) 球；下界连同\eqref{g3:eq:coordinate-exit-cost}说明足够小的 \(d\) 球包含于任意给定的小坐标球。两类邻域在 \(p\) 处互相细化，而 \(p\) 任意，故拓扑相同。
\end{proof}

我们分别用 \(B_d(p,r)\) 和 \(\overline B_d(p,r)=\{q:d(p,q)\leq r\}\) 表示开距离球和闭距离球；切空间中的球写为 \(B_p(0,r)=\{v\in\mathrm{T}_pM:|v|_g<r\}\)。三角不等式还给出
\[
	|d(p,q)-d(p',q')|\leq d(p,p')+d(q,q'),
\]
所以距离作为 \(M\times M\) 上的函数连续。这里的闭距离球仅按不等式定义；它是否紧以及是否等于某个指数映射的像，仍需进一步证明。

\subsection{短测地线与强凸邻域}
\label{g3:subsec:0202:02}

一个固定中心的法邻域来自 \(\exp_p\) 的局部可逆性。为了把中心也作为变量，考虑同时记录起点与终点的映射
\[
	E(x,v)=(x,\exp_xv).
\]
\cref{g3:prop:exponential-local}已经证明它在零截面附近局部可逆。下面把这一结论转成对一族中心都适用的半径，并用径向极小性辨认相应的距离球。

\begin{lemma}
	\label{g3:lem:uniform-normal-neighborhood}
	设 \((M,g)\) 为连通无边界光滑黎曼流形，\(p\in M\)。存在 \(p\) 的开邻域 \(V\) 和常数 \(r>0\)，使对每个 \(x\in V\)，指数映射在 \(B_x(0,4r)\) 上有定义并且是到其像的微分同胚。对 \(0<\rho\leq2r\)，有
	\begin{equation}
		\label{g3:eq:uniform-normal-balls}
		\begin{aligned}
			B_d(x,\rho)&=\exp_x B_x(0,\rho),\\
			\overline B_d(x,\rho)&=\exp_x\overline B_x(0,\rho).
		\end{aligned}
	\end{equation}
	这些闭距离球紧。每个 \(y\in\overline B_d(x,2r)\) 都由唯一的常速最短测地线 \([0,1]\to M\) 与 \(x\) 连接；任意其他最短曲线都是该线段的单调重参数化，允许停顿。

	还可使任意 \(x,y\in V\) 满足 \(d(x,y)<r\)，并使
	\[
		\log_x y=\exp_x^{-1}(y),\qquad
		\gamma_{x,y}(t)=\exp_x(t\log_xy)
	\]
	光滑依赖于 \((x,y)\) 及 \(0\leq t\leq1\)。对任意紧集 \(K\subset M\)，也存在一个 \(r_K>0\)，使上述关于指数映射、距离球与最短线的结论对所有中心 \(x\in K\) 同时成立。
\end{lemma}
\begin{proof}
	先给出统一的指数映射范围，再排除从外部绕行的短曲线。由\cref{g3:prop:exponential-local}，可取 \((p,0)\) 在 \(\mathrm{T}M\) 中的邻域 \(\mathcal U\)，使 \(E\) 在 \(\mathcal U\) 上为微分同胚。在局部切丛平凡化中，缩小一个相对紧的底邻域 \(V_0\)，再取足够小的正数 \(r\)，可使
	\[
		\{(x,v):x\in V_0,\ |v|_g<4r\}\subset\mathcal U.
	\]
	具体说，\(\mathcal U\) 含有底坐标小球与某个 Euclid 纤维小球的乘积，而\cref{g3:prop:metric-coordinate-bounds}在底坐标闭球上给出统一的范数比较；将 \(4r\) 选得足够小即可。由于 \(E\) 保持第一个分量，限制到每一纤维就得到 \(\exp_x\) 的单射性和光滑逆。

	固定 \(x\in V_0\)，记 \(U_x=\exp_x B_x(0,3r)\)。\(\exp_x\) 在更大的 \(4r\) 球上为微分同胚，所以 \(\overline U_x=\exp_x\overline B_x(0,3r)\)，边界是半径 \(3r\) 的切球面之像。由\cref{g3:prop:normal-radial-minimality}，在 \(U_x\) 中，从 \(x\) 到 \(\exp_xv\) 的任意连接曲线长度至少为 \(|v|_g\)，等号只对应单调径向曲线。任意从 \(x\) 离开 \(U_x\) 的曲线，第一次到达边界之前已经付出至少 \(3r\) 的长度；这一点可沿曲线先取趋于首次出界时刻的内点，应用径向下界，再令径向长度趋于 \(3r\)。

	因此，对 \(|v|_g\leq2r\)，径向测地线的长度 \(|v|_g\) 比出界代价小，它是流形中所有连接曲线中的最短线，且等号仍只能是单调径向曲线。反过来，若 \(d(x,y)\leq2r\)，按下确界定义可选长度小于 \(3r\) 的连接曲线。它不能离开 \(U_x\)，故 \(y\in U_x\)。对每条长度趋于 \(d(x,y)\) 的此类曲线使用径向下界，得
	\[
		|\exp_x^{-1}y|_g\leq d(x,y).
	\]
	径向线给出反向不等式，所以两者相等。开、闭球等式同时成立；闭切球在有限维向量空间中紧，它的连续像也紧。常速条件确定了径向线的参数，故常速最短线唯一；\(x=y\) 时唯一最短线为常值曲线。

	逆映射 \(E^{-1}\) 在其像上光滑，因而刚得到的 \(\log_xy\) 随两端点光滑变化。利用已经证明的距离拓扑一致性，将 \(V_0\) 缩小到 \(p\) 的邻域 \(V\)，使 \(V\subset B_d(p,r/3)\)。三角不等式给出任意 \(x,y\in V\) 的距离小于 \(2r/3\)，所以 \((x,y)\) 都落在该光滑逆的定义域内。\(t\log_xy\) 始终在相同的指数映射范围内，故 \(\gamma_{x,y}\) 也光滑依赖所有参数，包括对角线 \(x=y\)。

	最后，对紧集 \(K\) 中每点作上述构造，从所得底邻域选有限子覆盖 \(V_1,\ldots,V_m\)。若相应半径为 \(r_1,\ldots,r_m\)，取 \(r_K=\min_i r_i>0\)。每个 \(x\in K\) 至少属于一个 \(V_i\)，在其中已经证明的结论限制到较小半径后仍成立。
\end{proof}

这个引理同时说明黎曼距离局部具有紧闭球，并给出了最短线随端点变化的公式。它还没有说明连接 \(V\) 中两点的短线始终留在 \(V\) 中：法邻域控制的是从一个中心出发的径向线，而任意两点之间的线需要另外的凸性条件。

\begin{definition}
	\label{g3:def:strongly-convex-neighborhood}
	黎曼流形中的开集 \(U\) 若满足：任意 \(x,y\in U\) 在整个流形中有唯一的常速最短测地线连接，并且这条线的像包含于 \(U\)，则称 \(U\) 为\term[qiangtulinyu]{强凸邻域}[strongly convex neighborhood]。
\end{definition}

\begin{theorem}
	\label{g3:thm:strongly-convex-neighborhood}
	设 \((M,g)\) 为连通无边界光滑黎曼流形，\(p\in M\)。对 \(p\) 的任意开邻域 \(O\)，存在强凸邻域 \(U\)，使 \(p\in U\)、\(\overline U\) 紧且 \(\overline U\subset O\)。其中的最短线 \(\gamma_{x,y}(t)\) 光滑依赖两端点与参数，函数 \((x,y)\mapsto d(x,y)^2\) 在 \(U\times U\) 上光滑。
\end{theorem}
\begin{proof}
	先用\cref{g3:lem:uniform-normal-neighborhood}取得 \(p\) 的邻域 \(V\) 及其中两点间的唯一短线。我们将选择一个更小的坐标球：先证明这些短线不可能逃出一个较大的坐标球，再用一个沿测地线凸的函数把它们压回小球。

	在 \(V\cap O\) 中取以 \(p\) 为原点的坐标 \(z=(z^1,\ldots,z^n)\)，令 \(h(z)=|z|^2\)。对仿射参数测地线 \(\gamma\)，坐标方程给出
	\begin{equation}
		\label{g3:eq:coordinate-square-convexity}
		\frac{\mathrm{d}^2}{\mathrm{d}t^2}h(\gamma(t))
		=2\sum_k(\dot z^k)^2
		-2\sum_{i,j,k}z^k\Gamma^k_{ij}(z)\dot z^i\dot z^j.
	\end{equation}
	Christoffel 系数在 \(0\) 附近有界，而第二项还带有因子 \(z^k\)。故可取 \(R>0\)，使闭坐标球 \(\{|z|\leq R\}\) 紧含于 \(V\cap O\)，并且在该闭球内，对所有坐标向量 \(\xi\)，有
	\[
		\left|\sum_{i,j,k}z^k\Gamma^k_{ij}(z)\xi^i\xi^j\right|
		\leq\tfrac12|\xi|^2.
	\]
	于是任何留在这个球中的测地线都满足 \((h\circ\gamma)''\geq|\dot z|^2\geq0\)。再在闭球上取度量比较常数 \(0<c\leq C\)。

	选择 \(\delta>0\)，满足
	\[
		\delta<R/2,\qquad 2\sqrt C\,\delta<\sqrt c\,R/2,
	\]
	并令 \(U=\{|z|<\delta\}\)。对 \(x,y\in U\)，\(\gamma_{x,y}\) 是上述引理给出的全局最短测地线。坐标直线段留在 \(U\) 中，给出
	\[
		L_g(\gamma_{x,y})=d(x,y)\leq\sqrt C\,|z(y)-z(x)|
		<2\sqrt C\,\delta.
	\]
	如果 \(\gamma_{x,y}\) 离开较大的坐标球 \(\{|z|<R\}\)，那么从 \(x\) 到首次出界点的一段，坐标位移至少为 \(R-|z(x)|>R/2\)，几何长度便大于 \(\sqrt c\,R/2\)。这与刚得到的长度上界矛盾，所以整条短线留在大球中。

	现在可沿整条线使用\eqref{g3:eq:coordinate-square-convexity}。凸函数不超过其两端点之间的线性插值，故
	\[
		h(\gamma_{x,y}(t))\leq(1-t)h(x)+th(y)<\delta^2
		\qquad(0\leq t\leq1).
	\]
	因此它实际上始终留在 \(U\) 内。最短线在整个流形中的唯一性已由前一个引理给出，于是 \(U\) 强凸。其闭包为闭坐标球，紧且包含于 \(O\)。光滑端点依赖同样由前一个引理继承；最后
	\[
		d(x,y)^2=g_x(\log_xy,\log_xy)
	\]
	是光滑函数，包括 \(x=y\) 时也成立。
\end{proof}

证明使用的是固定坐标中的 \(|z|^2\)，并不要求先知道距离平方的二阶导数。它将两种作用分开：出界长度下界保证所选短线留在可估计的区域，沿线的凸性再保证它留在所需的小邻域。一个点的法邻域未必本身就具有这种两端点凸性。

\subsection{局部唯一性和整体多条最短线}
\label{g3:subsec:0202:03}

最短路的唯一性只在上述邻域范围内得到。另一个局部事实则没有端点接近的限制：如果一条最短路由光滑测地线段拼成，拼接处的方向必须相接。否则可以把角点附近替换成一条更短的小弦。

\begin{lemma}
	\label{g3:lem:corner-shortening}
	设 \((M,g)\) 为无边界光滑黎曼流形，两条单位速测地线段 \(\gamma_-:[-a,0]\to M\)、\(\gamma_+:[0,b]\to M\) 满足 \(a,b>0\) 及 \(\gamma_-(0)=\gamma_+(0)=z\)。若
	\[
		A=\dot\gamma_-(0)\ne\dot\gamma_+(0)=B,
	\]
	则它们的拼接不是两端点之间的最短曲线。更准确地说，在 \(z\) 的任意邻域内都能替换拼接曲线的一小段，使其余部分保持不变而总长度严格减小。
\end{lemma}
\begin{proof}
	在 \(\mathrm{T}_zM\) 中选 \(g_z\) 正交单位基，并用 \(\exp_z\) 给出法坐标。因反向测地线仍是测地线，由初值唯一性，对足够小的 \(s>0\)，有
	\[
		\gamma_-(-s)=\exp_z(-sA),\qquad
		\gamma_+(s)=\exp_z(sB).
	\]
	原拼接段的长度为 \(2s\)。用切空间中的直线段
	\[
		v_s(t)=-sA+ts(A+B),\qquad 0\leq t\leq1,
	\]
	经指数映射后的曲线代替它。\(|A|=|B|=1\) 且 \(A\ne B\)，故 \(|A+B|<2\)。法坐标度量在原点为单位矩阵，由连续性，在 \(|v|\leq s\) 上存在 \(\varepsilon_s\to0\)，使
	\[
		g_{ij}(v)\xi^i\xi^j\leq(1+\varepsilon_s)|\xi|^2.
	\]
	线段 \(v_s\) 留在该球内，因而替代曲线的长度至多为
	\[
		\sqrt{1+\varepsilon_s}\,s|A+B|<2s
	\]
	只要 \(s\) 足够小。若 \(A+B=0\)，两截断端点相同，替代曲线直接取常值即可。将 \(s\) 同时取得小于 \(a,b\)，并使整个坐标小球包含于指定邻域，便得到所需的局部替换。新曲线仍分片光滑，其在替换端点出现的角点不影响它作为长度竞争曲线的资格。
\end{proof}

\begin{corollary}
	\label{g3:cor:minimizing-curve-geodesic}
	在连通无边界黎曼流形中，每条按弧长参数化的分片光滑最短曲线都是光滑的单位速测地线。特别地，两条单位速测地线段若拼成一条最短曲线，则它们在接点的速度相同，因而由初值唯一性属于同一条测地线。
\end{corollary}
\begin{proof}
	最短曲线的每个子段仍最短，否则替换该子段即可缩短全曲线。对曲线任一内点，使用\cref{g3:lem:uniform-normal-neighborhood}选取邻域 \(V\) 与半径 \(r\)，再缩小该点两侧的参数区间，使两端点都在 \(V\) 中，子段长度小于 \(r\)。引理中的等号刻画说明，这个小段是唯一径向最短线的单调重参数化；弧长参数使重参数化在非零长度小段上为平移。因此曲线在每个内点附近都是光滑单位速测地线，在端点具有相应的单侧延拓。特别结论也可直接由\cref{g3:lem:corner-shortening}排除方向不接的情形，再用\cref{g3:thm:geodesic-ivp}连接两段。
\end{proof}

这个结果只说明达到距离的曲线必须满足测地线方程。逆命题需要长度范围：一条测地线可以在短时间内最短，继续走下去却变得更长；两点之间还可能有多条同样短的测地线。

\begin{example}
	\label{g3:exm:circle-distance}
	半径 \(a>0\) 的圆周写为 \(\theta\mapsto(a\cos\theta,a\sin\theta)\)，诱导度量为 \(a^2\mathrm{d}\theta^2\)。若两点的角度代表为 \(\alpha,\beta\)，则
	\begin{equation}
		\label{g3:eq:circle-distance}
		d(\alpha,\beta)=a\min_{k\in\mathbb Z}|\beta-\alpha+2\pi k|.
	\end{equation}
	为证明下界，沿任意连接曲线连续选取实值角度，使起始角为 \(\alpha\)。这种选择可先在一个角度坐标片中进行，再沿有限个相邻坐标片以加上 \(2\pi\) 的整数倍来接续；初始值固定后，交叠处连续性使接续唯一。终点角度必为 \(\beta+2\pi k\)，而曲线长度满足
	\[
		L=a\int_0^1|\dot\theta(t)|\,\mathrm{d}t
		\geq a|\beta-\alpha+2\pi k|.
	\]
	对使右侧最小的 \(k\)，用线性角度参数就取得等号，证明距离公式。

	若两点不对径，较短圆弧是唯一最短线；若两点对径，两个半圆都长 \(\pi a\)，给出两条最短线。再绕一周的常速圆弧仍是测地线，却不再最短。因此强凸邻域中的唯一性是最短线的唯一性，不排除连接同两点而绕行到邻域外的较长测地线。
\end{example}

\begin{example}
	\label{g3:exm:sphere-multiple-minimizers}
	在半径 \(a\) 的标准球面 \(S_a^n\) 上，\cref{g3:exm:round-sphere-exponential}给出 \(\exp_p\) 在 \(B_p(0,\pi a)\) 上的微分同胚，其像为去掉对径点 \(-p\) 的球面。对 \(q\ne-p\)，径向极小性因此给出唯一最短线，长度为
	\[
		d(p,q)=a\arccos\frac{\langle p,q\rangle}{a^2}.
	\]
	让 \(q\) 沿一个大圆趋于 \(-p\)。由\cref{g3:thm:riemannian-distance}及距离的连续性，得到 \(d(p,-p)=\pi a\)。每个单位切向量 \(u\in\mathrm{T}_pS_a^n\) 都给出一条达到这个距离的半大圆
	\[
		\gamma_u(t)=\cos(t/a)p+a\sin(t/a)u,
		\qquad 0\leq t\leq\pi a.
	\]
	当 \(n\geq2\) 时，\(u\) 可以在单位切球面上连续变化，所以对径点之间有无穷多条最短线；\(n=1\) 时恰好回到圆周的两条半圆。局部端点参数化在对径点处失去单值性，这为后面的割迹提供了第一个模型。
\end{example}

% 纲要条目排为编号小节标题，并进入目录与书签；说明段仍保留于纲要注释中。

% !TeX root = ../../Book3.tex
% 纲要标题：### 2.3 Hopf–Rinow 定理
% 纲要位置：几何学计划纲要.md，第 2072 行。
% 正文状态：未撰写；本文件目前是框架。

% 纲要内容开始
% 
% * 度量完备推出延拓
% * 从测地完备构造最短测地线
% * 有界闭集紧性
% * 等价条件的完整闭环
% * 非完备和一般长度空间的对照
% 
% 纲要内容结束

\section{Hopf–Rinow 定理}
\label{g3:sec:0203}

待撰写

\subsection{度量完备推出延拓}
\label{g3:subsec:0203:01}
待撰写

\subsection{从测地完备构造最短测地线}
\label{g3:subsec:0203:02}
待撰写

\subsection{有界闭集紧性}
\label{g3:subsec:0203:03}
待撰写

\subsection{等价条件的完整闭环}
\label{g3:subsec:0203:04}
待撰写

\subsection{非完备和一般长度空间的对照}
\label{g3:subsec:0203:05}
待撰写
% 纲要条目排为编号小节标题，并进入目录与书签；说明段仍保留于纲要注释中。

% !TeX root = ../../Book3.tex
% 纲要标题：### 2.4 割迹与注入半径
% 纲要位置：几何学计划纲要.md，第 2183 行。
% 正文状态：初稿；尚未正式审核。

% 纲要内容开始
% 
% * 最大极小区间与割点函数
% * 指数映射的单射范围
% * 注入半径及紧流形下界
% * 球面、平坦环面和商空间模型
% 
% ---
% 
% 纲要内容结束

\section{割迹与注入半径}
\label{g3:sec:0204}

在一个法邻域内，径向测地线既唯一又最短；把它继续延长，这两个性质都可能改变。圆周上的短弧到达对径点时出现另一条等长弧，越过对径点后原来的弧就不再最短。我们先沿每个方向记录这种变化，再讨论指数映射在哪些半径内仍能给出坐标。本节除另有说明外，假设 \((M,g)\) 连通、完备、无边界，且维数 \(n\ge1\)。

\subsection{最大极小区间与割点函数}
\label{g3:subsec:0204:01}
固定 \(p\in M\)，记 \(S_pM=\{u\in \mathrm{T}_pM:|u|_g=1\}\)，并令 \(\gamma_u(t)=\exp_p(tu)\)。完备性保证它对所有实数 \(t\) 有定义。沿这条单位速测地线，长度总是参数差，端点距离却可能更小。

\begin{definition}[割时与割迹]\label{g3:def:cut-time}
	方向 \(u\in S_pM\) 的\term[geshi]{割时}[cut time]定义为
	\[
		c_p(u)=\sup\{t\ge0:d(p,\gamma_u(t))=t\}\in(0,\infty].
	\]
	若 \(c_p(u)<\infty\)，称 \(\gamma_u(c_p(u))\) 为该方向的\term[gedian]{割点}[cut point]。所有这样的点组成 \(p\) 的\term[geji]{割迹}[cut locus]，记为
	\[
		\operatorname{Cut}(p)=\{\exp_p(c_p(u)u):u\in S_pM,\ c_p(u)<\infty\}.
	\]
	\symidx{\(c_p(u)\)}
	\symidx{\(\operatorname{Cut}(p)\)}
\end{definition}

定义中的集合确实是一个区间。若 \(d(p,\gamma_u(t))=t\) 且 \(0\le s\le t\)，三角不等式给出
\[
	t=d(p,\gamma_u(t))\le d(p,\gamma_u(s))+(t-s)\le s+(t-s)=t.
\]
因此两处不等式都取等号，\(d(p,\gamma_u(s))=s\)。\cref{g3:prop:normal-radial-minimality}又保证足够小的正数 \(s\) 满足这个等式，所以 \(c_p(u)>0\)。若 \(c_p(u)\) 有限，距离的连续性使等式在端点仍成立。于是最大极小区间是 \([0,c_p(u)]\)；若割时无穷，就是 \([0,\infty)\)。在有限割时之后，每个 \(t>c_p(u)\) 都满足严格不等式 \(d(p,\gamma_u(t))<t\)。

\ProseParagraphBreak

这里的“极小”表示这段曲线在所有连接端点的曲线中长度最小。每条测地线在每个参数值附近仍有足够短的极小子段；越过割点所失去的是从原起点 \(p\) 算起的整体极小性。

\begin{proposition}\label{g3:prop:cut-time-upper-semicontinuity}
	割时函数 \(c_p:S_pM\to(0,\infty]\) 是上半连续的：若 \(u_j\to u\)，则
	\[
		\limsup_{j\to\infty}c_p(u_j)\le c_p(u).
	\]
\end{proposition}
\begin{proof}
	设 \(t<c_p(u_j)\) 对某个子列的所有项成立。沿该子列有 \(d(p,\exp_p(tu_j))=t\)；由指数映射与距离的连续性，极限给出 \(d(p,\exp_p(tu))=t\)，所以 \(t\le c_p(u)\)。若上极限大于 \(c_p(u)\)，在二者之间取这样的 \(t\) 即得矛盾。\(c_p(u)=\infty\) 时不等式自动成立。
\end{proof}

上半连续性说明：当极小线段的初方向与长度同时收敛时，极限线段仍极小。它来自端点距离等式的封闭性，不需要曲率计算。

\subsection{指数映射的单射范围}
\label{g3:subsec:0204:02}
令
\[
	\mathcal D_p=\{0\}\cup\{tu:u\in S_pM,\ 0<t<c_p(u)\}\subset \mathrm{T}_pM.
\]
这保留每条射线上严格位于割时之前的向量。以下结论首先说明它在集合意义上如何参数化割迹的补集。

\begin{proposition}[割点之前的唯一性]\label{g3:prop:precut-injectivity}
	若 \(0<t<c_p(u)\)，则 \(\gamma_u|_{[0,t]}\) 是从 \(p\) 到 \(q=\gamma_u(t)\) 的唯一单位速最短测地线。指数映射限制给出双射
	\[
		\exp_p:\mathcal D_p\longrightarrow M\setminus\operatorname{Cut}(p).
	\]
\end{proposition}
\begin{proof}
	取 \(\varepsilon>0\)，使 \(t+\varepsilon<c_p(u)\)。若另有一条从 \(p\) 到 \(q\) 的单位速最短测地线 \(\sigma:[0,t]\to M\)，将它与 \(\gamma_u|_{[t,t+\varepsilon]}\) 拼接，所得曲线长 \(t+\varepsilon\)，仍等于两端距离。由\cref{g3:lem:corner-shortening}，接点不能有角点，所以 \(\dot\sigma(t)=\dot\gamma_u(t)\)。测地线初值唯一性向后应用，得到 \(\sigma=\gamma_u|_{[0,t]}\)。

	\(\mathcal D_p\) 中的非零向量长度等于其像到 \(p\) 的距离。同一点的两个原像因而有相同长度，由刚证的唯一性又有相同方向，故限制映射单射。若其像 \(q\) 同时为某个方向的割点，沿那个方向到 \(q\) 的线段也是最短线，只能与已有线段相同；这样同一方向的同一参数既小于割时又等于割时，矛盾。

	最后取 \(q\notin\operatorname{Cut}(p)\)、\(q\ne p\)。Hopf--Rinow 定理给出一条从 \(p\) 到 \(q\) 的单位速最短测地线，写成 \(\gamma_u|_{[0,t]}\)，其中 \(t=d(p,q)\le c_p(u)\)。等号会使 \(q\) 成为割点，所以 \(t<c_p(u)\)。点 \(p\) 则由零向量表示。
\end{proof}

这个证明还给出一个实用判据：若同一点由两条不同的最短测地线到达，则它必在起点的割迹中。反向陈述并不由唯一性证明得到。割时处指数映射的微分也可能发生退化；下面把这两种可能性分开。

\begin{proposition}[有限割时处的两种可能]\label{g3:prop:cut-endpoint-alternative}
	设 \(c=c_p(u)<\infty\)，\(q=\exp_p(cu)\)。则至少发生下列一件事：
	\begin{enumerate}
		\item 从 \(p\) 到 \(q\) 存在两条不同的最短测地线；
		\item \((\mathrm d\exp_p)_{cu}\) 不可逆。
	\end{enumerate}
\end{proposition}
\begin{proof}
	假设最短测地线唯一，且该微分可逆。取 \(t_j\downarrow c\)、\(t_j>c\)，令 \(q_j=\exp_p(t_ju)\)。由 Hopf--Rinow 定理，可取 \(v_j\in \mathrm{T}_pM\)，使
	\[
		\exp_p(v_j)=q_j,\qquad |v_j|=d(p,q_j)<t_j.
	\]
	距离连续性给出 \(|v_j|\to c\)。有限维切空间中的有界序列有收敛子列，记极限为 \(v\)。则 \(|v|=c\)、\(\exp_p(v)=q\)，相应径向线长等于端点距离，故是最短线。唯一性使 \(v=cu\)。逆函数定理给出 \(cu\) 的一个邻域，\(\exp_p\) 在其中单射。子列的 \(v_j\) 与 \(t_ju\) 最终都在此邻域内，像相同，却因长度不同而不相等，矛盾。
\end{proof}

微分退化如何由沿测地线的变分场检测，将在\cref{g3:sec:0602}建立。\cref{g3:sec:0603}再把它与割时联系起来。这里的双射与上述二择结论已经足以解释以下模型中指数坐标失效的位置。

\subsection{注入半径及紧流形下界}
\label{g3:subsec:0204:03}
即使不假设完备，\cref{g3:prop:exponential-local}也保证指数映射在零向量附近是微分同胚。以整个切球作为坐标范围所能达到的最大半径，给出一个可以逐点比较的尺度。

\begin{definition}[注入半径]\label{g3:def:injectivity-radius}
	对无边界黎曼流形上的 \(p\)，定义 \(p\) 的\term[zhurubanjing]{注入半径}[injectivity radius]为
	\[
		\operatorname{inj}_g(p)=\sup\bigl\{r>0:\exp_p\text{ 在 }B_r(0)\subset \mathrm{T}_pM
		\text{ 上有定义且为到其像的微分同胚}\bigr\}.
	\]
	流形的注入半径定义为 \(\operatorname{inj}_g(M)=\inf_{p\in M}\operatorname{inj}_g(p)\)。上确界允许为无穷，整体下确界则可能为零。
	\symidx{\(\operatorname{inj}_g(p)\)}
\end{definition}

\begin{proposition}\label{g3:prop:compact-injectivity-lower-bound}
	每点的注入半径为正，且在任意紧子集上有统一正下界。特别地，紧黎曼流形的注入半径为正。若流形还完备，则对 \(0<r<\operatorname{inj}_g(p)\)，
	\[
		\exp_p(B_r(0))=B_d(p,r),\qquad
		d(p,\exp_p v)=|v|\quad(|v|<r),
	\]
	并且 \(\operatorname{inj}_g(p)\le\inf_{u\in S_pM}c_p(u)\)。
\end{proposition}
\begin{proof}
	逐点正性来自逆函数定理。\cref{g3:lem:uniform-normal-neighborhood}进一步给出：每点 \(p\) 有邻域 \(V_p\) 和常数 \(r_p>0\)，使所有 \(x\in V_p\) 的指数映射在 \(B_{r_p}(0)\subset \mathrm{T}_xM\) 上均为微分同胚。有限个 \(V_p\) 覆盖给定紧集；取对应 \(r_p\) 的最小值即可。

	现在假设完备，并取 \(r<\operatorname{inj}_g(p)\)。对 \(|v|<r\)，径向线给出 \(d(p,\exp_pv)\le|v|\)。若严格小于，取一条最短测地线，把其初向量按终点参数一缩放为 \(w\)，便有 \(|w|=d(p,\exp_pv)<|v|<r\) 且 \(\exp_pw=\exp_pv\)，与切球内的单射性矛盾。因此距离等式成立。反过来，\(d(p,q)<r\) 时最短测地线给出长度小于 \(r\) 的原像，所以像正是距离球。对每个方向与所有 \(t<r\) 应用距离等式，得到 \(c_p(u)\ge r\)；让 \(r\) 趋向注入半径即得最后的不等式。
\end{proof}

这一正下界的证明只用局部指数映射和紧性。在非紧流形上，各点都有法邻域，却未必能选同一个半径。习题中的尖缩圆柱会给出完备但整体注入半径为零的例子。一般流形上注入半径与最早割时的相等关系，连同割时函数的进一步正则性，将随第6章的指数映射微分理论完成。

\subsection{球面、平坦环面和商空间模型}
\label{g3:subsec:0204:04}
\begin{example}[圆球面的割迹]\label{g3:exm:sphere-cut-locus}
	在半径 \(a\) 的圆球面 \(S_a^n\)、\(n\ge1\) 上，\cref{g3:exm:round-sphere-exponential}给出
	\[
		\exp_p(tu)=\cos(t/a)p+a\sin(t/a)u\qquad(|u|=1).
	\]
	这些曲线全时存在，所以球面完备。对 \(q\ne p\)，取到 \(q\) 的最短测地线，长度记为 \(\ell\)。显式公式中长度超过 \(\pi a\) 的弧总能换成同一大圆上长度至多 \(\pi a\) 的弧：先减去整圈，再在长弧和短弧之间选择。因此 \(0<\ell\le\pi a\)，并且
	\[
		d(p,q)=a\arccos\frac{\langle p,q\rangle}{a^2}.
	\]
	从而 \(c_p(u)=\pi a\) 对所有方向成立，\(\operatorname{Cut}(p)=\{-p\}\)。指数映射在 \(B_{\pi a}(0)\) 上是到去掉对径点的球面的微分同胚，而任意更大的切球内都有两个向量 \(\pi a u\)、\(-\pi a u\) 映到 \(-p\)。故
	\[
		\operatorname{inj}(S_a^n)=\pi a.
	\]
	当 \(n=1\)，到对径点有两条半圆最短弧；当 \(n\ge2\)，每个单位初方向都给出一条最短半大圆。对于 \(n=1\)，指数映射在割时处仍是局部微分同胚，失效的是单射性。
\end{example}

\begin{example}[平坦环面的最近格点区域]\label{g3:exm:flat-torus-cut-locus}
	设 \(\Lambda\subset\mathbb R^n\) 为由一组基生成的格，\(T=\mathbb R^n/\Lambda\) 取 Euclid 度量诱导的商度量。商映射在足够小的球上是等距坐标图，所以测地线就是直线的投影，且 \(T\) 完备。若 \(p=[x]\)、\(q=[y]\)，每条从 \(p\) 到 \(q\) 的曲线从 \(x\) 提升后终止于某个 \(y+\lambda\)，长度不变；其长度至少为 \(|y-x+\lambda|\)。直线段达到该下界。因此
	\begin{equation}\label{g3:eq:flat-torus-distance}
		d_T([x],[y])=\min_{\lambda\in\Lambda}|y-x+\lambda|.
	\end{equation}
	最小值确实取得：任何有界 Euclid 球只含有限个格点。所用路径提升来自第一卷第5.3节的道路提升定理；这里的平移商在小球上满足其均匀覆盖条件。局部逆光滑，所以分片光滑曲线的提升也分片光滑。

	将 \(\mathrm{T}_pT\) 与 \(\mathbb R^n\) 认同，则 \(\exp_pv=[x+v]\)。定义格的最近点区域
	\[
		V=\{v\in\mathbb R^n:|v|\le|v-\lambda|\text{ 对所有 }\lambda\in\Lambda\}
		=\bigcap_{\lambda\in\Lambda\setminus\{0\}}
		\{v:2\langle v,\lambda\rangle\le|\lambda|^2\}.
	\]
	它是闭凸集。取一个有界的基本平行多面体，可知任一点到格的距离有统一上界；\(v\in V\) 时最近格点是零，故 \(|v|\) 也有统一上界，所以 \(V\) 有界。在任意有界区域内，足够长的格向量自动满足严格不等式，因而只有有限个不等式影响边界。由此，\(V\) 的内部由所有严格不等式刻画，边界点至少满足一个非零格向量对应的等式。

	径向线段到 \([x+tu]\) 最短，当且仅当 \(tu\in V\)。每条射线与 \(V\) 相交为一个闭区间，所以
	\[
		c_p(u)=\min_{\substack{\lambda\in\Lambda\\\langle u,\lambda\rangle>0}}
		\frac{|\lambda|^2}{2\langle u,\lambda\rangle},\qquad
		\operatorname{Cut}(p)=[x+\partial V].
	\]
	式中的最小值存在：\(V\) 有界使射线在有限参数退出，而该边界点处有一个起约束作用的格向量。若 \(v\in\partial V\) 且 \(|v|=|v-\lambda|\)，向量 \(v\) 与 \(v-\lambda\) 给出同一终点的两条最短线。切空间中的相邻区域因此在割迹上发生识别。

	设 \(\ell_\Lambda=\min_{\lambda\ne0}|\lambda|\)。指数映射微分处处可逆；若 \(|v|,|w|<\ell_\Lambda/2\) 且像相同，则 \(v-w\in\Lambda\)，但 \(|v-w|<\ell_\Lambda\)，故 \(v=w\)。任意大于 \(\ell_\Lambda/2\) 的切球又包含一对 \(\lambda/2,-\lambda/2\)，其中 \(|\lambda|=\ell_\Lambda\)，它们像相同。于是
	\[
		\operatorname{inj}(T)=\tfrac12\ell_\Lambda.
	\]
\end{example}

\begin{figure}[htbp]
	\centering
	\input{Figures/Ch02-TorusCutLocus}
	\caption[矩形平坦环面的割迹]{矩形平坦环面的割迹。取格 \(L\mathbb Z\times H\mathbb Z\)，图中矩形为原点的最近点区域 \([-L/2,L/2]\times[-H/2,H/2]\)。相对边在商中识别，整条边界的像组成基点的割迹；边中点有两条最短线，角点有四条。}
	\label{g3:fig:torus-cut-locus}
\end{figure}

覆盖空间也能用来计算其他商度量。下面先把所需的距离公式和完备性传递写清楚。

\begin{proposition}[黎曼覆盖的距离公式]\label{g3:prop:riemannian-covering-distance}
	设 \(\pi:(\widetilde M,\widetilde g)\to(M,g)\) 是连通无边界黎曼流形之间的覆盖，且每点附近的限制是等距映射。若 \(\widetilde M\) 完备，则 \(M\) 完备。对固定提升 \(\widetilde p\in\pi^{-1}(p)\)，有
	\[
		d_M(p,q)=\min_{\widetilde q\in\pi^{-1}(q)}
		d_{\widetilde M}(\widetilde p,\widetilde q).
	\]
	若覆盖来自等距变换群 \(\Gamma\) 的自由作用且纤维正是其轨道，右边等于 \(\min_{\gamma\in\Gamma}d_{\widetilde M}(\widetilde p,\gamma\widetilde q_0)\)，其中 \(\widetilde q_0\) 是任一固定提升。
\end{proposition}
\begin{proof}
	局部等距保持 Levi-Civita 联络：将一个联络通过局部等距搬运，仍满足度量相容和无挠，由唯一性就是另一侧的 Levi-Civita 联络。因此它保持测地线与长度。给定 \(M\) 上的测地线初值，在 \(\widetilde M\) 选取对应初值，所得全时测地线投影后就是原初值的全时解。故 \(M\) 测地完备，再由 Hopf--Rinow 定理得到度量完备。

	在 \(M\) 上取从 \(p\) 到 \(q\) 的最短测地线，从 \(\widetilde p\) 将其提升，终点记为 \(\widetilde q\)。提升长度仍为 \(d_M(p,q)\)，所以 \(d_{\widetilde M}(\widetilde p,\widetilde q)\le d_M(p,q)\)。反向不等式对任意提升终点都成立，因为连接提升端点的曲线投影后长度不变。两者合并，既得到等式，也证明最小值取得。轨道表述只是把纤维中的点重新标记。
\end{proof}

\begin{example}[实射影空间]\label{g3:exm:projective-cut-locus}
	令 \(\mathbb{RP}^n=S_a^n/\{p\sim-p\}\)，取使商映射局部等距的度量。反极点作用自由，足够小的球与其反极点球不交，所以商映射是二重覆盖。令 \(\theta=\arccos(\langle p,q\rangle/a^2)\)，上一命题给出
	\[
		d([p],[q])=a\min\{\theta,\pi-\theta\}
		=a\arccos\frac{|\langle p,q\rangle|}{a^2}.
	\]
	从 \([p]\) 出发的单位速测地线提升为从 \(p\) 出发的大圆；它在 \(t\le\pi a/2\) 时最短，随后一段时间选取反极点提升更短。因此每个方向的割时为 \(\pi a/2\)，割迹为赤道 \(\{q:\langle p,q\rangle=0\}\) 在商中的像，即嵌入的 \(\mathbb{RP}^{n-1}\)。在半径 \(\pi a/2\) 的开切球上，球面指数映射的像是以 \(p\) 为中心的开半球，它与自身的反极点像不交，故商后的指数映射是微分同胚。任意更大的切球内，\((\pi a/2)u\) 与 \(-(\pi a/2)u\) 的像相同，故
	\[
		\operatorname{inj}(\mathbb{RP}^n)=\pi a/2.
	\]
\end{example}

球面在一个点处汇集所有割方向，矩形平坦环面把最近点区域的边界粘成图形，实射影空间则出现整个射影超平面。这些差别来自指数映射的全局行为；点附近总能建立法坐标，并不能决定较远处最短线的条数。
% 纲要条目排为编号小节标题，并进入目录与书签；说明段仍保留于纲要注释中。


\begin{chaptersummary}
	测地线由初始位置和速度唯一确定，速度长度恒定。指数映射把初速度变成时间一的终点；其零处微分为恒等映射，所以零向量附近有法坐标。Gauss 引理进一步说明径向与角向变化正交，从而给出径向线的长度下界及等号条件。出界曲线也必须先付出足够的长度，这使局部计算能够与流形中的所有竞争曲线比较。

	\ProseParagraphBreak

	曲线长度的下确界是与流形拓扑一致的距离。双端点指数映射提供统一的短线，坐标平方函数沿短测地线的凸性则给出强凸邻域。最短曲线不能有角点，按弧长参数化后必为测地线；反过来，一条测地线延长后可能不再最短。

	\ProseParagraphBreak

	Hopf--Rinow 定理把度量完备、所有测地线全时存在、某个基点的指数映射全域定义，以及闭有界集紧性连成等价闭环。完备性还保证任意两点之间存在最短测地线。割时记录这种极小性沿某个方向能保持多久；割时之前的最短线唯一，有限割时处则出现多条最短线或指数映射微分退化。注入半径给出指数坐标可用的球半径，并在紧集上有统一正下界。球面、平坦环面与实射影空间展示了不同的全局识别方式。
\end{chaptersummary}

\BookExercises

\begin{exercise}\label{g3:ex:exponential-line-incomplete}
	在 \(M=\mathbb R\) 上取度量 \(g=e^{2x}\mathrm{d}x^2\)。
	\begin{enumerate}
		\item 证明 \(F(x)=e^x\) 是到 Euclid 半直线 \((0,\infty)\) 的等距微分同胚，并求内在距离 \(d_g(p,q)\)。
		\item 对初值 \(\gamma(0)=p\)、\(\dot\gamma(0)=v\partial_x|_p\)，求最大测地线及其定义区间，进而求 \(\exp_p\) 的定义域和表达式。
		\item 按\cref{g3:def:injectivity-radius}计算 \(\operatorname{inj}_g(p)\)，并证明此流形不完备。
	\end{enumerate}
\end{exercise}
\begin{solution}
	1. \(F\) 是 \(\mathbb R\) 到 \((0,\infty)\) 的光滑双射，其逆为对数函数，且
	\[
		F^*(\mathrm{d}y^2)=(e^x\mathrm{d}x)^2=g.
	\]
	因此它保持曲线长度。半直线内任意连接 \(e^p,e^q\) 的曲线长度至少为 \(|e^q-e^p|\)，连接端点的直线段仍在半直线内并达到这个下界，故
	\[
		d_g(p,q)=|e^q-e^p|.
	\]

	2. 在坐标 \(y=F(x)\) 中，测地线方程是 \(\ddot y=0\)。初速度变为 \(\dot y(0)=e^pv\)，所以
	\[
		y(t)=e^p(1+vt),\qquad
		\gamma(t)=p+\log(1+vt).
	\]
	允许的最大时间区间由 \(1+vt>0\) 决定：
	\[
		I_{p,v}=
		\begin{cases}
			(-1/v,\infty),&v>0,\\
			\mathbb R,&v=0,\\
			(-\infty,-1/v),&v<0.
		\end{cases}
	\]
	有限端点处 \(y(t)\to0\)，对应 \(\gamma(t)\to-\infty\)，没有流形内的延拓点。时间一可取当且仅当 \(v>-1\)，故
	\[
		\mathcal E_p=\{v\partial_x|_p:v>-1\},\qquad
		\exp_p(v\partial_x|_p)=p+\log(1+v).
	\]

	3. 切向量的度量长度为 \(g_p(v\partial_x,v\partial_x)^{1/2}=e^p|v|\)，因而半径 \(r\) 的切球对应区间 \((-re^{-p},re^{-p})\)。当 \(0<r\le e^p\) 时，这个开区间包含于 \(v>-1\)，且指数映射的导数 \(1/(1+v)\) 为正，故指数映射在该球上为微分同胚。若 \(r>e^p\)，切球包含 \(v=-1\)，指数映射在此没有定义。因此
	\[
		\operatorname{inj}_g(p)=e^p.
	\]
	序列 \(p_j=-j\) 关于 \(d_g\) 为 Cauchy 序列，因为 \(F(p_j)=e^{-j}\to0\)；它在 \(M\) 中没有极限。也可以取 \(v=-1\)，得到在时间一终止的最大测地线。这分别检验了度量不完备和测地不完备。
\end{solution}

\begin{exercise}\label{g3:ex:complete-product-distance-injectivity}
	设 \((M,g)\)、\((N,h)\) 为连通、无边界的完备光滑黎曼流形，乘积上取度量 \(G=g\oplus h\)。证明
	\[
		d_G((p,q),(p',q'))^2=d_g(p,p')^2+d_h(q,q')^2,
	\]
	并证明 \(M\times N\) 完备。求乘积的指数映射，再按定义证明
	\[
		\operatorname{inj}_G(p,q)
		=\min\{\operatorname{inj}_g(p),\operatorname{inj}_h(q)\},
	\]
	其中允许半径为无穷。
\end{exercise}
\begin{solution}
	对乘积曲线 \(c=(\alpha,\beta):[0,1]\to M\times N\)，记两个分量的速度长度为 \(a(t),b(t)\)，曲线长度为 \(L_\alpha,L_\beta\)。积分的向量三角不等式给出
	\[
		L_G(c)=\int_0^1\sqrt{a(t)^2+b(t)^2}\,\mathrm{d}t
		\ge\sqrt{L_\alpha^2+L_\beta^2}
		\ge\sqrt{d_g(p,p')^2+d_h(q,q')^2}.
	\]
	第一处不等式也可直接检验：若 \((L_\alpha,L_\beta)\ne(0,0)\)，将 \((a(t),b(t))\) 与其积分向量的单位方向作内积，再逐点用内积不等式，积分即可。零向量情形直接成立。

	由\cref{g3:thm:hopf-rinow}，两个因子中均存在实现端点距离的测地线。将它们都参数化在 \([0,1]\) 上，使速度长度分别恒为 \(d_g(p,p')\)、\(d_h(q,q')\)；距离为零的因子取常曲线。它们组成的乘积曲线速度长度恒为这两个数的平方和开根，达到上述下界。因此距离公式成立。乘积中的 Cauchy 序列投影到两个因子后仍是 Cauchy 序列，各有极限；距离公式又说明乘积序列收敛到两极限组成的点，故乘积完备。

	由\cref{g3:prop:product-connection}，乘积测地线方程分成两个因子的测地线方程，初值唯一性给出
	\[
		\exp^G_{(p,q)}(v,w)
		=(\exp^g_pv,\exp^h_qw),\qquad
		|(v,w)|_G^2=|v|_g^2+|w|_h^2.
	\]
	完备性保证这些指数映射处处有定义。若 \(r\) 小于两个因子的注入半径，则各指数映射在半径 \(r\) 的切球上都是微分同胚，它们的乘积在两个切球的乘积上也是微分同胚。半径 \(r\) 的乘积切球是其中的开子集，故限制后仍为微分同胚，得到所求等式的“\(\ge\)”方向。

	反过来，若乘积指数映射在半径 \(r\) 的切球上为微分同胚，则限制到 \((v,0)\)、\(|v|_g<r\)，可知 \(\exp_p^g\) 在 \(B_p(0,r)\) 上单射。同时
	\[
		\mathrm{d}\exp^G_{(p,q)}\big|_{(v,0)}
		=\operatorname{diag}\bigl((\mathrm{d}\exp^g_p)_v,
			(\mathrm{d}\exp^h_q)_0\bigr)
	\]
	可逆，且第二块为恒等，所以第一块可逆。于是 \(\exp_p^g\) 是单射的局部微分同胚，其局部光滑逆拼成像上的光滑逆；它在整个 \(B_p(0,r)\) 上为微分同胚。交换两个因子，得到同样的结论。故任意乘积可用半径都是两个因子各自的可用半径，对上确界取值即得“\(\le\)”方向，也包括无穷半径的情形。
\end{solution}

\begin{exercise}\label{g3:ex:complete-metric-lower-bound}
	设 \(g,h\) 是同一连通、无边界光滑流形 \(M\) 上的黎曼度量，且 \((M,g)\) 完备。
	\begin{enumerate}
		\item 若存在常数 \(c>0\)，使 \(h(\xi,\xi)\ge c^2g(\xi,\xi)\) 对一切切向量成立，证明 \((M,h)\) 完备。说明仅有距离下界以后，还需怎样把 \(g\) 意义下的极限变成 \(h\) 意义下的极限。
		\item 证明全局上界 \(h\le C^2g\) 不能替代这一假设。可在 \(\mathbb R\) 上取
		\[
			g=\mathrm{d}x^2,\qquad
			h=e^{-2\sqrt{1+x^2}}\mathrm{d}x^2.
		\]
	\end{enumerate}
\end{exercise}
\begin{solution}
	1. 对每条分片光滑曲线积分速度下界，得到 \(L_h\ge cL_g\)，继而
	\[
		d_h(x,y)\ge c\,d_g(x,y).
	\]
	所以 \(h\)-Cauchy 序列 \((p_j)\) 也是 \(g\)-Cauchy 序列。由 \(g\) 完备，存在 \(p\in M\)，使 \(d_g(p_j,p)\to0\)。\cref{g3:thm:riemannian-distance}说明 \(p_j\) 在流形拓扑中也趋于 \(p\)。

	取以 \(p\) 为原点的坐标图，并在其中取闭包紧含于图域的坐标球。度量 \(h\) 的系数在这个闭球上有界，因此存在 \(A>0\)，使 \(h_{ab}(x)\xi^a\xi^b\le A^2|\xi|^2\)。充分大的 \(j\) 位于较小坐标球中，将 \(p_j\) 与 \(p\) 用坐标直线段连接，便有
	\[
		d_h(p_j,p)\le A\,|x(p_j)-x(p)|\longrightarrow0.
	\]
	这一步只在极限点附近使用度量上界；全局不要求 \(h\) 被 \(g\) 的常数倍控制。由此每个 \(h\)-Cauchy 序列都收敛，故 \(h\) 完备。

	2. 所给度量光滑正定，且 \(h\le g\)。令
	\[
		F(x)=\int_0^x e^{-\sqrt{1+s^2}}\,\mathrm{d}s,\qquad
		A=\int_0^\infty e^{-\sqrt{1+s^2}}\,\mathrm{d}s.
	\]
	因 \(\sqrt{1+s^2}\ge s\) 对 \(s\ge0\) 成立，有 \(0<A\le1\)。\(F'\) 处处为正，且被积函数为偶函数，所以 \(F\) 是到有限开区间 \((-A,A)\) 的微分同胚，满足 \(F^*(\mathrm{d}y^2)=h\)。因此
	\[
		d_h(x,y)=|F(x)-F(y)|.
	\]
	序列 \(x_j=j\) 映到趋于 \(A\) 的序列，是 \(h\)-Cauchy 序列，却没有流形内的极限。缩小远处的长度尺度可以使通往无穷远的路径具有有限长度，因此全局上界 \(h\le g\) 不能保证完备性。
\end{solution}

\begin{exercise}\label{g3:ex:squared-distance-at-center}
	设 \((M,g)\) 为无边界光滑黎曼流形，\(\exp_p:B_p(0,R)\to U_R\) 为微分同胚，令 \(r(q)=d_g(p,q)\)。证明 \(r\) 在 \(U_R\setminus\{p\}\) 上光滑，并求其梯度，证明 \(|\operatorname{grad}_g r|_g=1\)。再证明 \(f=r^2\) 在整个 \(U_R\) 上光滑，且
	\[
		(H_f)_p=2g_p.
	\]
	这里沿用协变二阶导数的定义
	\[
		H_f(X,Y)=X(Yf)-(\nabla_XY)f.
	\]
\end{exercise}
\begin{solution}
	由\cref{g3:prop:normal-radial-minimality}，若 \(q=\exp_p(v)\)，则 \(r(q)=|v|_g\)。因此 \(r\) 在去掉中心的法邻域中光滑。固定非零 \(v\)，令 \(u=v/|v|_g\)。对任意 \(w\in\mathrm{T}_pM\)，有
	\[
		\mathrm{d}r_q\bigl((\mathrm{d}\exp_p)_vw\bigr)
		=\frac{g_p(v,w)}{|v|_g}.
	\]
	Gauss 引理把右端改写为
	\[
		g_q\bigl((\mathrm{d}\exp_p)_vu,(\mathrm{d}\exp_p)_vw\bigr).
	\]
	由于 \((\mathrm{d}\exp_p)_v\) 可逆，这对 \(q\) 处的全部切向量确定了梯度：
	\[
		\operatorname{grad}_g r\big|_q=(\mathrm{d}\exp_p)_vu=\partial_r|_q,
		\qquad |\operatorname{grad}_g r|_g=1.
	\]

	取由 \(\mathrm{T}_pM\) 的正交单位基得到的法坐标，则
	\[
		f(x)=r(x)^2=\sum_i(x^i)^2.
	\]
	此表达在中心也光滑，且 \(\mathrm{d}f_p=0\)。由 \(H_f\) 的定义直接得到坐标分量
	\[
		(H_f)_{ij}=\partial_i\partial_jf-\Gamma^k_{ij}\partial_kf.
	\]
	在 \(p\) 处第二项为零，第一项是 \(2\delta_{ij}\)。又 \(g_{ij}(p)=\delta_{ij}\)，故 \((H_f)_p=2g_p\)。距离函数本身在中心有绝对值型行为，平方则消除了这一不可微点。
\end{solution}

\begin{exercise}\label{g3:ex:completeness-escaping-curves}
	设 \((M,g)\) 为连通、无边界光滑黎曼流形。称连续曲线 \(c:[0,b)\to M\)、\(0<b\le\infty\)，在 \(t\uparrow b\) 时离开所有紧集，是指对每个紧集 \(K\subset M\)，都存在 \(t_K<b\)，使 \(c(t)\notin K\) 对 \(t_K<t<b\) 成立。假设曲线在每个紧参数子区间上分片光滑，并定义
	\[
		L_g(c)=\sup_{0<T<b}L_g(c|_{[0,T]}).
	\]
	证明：\(M\) 完备，当且仅当每条离开所有紧集的这种曲线都有无限长度。反向证明中，可使用有限端点的最大测地线，但须证明其位置确实离开所有紧集。
\end{exercise}
\begin{solution}
	先设 \(M\) 完备。若某条曲线的总长度为有限数 \(L\)，则对所有 \(t<b\)，
	\[
		d_g(c(0),c(t))\le L.
	\]
	由 Hopf--Rinow 定理，闭距离球 \(\overline B_{d_g}(c(0),L)\) 紧，而整条曲线都在其中。因此它不能离开所有紧集，证明一个方向。

	反过来，若 \(M\) 不完备，Hopf--Rinow 定理给出一条不能全时延拓的最大测地线。必要时反向并作仿射重参数化，可取单位速测地线 \(\gamma:[0,b)\to M\)，其中 \(b<\infty\) 是其最大右端点。它的长度为 \(b\)。下面证明它离开所有紧集。

	若存在紧集 \(K\) 和时刻 \(t_j\uparrow b\)，使 \(\gamma(t_j)\in K\)，取子列可使 \(\gamma(t_j)\to q\in K\)。因速度为一，当 \(t<t_j\) 时有
	\[
		d_g(\gamma(t),\gamma(t_j))\le t_j-t.
	\]
	让 \(j\to\infty\)，得到 \(d_g(\gamma(t),q)\le b-t\)，所以整个曲线尾部趋于 \(q\)。这里从反复进入紧集到存在位置极限，使用了有限剩余长度；仅知道测地线的某些位置有收敛子列尚不足以直接给出二阶方程的初值。

	在 \(q\) 附近取闭包紧含于一个坐标图的邻域，曲线尾部最终留在其中。度量矩阵有统一正下界，因此单位速条件使坐标速度 \(\dot x(t)\) 有界；Christoffel 系数也有界，由
	\[
		\ddot x^k=-\Gamma^k_{ij}(x)\dot x^i\dot x^j
	\]
	可知坐标加速度有界。于是 \(\dot x(t)\) 在 \(t\uparrow b\) 时有极限 \(w\)，并且 \((x(t),\dot x(t))\to(x(q),w)\)。测地线一阶系统的积分形式延续到 \(b\)；以 \((q,w)\) 为初值应用\cref{g3:thm:geodesic-ivp}，便将原解延长到 \(b\) 之后，违背最大性。因此每个紧集都只能在某个末时刻以前被访问，\(\gamma\) 确实离开所有紧集。它具有有限长度，与题设性质矛盾，故该性质推出完备。
\end{solution}

\begin{exercise}\label{g3:ex:sphere-normal-volume}
	在半径 \(a>0\) 的圆球面 \(S_a^n\subset\mathbb R^{n+1}\)、\(n\ge1\) 上，固定 \(p\) 并使用由 \(\mathrm{T}_pS_a^n\) 的正交单位基产生的法坐标。设 \(v=ru\)、\(0<r<\pi a\)、\(|u|=1\)。直接对式\eqref{g3:eq:round-sphere-exponential}求微分，证明 \(\exp_p\) 在径向上的长度倍率为 \(1\)，在与 \(v\) 正交的方向上的长度倍率为
	\[
		s(r)=\frac{\sin(r/a)}{r/a}.
	\]
	由此求法坐标中的度量矩阵及体积密度，并说明其在中心的延拓。
\end{exercise}
\begin{solution}
	对任意 \(w\in\mathrm{T}_pS_a^n\)，写
	\[
		w=\lambda u+w_\perp,\qquad
		\lambda=\langle u,w\rangle,\quad w_\perp\perp u.
	\]
	因为 \(\mathrm{d}r_v(w)=\lambda\)、\(\mathrm{d}u_v(w)=w_\perp/r\)，对
	\[
		\exp_p(v)=\cos(r/a)p+a\sin(r/a)u
	\]
	求微分得到
	\[
		(\mathrm{d}\exp_p)_vw
		=\lambda E_r+s(r)w_\perp,\qquad
		E_r=-\frac{\sin(r/a)}a\,p+\cos(r/a)u.
	\]
	向量 \(E_r\) 的环境长度为一，并且与 \(w_\perp\) 正交；两者都切于终点球面。因此
	\[
		\left|(\mathrm{d}\exp_p)_vw\right|_g^2
		=\lambda^2+s(r)^2|w_\perp|^2,
	\]
	给出两种所求倍率。这个计算还直接验证了 Gauss 引理中的径向正交性。

	在固定正交单位基中写 \(v=(x^1,\ldots,x^n)\)，则 \(r^2=\sum_i(x^i)^2\)，拉回度量矩阵为
	\[
		g_{ij}(x)=s(r)^2\delta_{ij}
		+\bigl(1-s(r)^2\bigr)\frac{x^ix^j}{r^2}.
	\]
	径向特征值为 \(1\)，其余 \(n-1\) 个特征值为 \(s(r)^2\)。在 \(0<r<\pi a\) 内，\(s(r)>0\)，所以
	\[
		\sqrt{\det(g_{ij}(x))}
		=s(r)^{n-1},\qquad
		\mu_g=
		\left(\frac{\sin(r/a)}{r/a}\right)^{n-1}
		\left|\mathrm{d}x^1\cdots\mathrm{d}x^n\right|.
	\]
	当 \(n=1\) 时没有横向方向，密度系数为 \(1\)，与此公式一致。

	最后，\(\sin z/z\) 在 \(z=0\) 的值取 \(1\) 后，是 \(z^2\) 的光滑函数。因而 \(s(r)\) 是 \(r^2\) 的光滑函数，且
	\[
		s(r)^2=1-\frac{r^2}{3a^2}+O(r^4).
	\]
	于是 \((1-s(r)^2)/r^2\) 在零处光滑延拓，取值为 \(1/(3a^2)\)；度量矩阵在中心延拓为 \(\delta_{ij}\)，体积密度系数光滑延拓为 \(1\)。以上计算直接使用显式指数映射，给出了正曲率球面相对于切空间 Euclid 体积的局部变化。
\end{solution}

\begin{exercise}\label{g3:ex:rectangular-torus-distance}
	设 \(T=\mathbb R^2/(L\mathbb Z\times H\mathbb Z)\)，其中 \(L,H>0\)，取平坦商度量。计算其直径、注入半径，以及从原点类到任意一点的最短测地线条数。再考虑平坦圆柱 \(C=\mathbb R\times(\mathbb R/L\mathbb Z)\)，求原点类的割迹和单位方向 \(u=(u_1,u_2)\) 的割时。
\end{exercise}
\begin{solution}
	每个点类都有代表 \((x,y)\in[-L/2,L/2]\times[-H/2,H/2]\)。由式\eqref{g3:eq:flat-torus-distance}，到原点类的距离为 \(\sqrt{x^2+y^2}\)，故最大值在角点取得，为 \(\tfrac12\sqrt{L^2+H^2}\)。平移是等距映射，所以任意点对的距离也不超过此值，从而
	\[
		\operatorname{diam}(T)=\tfrac12\sqrt{L^2+H^2},\qquad
		\operatorname{inj}(T)=\tfrac12\min\{L,H\}.
	\]
	最近代表在每个坐标方向独立选择：若 \(|x|<L/2\)，第一个坐标唯一；若 \(|x|=L/2\)，两个符号都可选。第二个坐标同理。因此最近点区域内部的点有一条最短线，边的内部有两条，四角共同代表的点有四条。从基点到自身，最短线只有常值曲线。每个最近代表给出一个直线段，其投影的初速度不同，故这些线确实互不相同；任意最短线提升后又必须是到某个最近代表的直线段，所以没有遗漏。

	圆柱只在第二坐标上取商，距离公式为
	\[
		d_C([0,0],[x,y])=\min_{k\in\mathbb Z}\sqrt{x^2+(y+kL)^2}.
	\]
	最近点区域是无限条带 \(|y|\le L/2\)，两条边界在商中识别，故原点类的割迹为直线 \(\{[x,L/2]:x\in\mathbb R\}\)。沿方向 \(u\) 的线最短恰好到第一次抵达该边界，因此
	\[
		c_p(u)=
		\begin{cases}
			L/(2|u_2|),&u_2\ne0,\\
			\infty,&u_2=0.
		\end{cases}
	\]
	相应注入半径为 \(L/2\)。同一个完备流形上，可以有永远最短的方向，也可以有有限割时的方向。
\end{solution}

\begin{exercise}\label{g3:ex:complete-cusp-zero-injectivity}
	设 \((M,g)\) 连通、完备、无边界。
	\begin{enumerate}
		\item 若基于 \(p\) 的分片光滑环路 \(\alpha\) 不能在 \(M\) 中收缩为常值环路，证明 \(\operatorname{inj}_g(p)\le L_g(\alpha)/2\)。
		\item 在圆柱 \(M=\mathbb R\times(\mathbb R/2\pi\mathbb Z)\) 上取
		\[
			g=\mathrm dr^2+e^{-2r}\mathrm d\theta^2.
		\]
		证明它完备，但 \(\operatorname{inj}_g(M)=0\)。
	\end{enumerate}
\end{exercise}
\begin{solution}
	1. 记 \(\ell=L_g(\alpha)\)。沿环路的正向和反向两段都能连接 \(p\) 与任意一点 \(q\)，其中较短的一段长至多 \(\ell/2\)，所以整条环路位于 \(\overline B_d(p,\ell/2)\)。若 \(\ell/2<\operatorname{inj}_g(p)\)，取
	\[
		\ell/2<\rho<\operatorname{inj}_g(p).
	\]
	由\cref{g3:prop:compact-injectivity-lower-bound}，\(B_d(p,\rho)=\exp_p B_p(0,\rho)\)，且指数映射是微分同胚。于是
	\[
		H(s,t)=\exp_p\bigl((1-s)\exp_p^{-1}\alpha(t)\bigr)
	\]
	在固定基点的条件下把 \(\alpha\) 收缩为常值环路，矛盾。

	2. 对任意连接曲线，都有 \(L_g\ge\int|\dot r|\)，因此 \(|r-r'|\le d_g((r,\theta),(r',\theta'))\)。若 \((r_j,\theta_j)\) 为 Cauchy 序列，\(r_j\) 为实数 Cauchy 序列，特别地全部落在某个有界闭区间 \([A,B]\) 中。拓扑圆柱带 \([A,B]\times S^1\) 紧，故该序列有子列在流形拓扑中收敛；距离拓扑相同，所以此子列也在 \(d_g\) 下收敛。Cauchy 性再使整个序列收敛，证明完备。

	基于 \(p_r=(r,0)\) 的圆周环路 \(\alpha_r(t)=(r,2\pi t)\) 长为 \(2\pi e^{-r}\)，且不可收缩。事实上，若它可收缩，投影到第二因子就会把圆周绕行一周的环路收缩；在覆盖 \(\mathbb R\to S^1\)、\(s\mapsto e^{is}\) 下，该环路从零出发的提升终止于 \(2\pi\)，而常值环路的提升终止于零，与第一卷第5.3节的定端点同伦提升结论矛盾。由第一问，
	\[
		0<\operatorname{inj}_g(p_r)\le\pi e^{-r}.
	\]
	令 \(r\to+\infty\)，得到整体下确界为零。各点的注入半径仍为正，趋于零发生在离开所有紧集的过程中。
\end{solution}

\begin{exercise}\label{g3:ex:covering-completeness-converse}
	设 \(\pi:(\widetilde M,\widetilde g)\to(M,g)\) 是连通无边界黎曼流形之间的覆盖，且处处局部等距。证明 \(M\) 完备当且仅当 \(\widetilde M\) 完备。给出一个局部等距映射，说明删去覆盖条件后，从目标完备不能推出源完备。
\end{exercise}
\begin{solution}
	从源完备推出目标完备已由\cref{g3:prop:riemannian-covering-distance}证明。反过来设 \(M\) 完备，给定 \(\widetilde M\) 上的初值 \((\widetilde p,\widetilde v)\)。将其投影到 \(M\)，所得测地线 \(\gamma\) 定义在整个实轴上。对每个正整数 \(N\)，在 \([-N,N]\) 上从 \(t=0\) 分别向正负方向提升 \(\gamma\)，初点指定为 \(\widetilde p\)。道路提升唯一性保证不同 \(N\) 的提升在公共区间相等，因而拼成 \(\mathbb R\) 上的一条提升 \(\widetilde\gamma\)。局部覆盖逆为等距映射，所以提升光滑，并保持测地线方程；其初速度就是 \(\widetilde v\)。由此全部测地线初值全时可解，Hopf--Rinow 定理给出 \(\widetilde M\) 完备。

	带普通度量的包含映射 \((0,1)\hookrightarrow\mathbb R\) 局部等距，但目标完备而源不完备。有限时间内，源中的直线可以到达被删去的端点；该映射不是到 \(\mathbb R\) 的覆盖，不能对全时的目标直线作所需提升。
\end{solution}

\begin{exercise}\label{g3:ex:noncompact-complete-ray}
	设 \((M,g)\) 连通、完备、无边界且非紧。证明对任意 \(p\in M\)，存在单位速测地线 \(\gamma:[0,\infty)\to M\)，满足
	\[
		d(\gamma(s),\gamma(t))=t-s\qquad(0\le s\le t).
	\]
	这样的曲线称为从 \(p\) 出发的测地射线。由此证明至少存在一个 \(u\in S_pM\)，使 \(c_p(u)=\infty\)。
\end{exercise}
\begin{solution}
	由 Hopf--Rinow 定理，有限半径闭球紧。若 \(d(p,\cdot)\) 有上界，整个流形就是一个紧闭球，与非紧性矛盾。因此可选 \(q_j\)，使 \(\ell_j=d(p,q_j)\to\infty\)。取从 \(p\) 到 \(q_j\) 的单位速最短测地线，初方向记为 \(u_j\in S_pM\)。有限维单位球面紧，故取子列可使 \(u_j\to u\)。完备性保证 \(\gamma(t)=\exp_p(tu)\) 对所有 \(t\ge0\) 有定义。

	固定 \(0\le s\le t\)。充分大的 \(j\) 有 \(\ell_j\ge t\)，最短线的子段仍最短，故
	\[
		d(\exp_p(su_j),\exp_p(tu_j))=t-s.
	\]
	指数映射和距离均连续，令 \(j\to\infty\) 得到所求等式。取 \(s=0\)，可见从 \(p\) 到 \(\gamma(t)\) 的线段对每个 \(t\) 都最短，所以 \(c_p(u)=\infty\)。
\end{solution}

\begin{exercise}\label{g3:ex:isometry-first-jet}
	设 \(F:(M,g)\to(N,h)\) 为黎曼等距同构，即满足 \(F^*h=g\) 的微分同胚。证明在指数映射有定义处，
	\[
		F(\exp_pv)=\exp_{F(p)}((\mathrm dF)_pv).
	\]
	若 \(M\) 还连通、完备，证明 \(F\) 由一点处的值 \(F(p)\) 及微分 \((\mathrm dF)_p\) 唯一确定。并证明
	\[
		F(\operatorname{Cut}_M(p))=\operatorname{Cut}_N(F(p)),\qquad
		\operatorname{inj}_h(F(p))=\operatorname{inj}_g(p).
	\]
\end{exercise}
\begin{solution}
	将 \(M\) 上的联络通过 \(F\) 搬到 \(N\)，度量相容和无挠性质保持成立，所以它由 Levi-Civita 唯一性就是 \(N\) 上的联络。因此 \(F\circ\gamma_{p,v}\) 是初值 \((F(p),(\mathrm dF)_pv)\) 的测地线。初值唯一性给出全时间的对应，取时间一即得指数映射公式。

	若另一个等距同构 \(G:M\to N\) 在 \(p\) 处有同样的值与微分，则对每个 \(v\in\mathrm T_pM\)，有 \(F(\exp_pv)=G(\exp_pv)\)。Hopf--Rinow 定理使 \(\exp_p\) 满射，所以 \(F=G\) 于整个 \(M\)。

	等距同构保持全部曲线长度与端点距离，其微分把单位切球面等距地对应起来。因此
	\[
		c_{F(p)}((\mathrm dF)_pu)=c_p(u),
	\]
	并把有限割时的终点一一对应，得到割迹等式。指数映射公式还说明，两侧指数映射分别通过线性等距同构 \((\mathrm dF)_p\) 和微分同胚 \(F\) 共轭，所以使指数映射在切球上成为微分同胚的半径完全相同，注入半径相等。
\end{solution}

\begin{exercise}\label{g3:ex:scaling-cut-injectivity}
	设 \((M,g)\) 连通、完备、无边界，\(a>0\)，令 \(\widetilde g=a^2g\)。比较两种度量的距离、指数映射、割迹、割时及注入半径。写出割时公式时，明确两侧单位初方向的对应。
\end{exercise}
\begin{solution}
	每条曲线的长度乘以 \(a\)，取下确界得到 \(d_{\widetilde g}=a\,d_g\)，所以完备性保持不变。上一章的常数缩放计算说明两种度量的 Levi-Civita 联络相同，故同一个切向量给出同一个带参数测地线，指数映射也相同。

	若 \(|u|_g=1\)，则 \(|u/a|_{\widetilde g}=1\)，且
	\[
		d_{\widetilde g}\bigl(p,\exp_p(tu/a)\bigr)
		=a\,d_g\bigl(p,\exp_p((t/a)u)\bigr).
	\]
	左边等于 \(t\)，当且仅当右侧原距离等于 \(t/a\)。因此
	\[
		c_p^{\widetilde g}(u/a)=a\,c_p^g(u).
	\]
	达到有限割时的向量相同：\((a c_p^g(u))(u/a)=c_p^g(u)u\)，所以割迹作为点集不变。最后，\(\widetilde g\) 的半径 \(r\) 切球就是 \(g\) 的半径 \(r/a\) 切球，指数映射本身不变，故
	\[
		\operatorname{inj}_{\widetilde g}(p)=a\,\operatorname{inj}_g(p),\qquad
		\operatorname{inj}_{\widetilde g}(M)=a\,\operatorname{inj}_g(M).
	\]
\end{solution}

% !TeX root = ../../Book3.tex
% 纲要标题：## 第3章　曲率张量与局部几何的展开
% 纲要位置：几何学计划纲要.md，第 2192 行。
% 正文状态：初稿；尚未正式审核。

% 纲要内容开始
% 
% > 本章目标与证明范围：从协变微分的非交换性得到代数曲率张量，再把各次收缩和法坐标展开联系到可观测的长度、面积和体积变化。
% 
% 纲要内容结束

\chapter{曲率张量与局部几何的展开}
\label{g3:ch:03}
\BookChapterNumber{g3:ch:03}

\begin{chapterabstract}
	法坐标消去度量在中心的一阶变化，二阶变化则由曲率控制。本章从协变微分的非交换性证明曲率的张量性、代数对称性及两条 Bianchi 恒等式，建立曲率算子、截面曲率、Ricci 曲率与标量曲率之间的关系。随后推导联络、度量和体积密度的法坐标展开，把不同收缩解释为长度、面积和小球体积的低阶变化。空间形式、乘积、翘曲积与共形度量提供计算模型，Weyl 分解则分离标量、无迹 Ricci 和无迹四阶曲率三部分。
\end{chapterabstract}

\begin{learninggoals}
	\begin{itemize}
		\item 在向量曲率、四阶曲率张量和二向量曲率算子之间转换，并核对指标次序与球面的正曲率符号。
		\item 证明两条 Bianchi 恒等式及其收缩，说明截面、Ricci、标量和曲率算子正性之间的含义与差别。
		\item 从法坐标径向条件求出联络和度量的低阶项，推导体积密度、小球体积及面积的展开，并控制余项。
		\item 计算乘积、翘曲积及共形变化的曲率，使用 Weyl 分解解释低维确定性和高维信息损失。
	\end{itemize}
\end{learninggoals}

在平面上画一个半径很小的圆，它的周长是 \(2\pi r\)。在半径 \(a\) 的球面上，以沿球面的距离 \(r\) 为半径画圆，周长却是 \(2\pi a\sin(r/a)\)。当 \(r\) 很小时，两者的一阶项相同，差异从三次项开始：
\[
	2\pi a\sin(r/a)=2\pi r-\frac{\pi}{3a^2}r^3+O(r^5).
\]
这个差异可以由长度测量得到，不需要先把球面与环境空间中的某个法向比较。高维中，不同二维方向可以出现不同的差异；将它们沿某个方向相加，会得到 Ricci 曲率，再把方向全部平均，则出现标量曲率。

\ProseParagraphBreak

第二章的法坐标使 \(g_{ij}(p)=\delta_{ij}\)、\(\partial_kg_{ij}(p)=0\)。因此我们可以把度量与 Euclid 度量的首个差异集中在二阶项中。证明将沿两个方向进行：先由联络求导的交换差得到曲率，找出它必须满足的对称性；再由这些曲率分量恢复法坐标中的二次项，计算可观测的几何量。曲率约束之间的区别也将在实际局部度量中检验。

\ProseParagraphBreak

本章使用第一章的 Levi-Civita 联络、体积密度及非负 Laplace 算子 \(\Delta_g=-\operatorname{div}_g\operatorname{grad}_g\)，以及第二章的指数映射、Gauss 引理和法坐标。外积空间只需有限维线性代数：\(X\wedge Y\) 表示由两个向量张成的有向面积元素。流形及度量均光滑，通常假设无边界；紧性、连通性和维数限制在实际需要的结果中另行写明。法坐标展开由坐标恒等式和 Taylor 公式推出。

\ProseParagraphBreak

作为阅读对照，曲率、曲率算子及各次收缩可参看 Petersen 第3版第3.1.1--3.1.6节；球面、乘积及旋转型度量见第4.2.1--4.2.3节，共形表示见第4.3.3节，指数坐标性质见第5.5.3节\cite[第3.1.1--3.1.6节; 第4.2.1--4.2.3节; 第4.3.3节; 第5.5.3节]{Petersen2016}。本章固定的四阶张量次序是 \(\mathrm{Rm}(X,Y,Z,W)=g(R(X,Y)W,Z)\)，因此 \(\mathrm{Rm}(u,v,u,v)\) 是正交单位向量所张截面的曲率。

% !TeX root = ../../Book3.tex
% 纲要标题：### 3.1 曲率的代数结构
% 纲要位置：几何学计划纲要.md，第 2196 行。
% 正文状态：未撰写；本文件目前是框架。

% 纲要内容开始
% 
% * 曲率张量性及符号
% * 四阶张量的对称性
% * 第一、第二 Bianchi 恒等式
% * 曲率算子及二形式上的作用
% 
% 纲要内容结束

\section{曲率的代数结构}
\label{g3:sec:0301}

待撰写

\subsection{曲率张量性及符号}
\label{g3:subsec:0301:01}
待撰写

\subsection{四阶张量的对称性}
\label{g3:subsec:0301:02}
待撰写

\subsection{第一、第二 Bianchi 恒等式}
\label{g3:subsec:0301:03}
待撰写

\subsection{曲率算子及二形式上的作用}
\label{g3:subsec:0301:04}
待撰写
% 纲要条目排为编号小节标题，并进入目录与书签；说明段仍保留于纲要注释中。

% !TeX root = ../../Book3.tex
% 纲要标题：### 3.2 不同曲率约束
% 纲要位置：几何学计划纲要.md，第 2203 行。
% 正文状态：初稿；尚未正式审核。

% 纲要内容开始
% 
% * 截面曲率和二维截面
% * Ricci 曲率与体积效应
% * 标量曲率与收缩
% * 高维不同正性条件的严格区别
% 
% 纲要内容结束

\section{不同曲率约束}
\label{g3:sec:0302}

曲率张量同时接受四个方向，实际几何问题却常常只指定一个二维平面、一个径向方向，或者一小块区域的总体积。截面曲率、Ricci 曲率与标量曲率分别对应这些不同的观察方式。\footnote{Gregorio Ricci-Curbastro（1853--1925），意大利数学家，创立绝对微分学，并与 Levi-Civita 合作发展其几何应用。}它们都来自同一个曲率张量，但取值所保留的信息和“为正”的含义并不相同。本节先建立这些量之间的精确关系，再用可以直接核算的例子辨明各种正性条件。

\subsection{截面曲率和二维截面}
\label{g3:subsec:0302:01}

回顾\cref{g3:def:sectional-curvature}。对 \(\sigma\subset\mathrm{T}_pM\) 的任意一组基 \(X,Y\)，本章的四阶张量约定将截面曲率写为
\begin{equation}
	\label{g3:eq:sectional-bivector}
	K_p(\sigma)
	=\frac{\mathrm{Rm}(X,Y,X,Y)}{|X\wedge Y|_g^2}
	=\frac{\langle\mathcal R(X\wedge Y),X\wedge Y\rangle}{|X\wedge Y|_g^2},
\end{equation}
其中 \(|X\wedge Y|_g^2=g(X,X)g(Y,Y)-g(X,Y)^2\) 是平行四边形面积的平方。因此当 \(u,v\) 为 \(\sigma\) 的正交单位基时，\(K_p(\sigma)=\mathrm{Rm}(u,v,u,v)\)。

\ProseParagraphBreak

这个商与平面中的基无关。事实上，若 \(X'=aX+bY\)、\(Y'=cX+dY\)，则 \(X'\wedge Y'=(ad-bc)X\wedge Y\)。\eqref{g3:eq:sectional-bivector}的分子和分母同时乘以 \((ad-bc)^2\)。特别地，反转平面定向也不改变曲率。这个检验解释了为何必须除以面积平方：仅有四阶张量的分子，会随着所选基的长度和夹角变化。

\begin{definition}
	\label{g3:def:sectional-curvature-algebraic}
	设 \((V,g)\) 为维数 \(n\geq2\) 的实内积空间，\(A\) 为 \(V\) 上的代数曲率张量。对二维子空间 \(\sigma=\operatorname{span}\{X,Y\}\)，定义它的截面曲率为
	\[
		K_A(\sigma)=\frac{A(X,Y,X,Y)}{|X\wedge Y|_g^2}.
	\]
	同样的换基计算保证它只依赖 \(\sigma\)。
\end{definition}

截面曲率看似只读取了四阶张量中两个输入重复的值，却能恢复全部分量。下面的恢复过程除了利用双线性，还将用到第一 Bianchi 恒等式。

\begin{proposition}
	\label{g3:prop:sectional-determines-curvature}
	设 \((V,g)\) 为有限维实内积空间，\(\dim V\geq2\)，\(A,B\) 是其上的代数曲率张量。若它们在每个二维子空间上的截面曲率相同，则 \(A=B\)。
\end{proposition}
\begin{proof}
	令 \(D=A-B\)。截面曲率相等说明对线性无关的 \(x,y\)，有 \(D(x,y,x,y)=0\)；线性相关时，同一式子由反对称性成立。因此这个等式对所有 \(x,y\) 都成立。下面作两次极化，再用第一 Bianchi 恒等式消去剩余分量。

	将第一个重复向量替换为 \(x+z\)，利用交换两对输入的对称性，得到
	\[
		0=D(x+z,y,x+z,y)=2D(x,y,z,y).
	\]
	再在刚得到的等式中将 \(y\) 替换为 \(y+w\)，便有
	\[
		D(x,y,z,w)+D(x,w,z,y)=0.
	\]
	由最后两个输入的反对称性，这等价于
	\begin{equation}
		\label{g3:eq:sectional-zero-polarization}
		D(x,y,z,w)=D(x,w,y,z).
	\end{equation}
	于是交换两对输入再用该式，得到
	\[
		D(y,z,x,w)=D(x,w,y,z)=D(x,y,z,w).
	\]
	同样，
	\[
		\begin{aligned}
			D(z,x,y,w)&=D(y,w,z,x)=D(y,x,w,z)\\
			&=D(x,y,z,w),
		\end{aligned}
	\]
	最后一步同时交换第一对和第二对输入。第一 Bianchi 恒等式的三个循环项因而相等：
	\[
		0=D(x,y,z,w)+D(y,z,x,w)+D(z,x,y,w)
		=3D(x,y,z,w).
	\]
	故 \(D=0\)。
\end{proof}

\begin{corollary}
	\label{g3:cor:pointwise-isotropic-curvature}
	设 \((V,g)\) 为维数至少为二的实内积空间。代数曲率张量 \(A\) 在所有二维子空间上的截面曲率均为同一个实数 \(k\)，当且仅当
	\begin{equation}
		\label{g3:eq:constant-sectional-algebraic}
		A(X,Y,Z,W)
		=k\bigl(g(X,Z)g(Y,W)-g(X,W)g(Y,Z)\bigr).
	\end{equation}
\end{corollary}
\begin{proof}
	右端具有两对输入的反对称性与交换两对输入的对称性，第一 Bianchi 循环和中的六个乘积逐对抵消，故它是代数曲率张量。令 \(Z=X,W=Y\)，分子恰为 \(k|X\wedge Y|_g^2\)，所以它的截面曲率恒为 \(k\)。反过来，若 \(A\) 也有这些截面值，\cref{g3:prop:sectional-determines-curvature}使两个张量相同。
\end{proof}

这里的 \(k\) 只描述一个切空间中的各个平面。若在流形上逐点应用这一推论，首先得到的仍是可能随点变化的 \(k(p)\)；它何时必须为常数，将由本节稍后的微分恒等式决定。

\subsection{Ricci 曲率与体积效应}
\label{g3:subsec:0302:02}

固定一个方向 \(v\)，可以分别观察包含 \(v\) 的各个二维平面。若把正交于 \(v\) 的所有方向同时考虑，对相应曲率求和便比单个截面更适合描述横向体积的变化。这个求和应由张量的迹来定义，使结果不依赖所选的正交基。

\begin{definition}
	\label{g3:def:ricci-scalar}
	设 \((M^n,g)\) 为光滑黎曼流形，\(\mathrm{Rm}\) 采用\cref{g3:def:curvature-four-tensor}的约定。在任意正交单位基 \(e_1,\ldots,e_n\) 下，定义\term[ricciqulv]{Ricci 曲率}[Ricci curvature]为二阶协变张量
	\begin{equation}
		\label{g3:eq:ricci-definition}
		\operatorname{Ric}(Y,W)
		=\sum_{i=1}^n\mathrm{Rm}(e_i,Y,e_i,W)
		=\operatorname{tr}\bigl(X\mapsto R(X,Y)W\bigr).
	\end{equation}
	它的迹称为\term[biaoliangqulv]{标量曲率}[scalar curvature]，记为
	\[
		\operatorname{Scal}=\operatorname{tr}_g\operatorname{Ric}
		=\sum_{j=1}^n\operatorname{Ric}(e_j,e_j).
	\]
\end{definition}

定义中的 \(X\mapsto R(X,Y)W\) 是同一个切空间上的线性自映射，其迹不依赖基。一般坐标中的公式为
\begin{equation}
	\label{g3:eq:ricci-scalar-contractions}
	\operatorname{Ric}_{jl}=g^{ik}\mathrm{Rm}_{ijkl},\qquad
	\operatorname{Scal}=g^{jl}\operatorname{Ric}_{jl}.
\end{equation}
这些公式还说明两者随基点光滑变化。\(\operatorname{Ric}\) 对两个向量线性，\(\operatorname{Scal}\) 则是每一点处的实数。

\symidx{\(\operatorname{Ric}\)}
\symidx{\(\operatorname{Scal}\)}

\begin{proposition}
	\label{g3:prop:ricci-sectional-trace}
	在 \(n\) 维光滑黎曼流形上，\(\operatorname{Ric}\) 是对称二阶张量。若 \(n\geq2\)、\(|v|_g=1\)，并将 \(v=e_1\) 补成正交单位基，则
	\begin{equation}
		\label{g3:eq:ricci-sectional-sum}
		\operatorname{Ric}(v,v)=\sum_{a=2}^nK(\operatorname{span}\{v,e_a\}).
	\end{equation}
	特别地，若一点处所有截面曲率位于 \([a,b]\)，则在该点有二次型不等式
	\[
		(n-1)a\,g\leq\operatorname{Ric}\leq(n-1)b\,g.
	\]
\end{proposition}
\begin{proof}
	由交换两对输入的对称性，
	\[
		\mathrm{Rm}(e_i,Y,e_i,W)=\mathrm{Rm}(e_i,W,e_i,Y),
	\]
	求和即得 Ricci 对称性。将 \(Y=W=v=e_1\) 代入定义，\(i=1\) 的项因反对称性为零，其余每项是一个正交单位二维平面的截面曲率，得到\eqref{g3:eq:ricci-sectional-sum}。最后对单位向量应用逐项上下界，再利用二次齐次性恢复一般向量，即得二次型估计。
\end{proof}

因此 \(\operatorname{Ric}(v,v)/(n-1)\) 是沿单位方向 \(v\) 所见的这些正交截面的平均曲率。求和中的个别截面值可以变化，和却不受补全正交基的方式影响。在维数二，只有一个横向方向，所以 \(\operatorname{Ric}=Kg\)；维数增加后，正负截面可能在同一个方向的求和中抵消。

\ProseParagraphBreak

这个迹如何进入体积，可以先由下一节将证明的展开辨认。固定 \(p\)，在 \(\mathrm{T}_pM\) 中取正交单位基，并用指数映射得到法坐标。对单位向量 \(v\)，\cref{g3:prop:normal-volume-density}将证明
\begin{equation}
	\label{g3:eq:ricci-volume-direction}
	\sqrt{\det(g_{ij}(tv))}
	=1-\frac16\operatorname{Ric}_p(v,v)t^2+O(t^3).
\end{equation}
左侧是指数坐标下的体积密度与切空间 Euclid 体积密度之比。横向各方向的一阶偏差为零，二阶偏差在行列式中相加，因而出现截面曲率的迹。对某个 \(v\) 有 \(\operatorname{Ric}(v,v)>0\) 时，这个方向上足够近的体积比小于一；负值则使二阶修正朝相反方向变化。公式的推导放在法坐标二阶展开之后，当前 Ricci 张量的定义和对称性只用了代数收缩。

\subsection{标量曲率与收缩}
\label{g3:subsec:0302:03}

再对 Ricci 张量取迹，就把对某个方向的观察合成一个数。对任意正交单位基，\eqref{g3:eq:ricci-sectional-sum}给出
\begin{equation}
	\label{g3:eq:scalar-sectional-average}
	\operatorname{Scal}
	=2\sum_{i<j}K(\operatorname{span}\{e_i,e_j\})
	=2\operatorname{tr}\mathcal R.
\end{equation}
每个坐标二维平面在 Ricci 的总和中出现两次，这解释了系数 \(2\)。因此 \(\operatorname{Scal}/(n(n-1))\) 是这 \(n(n-1)/2\) 个截面曲率的平均；其数值与所选正交基无关。常截面曲率 \(k\) 的代数模型满足
\[
	\operatorname{Ric}=(n-1)kg,\qquad
	\operatorname{Scal}=n(n-1)k.
\]

\ProseParagraphBreak

也可以对所有单位径向方向作平均。记 \(\mathrm{d}\sigma\) 为 \(\mathrm{T}_pM\) 的单位球面上的面积测度。由坐标反射和坐标置换的对称性，\(\int v^iv^j\,\mathrm{d}\sigma=0\) 当 \(i\ne j\)，而各 \(\int(v^i)^2\,\mathrm{d}\sigma\) 相同；再用 \(\sum_i(v^i)^2=1\)，得到
\[
	\frac1{\operatorname{Area}(S^{n-1})}
	\int_{S^{n-1}}\operatorname{Ric}_p(v,v)\,\mathrm{d}\sigma(v)
	=\frac1n\operatorname{Scal}(p).
\]
这正是所有径向的体积偏差合成小球体积偏差时出现标量曲率的原因。\cref{g3:thm:small-ball-volume}将在下一节证明：对充分小的 \(r>0\)，
\begin{equation}
	\label{g3:eq:scalar-small-ball-preview}
	\operatorname{Vol}_g(B_d(p,r))
	=\omega_n r^n\left(1-\frac{\operatorname{Scal}(p)}{6(n+2)}r^2+O(r^4)\right),
\end{equation}
其中 \(\omega_n\) 是单位 Euclid \(n\) 球的体积。定向不参与这一公式，体积始终按密度积分。

\ProseParagraphBreak

曲率的这些收缩还满足微分约束。第二 Bianchi 恒等式将曲率在不同方向上的导数联系起来；取迹后，Ricci 的散度恰好由标量曲率的微分控制。

\begin{proposition}[缩并 Bianchi 恒等式]
	\label{g3:prop:contracted-bianchi}
	在光滑黎曼流形 \((M,g)\) 上，定义对称二阶张量 \(S\) 的散度为一形式
	\[
		(\operatorname{div}S)(Y)=\sum_i(\nabla_{e_i}S)(e_i,Y),
	\]
	其中 \(e_i\) 是任意正交单位基。则
	\begin{equation}
		\label{g3:eq:contracted-bianchi}
		\operatorname{div}\operatorname{Ric}=\tfrac12\mathrm{d}\operatorname{Scal}.
	\end{equation}
\end{proposition}
\begin{proof}
	在任意一点取正交法坐标，使该点处 \(g_{ij}=\delta_{ij}\)、\(\nabla_{e_i}e_j=0\)。因 \(\nabla g=0\)，协变微分与用度量作收缩可交换。\eqref{g3:eq:second-bianchi-components}给出
	\[
		\nabla_a\mathrm{Rm}_{ijkl}
		+\nabla_i\mathrm{Rm}_{jakl}
		+\nabla_j\mathrm{Rm}_{aikl}=0.
	\]
	先令 \(k=a\) 并求和。由第一对指标的反对称性及 Ricci 定义，得到
	\[
		\sum_a\nabla_a\mathrm{Rm}_{ijal}
		-\nabla_i\operatorname{Ric}_{jl}
		+\nabla_j\operatorname{Ric}_{il}=0.
	\]
	再令 \(l=i\) 并对 \(i\) 求和。由于 \(\sum_i\mathrm{Rm}_{ijai}=-\operatorname{Ric}_{ja}\)，上式变成
	\[
		-\sum_a\nabla_a\operatorname{Ric}_{ja}
		-\sum_i\nabla_i\operatorname{Ric}_{ji}
		+\nabla_j\operatorname{Scal}=0.
	\]
	Ricci 对称性使前两项相同，所以 \(2(\operatorname{div}\operatorname{Ric})_j=\partial_j\operatorname{Scal}\)。所选点任意，且两边都是一形式，故所得等式在全流形成立。
\end{proof}

例如 \(\operatorname{div}(\operatorname{Ric}-\tfrac12\operatorname{Scal}\,g)=0\)。若某类曲率条件使 Ricci 与度量成比例，这个微分恒等式就会限制比例因子。Schur 定理给出一个基本应用。\footnote{Friedrich Heinrich Schur（1856--1932），德国数学家，研究微分几何、变换群与几何基础；这里的 Schur 定理与群表示论中 Issai Schur 的引理由不同数学家得到。}

\begin{theorem}[Schur]
	\label{g3:thm:schur}
	设 \((M^n,g)\) 为连通光滑黎曼流形，\(n\geq3\)。若在每一点 \(p\)，所有二维子空间的截面曲率具有同一个值 \(k(p)\)，则 \(k\) 在 \(M\) 上为常数。因此 \(M\) 具有常截面曲率。
\end{theorem}
\begin{proof}
	由\cref{g3:cor:pointwise-isotropic-curvature}，在每一点都有\eqref{g3:eq:constant-sectional-algebraic}，从而
	\[
		\operatorname{Ric}=(n-1)kg,\qquad
		\operatorname{Scal}=n(n-1)k.
	\]
	第二式先说明 \(k=\operatorname{Scal}/(n(n-1))\) 是光滑函数。利用 \(\nabla g=0\)，第一式给出 \(\operatorname{div}\operatorname{Ric}=(n-1)\mathrm{d}k\)。代入\eqref{g3:eq:contracted-bianchi}，得
	\[
		(n-1)\mathrm{d}k=\tfrac12n(n-1)\mathrm{d}k,
		\qquad (n-1)(n-2)\mathrm{d}k=0.
	\]
	维数条件使 \(\mathrm{d}k=0\)，所以 \(k\) 局部为常数；连通性使这些局部常数在整个 \(M\) 上相同。
\end{proof}

维数二时，每个切空间本来就只有一个二维子空间，定理的逐点条件不施加额外限制。例如在 \(\mathbb R^2\) 上取
\[
	g=\mathrm{d}r^2+e^{2r^2}\mathrm{d}\theta^2.
\]
由\eqref{g3:eq:warped-surface-curvature}，令 \(f(r)=e^{r^2}\) 后得到
\[
	K(r,\theta)=-\frac{f''(r)}{f(r)}=-2-4r^2,
\]
它不是常数。在这个维数，\(\operatorname{Ric}=Kg\)、\(\operatorname{Scal}=2K\)，缩并 Bianchi 等式两边都等于 \(\mathrm{d}K\)，不再约束 \(K\) 的变化。

\subsection{高维不同正性条件的严格区别}
\label{g3:subsec:0302:04}

曲率算子非负，是说 \(\langle\mathcal R\omega,\omega\rangle\geq0\) 对所有二向量 \(\omega\) 成立；截面曲率非负只在 \(\omega=X\wedge Y\) 这类可分解二向量上作同样的检验。Ricci 非负是二次型 \(\operatorname{Ric}(v,v)\geq0\)，标量曲率非负则只检验一个实数。正曲率算子、正 Ricci 曲率分别指相应二次型正定；正截面曲率要求每个二维截面的值严格为正。

\begin{proposition}
	\label{g3:prop:curvature-positivity-implications}
	在维数 \(n\geq2\) 的黎曼流形的每一点，都有
	\begin{equation}
		\label{g3:eq:curvature-positivity-chain}
		\mathcal R\geq0\ \Longrightarrow\ K\geq0
		\ \Longrightarrow\ \operatorname{Ric}\geq0
		\ \Longrightarrow\ \operatorname{Scal}\geq0.
	\end{equation}
	将其中全部非负号换为严格正号，同样成立。维数二时，四种非负性彼此等价，四种严格正性也彼此等价；维数三时，曲率算子与截面曲率具有相同的非负性和严格正性。
\end{proposition}
\begin{proof}
	第一个箭头来自\eqref{g3:eq:sectional-bivector}，第二个来自\eqref{g3:eq:ricci-sectional-sum}，最后一个来自正交基下取迹。严格正性时，单位方向的 Ricci 和包含 \(n-1\geq1\) 个正数，标量曲率又是 \(n\) 个正数之和。

	维数二时，\(\Lambda^2\mathrm{T}_pM\) 是一维的，曲率算子就是乘以 \(K\)；同时 \(\operatorname{Ric}=Kg\)、\(\operatorname{Scal}=2K\)，故各条件等价。维数三时，每个非零二向量都可分解。具体说，对正交基 \(e_1,e_2,e_3\)，若
	\[
		\omega=a\,e_1\wedge e_2+b\,e_1\wedge e_3+c\,e_2\wedge e_3,
		\qquad a\ne0,
	\]
	则
	\[
		\omega=\left(e_1-\frac ca e_3\right)\wedge(ae_2+be_3).
	\]
	若 \(a=0\)，可重新排列基，使某个非零系数占据该位置。因此在三维，截面曲率的检验已经遍历全部非零二向量，得到所述等价性。
\end{proof}

截面曲率决定整个张量，并不意味着截面上的正性自动成为整个二向量空间上的正性。\cref{g3:prop:sectional-determines-curvature}的恢复过程包含极化求差，求差不会保持正号。下面先在单个内积空间上展示这些区别，再由下一节的局部实现将它们用于度量。

\ProseParagraphBreak

构造代数例子时需要核对第一 Bianchi 恒等式。任取 \(\Lambda^2V\) 上的自伴算子 \(A\)，令 \(T_{ijkl}=\langle A(e_i\wedge e_j),e_k\wedge e_l\rangle\)。它自动具有两对指标的反对称性与交换两对的对称性。第一 Bianchi 恒等式等价于
\begin{equation}
	\label{g3:eq:operator-bianchi-test}
	T_{ijkl}-T_{ikjl}+T_{iljk}=0.
\end{equation}
左侧交换任意相邻两个指标都会变号，这可直接由上述三种对称性核对；例如交换 \(i,j\) 后得到 \(-T_{ijkl}-T_{iljk}+T_{ikjl}\)。所以有重复指标时它自动为零。三维没有四个互异的指标，任意自伴算子因而都通过这个检验；四维则只剩下一个独立条件 \(T_{1234}-T_{1324}+T_{1423}=0\)。

\begin{example}
	\label{g3:exm:ricci-positive-sectional-negative}
	在 \(\mathbb R^3\) 的标准正交基下，令曲率算子在 \(e_1\wedge e_2,e_1\wedge e_3,e_2\wedge e_3\) 上为对角算子，相应特征值记为 \(\kappa_{12},\kappa_{13},\kappa_{23}\)。上一段已经验证它给出代数曲率张量。直接收缩可得
	\[
		\operatorname{Ric}
		=\operatorname{diag}(\kappa_{12}+\kappa_{13},
		\kappa_{12}+\kappa_{23},\kappa_{13}+\kappa_{23}),
	\]
	且 \(\operatorname{Scal}=2(\kappa_{12}+\kappa_{13}+\kappa_{23})\)。取
	\[
		(\kappa_{12},\kappa_{13},\kappa_{23})=(-1,2,2),
	\]
	则 \(\operatorname{Ric}=\operatorname{diag}(1,1,4)\) 正定，但 \(K(e_1,e_2)=-1\)。因此 Ricci 严格为正不能推出截面曲率非负。
\end{example}

\begin{example}
	\label{g3:exm:scalar-positive-ricci-negative}
	在同一三维构造中取
	\[
		(\kappa_{12},\kappa_{13},\kappa_{23})=(-2,3,0).
	\]
	则 \(\operatorname{Ric}=\operatorname{diag}(1,-2,3)\)，而 \(\operatorname{Scal}=2>0\)。标量曲率严格为正仍允许 Ricci 张量有负方向。它记录总和的符号，不能判断每一项的符号。
\end{example}

\begin{example}
	\label{g3:exm:sectional-positive-operator-negative}
	在定向 Euclid 空间 \(\mathbb R^4\) 中，记 \(e_{ij}=e_i\wedge e_j\)，并取下列六个正交单位二向量：
	\[
		\begin{aligned}
			\sigma_1&=(e_{12}+e_{34})/\sqrt2,&
			\tau_1&=(e_{12}-e_{34})/\sqrt2,\\
			\sigma_2&=(e_{13}-e_{24})/\sqrt2,&
			\tau_2&=(e_{13}+e_{24})/\sqrt2,\\
			\sigma_3&=(e_{14}+e_{23})/\sqrt2,&
			\tau_3&=(e_{14}-e_{23})/\sqrt2.
		\end{aligned}
	\]
	定义自伴算子 \(\mathcal R\)，使 \(\sigma_1,\sigma_2,\sigma_3\) 的特征值分别为 \(-1,3,3\)，而三个 \(\tau_i\) 的特征值都为 \(5/3\)。在原二向量基中，相应三个交叉分量为
	\[
		T_{1234}=-\tfrac43,\qquad
		T_{1324}=-\tfrac23,\qquad
		T_{1423}=\tfrac23.
	\]
	它们满足\eqref{g3:eq:operator-bianchi-test}，所以这确实是一个代数曲率张量的算子。

	现在检查所有截面，而不只检查六个坐标截面。记 \(\nu=e_1\wedge e_2\wedge e_3\wedge e_4\)。直接作外积得到
	\[
		\sigma_i\wedge\sigma_j=\delta_{ij}\nu,\qquad
		\tau_i\wedge\tau_j=-\delta_{ij}\nu,\qquad
		\sigma_i\wedge\tau_j=0.
	\]
	若 \(\omega=X\wedge Y\) 是单位可分解二向量，写成 \(\omega=\sum_i a_i\sigma_i+\sum_i b_i\tau_i\)。因为 \(\omega\wedge\omega=0\) 且 \(|\omega|=1\)，有
	\[
		\sum_i a_i^2-\sum_i b_i^2=0,\qquad
		\sum_i a_i^2+\sum_i b_i^2=1.
	\]
	两组平方和因此都为 \(1/2\)，从而
	\[
		\langle\mathcal R\omega,\omega\rangle
		\geq-\sum_i a_i^2+\tfrac53\sum_i b_i^2
		=\tfrac13>0.
	\]
	所有二维截面曲率都严格为正；然而 \(\langle\mathcal R\sigma_1,\sigma_1\rangle=-1\)，所以曲率算子具有负特征值。这里 \(\sigma_1\wedge\sigma_1=\nu\ne0\)，它恰好不是某个二维平面的可分解二向量。
\end{example}

这三个例子首先是单个切空间中的代数数据。\cref{g3:prop:curvature-local-realization}将在下一节证明：任意这样的代数曲率张量都能实现为一个光滑黎曼度量在原点的曲率。上述严格正性与所展示的负方向在充分小的邻域内还会保持。对 Ricci 张量和曲率算子，这来自矩阵分量的连续性；对四维例子的截面曲率，取光滑正交标架后，单位可分解二向量构成固定的紧集，原点处统一的 \(1/3\) 下界可在邻近点保留为某个正下界。因此这些数据给出实际的小开球上的度量反例，说明\eqref{g3:eq:curvature-positivity-chain}的逆向推理在相应维数确实失败；附加紧性或完备性后的整体存在问题需要另外的构造。

\ProseParagraphBreak

相反，若一点处 \(K\geq0\) 且 \(\operatorname{Scal}=0\)，则任一正交基下的非负截面曲率之和为零。将任意指定的二维平面补入一组正交基，就知该平面的曲率也为零；\cref{g3:prop:sectional-determines-curvature}遂给出该点处 \(\mathrm{Rm}=0\)。这个结论同时使用了非负性与零总和。单独的标量曲率为零仍允许不同方向的正负曲率相互抵消。

% 纲要条目排为编号小节标题，并进入目录与书签；说明段仍保留于纲要注释中。

% !TeX root = ../../Book3.tex
% 纲要标题：### 3.3 法坐标中的二阶几何
% 纲要位置：几何学计划纲要.md，第 2210 行。
% 正文状态：初稿；尚未正式审核。

% 纲要内容开始
% 
% * 联络的一阶展开
% * 度量的二阶展开
% * 体积密度和小球体积
% * 展开余项及正则性
% 
% 纲要内容结束

\section{法坐标中的二阶几何}
\label{g3:sec:0303}

在法坐标中心，度量矩阵是单位矩阵，一阶偏导也全部为零。因而仅看一点处的长度和一阶变化，还分辨不出它附近是否弯曲。现在要找的是二阶变化中无法通过这种坐标选择消去的部分。曲率将决定度量的二次项；取矩阵行列式后，这些方向性信息收缩为 Ricci 曲率，再对整个小球积分，就出现标量曲率。

\ProseParagraphBreak

本节的度量始终光滑。固定无边界黎曼流形 \((M^n,g)\) 中一点 \(p\)，在 \(\mathrm T_pM\) 选正交单位基 \((E_1,\ldots,E_n)\)，并用
\[
	q=\exp_p(x^iE_i)
\]
定义足够小邻域内的法坐标。记 \(|x|^2=\sum_i(x^i)^2\)，原点对应 \(p\)。曲率分量沿用本章约定
\[
	R(\partial_i,\partial_j)\partial_b
	=R^a{}_{b i j}\partial_a,\qquad
	R_{ijkl}=g(R(\partial_i,\partial_j)\partial_l,\partial_k).
\]
在原点，因 \(g_{ij}=\delta_{ij}\)，有 \(R^a{}_{b i j}=R_{ijab}\)。下文展开式右侧未注明位置的曲率分量均在 \(p\) 取值；\(O(|x|^m)\) 表示各分量在原点附近的绝对值不超过某个常数乘以 \(|x|^m\)。

\subsection{联络的一阶展开}
\label{g3:subsec:0303:01}

\cref{g3:prop:normal-coordinate-first-order}已经给出 \(\Gamma^a_{ij}(0)=0\)。要确定其一阶项，需要同时使用两类关系：曲率给出联络偏导的反对称部分，径向直线是测地线则给出其完全对称部分。两者合起来恰好确定所有一阶系数。

\begin{proposition}\label{g3:prop:normal-connection-expansion}
	设 \((M^n,g)\) 为光滑无边界黎曼流形，\(x\) 为在 \(p\) 处由正交单位基产生的法坐标。则
	\begin{equation}\label{g3:eq:normal-connection-derivative}
		\partial_k\Gamma^a_{ij}(0)
		=-\frac13\bigl(R^a{}_{i j k}+R^a{}_{j i k}\bigr),
	\end{equation}
	因而
	\begin{equation}\label{g3:eq:normal-connection-expansion}
		\Gamma^a_{ij}(x)
		=-\frac13\bigl(R^a{}_{i j k}+R^a{}_{j i k}\bigr)x^k
		+O(|x|^2).
	\end{equation}
\end{proposition}
\begin{proof}
	在法坐标中，\(t\mapsto tx\) 表示径向测地线。其坐标加速度为零，在时间一代入测地线方程，得到原点附近的恒等式
	\[
		\Gamma^a_{ij}(x)x^ix^j=0.
	\]
	令 \(C^a{}_{ij k}=\partial_k\Gamma^a_{ij}(0)\)。无挠性给出 \(C^a{}_{ij k}=C^a{}_{ji k}\)。把 \(\Gamma\) 的一阶 Taylor 展开代入上式，比较三次齐次项，得到
	\[
		C^a{}_{ij k}x^ix^jx^k=0\qquad\text{对所有 }x.
	\]
	对这个三次多项式取三次偏导，或对其作极化，得
	\begin{equation}\label{g3:eq:normal-connection-symmetric-part}
		C^a{}_{ij k}+C^a{}_{jk i}+C^a{}_{ki j}=0.
	\end{equation}
	这里三项已经包含了完全对称化，因为前两个下指标对称。

	另一方面，原点处 \(\Gamma=0\)，所以\eqref{g3:eq:curvature-coordinate-components}中的二次联络项消失，留下
	\[
		R^a{}_{b i j}=C^a{}_{jb i}-C^a{}_{ib j}.
	\]
	特别地，
	\[
		\begin{aligned}
			R^a{}_{i j k}+R^a{}_{j i k}
			&=C^a{}_{ki j}-C^a{}_{ji k}
			+C^a{}_{kj i}-C^a{}_{ij k}\\
			&=C^a{}_{ki j}+C^a{}_{jk i}-2C^a{}_{ij k}
			=-3C^a{}_{ij k}.
		\end{aligned}
	\]
	最后一步使用\eqref{g3:eq:normal-connection-symmetric-part}。这证明了偏导公式，再对光滑函数 \(\Gamma^a_{ij}\) 使用一阶 Taylor 公式，得到所述展开。
\end{proof}

联络的一阶系数本身依赖坐标，但法坐标的径向条件已经固定了其中可以自由改变的部分。因此在这种坐标中，一阶联络变化可以全部由原点的曲率写出。

\subsection{度量的二阶展开}
\label{g3:subsec:0303:02}

联络与度量之间的联系是度量相容性。它将度量的一阶偏导写成联络系数；再求一次导数，刚得到的一阶联络公式便转化为二阶度量公式。

\begin{theorem}\label{g3:thm:normal-metric-expansion}
	设 \((M^n,g)\) 为光滑无边界黎曼流形，\(x\) 为在 \(p\) 处由正交单位基产生的法坐标。则
	\begin{equation}\label{g3:eq:normal-metric-second-derivative}
		\partial_k\partial_lg_{ij}(0)
		=-\frac13\bigl(R_{ikjl}+R_{iljk}\bigr),
	\end{equation}
	并且
	\begin{equation}\label{g3:eq:normal-metric-expansion}
		g_{ij}(x)=\delta_{ij}-\frac13R_{ikjl}x^kx^l+O(|x|^3).
	\end{equation}
\end{theorem}
\begin{proof}
	度量相容性给出
	\[
		\partial_lg_{ij}=\Gamma^a_{li}g_{aj}+\Gamma^a_{lj}g_{ia}.
	\]
	再对 \(x^k\) 求导，并利用原点处 \(\Gamma=0\)、\(g_{ij}=\delta_{ij}\)，得到
	\[
		\partial_k\partial_lg_{ij}(0)
		=\partial_k\Gamma^j_{li}(0)+\partial_k\Gamma^i_{lj}(0).
	\]
	代入\eqref{g3:eq:normal-connection-derivative}，再将 \(R^a{}_{b i j}\) 改写成四阶分量，右侧变为
	\[
		-\frac13\bigl(R_{ikjl}+R_{lkji}+R_{jkil}+R_{lkij}\bigr).
	\]
	中间的 \(R_{lkji}\) 与最后的 \(R_{lkij}\) 因末两个指标反对称而抵消。成对交换对称性又给出 \(R_{jkil}=R_{iljk}\)，所以得到\eqref{g3:eq:normal-metric-second-derivative}。

	度量的一阶偏导在原点为零，Taylor 公式于是给出
	\[
		g_{ij}(x)=\delta_{ij}
		-\frac16\bigl(R_{ikjl}+R_{iljk}\bigr)x^kx^l+O(|x|^3).
	\]
	第二项中的两个求和式在交换哑指标 \(k,l\) 后相同，便得到\eqref{g3:eq:normal-metric-expansion}。
\end{proof}

这个二次项具有直接的长度含义。在 \(p\) 的切空间中把 \(a=(a^i)\)、\(x=(x^i)\) 看成向量，并在坐标点 \(x\) 考察具有相同分量 \(a^i\) 的坐标切向量，则
\[
	g_{ij}(x)a^ia^j
	=|a|^2-\frac13\mathrm{Rm}_p(a,x,a,x)
	+O(|x|^3)|a|^2.
\]
若 \(a\) 与 \(x\) 垂直，二次修正由它们张成平面的截面曲率决定。若 \(a\) 与 \(x\) 平行，这个曲率项因反对称性消失，与 Gauss 引理所给的径向长度保持相符。

\begin{example}\label{g3:exm:constant-curvature-normal-quadratic}
	设 \(n\ge2\)。若 \(p\) 处所有截面曲率等于常数 \(\kappa\)，由\cref{g3:cor:pointwise-isotropic-curvature}有
	\[
		R_{ikjl}=\kappa(\delta_{ij}\delta_{kl}-\delta_{il}\delta_{kj}),
	\]
	故二阶度量展开为
	\[
		g_{ij}(x)=\delta_{ij}
		-\frac{\kappa}{3}\bigl(|x|^2\delta_{ij}-x^ix^j\bigr)+O(|x|^3).
	\]
	对垂直于 \(x\) 的单位向量 \(a\)，其在坐标点 \(x\) 的长度平方为 \(1-\kappa|x|^2/3+O(|x|^3)\)。正曲率使这一横向长度在二阶上小于平坦模型，负曲率则使它大于平坦模型。径向与横向表现不同，因此不能把二阶修正理解为对所有方向的统一缩放。
\end{example}

固定一个二维切平面，再把其中的小圆盘经指数映射送到流形，可以直接观测这个截面的曲率。得到的是嵌入二维圆盘及其诱导度量；其边界是参数圆周的像。

\begin{proposition}\label{g3:prop:normal-plane-disk-geometry}
	设 \((M^n,g)\) 为光滑无边界黎曼流形，\(n\ge2\)，\(\sigma\subset\mathrm T_pM\) 为二维线性子空间。对足够小的 \(r>0\)，令
	\[
		D_r=\exp_p\{v\in\sigma:|v|_g\le r\},\qquad
		\partial D_r=\exp_p\{v\in\sigma:|v|_g=r\}.
	\]
	这里 \(\partial D_r\) 表示这个二维圆盘的边界。其诱导面积与边界长度满足
	\begin{equation}\label{g3:eq:normal-plane-disk-geometry}
		\begin{aligned}
			\operatorname{Area}_g(D_r)
			&=\pi r^2\left(1-\frac{K_p(\sigma)}{12}r^2+O(r^4)\right),\\
			\operatorname{Length}_g(\partial D_r)
			&=2\pi r\left(1-\frac{K_p(\sigma)}6r^2+O(r^4)\right).
		\end{aligned}
	\end{equation}
\end{proposition}
\begin{proof}
	选取正交单位基，使 \(\sigma\) 由 \(E_1,E_2\) 张成。令 \(z=(x^1,x^2)\)，其余坐标为零。指数映射在一个较大的切球上为微分同胚，所以限制后得到嵌入圆盘；其诱导度量矩阵为 \(h_{ab}(z)=g_{ab}(z,0)\)，\(1\le a,b\le2\)。记 \(K=K_p(\sigma)\)。曲率限制到这个二维空间后，只有 \(R_{1212}=K\) 一个独立分量，故\eqref{g3:eq:normal-metric-expansion}给出
	\[
		h_{ab}(z)=\delta_{ab}
		-\frac K3\bigl(|z|^2\delta_{ab}-z^az^b\bigr)
		+P_{ab,3}(z)+O(|z|^4),
	\]
	其中 \(P_{ab,3}\) 为光滑系数的三次齐次 Taylor 项。对单位矩阵附近的矩阵 \(I+H\)，有 \(\sqrt{\det(I+H)}=1+\tfrac12\operatorname{tr}H+O(\|H\|^2)\)。这里 \(H=O(|z|^2)\)，所以诱导面积密度为
	\[
		J_\sigma(z)=\sqrt{\det h(z)}
		=1-\frac K6|z|^2+P_3(z)+O(|z|^4),
	\]
	其中 \(P_3\) 仍为三次齐次多项式。在半径 \(r\) 的平面圆盘上积分，奇函数 \(P_3\) 的积分为零，而 \(\int_{|z|\le r}|z|^2\,\mathrm dz=\pi r^4/2\)。四阶余项积分为 \(O(r^6)\)，得到面积公式。

	对边界令 \(u(\theta)=(\cos\theta,\sin\theta)\)、\(v(\theta)=(-\sin\theta,\cos\theta)\)。其坐标速度为 \(rv(\theta)\)，由 \(u,v\) 为正交单位向量，有
	\[
		|\dot c_r(\theta)|_g^2
		=r^2\left(1-\frac K3r^2+O(r^3)\right).
	\]
	为同时控制积分余项，使用 Gauss 引理：圆盘上的径向方向长度为一，并与角向方向正交。因此极坐标中的面积密度既等于 \(|\dot c_r(\theta)|_g\,\mathrm dr\,\mathrm d\theta\)，也等于 \(rJ_\sigma(ru(\theta))\,\mathrm dr\,\mathrm d\theta\)，从而
	\[
		\operatorname{Length}_g(\partial D_r)
		=r\int_0^{2\pi}J_\sigma(ru(\theta))\,\mathrm d\theta.
	\]
	在 \(\theta\mapsto\theta+\pi\) 下，三次项变号，积分为零；余下四阶项贡献 \(O(r^5)\)。代入二次项便得长度公式。
\end{proof}

这些圆盘由同一个点的径向测地线组成；盘内其余方向的内在测地线未必是环境测地线，因此不能据此把圆盘视为全测地子流形。面积和长度是对这个指定二维指数像作出的测量。当 \(n>2\) 时，它也不同于环境中的 \(n\) 维测地球。截面曲率正是这类二维测量的首个非平坦系数。

\ProseParagraphBreak
二阶公式还有一个反向用途。前两节讨论的代数曲率张量是否真能出现在某个度量的一点处，可以通过直接指定度量的二次项来回答。

\begin{proposition}\label{g3:prop:curvature-local-realization}
	设 \(A\) 为 \(n\) 维实 Euclid 空间上的代数曲率张量，即满足\cref{g3:def:algebraic-curvature-tensor}中的两对反对称性、成对交换对称性和第一 Bianchi 恒等式。则存在原点的一个邻域及其上的光滑黎曼度量，使原点处度量为 Euclid 内积，曲率张量为 \(A\)。具体可取
	\begin{equation}\label{g3:eq:curvature-local-realization-metric}
		g_{ij}(x)=\delta_{ij}-\frac13A_{ikjl}x^kx^l
	\end{equation}
	并将定义域限制在足够小的开球内。
\end{proposition}
\begin{proof}
	先核对它确实是度量。由成对交换对称性和哑指标互换，
	\[
		A_{jkil}x^kx^l=A_{iljk}x^kx^l=A_{ikjl}x^kx^l,
	\]
	故 \(g_{ij}=g_{ji}\)。因为 \(A\) 固定，存在常数 \(C\ge0\)，使任意 \(v\in\mathbb R^n\) 满足
	\[
		\left|\frac13A_{ikjl}v^iv^jx^kx^l\right|
		\le C|x|^2|v|^2.
	\]
	在 \(C|x|^2<1/2\) 的小球内，\(g(v,v)\ge|v|^2/2\)，所以正定。系数是多项式，因而光滑。原点处 \(g_{ij}=\delta_{ij}\)、\(\partial_kg_{ij}=0\)，故 \(\Gamma^a_{ij}(0)=0\)，且
	\[
		\partial_k\partial_lg_{ij}(0)
		=-\frac13\bigl(A_{ikjl}+A_{iljk}\bigr).
	\]
	此处的 \(g(0)=I\)、\(\partial g(0)=0\) 正好满足\eqref{g3:eq:curvature-normal-second-derivatives}推导时所用的点处条件。将二阶偏导代入该曲率公式，整理两对反对称性与成对交换对称性，得到
	\[
		R^g_{ijkl}(0)
		=\frac13\bigl(2A_{ijkl}+A_{ikjl}-A_{iljk}\bigr).
	\]
	第一 Bianchi 恒等式为
	\[
		A_{ijkl}+A_{jkil}+A_{kijl}=0.
	\]
	其中 \(A_{jkil}=A_{iljk}\)、\(A_{kijl}=-A_{ikjl}\)，所以 \(A_{ikjl}-A_{iljk}=A_{ijkl}\)。代回即得 \(R^g_{ijkl}(0)=A_{ijkl}\)，完成构造。
\end{proof}

这个构造实现的是指定点的曲率数据。若某个曲率不等式在 \(A\) 处严格成立，则由曲率的连续性，将邻域再缩小后，该严格不等式仍成立；涉及所有单位向量或单位二维截面时，所需的统一余量来自相应单位参数空间的紧性。因此，严格的代数正性反例也能给出局部黎曼度量的反例。

\subsection{体积密度和小球体积}
\label{g3:subsec:0303:03}

度量矩阵决定每个方向的长度，体积密度则把这些方向的伸缩合成为一个标量。在法坐标中，令
\[
	J_p(x)=\sqrt{\det(g_{ij}(x))},\qquad
	\mu_g=J_p(x)|\mathrm dx^1\cdots\mathrm dx^n|.
\]
这里的密度定义不需要流形可定向。

\begin{proposition}\label{g3:prop:normal-volume-density}
	设 \((M^n,g)\) 为光滑无边界黎曼流形，在 \(p\) 处采用正交单位基产生的法坐标。则
	\begin{equation}\label{g3:eq:normal-volume-density}
		J_p(x)=1-\frac16\operatorname{Ric}_{kl}(p)x^kx^l+O(|x|^3).
	\end{equation}
\end{proposition}
\begin{proof}
	写 \(G=I+H\)。由\eqref{g3:eq:normal-metric-expansion}，有
	\[
		H_{ij}=-\frac13R_{ikjl}x^kx^l+O(|x|^3),
		\qquad H=O(|x|^2).
	\]
	行列式是矩阵元的多项式，故在单位矩阵附近，
	\[
		\det(I+H)=1+\operatorname{tr}H+O(\|H\|^2),
		\qquad
		\sqrt{\det(I+H)}=1+\tfrac12\operatorname{tr}H+O(\|H\|^2).
	\]
	另一方面，按本章的 Ricci 收缩约定，
	\[
		\operatorname{tr}H
		=-\frac13\sum_iR_{ikil}x^kx^l+O(|x|^3)
		=-\frac13\operatorname{Ric}_{kl}(p)x^kx^l+O(|x|^3).
	\]
	因为 \(\|H\|^2=O(|x|^4)\)，代回便得所述展开。
\end{proof}

沿固定单位方向 \(u\) 取 \(x=ru\)，则密度的二阶项为 \(-\operatorname{Ric}_p(u,u)r^2/6\)。这说明 Ricci 曲率测量的是沿该方向向外展开时，横向体积相对于 Euclid 模型的二阶变化。将所有方向平均，才得到标量曲率的作用。

\begin{theorem}\label{g3:thm:small-ball-volume}
	设 \((M^n,g)\) 为光滑无边界黎曼流形，\(p\in M\)，\(d\) 为 \(p\) 所在连通分支的内在距离，\(\omega_n\) 表示 \(\mathbb R^n\) 中单位球的 Euclid 体积。则当 \(r\downarrow0\) 时，内在距离球的体积满足
	\begin{equation}\label{g3:eq:small-ball-volume}
		\operatorname{Vol}_g(B_d(p,r))
		=\omega_n r^n\left(1-\frac{\operatorname{Scal}(p)}{6(n+2)}r^2+O(r^4)\right).
	\end{equation}
	若 \(n\ge2\)，足够小的距离球边界是光滑超曲面，其面积满足
	\begin{equation}\label{g3:eq:small-sphere-area}
		\operatorname{Area}_g(\partial B_d(p,r))
		=n\omega_n r^{n-1}\left(1-\frac{\operatorname{Scal}(p)}{6n}r^2+O(r^4)\right).
	\end{equation}
\end{theorem}
\begin{proof}
	由\cref{g3:lem:uniform-normal-neighborhood}，对足够小的 \(r\)，\(B_d(p,r)=\exp_p B_p(0,r)\)，且指数映射在更大的切球上为微分同胚。因此体积可以在普通 Euclid 球上积分：
	\[
		\operatorname{Vol}_g(B_d(p,r))=\int_{|x|<r}J_p(x)\,\mathrm dx.
	\]
	为了确定积分后的余项阶数，需把密度再展开一阶。\(J_p\) 光滑，一次项为零，二次项已经由\cref{g3:prop:normal-volume-density}确定，所以存在三次齐次多项式 \(P_{3,p}\)，使
	\begin{equation}\label{g3:eq:normal-density-cubic-remainder}
		J_p(x)=1-\frac16\operatorname{Ric}_{kl}(p)x^kx^l
		+P_{3,p}(x)+O(|x|^4).
	\end{equation}
	三次齐次项满足 \(P_{3,p}(-x)=-P_{3,p}(x)\)，在中心对称的球上积分为零。对二次项，由反射和坐标置换对称性，非对角积分为零，各对角积分相等；再把它们相加，得到
	\[
		\int_{|x|<r}x^kx^l\,\mathrm dx
		=\delta_{kl}\frac{\omega_n}{n+2}r^{n+2}.
	\]
	故二次项的积分为 \(-\omega_n\operatorname{Scal}(p)r^{n+2}/(6(n+2))\)。四阶余项的绝对值积分至多为常数乘以 \(r^{n+4}\)，除去主项 \(\omega_n r^n\) 后即得\eqref{g3:eq:small-ball-volume}。

	再求球面面积。令 \(x=ru\)、\(u\in S^{n-1}\)，并以 \(\mathrm d\sigma(u)\) 表示单位 Euclid 球面的面积密度。Gauss 引理给出径向单位向量与距离球面正交，因而在这些极坐标中，体积密度等于 \(\mathrm dr\) 与各距离球面面积密度的乘积。同时 Euclid 极坐标变换给出 \(\mathrm dx=r^{n-1}\mathrm dr\,\mathrm d\sigma\)，所以
	\[
		\operatorname{Area}_g(\partial B_d(p,r))
		=r^{n-1}\int_{S^{n-1}}J_p(ru)\,\mathrm d\sigma(u).
	\]
	将\eqref{g3:eq:normal-density-cubic-remainder}代入，三次项仍因 \(u\mapsto-u\) 对称而积分为零；又有
	\[
		\int_{S^{n-1}}1\,\mathrm d\sigma=n\omega_n,
		\qquad
		\int_{S^{n-1}}u^ku^l\,\mathrm d\sigma=\delta_{kl}\omega_n.
	\]
	这便给出\eqref{g3:eq:small-sphere-area}。面积公式直接来自极坐标密度积分，所需的余项估计也在该积分中取得。
\end{proof}

若 \(\operatorname{Scal}(p)>0\)，足够小的测地球体积小于同半径 Euclid 球；若 \(\operatorname{Scal}(p)<0\)，则大于它。标量曲率为零时，二次修正消失，体积差异由更高阶项决定。这个平均结果并不决定每个方向的 Ricci 曲率符号。

\begin{example}\label{g3:exm:surface-small-ball-check}
	二维时，\(\operatorname{Scal}=2K\)、\(\omega_2=\pi\)，所以
	\[
		\begin{aligned}
			\operatorname{Area}_g(B_d(p,r))
			&=\pi r^2\left(1-\frac{K(p)}{12}r^2+O(r^4)\right),\\
			\operatorname{Length}_g(\partial B_d(p,r))
			&=2\pi r\left(1-\frac{K(p)}6r^2+O(r^4)\right).
		\end{aligned}
	\]
	可以用刚构造的二次度量作具体核验。对实常数 \(\kappa\)，在原点附近令
	\[
		g_{ij}(x)=\delta_{ij}-\frac{\kappa}{3}
		\bigl(|x|^2\delta_{ij}-x^ix^j\bigr).
	\]
	它是\cref{g3:prop:curvature-local-realization}中的二维模型，故 \(K(0)=\kappa\)。改写成极坐标得
	\[
		g=\mathrm dr^2+r^2\left(1-\frac{\kappa r^2}{3}\right)\mathrm d\theta^2.
	\]
	在足够小的范围内括号内为正。翘曲积联络公式说明径向曲线以 \(r\) 为单位速参数是测地线；原来的直角坐标系数在原点光滑，故这些径向线与原点的测地线初值相接，\((x^1,x^2)\) 就是原点的法坐标。此时
	\[
		J_0(x)=\sqrt{1-\frac{\kappa|x|^2}{3}}
		=1-\frac{\kappa}{6}|x|^2-\frac{\kappa^2}{72}|x|^4+O(|x|^6).
	\]
	沿半径积分或沿圆周积分，分别得到
	\[
		\begin{aligned}
			\operatorname{Area}_g(B_d(0,r))
			&=\pi r^2\left(1-\frac{\kappa}{12}r^2
			-\frac{\kappa^2}{216}r^4+O(r^6)\right),\\
			\operatorname{Length}_g(\partial B_d(0,r))
			&=2\pi r\left(1-\frac{\kappa}{6}r^2
			-\frac{\kappa^2}{72}r^4+O(r^6)\right).
		\end{aligned}
	\]
	两式的二阶相对修正只用 \(K(0)=\kappa\) 即已确定。球的面积与边界长度采用不同系数，因为前者还对径向变量作了一次积分；再高阶的系数则取决于这个具体度量。
\end{example}

\subsection{展开余项及正则性}
\label{g3:subsec:0303:04}

上述公式同时记录了系数和误差，两者承担不同职责。曲率决定二阶系数；误差界还依赖附近度量系数的更高阶导数，不能只由 \(R(p)\) 的值控制。

\ProseParagraphBreak

在一个固定法坐标图中，令
\[
	E_{ij}(x)=g_{ij}(x)-\delta_{ij}+\frac13R_{ikjl}(p)x^kx^l.
\]
它在原点的零阶、一阶和二阶导数全部为零。对每个满足 \(|\alpha|\le2\) 的多重指标，分别对 \(\partial^\alpha E_{ij}\) 使用 Taylor 公式，得到
\begin{equation}\label{g3:eq:normal-metric-remainder-derivatives}
	\partial^\alpha E_{ij}(x)=O(|x|^{3-|\alpha|}).
\end{equation}
这些导数估计来自光滑函数的 Taylor 余项，而非对一个任意 \(O\) 记号形式求导。同样，若要在小球体积公式中保留相对 \(O(r^4)\)，就应使用\eqref{g3:eq:normal-density-cubic-remainder}的三次项和四阶余项；仅对密度中的 \(O(|x|^3)\) 作绝对值估计，只能得到相对 \(O(r^3)\)。

\begin{proposition}\label{g3:prop:normal-expansion-uniform-remainders}
	设 \((M,g)\) 为光滑无边界黎曼流形，\(K\subset M\) 紧。则存在共同半径 \(r_0>0\)，使本节的联络、度量、密度、小球体积展开，以及在 \(n\ge2\) 时的小球面面积展开，对所有中心 \(p\in K\)、\(0<r<r_0\) 都成立，且各余项中的常数可与 \(p\) 无关。
\end{proposition}
\begin{proof}
	由\cref{g3:lem:uniform-normal-neighborhood}，\(K\) 上存在统一可用的法坐标半径。取有限个较小的坐标邻域覆盖 \(K\)，使各自闭包紧含于一个有光滑正交单位标架 \((E_1(p),\ldots,E_n(p))\) 的较大邻域。在每个这样的邻域中，
	\[
		F(p,x)=\exp_p(x^iE_i(p))
	\]
	在共同的小 \(x\) 球上光滑依赖 \((p,x)\)。这由\cref{g3:prop:exponential-local}的指数映射联合光滑性和标架光滑性得到。拉回度量系数
	\[
		G_{ij}(p,x)=g_{F(p,x)}
		\bigl(\mathrm d_xF(p,x)e_i,\mathrm d_xF(p,x)e_j\bigr)
	\]
	因此也联合光滑。将 \(x\) 球再缩小一层后，各底邻域的闭包与闭 \(x\) 球的乘积紧，所以证明中使用的有限阶导数均有统一上界，度量正定性也有统一下界。Taylor 公式、行列式及平方根的导数估计便给出统一余项界。

	有限个邻域取最小半径、最大误差常数，即得对 \(K\) 一致的结论。改用同一点处的其他正交单位基，只对 \(x\) 作正交线性变换，故不改变 \(|x|\) 的大小；相应张量范数也保持不变。体积和面积的积分估计由这些统一界直接得到。
\end{proof}

这里始终假设原度量光滑，从而指数映射与法坐标中的度量系数也光滑。若降低原度量的正则性，需要分别核对测地线方程、指数映射和换坐标后的系数具有多少阶导数。就 Taylor 估计本身而言，已给定坐标中的 \(C^2\) 系数只能直接保证二阶展开的 \(o(|x|^2)\) 余项；三阶导数有界给出 \(O(|x|^3)\)，而密度展开至三次项后控制 \(O(|x|^4)\) 需要相应的四阶导数界。原坐标中的可微阶数不能未经检查就沿用到法坐标中。

% !TeX root = ../../Book3.tex
% 纲要标题：### 3.4 模型与代数分解
% 纲要位置：几何学计划纲要.md，第 2114 行。
% 正文状态：未撰写；本文件目前是框架。

% 纲要内容开始
% 
% * 空间形式、乘积和翘曲积
% * 共形变化和标量曲率
% * Weyl 张量及低维特例
% * 由模型检验符号与等号
% 
% ---
% 
% 纲要内容结束

\section{模型与代数分解}
\label{g3:sec:0304}

待撰写

\subsection{空间形式、乘积和翘曲积}
\label{g3:subsec:0304:01}
待撰写

\subsection{共形变化和标量曲率}
\label{g3:subsec:0304:02}
待撰写

\subsection{Weyl 张量及低维特例}
\label{g3:subsec:0304:03}
待撰写

\subsection{由模型检验符号与等号}
\label{g3:subsec:0304:04}
待撰写
% 纲要条目排为编号小节标题，并进入目录与书签；说明段仍保留于纲要注释中。


\begin{chaptersummary}
	曲率是协变求导的交换差扣除 Lie 括号方向导数后留下的张量。度量相容性给出反对称性，无挠性给出第一 Bianchi 恒等式；第二 Bianchi 恒等式进一步限制曲率的变化，其收缩产生 Ricci 散度与标量曲率微分之间的关系。把四阶张量视为二向量空间上的自伴算子后，截面曲率只检测可分解二向量；这一点解释了高维中截面正性与算子正性之间的差别。

	\ProseParagraphBreak

	法坐标的径向条件与曲率公式共同确定联络的一阶项，度量相容性再确定度量的二阶项。二维切平面指数像的周长和面积感受到相应的截面曲率，法坐标体积密度的二次项由 Ricci 曲率控制，小球体积的方向平均则由标量曲率控制。Taylor 余项与对称积分使这些解释成为带有误差阶数的局部公式；紧中心集上的共同法坐标范围还给出一致估计。

	\ProseParagraphBreak

	乘积与翘曲积使曲率的不同来源可以逐项计算，共形变化把长度尺度的导数转化为曲率修正。Weyl 分解将曲率拆成相互正交的三部分：标量部分、无迹 Ricci 部分和无迹四阶部分。二维曲率由一个数决定，三维由 Ricci 决定，四维起则可以有 Ricci 收缩看不见的曲率。标准模型既核验了符号，也说明局部体积或某一次收缩不能单独恢复全部几何。
\end{chaptersummary}

\BookExercises


\begin{exercise}\label{g3:ex:constant-metric-curvature-scaling}
	设 \((M^n,g)\) 为光滑黎曼流形，\(a>0\) 为常数，\(\widetilde g=a^2g\)。将两种度量下的切空间、张量空间和二向量空间按恒等映射识别。
	\begin{enumerate}
		\item 求 \((1,3)\) 型曲率 \(R\)、四阶协变曲率 \(\mathrm{Rm}\)、协变 Ricci 张量及标量曲率在此次缩放下的变化。
		\item 求曲率算子 \(\mathcal R\) 与曲率平方范数 \(|\mathrm{Rm}|^2\) 的变化。说明在固定基与分别单位化的基中，曲率分量的比较为何不同。
	\end{enumerate}
\end{exercise}
\begin{solution}
	1. 原联络 \(\nabla\) 仍然无挠，而且
	\[
		X\widetilde g(Y,Z)
		=a^2Xg(Y,Z)
		=\widetilde g(\nabla_XY,Z)+\widetilde g(Y,\nabla_XZ).
	\]
	由 \cref{g3:thm:levi-civita} 的唯一性，\(\widetilde\nabla=\nabla\)。曲率作为由联络定义的 \((1,3)\) 型张量因而不变：
	\[
		\widetilde R(X,Y)Z=R(X,Y)Z.
	\]
	四阶张量还含一次度量配对，故
	\[
		\widetilde{\mathrm{Rm}}(X,Y,Z,W)
		=\widetilde g(\widetilde R(X,Y)W,Z)
		=a^2\mathrm{Rm}(X,Y,Z,W).
	\]
	在同一坐标基中，\(\widetilde g^{ik}=a^{-2}g^{ik}\)，所以
	\[
		\widetilde{\operatorname{Ric}}_{jl}
		=\widetilde g^{ik}\widetilde R_{ijkl}
		=g^{ik}R_{ijkl}=\operatorname{Ric}_{jl},
		\qquad
		\widetilde{\operatorname{Scal}}
		=a^{-2}\operatorname{Scal}.
	\]
	这里不变的是协变二阶 Ricci 张量；将它用度量升起一个指标后，所得自伴算子的特征值也乘以 \(a^{-2}\)。

	2. 二向量空间上的诱导内积含两次度量，满足
	\[
		\langle\xi,\eta\rangle_{\Lambda^2,\widetilde g}
		=a^4\langle\xi,\eta\rangle_{\Lambda^2,g}.
	\]
	由曲率算子的定义，对可分解二向量 \(\xi=X\wedge Y\)、\(\eta=Z\wedge W\)，
	\[
		a^4\langle\widetilde{\mathcal R}\xi,\eta\rangle_{\Lambda^2,g}
		=a^2\langle\mathcal R\xi,\eta\rangle_{\Lambda^2,g}.
	\]
	这类二向量张成整个空间，内积又非退化，因此
	\[
		\widetilde{\mathcal R}=a^{-2}\mathcal R.
	\]
	在二形式空间上用各自的度量同构定义曲率算子时，结论相同：\(\widetilde\flat_2=a^4\flat_2\)、\(\widetilde\sharp_2=a^{-4}\sharp_2\)，两个同构的尺度在共轭时抵消。

	设 \((e_i)\) 为 \(g\) 正交单位基，则 \((\widetilde e_i)=(a^{-1}e_i)\) 为 \(\widetilde g\) 正交单位基。在同一组固定向量 \(e_i\) 上，\(\widetilde{\mathrm{Rm}}\) 的分量是原来的 \(a^2\) 倍；在各自的单位基上则有
	\[
		\widetilde{\mathrm{Rm}}
			(\widetilde e_i,\widetilde e_j,\widetilde e_k,\widetilde e_l)
		=a^{-2}\mathrm{Rm}(e_i,e_j,e_k,e_l).
	\]
	因此
	\[
		|\widetilde{\mathrm{Rm}}|_{\widetilde g}^2
		=a^{-4}|\mathrm{Rm}|_g^2.
	\]
	同样，每个固定二维平面的截面曲率变为原来的 \(a^{-2}\) 倍。张量本身的缩放与分量比较中基向量的缩放共同决定这些公式。
\end{solution}

\begin{exercise}\label{g3:ex:algebraic-curvature-dimension}
	设 \(V\) 为 \(n\) 维实向量空间，\(n\ge2\)。将 \(\mathcal S=\operatorname{Sym}^2(\Lambda^2V^*)\) 视为满足两对槽位内反对称及成对交换对称的四阶协变张量空间。对 \(A\in\mathcal S\) 定义
	\[
		(\mathfrak bA)(X,Y,Z,W)
		=A(X,Y,Z,W)+A(Y,Z,X,W)+A(Z,X,Y,W).
	\]
	\begin{enumerate}
		\item 证明 \(\mathfrak bA\) 是交替四形式，且线性映射
		\(\mathfrak b:\mathcal S\to\Lambda^4V^*\) 满射。
		\item 证明代数曲率张量空间就是 \(\ker\mathfrak b\)，并求其维数。特别求 \(n=2,3,4\) 时的数值。
	\end{enumerate}
\end{exercise}
\begin{solution}
	1. 利用成对交换与第一对的反对称性，可将定义写成
	\[
		(\mathfrak bA)(X,Y,Z,W)
		=A(X,Y,Z,W)-A(X,Z,Y,W)+A(X,W,Y,Z).
	\]
	记右侧为 \(B(X,Y,Z,W)\)。例如，交换前两个输入并再次使用成对交换，得到
	\[
		\begin{aligned}
			B(Y,X,Z,W)
			&=-A(X,Y,Z,W)-A(Y,Z,X,W)+A(Y,W,X,Z)\\
			&=-A(X,Y,Z,W)-A(X,W,Y,Z)+A(X,Z,Y,W)\\
			&=-B(X,Y,Z,W).
		\end{aligned}
	\]
	其余两个相邻换位也可直接展开：
	\[
		\begin{aligned}
			B(X,Z,Y,W)
			&=A(X,Z,Y,W)-A(X,Y,Z,W)+A(X,W,Z,Y)\\
			&=-B(X,Y,Z,W),\\
			B(X,Y,W,Z)
			&=A(X,Y,W,Z)-A(X,W,Y,Z)+A(X,Z,Y,W)\\
			&=-B(X,Y,Z,W).
		\end{aligned}
	\]
	相邻换位生成全部置换，故 \(B\) 是交替四形式。

	反过来，任取 \(\eta\in\Lambda^4V^*\)，也把它看作四阶张量。交换每一对内的输入使它变号，交换两对输入是一个偶置换，所以 \(\eta\in\mathcal S\)。循环前三个槽位也是偶置换，因此
	\[
		\mathfrak b\eta=3\eta,\qquad
		\mathfrak b\left(\frac13\eta\right)=\eta.
	\]
	这证明满射。当 \(n<4\) 时，目标空间为零，上述结论仍成立。

	2. \(\mathcal S\) 已包含定义中的两对反对称性和成对交换对称性，剩下的第一 Bianchi 恒等式恰为 \(\mathfrak bA=0\)。所以其核正是代数曲率张量空间。

	令 \(N=\dim\Lambda^2V^*=n(n-1)/2\)。在一组二形式基中，\(\mathcal S\) 对应实对称 \(N\times N\) 矩阵，因此维数为 \(N(N+1)/2\)。满射的秩为 \(\binom n4\)，秩与零度关系给出
	\[
		\begin{aligned}
			\dim\ker\mathfrak b
			&=\frac{N(N+1)}2-\binom n4\\
			&=\frac{n(n-1)\bigl(n(n-1)+2\bigr)}8
				-\frac{n(n-1)(n-2)(n-3)}{24}\\
			&=\frac{n^2(n^2-1)}{12}.
		\end{aligned}
	\]
	这里约定 \(n<4\) 时 \(\binom n4=0\)。于是二维、三维、四维的维数依次为
	\[
		1,\qquad6,\qquad20.
	\]
	特别地，三维中每个对称的二向量双线性型都满足第一 Bianchi 恒等式；四维中则恰有一个独立的线性约束。
\end{solution}

\begin{exercise}\label{g3:ex:curvature-covariant-tensors}
	设 \((M,g)\) 为光滑黎曼流形，\(\nabla\) 为其 Levi-Civita 联络，并在协变张量上使用诱导联络
	\[
		(\nabla_XT)(Z_1,\ldots,Z_k)
		=X\bigl(T(Z_1,\ldots,Z_k)\bigr)
		-\sum_{a=1}^kT(Z_1,\ldots,\nabla_XZ_a,\ldots,Z_k).
	\]
	对函数或协变张量场 \(T\)，记
	\[
		\mathcal D_{X,Y}T
		=\nabla_X\nabla_YT-\nabla_Y\nabla_XT-\nabla_{[X,Y]}T.
	\]
	\begin{enumerate}
		\item 证明对任意一形式 \(\alpha\)，
		\[
			(\mathcal D_{X,Y}\alpha)(Z)=-\alpha(R(X,Y)Z).
		\]
		\item 对一般协变 \(k\) 阶张量，证明
		\[
			(\mathcal D_{X,Y}T)(Z_1,\ldots,Z_k)
			=-\sum_{a=1}^k
				T(Z_1,\ldots,R(X,Y)Z_a,\ldots,Z_k).
		\]
		\item 说明函数的协变二阶导数可以交换，但一形式的协变二阶导数一般不能交换，并在正曲率圆球面上给出一个具体的非零值。
	\end{enumerate}
\end{exercise}
\begin{solution}
	1. 把诱导联络的定义使用两次，得到
	\[
		\begin{aligned}
			(\nabla_X\nabla_Y\alpha)(Z)
			&=X\bigl((\nabla_Y\alpha)(Z)\bigr)
				-(\nabla_Y\alpha)(\nabla_XZ)\\
			&=XY\bigl(\alpha(Z)\bigr)
				-X\bigl(\alpha(\nabla_YZ)\bigr)
				-Y\bigl(\alpha(\nabla_XZ)\bigr)
				+\alpha(\nabla_Y\nabla_XZ).
		\end{aligned}
	\]
	交换 \(X,Y\) 后，两个中间项相消。再减去
	\[
		(\nabla_{[X,Y]}\alpha)(Z)
		=[X,Y]\bigl(\alpha(Z)\bigr)-\alpha(\nabla_{[X,Y]}Z),
	\]
	作用在标量函数 \(\alpha(Z)\) 上的三项也相消，余下
	\[
		(\mathcal D_{X,Y}\alpha)(Z)
		=-\alpha\bigl(
			\nabla_X\nabla_YZ-\nabla_Y\nabla_XZ-\nabla_{[X,Y]}Z
			\bigr),
	\]
	即所求公式。负号来自协变张量对输入向量变化的扣除。

	2. 对函数 \(f\)，有
	\[
		\mathcal D_{X,Y}f=XYf-YXf-[X,Y]f=0.
	\]
	诱导联络满足张量积的乘积规则。为检查两次求导后的规则，展开
	\[
		\begin{aligned}
			\nabla_X\nabla_Y(S\otimes T)
			={}&(\nabla_X\nabla_YS)\otimes T
				+(\nabla_YS)\otimes(\nabla_XT)\\
			&+(\nabla_XS)\otimes(\nabla_YT)
				+S\otimes(\nabla_X\nabla_YT).
		\end{aligned}
	\]
	交换 \(X,Y\) 时，两个交叉项相消；再减去括号方向的导数，得到
	\[
		\mathcal D_{X,Y}(S\otimes T)
		=(\mathcal D_{X,Y}S)\otimes T
			+S\otimes(\mathcal D_{X,Y}T).
	\]
	对 \(S=\alpha^1\otimes\cdots\otimes\alpha^k\) 逐次使用此式，再代入第 1 问的一形式公式，恰好得到每个槽位各插入一次 \(-R(X,Y)\) 的和。任意协变张量在局部余标架中都是这种纯张量的光滑函数线性组合；函数上的 \(\mathcal D_{X,Y}\) 为零，故同一公式对一般 \(T\) 成立。

	3. 函数 \(f\) 的协变二阶导数为
	\[
		H_f(X,Y)=X(Yf)-(\nabla_XY)f.
	\]
	利用无挠性，
	\[
		H_f(X,Y)-H_f(Y,X)
		=[X,Y]f-(\nabla_XY-\nabla_YX)f=0.
	\]
	这里需要区分普通迭代方向导数与张量二阶导数：\(XYf-YXf=[X,Y]f\)，对交换的坐标向量场这就是普通混合偏导相等。

	对一形式定义张量二阶导数
	\[
		(\nabla^2\alpha)(X,Y;Z)
		=\bigl(\nabla_X\nabla_Y\alpha
			-\nabla_{\nabla_XY}\alpha\bigr)(Z).
	\]
	再次利用无挠性及第 1 问，
	\[
		(\nabla^2\alpha)(X,Y;Z)-(\nabla^2\alpha)(Y,X;Z)
		=-\alpha(R(X,Y)Z).
	\]
	在半径 \(a\) 的圆球面取一点处的正交单位向量 \(e_1,e_2\)，并取局部一形式使 \(\alpha_p=e_1^\flat\)。该点有
	\[
		R(e_1,e_2)e_2=a^{-2}e_1,
	\]
	所以当 \((X,Y;Z)=(e_1,e_2;e_2)\) 时，上述交换差为 \(-a^{-2}\ne0\)。公式的两侧都是张量，任意光滑局部延拓都给出相同的点处数值。
\end{solution}

\begin{exercise}\label{g3:ex:three-dimensional-ricci-positivity}
	在三维黎曼流形的一点 \(p\)，设 Ricci 张量对应的自伴算子具有特征值 \(r_1,r_2,r_3\)。
	\begin{enumerate}
		\item 利用 \cref{g3:prop:low-dimensional-curvature}，求曲率算子的三个特征值及一组相应的特征二向量。
		\item 证明该点所有截面曲率严格为正，当且仅当
		\[
			r_1+r_2>r_3,\qquad
			r_1+r_3>r_2,\qquad
			r_2+r_3>r_1.
		\]
		\item 给出三个严格正的 \(r_i\)，使某个截面曲率为负，并说明这些数据可以由某个局部黎曼度量实现。
	\end{enumerate}
\end{exercise}
\begin{solution}
	1. 取 \(g_p\) 正交单位基 \(e_1,e_2,e_3\)，使 \(\operatorname{Ric}(e_i,e_j)=r_i\delta_{ij}\)，并记 \(s=r_1+r_2+r_3\)。三维公式为
	\[
		\mathrm{Rm}=g\wedge\operatorname{Ric}-\frac{s}{4}g\wedge g.
	\]
	展开乘积，在此基中得到
	\[
		R_{ijkl}
		=\left(r_i+r_j-\frac{s}{2}\right)
			(\delta_{ik}\delta_{jl}-\delta_{il}\delta_{jk}).
	\]
	上式中的 \(i,j\) 是自由指标，不对它们求和。因此在正交单位二向量基
	\(e_1\wedge e_2,e_1\wedge e_3,e_2\wedge e_3\) 中，曲率算子为对角矩阵，三个特征值分别为
	\[
		\lambda_{12}=\frac{r_1+r_2-r_3}{2},\qquad
		\lambda_{13}=\frac{r_1+r_3-r_2}{2},\qquad
		\lambda_{23}=\frac{r_2+r_3-r_1}{2}.
	\]

	2. 若所有截面曲率为正，对由 \(e_i,e_j\) 张成的三个截面计算，便有 \(\lambda_{ij}>0\)，这给出所列三个不等式。反过来，若它们成立，则任意二维平面的单位二向量可写成
	\[
		\xi=b_{12}e_1\wedge e_2+b_{13}e_1\wedge e_3
			+b_{23}e_2\wedge e_3,\qquad
		b_{12}^2+b_{13}^2+b_{23}^2=1.
	\]
	于是该平面的截面曲率为
	\[
		\langle\mathcal R\xi,\xi\rangle
		=\lambda_{12}b_{12}^2+\lambda_{13}b_{13}^2
			+\lambda_{23}b_{23}^2>0.
	\]
	这证明充分性。三个不等式两两相加还给出每个 \(r_i>0\)，符合正截面曲率蕴含正 Ricci 曲率的收缩关系。

	3. 取 \((r_1,r_2,r_3)=(1,1,3)\)，则
	\[
		(\lambda_{12},\lambda_{13},\lambda_{23})
		=\left(-\frac12,\frac32,\frac32\right).
	\]
	Ricci 张量正定，但 \(e_1,e_2\) 张成的截面曲率为负。为了实现这些数值，在 \(\mathbb R^3\) 的 Euclid 内积 \(g_0\) 下取对称张量 \(h=\operatorname{diag}(1,1,3)\)，令
	\[
		A=g_0\wedge h-\frac54g_0\wedge g_0.
	\]
	两个乘积均有代数曲率对称性，故 \(A\) 是代数曲率张量。由收缩公式，
	\[
		\operatorname{Ric}(A)
		=\bigl(h+5g_0\bigr)-\frac54(4g_0)=h.
	\]
	因此 \cref{g3:prop:curvature-local-realization} 给出一个原点邻域上的光滑度量，其原点曲率正是 \(A\)。由连续性再缩小邻域，可使 Ricci 张量在整个邻域保持正定，同时坐标前两个方向张成的截面仍有负曲率。这里 Ricci 的正定性有严格的正下界，负截面值也有严格的负余量，所以二者都在小邻域内保留。
\end{solution}

\begin{exercise}\label{g3:ex:curvature-tensor-operator-norm}
	设 \((M^n,g)\) 为光滑黎曼流形，\(n\ge2\)。曲率张量范数按四个有序指标的全部分量平方和定义，曲率算子的平方范数记为 \(|\mathcal R|_{\mathrm{HS}}^2\)，定义为 \(\Lambda^2\mathrm{T}_pM\) 的一组正交单位基在该算子下的像向量长度平方和。
	\begin{enumerate}
		\item 证明
		\[
			|\mathrm{Rm}|^2=4|\mathcal R|_{\mathrm{HS}}^2.
		\]
		\item 结合式\eqref{g3:eq:curvature-norm-decomposition}与低维公式，证明
		\[
			n=3:\quad|\mathrm{Rm}|^2=4|\operatorname{Ric}|^2
				-\operatorname{Scal}^2,
			\qquad
			n=2:\quad|\mathrm{Rm}|^2=\operatorname{Scal}^2.
		\]
		解释第 1 问中系数 \(4\) 的来源。
	\end{enumerate}
\end{exercise}
\begin{solution}
	1. 取一点处的正交单位基 \(e_1,\ldots,e_n\)。由曲率算子的分量表达，
	\[
		|\mathcal R|_{\mathrm{HS}}^2
		=\sum_{i<j}|\mathcal R(e_i\wedge e_j)|^2
		=\sum_{i<j,\ k<l}R_{ijkl}^2.
	\]
	另一方面，
	\[
		|\mathrm{Rm}|^2=\sum_{i,j,k,l}R_{ijkl}^2.
	\]
	第一对或第二对指标相等时，分量为零。对每个 \(i<j,k<l\)，独立交换前一对或后一对，只改变分量的符号，故其平方在全部有序指标的和中出现四次。因此
	\[
		\sum_{i,j,k,l}R_{ijkl}^2
		=4\sum_{i<j,\ k<l}R_{ijkl}^2,
	\]
	得到所求等式。

	2. 三维时 \(W=0\)。记 \(s=\operatorname{Scal}\)、\(E=\operatorname{Ric}-\tfrac{s}{3}g\)。由于 \(\operatorname{tr}_gE=0\)、\(|g|^2=3\)，有
	\[
		|\operatorname{Ric}|^2=|E|^2+\frac{s^2}{3}.
	\]
	把 \(n=3\) 代入范数分解式，得到
	\[
		|\mathrm{Rm}|^2
		=4|E|^2+\frac{s^2}{3}
		=4|\operatorname{Ric}|^2-s^2.
	\]
	二维时曲率算子作用在一维空间上，等于 \(K\operatorname{id}\)，故
	\[
		|\mathrm{Rm}|^2=4K^2=(2K)^2=\operatorname{Scal}^2.
	\]
	也可直接列出四个可能非零的分量：
	\(R_{1212}=R_{2121}=K\)、\(R_{1221}=R_{2112}=-K\)。

	系数 \(4\) 来自两对反对称指标各自的两次计数。曲率算子使用的单位基是 \(e_i\wedge e_j\)、\(i<j\)，每个方向只列一次；张量范数则对所有四个有序指标求和。固定这一规范后，同一几何对象的张量表达与算子表达具有上述确定的范数换算。
\end{solution}

\begin{exercise}\label{g3:ex:flat-nonconstant-conformal-metric}
	设 \(U\subset\mathbb R^2\) 为开集，\(u\in C^\infty(U)\)，取度量
	\[
		g=e^{2u(x,y)}(\mathrm{d}x^2+\mathrm{d}y^2).
	\]
	\begin{enumerate}
		\item 用正交余标架的结构方程计算 Gauss 曲率，证明
		\[
			K=-e^{-2u}(u_{xx}+u_{yy}),
		\]
		并核对它与式\eqref{g3:eq:conformal-scalar}一致。
		\item 在整个 \(\mathbb R^2\) 取 \(u=ax\)，其中 \(a\ne0\) 为常数。证明度量平坦，并证明
		\[
			F(x,y)=\frac{e^{ax}}{a}\bigl(\cos(ay),\sin(ay)\bigr)
		\]
		给出每点附近的平直坐标。说明 \(F\) 为局部等距映射，但在整个 \(\mathbb R^2\) 上不是单射。
	\end{enumerate}
\end{exercise}
\begin{solution}
	1. 取正交单位余标架
	\[
		\theta^1=e^u\mathrm{d}x,\qquad
		\theta^2=e^u\mathrm{d}y.
	\]
	直接求外微分，
	\[
		\mathrm{d}\theta^1=-e^uu_y\,\mathrm{d}x\wedge\mathrm{d}y,
		\qquad
		\mathrm{d}\theta^2=e^uu_x\,\mathrm{d}x\wedge\mathrm{d}y.
	\]
	记 \(\alpha=\omega^2{}_1=-\omega^1{}_2\)。无挠结构方程为
	\[
		\mathrm{d}\theta^1-\alpha\wedge\theta^2=0,\qquad
		\mathrm{d}\theta^2+\alpha\wedge\theta^1=0.
	\]
	设 \(\alpha=A\,\mathrm{d}x+B\,\mathrm{d}y\)，第一式给出 \(A=-u_y\)，第二式给出 \(B=u_x\)。所以
	\[
		\alpha=-u_y\,\mathrm{d}x+u_x\,\mathrm{d}y,\qquad
		\mathrm{d}\alpha=(u_{xx}+u_{yy})\,\mathrm{d}x\wedge\mathrm{d}y.
	\]
	由式\eqref{g3:eq:surface-curvature-form}，
	\[
		\mathrm{d}\alpha=-K\theta^1\wedge\theta^2
		=-Ke^{2u}\,\mathrm{d}x\wedge\mathrm{d}y,
	\]
	比较系数即得 \(K=-e^{-2u}(u_{xx}+u_{yy})\)。

	为核对共形标量公式，底度量取 Euclid 度量 \(g_0\)。本书符号约定给出
	\[
		\Delta_{g_0}u=-(u_{xx}+u_{yy}),\qquad
		\operatorname{Scal}_{g_0}=0.
	\]
	二维公式因而为
	\[
		\operatorname{Scal}_g
		=2e^{-2u}\Delta_{g_0}u
		=-2e^{-2u}(u_{xx}+u_{yy})=2K,
	\]
	与结构方程计算相同。

	2. \(u=ax\) 的二阶偏导均为零，故 \(K=0\)。二维的全部曲率由 \(K\) 决定，所以此度量平坦，而 \(e^{2ax}\) 随 \(x\) 非恒定。

	记 \(F=(F^1,F^2)\)。对给定表达式求导得到
	\[
		\begin{aligned}
			\mathrm{d}F^1
			&=e^{ax}\bigl(\cos(ay)\,\mathrm{d}x-\sin(ay)\,\mathrm{d}y\bigr),\\
			\mathrm{d}F^2
			&=e^{ax}\bigl(\sin(ay)\,\mathrm{d}x+\cos(ay)\,\mathrm{d}y\bigr).
		\end{aligned}
	\]
	因此
	\[
		F^*(\mathrm{d}X^2+\mathrm{d}Y^2)
		=(\mathrm{d}F^1)^2+(\mathrm{d}F^2)^2
		=e^{2ax}(\mathrm{d}x^2+\mathrm{d}y^2)=g.
	\]
	导数矩阵的行列式为 \(e^{2ax}>0\)，所以每一点附近 \(F\) 都是到一个平面开集的微分同胚。将邻域限制到这个可逆范围后，\((X,Y)=(F^1,F^2)\) 正是度量系数为单位矩阵的平直坐标。

	但周期性给出
	\[
		F\left(x,y+\frac{2\pi}{|a|}\right)=F(x,y).
	\]
	右边的两个原像点不同，所以整个映射不单射。这一例子同时给出了非恒定共形因子的平坦度量及显式局部等距坐标，并表明这些坐标不能由同一映射直接拼成全局单射坐标。
\end{solution}

\begin{exercise}\label{g3:ex:scalar-flat-warped-model}
	固定 \(b>0\)，令
	\[
		r(t)=b(t+\sinh t\cosh t),\qquad f(r(t))=b\cosh^2t.
	\]
	证明这在整个实轴上定义了光滑正函数 \(f\)。在 \(M=\mathbb R\times S^2\) 上取 \(g=\mathrm dr^2+f(r)^2g_{S^2}\)，其中 \(S^2\) 的截面曲率为一。计算径向与纤维截面曲率、Ricci 特征值及标量曲率，并证明 \(g\) 完备。这个模型是否平坦？
\end{exercise}
\begin{solution}
	因为 \(r'(t)=2b\cosh^2t>0\)，且 \(r(t)\to\pm\infty\) 当 \(t\to\pm\infty\)，映射 \(r\) 是实轴的光滑微分同胚，所定义的 \(f\) 光滑且 \(f\ge b\)。按 \(r\) 求导得
	\[
		f'=\tanh t,\qquad
		f''=\frac{1}{2b\cosh^4t}=\frac{b}{2f^2},\qquad
		1-(f')^2=\frac bf.
	\]
	由式\eqref{g3:eq:warped-sectional-curvatures}，包含单位径向量 \(T=\partial_r\) 的截面与纯纤维截面分别满足
	\[
		K_{\mathrm{rad}}=-\frac{b}{2f^3},\qquad
		K_{\mathrm{tan}}=\frac{b}{f^3}.
	\]
	因此在 \(g\) 正交单位基 \((T,E_1,E_2)\) 中，Ricci 的特征值为
	\[
		-\frac{b}{f^3},\quad \frac{b}{2f^3},\quad \frac{b}{2f^3}.
	\]
	它们之和为零，故 \(\operatorname{Scal}_g=0\)，而径向截面处处具有负曲率，度量并不平坦。这里 \(f(1-(f')^2)=b\) 也可直接作为标量曲率为零的方程 \(2ff''=1-(f')^2\) 的第一积分核对；求导时不需要除以可能为零的 \(f'\)。

	最后证明完备。任意连接 \((r,x)\) 与 \((s,y)\) 的曲线满足
	\[
		L_g\ge |r-s|,\qquad L_g\ge b\,d_{S^2}(x,y).
	\]
	所以一个 \(d_g\)-Cauchy 序列的两个坐标分别在 \(\mathbb R\) 和 \(S^2\) 中是 Cauchy 序列，收敛到某个 \((r_\infty,x_\infty)\)。在一个包含所有充分靠后实坐标的紧区间上，\(f\) 有上界 \(C\)。先沿实方向、再沿球面方向连接，得到
	\[
		d_g((r,x),(r_\infty,x_\infty))
		\le |r-r_\infty|+C\,d_{S^2}(x,x_\infty).
	\]
	故原序列在 \(d_g\) 中收敛，证明了度量完备。这给出一个完备、标量曲率恒零而截面曲率同时有正有负的三维模型。
\end{solution}

\begin{exercise}\label{g3:ex:product-weyl-criterion}
	设 \((M^p,g_1)\)、\((N^q,g_2)\) 的截面曲率分别恒为 \(k\)、\(l\)，其中 \(p,q\ge2\)。在乘积上取 \(g=g_1\oplus g_2\)，记 \(n=p+q\)。
	\begin{enumerate}
		\item 求 \(\operatorname{Ric}_g=\lambda g\) 对 \(k,l\) 的充要条件。
		\item 证明 \(W_g=0\) 当且仅当 \(k+l=0\)。若同时满足第一问的条件，能得到什么结论？
	\end{enumerate}
\end{exercise}
\begin{solution}
	1. 由乘积曲率公式，两因子切方向上的 Ricci 分别为 \((p-1)k g_1\)、\((q-1)l g_2\)，混合分量为零。因此所求条件是
	\[
		(p-1)k=(q-1)l=\lambda.
	\]

	2. \(n\ge4\)，可使用 \cref{g3:thm:weyl-decomposition}。记标量曲率为
	\[
		s=p(p-1)k+q(q-1)l,
	\]
	则 \(A=(\operatorname{Ric}-s g/(2(n-1)))/(n-2)\) 在两个因子上分别是 \(a g_1\)、\(c g_2\)，其中
	\[
		a=\frac{(p-1)k-s/(2(n-1))}{n-2},\qquad
		c=\frac{(q-1)l-s/(2(n-1))}{n-2}.
	\]
	在第一因子的单位二维截面、第二因子的单位二维截面及混合单位截面上，\(g\wedge A\) 的值分别是 \(2a,2c,a+c\)。若 \(W=0\)，与乘积曲率比较便有
	\[
		k=2a,\qquad l=2c,\qquad a+c=0,
	\]
	从而 \(k+l=0\)。条件 \(p,q\ge2\) 保证前两个截面确实存在。

	反过来，若 \(l=-k\)，则
	\[
		s=(p-q)(n-1)k,\qquad a=\frac k2,\quad c=-\frac k2.
	\]
	在适合乘积分解的正交单位基中，\(\mathrm{Rm}\) 和 \(g\wedge A\) 的纯因子分量分别相等，所有混合分量也同时为零。因此四阶张量相等，\(W=0\)。若还要求第一问成立，则 \((p-1)k=-(q-1)k\)，即 \((n-2)k=0\)。故 \(k=l=0\)，乘积平坦。
\end{solution}

\begin{exercise}\label{g3:ex:conformal-scalar-sign}
	设 \((M^n,g)\) 紧、连通、无边界，\(n\ge3\)，且 \(\operatorname{Scal}_g=0\)。对正光滑函数 \(v\) 取 \(\widetilde g=v^{4/(n-2)}g\)。证明：若 \(\operatorname{Scal}_{\widetilde g}\) 处处非负或处处非正，则 \(v\) 为常数；因此任意非常数的这种共形变化都产生既取正值又取负值的标量曲率。
\end{exercise}
\begin{solution}
	由式\eqref{g3:eq:conformal-scalar-power}，
	\[
		\frac{4(n-1)}{n-2}\Delta_gv
		=\operatorname{Scal}_{\widetilde g}\,v^{(n+2)/(n-2)}.
	\]
	散度定理 \cref{g3:thm:divergence-theorem} 给出 \(\int_M\Delta_gv\,\mu_g=0\)。右端的权 \(v^{(n+2)/(n-2)}\) 严格为正；若标量曲率处处具有同一弱符号，其积分为零便迫使它处处为零。于是 \(\Delta_gv=0\)。再由第一章的能量公式，
	\[
		0=\int_Mv\Delta_gv\,\mu_g
		=\int_M|\operatorname{grad}_gv|_g^2\mu_g.
	\]
	所以 \(\mathrm dv=0\)，连通性使 \(v\) 为常数。若 \(v\) 非常数，标量曲率既不可能处处非负，也不可能处处非正，因而必有正值点和负值点。紧性和无边界条件在两次积分中都已使用。
\end{solution}

\begin{exercise}\label{g3:ex:conformal-halfspace-scalar-flat}
	在半空间 \(\{(x^1,\ldots,x^{n-1},y):y>0\}\subset\mathbb R^n\) 上，设 \(n\ge3\)，考虑
	\[
		\widetilde g=y^{4/(n-2)}\left(\sum_{\alpha=1}^{n-1}(\mathrm dx^\alpha)^2+\mathrm dy^2\right).
	\]
	计算其标量曲率、协变 Ricci 张量，以及径向和纯水平截面曲率。证明此度量共形平坦但不平坦，并判断它是否完备。
\end{exercise}
\begin{solution}
	记括号内的 Euclid 度量为 \(g_0\)，令 \(\alpha=2/(n-2)\)、\(u=\alpha\log y\)，则 \(\widetilde g=e^{2u}g_0\)。由于 \(v=y\) 满足 \(\Delta_{g_0}v=0\)，幂次形式的共形公式立即给出 \(\operatorname{Scal}_{\widetilde g}=0\)。又有
	\[
		\mathrm du=\frac\alpha y\mathrm dy,\qquad
		H_u=-\frac\alpha{y^2}\mathrm dy\otimes\mathrm dy,\qquad
		\Delta_{g_0}u=\frac\alpha{y^2}.
	\]
	代入 Ricci 变换公式并使用 \((n-2)\alpha=2\)，得到
	\[
		\widetilde{\operatorname{Ric}}
		=\frac{2}{(n-2)y^2}
		\bigl(n\,\mathrm dy\otimes\mathrm dy-g_0\bigr).
	\]
	这是坐标基中的协变张量，取迹时须使用 \(\widetilde g^{-1}=y^{-2\alpha}g_0^{-1}\)，不能直接把坐标系数当成曲率特征值。

	进一步，\(B_u=H_u-\mathrm du\otimes\mathrm du+\tfrac12|\mathrm du|^2g_0\) 在 Euclid 单位水平向量和 \(\partial_y\) 上的值分别为
	\[
		\frac{\alpha^2}{2y^2},\qquad
		-\frac{\alpha+\alpha^2/2}{y^2}.
	\]
	共形四阶公式给出：对一对 Euclid 正交单位向量，将它们再乘 \(e^{-u}\) 单位化后，新截面曲率是 \(-e^{-2u}\) 乘以这两个 \(B_u\) 对角值之和。因此
	\[
		K_{\mathrm{rad}}=\frac{2}{n-2}\,y^{-2n/(n-2)},\qquad
		K_{\mathrm{hor}}=-\frac{4}{(n-2)^2}\,y^{-2n/(n-2)}.
	\]
	这里径向截面包含 \(\partial_y\)，纯水平截面位于 \(y\) 的等值超平面；两类截面都存在。曲率非零且异号，而该度量按定义与平直度量共形，所以共形平坦。

	它并不完备。固定所有 \(x^\alpha=0\)，沿竖直方向向 \(y=0\) 行进的总长度为
	\[
		\int_0^1 y^{2/(n-2)}\,\mathrm dy=\frac{n-2}{n}<\infty.
	\]
	尤其点列 \((0,1/j)\) 的尾段距离被这段积分的尾部控制，因而是 Cauchy 序列。若它在流形中收敛，距离拓扑与流形拓扑一致便使其 \(y\) 坐标收敛到一个正数，与 \(1/j\to0\) 矛盾。故此例还说明，局部曲率和共形公式本身不保证完备性。
\end{solution}

\begin{exercise}\label{g3:ex:normalized-total-scalar-variation}
	设 \((M^n,g)\) 紧、连通、无边界，\(n\ge3\)。记
	\[
		V(g)=\int_M\mu_g,\qquad
		I(g)=\int_M\operatorname{Scal}_g\,\mu_g,\qquad
		\mathcal F(g)=V(g)^{-(n-2)/n}I(g).
	\]
	对 \(g_t=e^{2tu}g\)、\(u\in C^\infty(M)\)，计算 \(\left.\frac{\mathrm d}{\mathrm dt}\right|_{t=0}\mathcal F(g_t)\)。证明它对每个 \(u\) 都为零，当且仅当 \(g\) 的标量曲率为常数。
\end{exercise}
\begin{solution}
	记 \(s=\operatorname{Scal}_g\)。由共形标量曲率公式及密度的行列式公式，
	\[
		\left.\frac{\mathrm d}{\mathrm dt}\right|_0s_{g_t}
		=2(n-1)\Delta_gu-2us,\qquad
		\left.\frac{\mathrm d}{\mathrm dt}\right|_0\mu_{g_t}=nu\mu_g.
	\]
	紧性允许在积分号下求导，散度定理又使 Laplace 项的积分为零，故
	\[
		V'(0)=n\int_Mu\mu_g,\qquad
		I'(0)=(n-2)\int_Mus\mu_g.
	\]
	令 \(\bar s=I(g)/V(g)\)，乘积求导得到
	\[
		\mathcal F'(0)
		=(n-2)V(g)^{-(n-2)/n}
		\int_Mu(s-\bar s)\mu_g.
	\]
	若 \(s\) 为常数，则 \(s=\bar s\)，所有这些导数都为零。反过来，若对每个光滑 \(u\) 都为零，可取 \(u=s-\bar s\)。由于前面的系数为正，便有
	\[
		\int_M(s-\bar s)^2\mu_g=0,
	\]
	从而 \(s=\bar s\)。这一判据把常标量曲率条件表示成归一化总标量曲率沿全部共形方向的驻定条件。
\end{solution}

\begin{exercise}\label{g3:ex:ricci-invisible-local-curvature}
	在四维 Euclid 空间的正交单位基 \((e_1,e_2,e_3,e_4)\) 中，定义二向量空间上的对角自伴算子，其在基 \((e_i\wedge e_j)_{i<j}\) 上的特征值为
	\[
		k_{12}=1,\quad k_{13}=-1,\quad k_{14}=0,\quad
		k_{23}=0,\quad k_{24}=-1,\quad k_{34}=1.
	\]
	证明它对应一个非零的代数曲率张量 \(A\)，但 \(\operatorname{Ric}_A=0\)。将它实现为某个局部黎曼度量在原点的曲率，并说明该度量在原点的法坐标体积密度和小球体积为何在本章所求的首个曲率修正阶数上都与 Euclid 情形相同。
\end{exercise}
\begin{solution}
	令 \(A_{ijij}=k_{ij}\) 并按两对反对称性及成对交换对称性延拓，其余独立分量取零。这给出题中算子对应的四阶张量。第一 Bianchi 恒等式在四个指标互不相同时的各项都为零；有重复指标时，各项则由反对称性两两抵消。因此它确是代数曲率张量。

	Ricci 的非对角分量全为零，对角分量是每个方向相邻的三个截面值之和：
	\[
		\begin{aligned}
		(\operatorname{Ric}_A)_{11}&=1-1+0=0,&
		(\operatorname{Ric}_A)_{22}&=1+0-1=0,\\
		(\operatorname{Ric}_A)_{33}&=-1+0+1=0,&
		(\operatorname{Ric}_A)_{44}&=0-1+1=0.
		\end{aligned}
	\]
	但是 \(A_{1212}=1\)，所以 \(A\ne0\)。由 \cref{g3:prop:curvature-local-realization}，原点邻域中的正定度量
	\[
		g_{ij}(x)=\delta_{ij}-\frac13A_{ikjl}x^kx^l
	\]
	在原点的曲率恰是 \(A\)，因而那里 \(\operatorname{Ric}=0\)、\(\operatorname{Scal}=0\)，但 \(K(e_1,e_2)=1\)、\(K(e_1,e_3)=-1\)。

	对这个度量重新取原点的法坐标，由体积展开定理可得
	\[
		J_0(x)=1+O(|x|^3),\qquad
		\operatorname{Vol}_g(B_g(0,r))
		=\omega_4r^4\bigl(1+O(r^4)\bigr).
	\]
	因此密度的二次项及小球体积相对于 \(\omega_4r^4\) 的二次修正都消失。它们分别读取 Ricci 与标量曲率，不能据此推出整个曲率张量为零。本题只比较这两个指定阶数，没有断言更高阶项或精确体积也与平直情形相同。
\end{solution}

% !TeX root = ../../Book3.tex
% 纲要标题：## 第4章　子流形、法丛与相容方程
% 纲要位置：几何学计划纲要.md，第 2226 行。
% 正文状态：初稿；尚未正式审核。

% 纲要内容开始
% 
% > 本章目标与证明范围：把第一卷曲面的两个基本形式推广到高余维。完整证明 Gauss、Codazzi、Ricci 方程及所选空间形式中的局部实现；法联络不作为附带记号。
% 
% 纲要内容结束

\chapter{子流形、法丛与相容方程}
\label{g3:ch:04}
\BookChapterNumber{g3:ch:04}

\begin{chapterabstract}
	浸入的诱导度量描述切向长度，第二基本形式和法联络描述切空间与法空间在环境中的变化。本章建立任意余维的切法分解，完整推导 Gauss、Codazzi 和 Ricci 方程，并将空间形式中的三组相容条件组织成一个平坦联络。通过局部平行标架与位置一形式的积分，构造所给数据的等距浸入，证明其在环境等距意义下的唯一性，并辨明整体延拓中的环路返回与周期障碍。球面、乘积和管状模型用于计算主曲率与内在曲率，面积一阶变分则解释平均曲率向量的几何意义。
\end{chapterabstract}

\begin{learninggoals}
	\begin{itemize}
		\item 从环境导数的切向与法向投影构造第二基本形式、形算子和法联络，核对它们的对应关系与符号。
		\item 完整证明三组相容方程，说明高余维中法曲率方程所增加的独立条件。
		\item 从相容数据构造 Euclid、球面及双曲模型中的局部浸入，解释唯一性的完整数据条件及整体单值障碍。
		\item 计算主曲率、管状体积和面积的一阶变化，将极小性与平均曲率、位置向量的 Laplace 公式相联系。
	\end{itemize}
\end{learninggoals}

把一个平面条带卷成圆柱，可以保持条带内的长度和夹角，却改变它在三维空间中的形状。平面与圆柱的 Gauss 曲率同为零；沿圆柱绕轴方向移动时，切平面不断转动，而沿平面移动时没有这种变化。内在曲率已经说明度量自身如何偏离平直，但要恢复浸入，还需要记录环境方向中的变化。

\ProseParagraphBreak

曲面在三维空间中只有一条法线方向，形算子便足以把法向变化写成切向量。高余维时，法空间本身也可以在各法向方向之间转动。它的导数和曲率分别由法联络及法曲率记录；第二基本形式将切方向与法方向联系起来。把环境曲率按切、法两部分展开，会得到必须同时满足的 Gauss、Codazzi 与 Ricci 方程。\footnote{Delfino Codazzi（1824--1873），意大利数学家，研究曲面理论与曲线坐标。}反过来，这些方程是否足以产生一个实际浸入，是本章的主要存在问题。

\ProseParagraphBreak

本章使用前三章的 Levi-Civita 联络、曲率张量、指数映射及散度定理，并回顾向量丛局部标架所需的记号。局部实现使用光滑线性常微分方程的初值唯一性与参数依赖；相关结论在第一卷第13.1节已证明，本章给出将它们用于多变量相容系统的步骤。第一卷曲面理论中的两个基本形式将在任意余维重新建立。各几何对象均光滑；边界、紧性、单连通性以及法向的整体可选性按具体结果说明。

\ProseParagraphBreak

本章固定 \(B(X,Y)=(\bar\nabla_X\mathrm df(Y))^\perp\)、\(g(A_\xi X,Y)=\langle B(X,Y),\xi\rangle\)，并取平均曲率向量 \(H=n^{-1}\operatorname{tr}_g B\)。因此球面的外法向形算子为负，而球面的内在截面曲率为正。曲率方向及四阶张量次序沿用第三章。球面与球面乘积的内在度量模型可对照 Petersen 第3版第4.2.1--4.2.2节\cite[第4.2.1--4.2.2节]{Petersen2016}；本章进一步计算它们在不同环境中的法向数据。

% !TeX root = ../../Book3.tex
% 纲要标题：### 4.1 切向与法向分解
% 纲要位置：几何学计划纲要.md，第 2230 行。
% 正文状态：初稿；尚未正式审核。

% 纲要内容开始
% 
% * 第二基本形式和形算子
% * 法联络及其曲率
% * 平均曲率向量
% * 全测地和全脐子流形
% 
% 纲要内容结束

\section{切向与法向分解}
\label{g3:sec:0401}

沿子流形移动一个切向量，环境中的导数可能离开切空间。第一章将导数的切向部分识别为诱导度量的 Levi-Civita 联络，现在保留下来的法向部分记录浸入怎样弯曲。法向量自身也会变化：它的切向导数由同一弯曲数据决定，法向导数则构成法丛上的联络。余维大于一时，法向空间中还有独立的转动，必须把这部分一起记录。

\subsection{第二基本形式和形算子}
\label{g3:subsec:0401:01}

设 \(f:(M^n,g)\to(\bar M^{n+r},\bar g)\) 为光滑等距浸入，其中 \(n\ge1\)，两流形均无边界，且 \(g=f^*\bar g\)。沿 \(f\) 的向量场，是给每个源点 \(p\) 指定一个向量 \(W(p)\in\mathrm{T}_{f(p)}\bar M\) 的光滑截面。它们所在的丛记为 \(f^*\mathrm{T}\bar M\)。即使两个不同源点有同一个像，这两个源点上的纤维仍分别记录。

\begin{definition}\label{g3:def:normal-bundle}
	设 \(f:(M^n,g)\to(\bar M^{n+r},\bar g)\) 为光滑等距浸入。其\term[facong]{法丛}[normal bundle]为
	\[
		N_fM\coloneqq
		\{(p,\xi):p\in M,\ \xi\in\mathrm{T}_{f(p)}\bar M,\
			\bar g(\xi,\mathrm df_p(v))=0
			\text{ 对所有 }v\in\mathrm{T}_pM\}.
	\]
	其截面称为法向量场。正交分解记为
	\[
		f^*\mathrm{T}\bar M=\mathrm df(\mathrm{T}M)\oplus N_fM,
		\qquad W=W^\top+W^\perp.
	\]
\end{definition}
\symidx{\(N_fM\)}

这个分解随点光滑。事实上，在源坐标中令 \(f_i=\mathrm df(\partial_i)\)，则切向投影为
\[
	W^\top=\sum_{i,j}g^{ij}\bar g(W,f_i)f_j.
\]
矩阵 \((g^{ij})\) 光滑，故两个投影均光滑，法向投影具有常秩 \(r\)。局部选择线性无关的投影像，再作光滑正交化，就得到法丛的局部正交单位标架。

\ProseParagraphBreak

沿 \(f\) 的环境协变导数采用\cref{g3:prop:induced-connection-projection}中的定义。若 \(f^a=y^a\circ f\)、\(W=W^a\partial_{y^a}|_f\)，则
\[
	\bar\nabla^f_XW
	=\left(X(W^a)+\bar\Gamma^a_{bc}(f)\,X(f^b)W^c\right)
		\partial_{y^a}|_f.
\]
它只使用沿 \(f\) 的数据。以下将 \(\bar\nabla^f_X\) 简写为 \(\bar\nabla_X\)，其中 \(X\) 始终是源流形上的向量场。这一联络满足环境度量的乘积规则，而上述命题给出
\[
	(\bar\nabla_X\mathrm df(Y))^\top=\mathrm df(\nabla_XY).
\]

\begin{definition}\label{g3:def:second-fundamental-form}
	设 \(f:(M^n,g)\to(\bar M,\bar g)\) 为光滑等距浸入。定义它的\term[dierjibenxingshi]{第二基本形式}[second fundamental form]为法丛值的二阶协变张量
	\[
		B(X,Y)\coloneqq(\bar\nabla_X\mathrm df(Y))^\perp.
	\]
\end{definition}
\symidx{\(B\)}

\begin{proposition}\label{g3:prop:gauss-decomposition}
	光滑等距浸入 \(f:(M,g)\to(\bar M,\bar g)\) 的第二基本形式光滑、对称，并对两个输入都为 \(C^\infty(M)\) 线性。切向与法向分解给出
	\begin{equation}\label{g3:eq:gauss-decomposition}
		\bar\nabla_X\mathrm df(Y)=\mathrm df(\nabla_XY)+B(X,Y).
	\end{equation}
\end{proposition}
\begin{proof}
	对光滑函数 \(h\)，第一输入的函数线性直接来自联络；第二输入则满足
	\[
		\bar\nabla_X\mathrm df(hY)
		=X(h)\mathrm df(Y)+h\bar\nabla_X\mathrm df(Y).
	\]
	第一项切向，取法向投影后为零，故 \(B(X,hY)=hB(X,Y)\)。各投影和联络系数光滑，因而 \(B\) 光滑。

	环境联络的无挠性沿浸入表现为
	\[
		\bar\nabla_X\mathrm df(Y)-\bar\nabla_Y\mathrm df(X)
		=\mathrm df([X,Y]).
	\]
	可由上面的坐标公式直接核对：含 \(\bar\Gamma^a_{bc}\) 的两项因下指标对称而相消，余下分量为 \(X(Yf^a)-Y(Xf^a)=[X,Y]f^a\)。右侧切向，取法向投影即得 \(B(X,Y)=B(Y,X)\)。最后将环境导数分成两个正交分量，切向部分由诱导联络公式给出，得到式\eqref{g3:eq:gauss-decomposition}。
\end{proof}

式\eqref{g3:eq:gauss-decomposition}称为 Gauss 公式。它把内在联络与浸入的额外二阶数据分开。反过来，对法向场求导，也会出现切向分量；由于切向场与法向场的内积恒为零，这个分量已由 \(B\) 决定。

\begin{definition}\label{g3:def:shape-operator}
	设 \(f:(M,g)\to(\bar M,\bar g)\) 为光滑等距浸入。对法向量 \(\xi\in N_{f,p}M\)，定义\term[xingsuanzi]{形算子}[shape operator]
	\(A_\xi:\mathrm{T}_pM\to\mathrm{T}_pM\)，使
	\begin{equation}\label{g3:eq:shape-operator-pairing}
		g(A_\xi X,Y)=\bar g(B(X,Y),\xi)
		\qquad(X,Y\in\mathrm{T}_pM).
	\end{equation}
	正定内积保证这个线性算子存在且唯一。
\end{definition}
\symidx{\(A_\xi\)}

\begin{proposition}\label{g3:prop:weingarten-decomposition}
	设 \(f:(M,g)\to(\bar M,\bar g)\) 为光滑等距浸入。形算子随 \(\xi\) 线性变化，且每个 \(A_\xi\) 都自伴。对光滑法向场 \(\xi\)，令
	\(\nabla^\perp_X\xi=(\bar\nabla_X\xi)^\perp\)，则
	\begin{equation}\label{g3:eq:weingarten-formula}
		\bar\nabla_X\xi=-\mathrm df(A_\xi X)+\nabla^\perp_X\xi.
	\end{equation}
\end{proposition}
\begin{proof}
	关于 \(\xi\) 的线性由式\eqref{g3:eq:shape-operator-pairing}及唯一性得到。在坐标中，\((A_\xi)^i{}_j=g^{ik}\bar g(B(\partial_j,\partial_k),\xi)\)，故对光滑法向场，形算子场也光滑。由 \(B\) 对称，
	\[
		g(A_\xi X,Y)=\bar g(B(X,Y),\xi)
		=\bar g(B(Y,X),\xi)=g(X,A_\xi Y),
	\]
	所以 \(A_\xi\) 自伴。对恒等式 \(\bar g(\xi,\mathrm df(Y))=0\) 沿 \(X\) 求导，得到
	\[
		\bar g(\bar\nabla_X\xi,\mathrm df(Y))
		=-\bar g(\xi,\bar\nabla_X\mathrm df(Y))
		=-g(A_\xi X,Y).
	\]
	这对所有 \(Y\) 成立，故切向分量是 \(-\mathrm df(A_\xi X)\)；余下分量正是所定义的 \(\nabla^\perp_X\xi\)。
\end{proof}

式\eqref{g3:eq:weingarten-formula}称为\term[weingartengongshi]{Weingarten 公式}[Weingarten formula]。\footnote{Julius Weingarten（1836--1910），德国数学家，研究曲面等距变形与曲率线。}其负号来自正交关系的求导。若余维为一，取局部单位法向 \(\nu\)，则 \(\bar\nabla_X\nu\) 与 \(\nu\) 正交，已没有法向分量，所以 \(\bar\nabla_X\nu=-\mathrm df(A_\nu X)\)。

\begin{example}\label{g3:exm:sphere-extrinsic-sign}
	对半径 \(a>0\) 的圆球面 \(S^n(a)\subset\mathbb R^{n+1}\)，取外向单位法向 \(\nu(p)=p/a\)。环境联络为普通导数，故对切向量 \(X\)，
	\[
		\mathrm D_X\nu=\frac1aX,\qquad
		A_\nu=-\frac1a\operatorname{id},\qquad
		B(X,Y)=-\frac1a g(X,Y)\nu.
	\]
	这些等式直接来自法向的变化，与第一章球面诱导联络的符号相容。把 \(\nu\) 改为 \(-\nu\) 会改变 \(A_\nu\) 的符号，但不改变法丛值张量 \(B\)。
\end{example}

\subsection{法联络及其曲率}
\label{g3:subsec:0401:02}

法向导数中的切向部分已经由 \(B\) 决定，余下部分比较的是相邻法向空间中的向量。余维为一时，单位法向的这部分导数为零；余维较高时，保持长度的法向量还可以在法向空间内部转动。

\begin{definition}\label{g3:def:normal-connection}
	设 \(f:(M,g)\to(\bar M,\bar g)\) 为光滑等距浸入。法丛上的\term[falianluo]{法联络}[normal connection]定义为
	\[
		\nabla^\perp_X\xi=(\bar\nabla_X\xi)^\perp,
		\qquad \xi\in\Gamma(N_fM).
	\]
	法向场 \(\xi\) 若满足 \(\nabla^\perp_X\xi=0\) 对所有 \(X\) 成立，则称为法平行的。
\end{definition}
\symidx{\(\nabla^\perp\)}

\begin{proposition}\label{g3:prop:normal-connection-metric}
	光滑等距浸入的 \(\nabla^\perp\) 是法丛上的光滑联络，并与法丛上的诱导内积相容，即
	\begin{equation}\label{g3:eq:normal-connection-metric}
		X\bar g(\xi,\zeta)
		=\bar g(\nabla^\perp_X\xi,\zeta)
			+\bar g(\xi,\nabla^\perp_X\zeta).
	\end{equation}
\end{proposition}
\begin{proof}
	正交投影为光滑丛映射，环境联络对求导方向为函数线性，所以
	\[
		\nabla^\perp_{hX}\xi=h\nabla^\perp_X\xi,\qquad
		\nabla^\perp_X(h\xi)=X(h)\xi+h\nabla^\perp_X\xi.
	\]
	与实线性一起，这正是丛联络的各项公理。环境度量相容性给出
	\[
		X\bar g(\xi,\zeta)
		=\bar g(\bar\nabla_X\xi,\zeta)+\bar g(\xi,\bar\nabla_X\zeta).
	\]
	这里 \(\xi,\zeta\) 都法向，另一个输入的切向分量不贡献内积，因此可将两个环境导数都替换为其法向投影，得到式\eqref{g3:eq:normal-connection-metric}。
\end{proof}

\begin{definition}\label{g3:def:normal-curvature}
	光滑等距浸入的\term[faqulv]{法曲率}[normal curvature]为法丛联络的曲率
	\begin{equation}\label{g3:eq:normal-curvature-definition}
		R^\perp(X,Y)\xi
		\coloneqq\nabla^\perp_X\nabla^\perp_Y\xi
			-\nabla^\perp_Y\nabla^\perp_X\xi
			-\nabla^\perp_{[X,Y]}\xi.
	\end{equation}
	若 \(R^\perp=0\)，称法联络平坦。
\end{definition}
\symidx{\(R^\perp\)}

\begin{proposition}\label{g3:prop:normal-connection-curvature}
	设 \(f:(M,g)\to(\bar M,\bar g)\) 为光滑等距浸入。\(R^\perp(X,Y)\xi\) 对 \(X,Y,\xi\) 均为函数线性，对 \(X,Y\) 反对称，且在法向输入上满足
	\begin{equation}\label{g3:eq:normal-curvature-skew}
		\bar g(R^\perp(X,Y)\xi,\zeta)
		=-\bar g(\xi,R^\perp(X,Y)\zeta).
	\end{equation}
	因而 \(R^\perp\) 是取值于法丛斜自伴算子的二形式。
\end{proposition}
\begin{proof}
	反对称性来自定义。对光滑函数 \(h\)，利用
	\([hX,Y]=h[X,Y]-Y(h)X\)，展开第一输入得
	\[
		\begin{aligned}
			R^\perp(hX,Y)\xi
			={}&h\nabla^\perp_X\nabla^\perp_Y\xi
				-\nabla^\perp_Y(h\nabla^\perp_X\xi)\\
			&-h\nabla^\perp_{[X,Y]}\xi
				+Y(h)\nabla^\perp_X\xi
			=hR^\perp(X,Y)\xi.
		\end{aligned}
	\]
	第二输入的函数线性由反对称性得到。对于法向输入，
	\[
		\nabla^\perp_X\nabla^\perp_Y(h\xi)
		=X(Yh)\xi+(Yh)\nabla^\perp_X\xi
			+(Xh)\nabla^\perp_Y\xi
			+h\nabla^\perp_X\nabla^\perp_Y\xi.
	\]
	交换 \(X,Y\) 并减去括号方向的导数后，函数的二阶导数合成为
	\(XYh-YXh-[X,Y]h=0\)，两个交叉项也抵消，因此第三输入函数线性。局部标架展开说明它只取决于各输入在一点的值，并且光滑。

	最后将度量相容式连续使用两次：
	\[
		\begin{aligned}
			XY\bar g(\xi,\zeta)
			={}&\bar g(\nabla^\perp_X\nabla^\perp_Y\xi,\zeta)
				+\bar g(\nabla^\perp_Y\xi,\nabla^\perp_X\zeta)\\
			&+\bar g(\nabla^\perp_X\xi,\nabla^\perp_Y\zeta)
				+\bar g(\xi,\nabla^\perp_X\nabla^\perp_Y\zeta).
		\end{aligned}
	\]
	交换 \(X,Y\) 后中间两项相消，再减去 \([X,Y]\bar g(\xi,\zeta)\)，左侧为零，右侧为式\eqref{g3:eq:normal-curvature-skew}两边移到一起的和。
\end{proof}

在局部正交单位法标架 \(\nu=(\nu_1,\ldots,\nu_r)\) 中，定义一形式矩阵 \(\eta=(\eta^\alpha{}_\beta)\)，使
\begin{equation}\label{g3:eq:normal-connection-matrix}
	\nabla^\perp\nu_\beta=\sum_\alpha\nu_\alpha\eta^\alpha{}_\beta,
	\qquad \eta^\alpha{}_\beta(X)=\bar g(\nabla^\perp_X\nu_\beta,\nu_\alpha).
\end{equation}
度量相容性给出 \(\eta^{\mathsf T}=-\eta\)。若法向场的列分量为 \(u\)，即 \(\xi=\nu u\)，则
\[
	\nabla^\perp_X\xi=\nu\bigl(Xu+\eta(X)u\bigr).
\]
因此平行性是关于各方向导数的方程 \(Xu=-\eta(X)u\)。

\begin{proposition}\label{g3:prop:normal-curvature-matrix}
	光滑等距浸入在局部正交单位法标架中的法曲率矩阵为
	\begin{equation}\label{g3:eq:normal-curvature-matrix}
		F^\perp=\mathrm d\eta+\eta\wedge\eta,\qquad
		R^\perp(X,Y)\nu_\beta
			=\sum_\alpha\nu_\alpha(F^\perp)^\alpha{}_\beta(X,Y).
	\end{equation}
	这里 \((\eta\wedge\eta)^\alpha{}_\beta
	=\sum_\gamma\eta^\alpha{}_\gamma\wedge\eta^\gamma{}_\beta\)。
\end{proposition}
\begin{proof}
	将式\eqref{g3:eq:normal-connection-matrix}代入曲率定义，\(\nu_\alpha\) 的系数为
	\[
		\begin{aligned}
			&X(\eta^\alpha{}_\beta(Y))
				-Y(\eta^\alpha{}_\beta(X))
				-\eta^\alpha{}_\beta([X,Y])\\
			&\qquad+\sum_\gamma
				\bigl(\eta^\alpha{}_\gamma(X)\eta^\gamma{}_\beta(Y)
				-\eta^\alpha{}_\gamma(Y)\eta^\gamma{}_\beta(X)\bigr).
		\end{aligned}
	\]
	第一行是 \(\mathrm d\eta^\alpha{}_\beta(X,Y)\)，第二行是矩阵外积的对应分量，故得到所述公式。
\end{proof}

局部正交单位法标架只要求每点的内积矩阵为单位矩阵。若进一步要求它在一个开邻域内法平行，就要求 \(\eta\) 在整个邻域为零，由式\eqref{g3:eq:normal-curvature-matrix}必有 \(R^\perp=0\)。下面的直接计算说明，这一要求对一般浸入不成立。

\begin{example}\label{g3:exm:nonflat-normal-connection}
	在 \(\mathbb R^4\) 中考虑图像浸入
	\[
		f(x,y)=\left(x,y,\frac{x^2}{2},xy\right).
	\]
	其坐标切向量为 \(f_x=(1,0,x,y)\)、\(f_y=(0,1,0,x)\)，以下两个场处处法向且线性无关：
	\[
		\xi_1=(-x,0,1,0),\qquad
		\xi_2=(-y,-x,0,1).
	\]
	它们不要求在整个邻域正交，在原点则为正交单位基。法向 Gram 矩阵及其行列式为
	\[
		G_N=
		\begin{pmatrix}1+x^2&xy\\xy&1+x^2+y^2\end{pmatrix},
		\qquad
		d_N=(1+x^2)(1+x^2+y^2)-x^2y^2>0.
	\]
	由于 \(\partial_y\xi_1=0\)，有 \(\nabla^\perp_{\partial_y}\xi_1=0\)。写
	\(\nabla^\perp_{\partial_x}\xi_1=a\xi_1+b\xi_2\)，对两法向场取内积便得
	\[
		G_N\binom{a}{b}=\binom{x}{y},\qquad
		a=\frac{x(1+x^2)}{d_N},\quad b=\frac{y}{d_N}.
	\]
	在原点 \(a=b=0\)、\(\partial_ya=0\)、\(\partial_yb=1\)，所以
	\[
		R^\perp(\partial_x,\partial_y)\xi_1\big|_{(0,0)}
		=-\nabla^\perp_{\partial_y}(a\xi_1+b\xi_2)\big|_{(0,0)}
		=-\xi_2(0,0)\ne0.
	\]
	因此任何包含原点的邻域都不存在处处法平行的法标架。环境空间本身平坦，并不使这个法联络平坦。
\end{example}

若法丛秩为一，任一局部单位法向 \(\nu\) 满足
\(\bar g(\nabla^\perp_X\nu,\nu)=0\)，而法向空间由 \(\nu\) 张成，故 \(\nabla^\perp\nu=0\)、\(R^\perp=0\)。这个局部结论不需要一个全局单位法向；是否能在整个流形选取这样的法向，是法线丛的另一项整体性质。

\subsection{平均曲率向量}
\label{g3:subsec:0401:03}

第二基本形式为每个切向方向给出一个法向弯曲向量。对正交方向求平均，应当得到不依赖所选切向基的法向量；这正是对 \(B\) 作度量迹的作用。

\begin{definition}\label{g3:def:mean-curvature-vector}
	设 \(f:(M^n,g)\to(\bar M,\bar g)\) 为光滑等距浸入，\(n\ge1\)。其\term[pingjunqulvxiangliang]{平均曲率向量}[mean curvature vector]为
	\begin{equation}\label{g3:eq:mean-curvature-vector}
		H\coloneqq\frac1n\operatorname{tr}_gB
			=\frac1n\sum_{i=1}^nB(e_i,e_i),
	\end{equation}
	其中 \(e_1,\ldots,e_n\) 为任意正交单位基。若 \(H\) 恒为零，称 \(f\) 为\term[jixiaojinru]{极小浸入}[minimal immersion]，相应浸入子流形称为极小子流形。
\end{definition}
\symidx{\(H\)}

若 \(e'_i=\sum_j e_jQ_{ji}\)，其中 \(Q\) 为正交矩阵，则
\[
	\sum_iB(e'_i,e'_i)
	=\sum_{j,k}\left(\sum_iQ_{ji}Q_{ki}\right)B(e_j,e_k)
	=\sum_jB(e_j,e_j).
\]
这既证明定义与正交基无关，也说明 \(H\) 是光滑法向场。在一般坐标中，
\[
	H=\frac1n g^{ij}B(\partial_i,\partial_j).
\]
与任意法向量 \(\xi\) 配对，式\eqref{g3:eq:shape-operator-pairing}又给出
\begin{equation}\label{g3:eq:mean-curvature-shape-trace}
	\bar g(H,\xi)=\frac1n\operatorname{tr}A_\xi.
\end{equation}
因此 \(H=0\) 等价于每个法向形算子都无迹。

\ProseParagraphBreak

对超曲面，选择单位法向 \(\nu\) 后可写 \(H=h\nu\)，其中 \(h=\operatorname{tr}A_\nu/n\) 是带符号的平均曲率。换成 \(-\nu\) 后，\(A_\nu\) 与 \(h\) 同时变号，向量 \(H\) 不变。由\cref{g3:exm:sphere-extrinsic-sign}，外向法向的圆球面满足
\[
	h=-\frac1a,\qquad H=-\frac1a\nu.
\]
平均曲率向量指向球心；它在这里采用平均值归一化，而非未除以维数的迹。

\begin{example}\label{g3:exm:curve-mean-curvature-vector}
	设 \(\gamma:I\to\bar M\) 为单位速光滑曲线，并给 \(I\) 度量 \(\mathrm ds^2\)。记 \(T=\dot\gamma\)。一维域中的单位坐标向量满足 \(\nabla_{\partial_s}\partial_s=0\)，所以 Gauss 公式给出
	\[
		B(\partial_s,\partial_s)=\bar{\mathrm D}_sT,\qquad
		H=\bar{\mathrm D}_sT.
	\]
	单位速条件使右侧与 \(T\) 正交。于是曲线的平均曲率向量就是其曲率向量，而单位速曲线是极小浸入当且仅当它是环境测地线。
\end{example}

\subsection{全测地和全脐子流形}
\label{g3:subsec:0401:04}

两种特殊弯曲情形可以直接用 \(B\) 描述。一种是所有方向都没有法向弯曲；另一种是所有单位切向方向的法向弯曲完全相同。前者对应测地线的保持，后者对应每个法向形算子都是数量算子。

\begin{definition}\label{g3:def:totally-geodesic-umbilical}
	设 \(f:(M^n,g)\to(\bar M,\bar g)\) 为光滑等距浸入，\(n\ge1\)。若 \(B\equiv0\)，称 \(f\) 为\term[quancedijinru]{全测地浸入}[totally geodesic immersion]。若每点 \(p\) 存在法向量 \(\zeta_p\)，使
	\[
		B_p(v,v)=g_p(v,v)\zeta_p\qquad
		\text{对所有 }v\in\mathrm{T}_pM,
	\]
	则称 \(f\) 为\term[quanqijinru]{全脐浸入}[totally umbilical immersion]。
\end{definition}

\begin{proposition}\label{g3:prop:totally-geodesic-characterization}
	设 \(f:(M,g)\to(\bar M,\bar g)\) 为光滑等距浸入。则 \(f\) 全测地，当且仅当每条 \(M\) 中的测地线经 \(f\) 映射后仍是保持原参数的环境测地线。此时，对每个 \(p\) 及足够小的 \(v\in\mathrm{T}_pM\)，有
	\begin{equation}\label{g3:eq:totally-geodesic-exponential}
		f(\exp_p^M v)=\exp_{f(p)}^{\bar M}(\mathrm df_pv).
	\end{equation}
\end{proposition}
\begin{proof}
	对任意光滑曲线 \(\gamma\)，将 Gauss 公式沿曲线使用，得到加速度分解
	\begin{equation}\label{g3:eq:immersion-curve-acceleration}
		\bar{\mathrm D}_t(\mathrm df(\dot\gamma))
		=\mathrm df(\mathrm D_t\dot\gamma)+B(\dot\gamma,\dot\gamma).
	\end{equation}
	具体地，在坐标中写 \(\dot\gamma=\dot x^i\partial_i\)，左侧为
	\(\ddot x^i f_i+\dot x^i\dot x^j\bar\nabla_{\partial_i}f_j\)；
	代入式\eqref{g3:eq:gauss-decomposition}，切向部分合成为
	\(\mathrm df((\ddot x^k+\Gamma^k_{ij}\dot x^i\dot x^j)\partial_k)\)，
	法向部分即上式所列的 \(B\)。

	若 \(B=0\)，当 \(\gamma\) 为测地线时，两项都为零，故 \(f\circ\gamma\) 是环境测地线。反过来，固定 \(p\) 及任意 \(v\in\mathrm{T}_pM\)。由\cref{g3:thm:geodesic-ivp}，存在以 \((p,v)\) 为初值的局部测地线。若其像也是环境测地线，在时间零代入加速度分解，得到 \(B_p(v,v)=0\)。由 \(B\) 对称，
	\[
		2B_p(v,w)=B_p(v+w,v+w)-B_p(v,v)-B_p(w,w)=0.
	\]
	所以 \(B=0\)。

	最后，在全测地情形，\(f\circ\gamma_{p,v}\) 与环境中以
	\((f(p),\mathrm df_pv)\) 为初值的测地线相同，这由同一初值唯一性得到。当 \(v\) 足够小，使两条曲线均定义到时间一时，比较时间一的值即得式\eqref{g3:eq:totally-geodesic-exponential}。
\end{proof}

\begin{proposition}\label{g3:prop:umbilical-characterization}
	设 \(f:(M^n,g)\to(\bar M,\bar g)\) 为光滑等距浸入，\(n\ge1\)。它全脐，当且仅当
	\begin{equation}\label{g3:eq:umbilical-second-fundamental-form}
		B(X,Y)=g(X,Y)H.
	\end{equation}
	这又等价于对每个法向量 \(\xi\)，有
	\[
		A_\xi=\bar g(H,\xi)\operatorname{id}.
	\]
	全脐浸入是全测地的，当且仅当其平均曲率向量为零。
\end{proposition}
\begin{proof}
	若 \(B_p(v,v)=g_p(v,v)\zeta_p\)，对 \(v+w\) 作同样计算，再减去 \(v,w\) 两式，得到
	\(B_p(v,w)=g_p(v,w)\zeta_p\)。在正交单位基中取迹，便有
	\(nH_p=n\zeta_p\)，所以 \(\zeta_p=H_p\)，得到式\eqref{g3:eq:umbilical-second-fundamental-form}。反向蕴含由令 \(Y=X\) 得到。

	将该式与任意 \(\xi\) 配对，形算子的定义给出
	\[
		g(A_\xi X,Y)=g(X,Y)\bar g(H,\xi),
	\]
	由度量非退化性得到数量算子表达。反过来，若所有 \(A_\xi\) 都具有这一形式，则 \(B(X,Y)-g(X,Y)H\) 与每个法向量正交；它本身法向，故为零。最后在全脐公式中令 \(H=0\) 即得 \(B=0\)，而 \(B=0\) 总使其迹 \(H\) 为零。
\end{proof}

Euclid 空间中的仿射子空间有常值切空间，环境导数的法向分量恒为零，因此全测地。圆球面由\cref{g3:exm:sphere-extrinsic-sign}满足 \(B=gH\)，故全脐，但它的 \(H\ne0\)，所以并非全测地。维数一时，每个单位切向量只差一个符号，任何曲线浸入都自动全脐；此时全测地与极小的条件则都要求曲率向量为零。
% 纲要条目排为编号小节标题，并进入目录与书签；说明段仍保留于纲要注释中。

% !TeX root = ../../Book3.tex
% 纲要标题：### 4.2 三组相容方程
% 纲要位置：几何学计划纲要.md，第 2134 行。
% 正文状态：未撰写；本文件目前是框架。

% 纲要内容开始
% 
% * Gauss 方程
% * Codazzi 方程
% * Ricci 法曲率方程
% * 高余维中不能删除的独立条件
% 
% 纲要内容结束

\section{三组相容方程}
\label{g3:sec:0402}

待撰写

\subsection{Gauss 方程}
\label{g3:subsec:0402:01}
待撰写

\subsection{Codazzi 方程}
\label{g3:subsec:0402:02}
待撰写

\subsection{Ricci 法曲率方程}
\label{g3:subsec:0402:03}
待撰写

\subsection{高余维中不能删除的独立条件}
\label{g3:subsec:0402:04}
待撰写
% 纲要条目排为编号小节标题，并进入目录与书签；说明段仍保留于纲要注释中。

% !TeX root = ../../Book3.tex
% 纲要标题：### 4.3 局部实现的结构方法
% 纲要位置：几何学计划纲要.md，第 2244 行。
% 正文状态：初稿；尚未正式审核。

% 纲要内容开始
% 
% * 给定度量、法丛及对称型
% * 标架系统及相容性
% * 空间形式中的局部基本定理
% * 唯一性及整体单值障碍
% 
% 纲要内容结束

\section{局部实现的结构方法}
\label{g3:sec:0403}

前三组方程是浸入所必须满足的条件，现在把问题反过来：若只给出长度、法方向及第二基本形式，能否把这些数据放进一个实际的环境空间？困难不在某一个点，而在相邻点的标架怎样同时移动。沿一条曲线可以解常微分方程；沿两个方向先后移动时，结果相容的条件正是 Gauss、Codazzi 与 Ricci 方程。下面将它们合成一个平坦联络，再由平行标架恢复位置。

\subsection{给定度量、法丛及对称型}
\label{g3:subsec:0403:01}

设 \(M^n\) 为无边界光滑流形，\(n\geq1\)。尚未有浸入时，“法丛”先只是一个抽象的实秩 \(r\geq0\) 光滑向量丛 \(E\to M\)，带正定纤维内积 \(\langle\ ,\ \rangle_E\) 和保持该内积的联络 \(\nabla^E\)。其余数据是光滑黎曼度量 \(g\) 和光滑对称双线性型
\[
	B\in\Gamma(\operatorname{Sym}^2\mathrm T^*M\otimes E).
\]
这时 \(B(X,Y)\) 已经指定“向哪个法方向弯曲”，但尚未把该法方向识别为某个环境切空间中的向量。记 \(\nabla\) 为 \(g\) 的 Levi-Civita 联络，按
\[
	g(A_\xi X,Y)=\langle B(X,Y),\xi\rangle_E
\]
定义自伴算子 \(A_\xi\)。秩 \(r=0\) 时，\(E,B,A\) 的各项均按零理解。

\ProseParagraphBreak

固定一个实数 \(c\)，表示待实现的环境截面曲率。上一节的空间形式相容方程在这些抽象数据上仍有意义：
\begin{align}
	\mathrm{Rm}(X,Y,Z,W)
	&=c\bigl(g(X,Z)g(Y,W)-g(X,W)g(Y,Z)\bigr)\notag\\
	&\quad+\langle B(X,Z),B(Y,W)\rangle_E
	-\langle B(X,W),B(Y,Z)\rangle_E,\label{g3:eq:abstract-gauss-data}\\
	(\nabla B)(X;Y,Z)&=(\nabla B)(Y;X,Z),\label{g3:eq:abstract-codazzi-data}\\
	\langle R^E(X,Y)\xi,\eta\rangle_E
	&=g([A_\xi,A_\eta]X,Y).\label{g3:eq:abstract-ricci-data}
\end{align}
这里
\[
	(\nabla B)(X;Y,Z)
	=\nabla^E_X(B(Y,Z))-B(\nabla_XY,Z)-B(Y,\nabla_XZ),
\]
而 \(R^E(X,Y)=\nabla^E_X\nabla^E_Y-\nabla^E_Y\nabla^E_X-\nabla^E_{[X,Y]}\)。三式分别对应\cref{g3:cor:space-form-compatibility}中的切向、混合与法向相容条件。

\ProseParagraphBreak

为明确实现的目标，令 \(N=n+r\)，并记 \(Q_c^N\) 为以下标准模型：\(c=0\) 时取 \(\mathbb R^N\)；\(c>0\) 时取半径 \(1/\sqrt c\) 的圆球面。若 \(c<0\)，令 \(a=\sqrt{|c|}\)，在 \(\mathbb R^{N+1}\) 上取双线性型
\[
	\langle z,z\rangle_-=\sum_{i=1}^N(z^i)^2-(z^{N+1})^2,
\]
并取双曲面的一片
\[
	Q_c^N=\{z:\langle z,z\rangle_-=1/c,\ z^{N+1}>0\}.
\]
它的诱导度量正定：写 \(z=(x,\sqrt{a^{-2}+|x|^2})\)，切向量 \(v\) 的平方长度为
\[
	|v|^2-\frac{\langle x,v\rangle^2}{a^{-2}+|x|^2}>0
	\qquad(v\ne0).
\]
因此这里虽然借助一个有负方向的线性空间，所构造的环境 \(Q_c^N\) 仍是黎曼流形。

\ProseParagraphBreak

对 \(c\ne0\)，统一记 \(\varepsilon=\operatorname{sign}c\)、\(a=\sqrt{|c|}\)，把环境线性空间最后一维的平方长度取为 \(\varepsilon\)。二次曲面满足 \(\langle f,f\rangle_\varepsilon=1/c\)，其单位法向 \(af\) 的平方长度为 \(\varepsilon\)。对曲面切向场 \(U,V\)，普通求导的法分量由 \(\langle V,f\rangle_\varepsilon=0\) 得到，因而
\begin{equation}\label{g3:eq:quadric-ambient-derivative}
	\partial_UV=\nabla^{Q}_UV-c\langle U,V\rangle_\varepsilon f.
\end{equation}
切向投影 \(\nabla^Q\) 保持诱导度量且无挠，所以是其 Levi-Civita 联络。将上式代入环境平直联络的曲率恒等式，取切向分量便得
\[
	R^{Q}(U,V)W=c\bigl(\langle V,W\rangle U-\langle U,W\rangle V\bigr).
\]
故这些模型的曲率与式\eqref{g3:eq:abstract-gauss-data}中的 \(c\) 一致。
\ProseParagraphBreak

这些模型也满足空间形式所要求的完备性。Euclid 空间的通常距离完备；圆球面由紧性及\cref{g3:cor:compact-riemannian-complete}完备。负曲率模型可与上一章的上半空间直接比较。对 \(x\in\mathbb R^{N-1}\)、\(y>0\)，令
\[
	\Theta(x,y)=\left(\frac{x}{ay},
	\frac{|x|^2+y^2-1}{2ay},
	\frac{|x|^2+y^2+1}{2ay}\right).
\]
直接代入得到 \(\langle\Theta,\Theta\rangle_-=-a^{-2}\)，且最后一分量为正。反过来，若双曲面点记为 \((z',z^N,t)\)，则 \(u=t-z^N>0\)，其唯一原像为
\[
	y=\frac1{au},\qquad x=\frac{z'}u.
\]
因此 \(\Theta\) 是光滑双射且逆光滑。记 \(v=t+z^N\)，环境内积写成 \(|\mathrm dz'|^2-\mathrm du\,\mathrm dv\)。将 \(z'=x/(ay)\)、\(u=1/(ay)\)、\(v=(|x|^2+y^2)/(ay)\) 代入并求导，混合项抵消，得到
\[
	\Theta^*g_{Q_c}=\frac1{a^2y^2}\bigl(|\mathrm dx|^2+\mathrm dy^2\bigr).
\]
当 \(N\geq2\) 时，\cref{g3:exm:hyperbolic-warped-halfspace}已经证明右侧未缩放的半空间度量完备，乘以正数 \(a^{-2}\) 后仍完备。当 \(N=1\) 时省去 \(x\)，用 \(r=a^{-1}\log y\) 把它写成整条实线的 \(\mathrm dr^2\)，同样完备。

\subsection{标架系统及相容性}
\label{g3:subsec:0403:02}

先把切空间和法空间放在同一个丛中。在 \(F_0=\mathrm TM\oplus E\) 上取直和内积，并定义
\begin{equation}\label{g3:eq:submanifold-data-connection}
	D_X(Y+\xi)
	=\nabla_XY-A_\xi X+B(X,Y)+\nabla^E_X\xi.
\end{equation}
若浸入已经存在，这就是环境联络分别写出切向和法向分量后的表达。它确为保持直和内积的联络：两个纯切向或纯法向输入由各自的度量相容性处理；一个切向、一个法向输入所产生的交叉项为
\[
	-g(A_\xi X,Y)+\langle\xi,B(X,Y)\rangle_E=0.
\]
因此这个构造把 \(B\) 与 \(A\) 放在彼此相反的两个非对角块中。

\begin{proposition}\label{g3:prop:compatibility-flat-connection}
	设 \(M^n,g,E,\nabla^E,B\) 为上述光滑数据，\(n\geq1\)，\(c\in\mathbb R\)。当 \(c=0\) 时，三式\eqref{g3:eq:abstract-gauss-data}--\eqref{g3:eq:abstract-ricci-data}等价于 \(D\) 平坦。当 \(c\ne0\) 时，在
	\[
		F_c=\mathrm TM\oplus E\oplus\mathbb R s,
		\qquad\langle s,s\rangle=\varepsilon
	\]
	上取各项正交，并定义
	\begin{equation}\label{g3:eq:augmented-submanifold-connection}
		\widehat D_X(Y+\xi+t s)
		=\nabla_XY-A_\xi X+a tX
		+B(X,Y)+\nabla^E_X\xi
		+\bigl(Xt-\varepsilon a g(X,Y)\bigr)s.
	\end{equation}
	则 \(\widehat D\) 保持该非退化内积，而且三组相容方程等价于 \(\widehat D\) 平坦。
\end{proposition}
\begin{proof}
	记
	\[
		C(X,Y,Z)=(\nabla B)(X;Y,Z)-(\nabla B)(Y;X,Z).
	\]
	直接在式\eqref{g3:eq:submanifold-data-connection}中求两次导数，切向输入 \(Z\) 给出
	\begin{equation}\label{g3:eq:data-connection-curvature-tangent}
		R^D(X,Y)Z
		=R(X,Y)Z-A_{B(Y,Z)}X+A_{B(X,Z)}Y+C(X,Y,Z).
	\end{equation}
	其中前三项在 \(\mathrm TM\)，最后一项在 \(E\)。例如法向部分展开为
	\[
		\begin{aligned}
		&B(X,\nabla_YZ)-B(Y,\nabla_XZ)
		+\nabla^E_XB(Y,Z)-\nabla^E_YB(X,Z)\\
		&\hspace{30mm}-B([X,Y],Z),
		\end{aligned}
	\]
	用 \([X,Y]=\nabla_XY-\nabla_YX\) 整理后正是 \(C\)。对法向输入 \(\xi\)，其法向曲率分量为
	\[
		\bigl(R^D(X,Y)\xi\bigr)^E
		=R^E(X,Y)\xi-B(X,A_\xi Y)+B(Y,A_\xi X).
	\]
	与 \(\eta\) 配对，便得到
	\[
		\left\langle\bigl(R^D(X,Y)\xi\bigr)^E,\eta\right\rangle_E
		=\langle R^E(X,Y)\xi,\eta\rangle_E-g([A_\xi,A_\eta]X,Y).
	\]
	由于 \(D\) 保持内积，其曲率算子对直和内积反对称；因此余下的法向输入、切向输出块由 \(C\) 的负转置确定。于是所有曲率块为零，恰等价于 \(c=0\) 的 Gauss、Codazzi、Ricci 三式。

	现在设 \(c\ne0\)。新增项也保持内积，因为
	\[
		\langle\widehat D_Xs,Y\rangle
		+\langle s,\widehat D_XY\rangle
		=a g(X,Y)-\varepsilon^2a g(X,Y)=0.
	\]
	其余新增配对立即为零，故 \(\widehat D\) 是度量联络。又有 \(\widehat D_Xs=aX\)，所以
	\[
		R^{\widehat D}(X,Y)s
		=a\bigl(\nabla_XY-\nabla_YX-[X,Y]+B(X,Y)-B(Y,X)\bigr)=0.
	\]
	这里也用到 \(g(X,Y)=g(Y,X)\)，使 \(s\) 分量抵消。由曲率的反对称性，任意输入的曲率输出都没有 \(s\) 分量。

	对切向输入 \(Z\)，新增的切向曲率项只有
	\[
		-\varepsilon a^2g(Y,Z)X+\varepsilon a^2g(X,Z)Y
		=-c\bigl(g(Y,Z)X-g(X,Z)Y\bigr).
	\]
	因此
	\[
		\begin{aligned}
		\bigl(R^{\widehat D}(X,Y)Z\bigr)^{\mathrm TM}
		={}&R(X,Y)Z-A_{B(Y,Z)}X+A_{B(X,Z)}Y\\
		&-c\bigl(g(Y,Z)X-g(X,Z)Y\bigr),
		\end{aligned}
	\]
	而法向部分仍为 \(C(X,Y,Z)\)。法向输入、法向输出块也没有改变。将切向公式的输入换为 \(W\) 再与 \(Z\) 配对，其为零恰好给出式\eqref{g3:eq:abstract-gauss-data}；其余两个块分别给出另外两式。因此三式与增广联络平坦等价。
\end{proof}

把曲率变成零以后，还要真正解出标架，不能只把“平坦”当作可积的同义词。下一引理证明所需的局部构造。这里用到的分析输入是第一卷第13.1节的“局部初值问题”与“初值与参数的光滑依赖”：光滑常微分方程的解局部存在唯一，并在共同小区间上光滑依赖初值与有限维参数。下面将矩阵条目作为未知量、尚未行进的坐标作为参数使用这两个结果。

\begin{lemma}[平坦联络的局部标架]\label{g3:lem:flat-connection-local-frame}
	设 \(F\to M\) 为有限秩实光滑向量丛，\(M\) 无边界，\(\mathscr D\) 为曲率为零的光滑联络。给定 \(p\in M\) 及一个线性同构 \(\Phi_p:F_p\to V\)，其中 \(V\) 为同维实向量空间，则在 \(p\) 的某个连通坐标盒 \(U\) 上，存在唯一的光滑丛同构 \(\Phi:F|_U\to U\times V\)，使它在 \(p\) 等于指定同构，且
	\begin{equation}\label{g3:eq:parallel-trivialization-derivative}
		X(\Phi\sigma)=\Phi(\mathscr D_X\sigma)
	\end{equation}
	对每个光滑截面 \(\sigma\) 成立。若联络保持一个非退化内积，且 \(\Phi_p\) 是到 \(V\) 的线性等距，则每个 \(\Phi_x\) 都是线性等距。
\end{lemma}
\begin{proof}
	在 \(p\) 附近取行标架 \(e\)，使 \(p\) 的坐标为零，并写
	\[
		\mathscr D(eu)=e\left(\mathrm du+\sum_i A_i u\,\mathrm dx^i\right).
	\]
	取 \(V\) 中与 \(e(p)\) 对应的基后，所求平行标架写为 \(b=eG\)，其方程为
	\begin{equation}\label{g3:eq:parallel-frame-matrix-system}
		\partial_iG=-A_iG,\qquad G(0)=I.
	\end{equation}
	平坦条件则是
	\[
		\partial_iA_j-\partial_jA_i+A_iA_j-A_jA_i=0.
	\]

	先沿第一坐标轴解式\eqref{g3:eq:parallel-frame-matrix-system}的 \(i=1\) 方程。假定已经在前 \(k-1\) 个坐标组成的盒上得到解，就以该解为 \(x^k=0\) 处的初值，沿第 \(k\) 个坐标解 \(\partial_kG=-A_kG\)，将前 \(k-1\) 个坐标作为参数，其余坐标仍取零。所引常微分方程定理保证在缩小的共同盒上得到光滑解。有限步构造只需缩小有限次。

	关键在于新一步不破坏先前方程。对 \(i<k\)，令误差矩阵 \(E_i=\partial_iG+A_iG\)。由刚解出的方程，
	\[
		\partial_kE_i
		=-A_kE_i+
		\bigl(\partial_kA_i-\partial_iA_k+A_kA_i-A_iA_k\bigr)G
		=-A_kE_i.
	\]
	在 \(x^k=0\) 上误差为零，唯一性便给出 \(E_i=0\)。归纳至 \(k=n\)，即解出了所有坐标方向的方程。每一步中，若初值矩阵可逆，另解 \(\partial_kK=KA_k\) 且初值为其逆，则 \(\partial_k(KG)=0\)，故 \(KG=I\)。因此所得 \(G\) 始终可逆，\(b\) 确为标架。

	这个标架满足 \(\mathscr Db=0\)，称为\term[pingxingbiaojia]{平行标架}[parallel frame]。令 \(\Phi\) 把 \(b\) 的各列送到 \(V\) 中固定基，便得到式\eqref{g3:eq:parallel-trivialization-derivative}。两个这样的标架之间的过渡矩阵所有偏导均为零，所以在连通盒中为常矩阵；初值相同则矩阵为恒等，证明唯一性。若 \(\mathscr D\) 保持内积，则 \(X\langle b_i,b_j\rangle=0\)，各配对恒等于初值配对。故初始线性等距延拓后仍为线性等距，这也适用于有一个负方向的内积。
\end{proof}

在平行标架中，丛中的变化已经变成固定向量空间中的普通求导。Euclid 情形还需把速度一形式积分成位置。为补齐这一步，设坐标域关于原点星形，\(\alpha=\sum_i\alpha_i\,\mathrm dx^i\) 是取值于有限维向量空间的光滑闭一形式。定义
\begin{equation}\label{g3:eq:star-shaped-one-form-primitive}
	f(x)=b+\int_0^1\sum_i\alpha_i(tx)x^i\,\mathrm dt.
\end{equation}
由闭性 \(\partial_j\alpha_i=\partial_i\alpha_j\)，在积分号内求导得到
\[
	\partial_jf(x)=\int_0^1\left(\alpha_j(tx)
	+t\sum_i x^i\partial_i\alpha_j(tx)\right)\mathrm dt
	=\int_0^1\frac{\mathrm d}{\mathrm dt}\bigl(t\alpha_j(tx)\bigr)\mathrm dt
	=\alpha_j(x).
\]
因此 \(\mathrm df=\alpha\)，且 \(f(0)=b\)。同一连通域中的两个原函数只差常向量。

\subsection{空间形式中的局部基本定理}
\label{g3:subsec:0403:03}

现在可以把相容性转成存在性。所谓实现，还必须指定抽象法丛与实际法丛怎样对应；只要求诱导度量相同，会丢掉本节给定的一部分数据。

\begin{theorem}[子流形的局部基本定理]\label{g3:thm:submanifold-fundamental-local}
	设 \(M^n\) 为无边界光滑流形，\(n\geq1\)，\(g\) 为光滑黎曼度量，\(E\to M\) 为秩 \(r\geq0\) 的光滑 Euclid 向量丛，\(\nabla^E\) 为其度量联络，\(B\in\Gamma(\operatorname{Sym}^2\mathrm T^*M\otimes E)\) 光滑。固定 \(c\in\mathbb R\)。若这些数据满足式\eqref{g3:eq:abstract-gauss-data}--\eqref{g3:eq:abstract-ricci-data}，则每个 \(p\in M\) 都有一个连通邻域 \(U\)，其上存在光滑等距浸入
	\[
		f:(U,g)\longrightarrow Q_c^{n+r}
	\]
	及覆盖恒等映射的光滑纤维等距 \(\iota:E|_U\to N_fU\)，使
	\[
		B_f(X,Y)=\iota(B(X,Y)),\qquad
		\nabla^{\perp,f}_X(\iota\xi)=\iota(\nabla^E_X\xi).
	\]
	反之，具有这样法丛识别的等距浸入必满足三组相容方程。
\end{theorem}

这是\term[ziliuxingjubujibendingli]{子流形的局部基本定理}[local fundamental theorem of submanifolds]。它给出浸入及完整法丛数据。

\begin{proof}
	先设 \(c=0\)。由\cref{g3:prop:compatibility-flat-connection}，\(F_0\) 上的 \(D\) 平坦。取小坐标盒，并由\cref{g3:lem:flat-connection-local-frame}构造保持内积的 \(\Phi:F_0|_U\to U\times\mathbb R^{n+r}\)。把切向量包含到 \(F_0\) 后定义向量值一形式
	\[
		\alpha(X)=\Phi(X).
	\]
	式\eqref{g3:eq:parallel-trivialization-derivative}给出
	\[
		\begin{aligned}
		\mathrm d\alpha(X,Y)
		&=\Phi(D_XY-D_YX-[X,Y])\\
		&=\Phi\bigl(\nabla_XY-\nabla_YX-[X,Y]+B(X,Y)-B(Y,X)\bigr)=0.
		\end{aligned}
	\]
	这里无挠性与 \(B\) 的对称性共同保证位置可积。由式\eqref{g3:eq:star-shaped-one-form-primitive}取得 \(f\)，满足 \(\mathrm df=\alpha\)。因为 \(\Phi\) 保持内积，
	\[
		\langle\mathrm df(X),\mathrm df(Y)\rangle=g(X,Y).
	\]
	故 \(\mathrm df\) 单射，\(f\) 是等距浸入。

	定义 \(\iota\xi=\Phi(\xi)\)。它与 \(\mathrm df(\mathrm TM)=\Phi(\mathrm TM)\) 正交，维数又恰为 \(r\)，所以识别了全部法丛。对切向与法向截面分别求导，得到
	\[
		\partial_X(\mathrm df(Y))=\mathrm df(\nabla_XY)+\iota B(X,Y),
		\qquad
		\partial_X(\iota\xi)=-\mathrm df(A_\xi X)+\iota(\nabla^E_X\xi).
	\]
	这两式正是所需的第二基本形式与法联络。

	再设 \(c\ne0\)。此时对平坦的 \(\widehat D\) 取内积保持的平凡化
	\[
		\Phi:F_c|_U\longrightarrow U\times\mathbb R^{n+r+1},
	\]
	目标的内积按前述 \(\varepsilon\) 选择。定义
	\begin{equation}\label{g3:eq:quadric-immersion-from-frame}
		f=\frac1a\Phi(s).
	\end{equation}
	由于 \(\widehat D_Xs=aX\)，有
	\[
		\mathrm df(X)=\Phi(X),\qquad
		\langle f,f\rangle_\varepsilon=\frac\varepsilon{a^2}=\frac1c.
	\]
	于是 \(f\) 落在所需二次曲面上，并且仍为等距浸入。当 \(c<0\) 时，在初始同构中把 \(s_p\) 送到所选正时间方向；由连续性及 \(U\) 连通，\(f\) 始终落在双曲面的同一片。

	仍令 \(\iota\xi=\Phi(\xi)\)。它既与 \(\Phi(\mathrm TM)\) 正交，也与 \(\Phi(s)=af\) 正交，所以属于 \(Q_c^{n+r}\) 内的法空间。对 \(\mathrm df(Y)\) 作普通求导，得到
	\[
		\partial_X(\mathrm df(Y))
		=\mathrm df(\nabla_XY)+\iota B(X,Y)-c\,g(X,Y)f.
	\]
	最后一项是式\eqref{g3:eq:quadric-ambient-derivative}中二次曲面本身的法分量。去掉该项，\(Q_c^{n+r}\) 内的第二基本形式恰为 \(\iota B\)。同时
	\[
		\partial_X(\iota\xi)
		=-\mathrm df(A_\xi X)+\iota(\nabla^E_X\xi).
	\]
	这里没有径向分量，因为 \(g(\mathrm df(X),\iota\xi)=0\)，所以它也直接给出 \(Q_c^{n+r}\) 中的法联络。至此存在性证明完毕。逆向结论正是\cref{g3:cor:space-form-compatibility}。
\end{proof}

当 \(r=0\) 时，这个构造给出常截面曲率度量到相应模型的局部等距。余维为正时，\(\iota\) 则使抽象法向量成为实际环境向量；平均曲率也随之实现为 \(H_f=\iota(\operatorname{tr}_gB/n)\)。高余维中若删去法联络或 Ricci 方程，前面的联络就未必平坦，因而不能继续这个构造。

\subsection{唯一性及整体单值障碍}
\label{g3:subsec:0403:04}

初始标架和初始位置的选择改变了所得映射，但没有改变实现的数据。以下唯一性始终固定源流形中的点对应，并同时比较法丛识别。

\begin{proposition}[保持法丛识别的唯一性]\label{g3:prop:submanifold-fundamental-uniqueness}
	设 \(U\) 连通，\((f_1,\iota_1)\)、\((f_2,\iota_2)\) 是同一组光滑数据 \((g,E,\nabla^E,B)\) 到标准模型 \(Q_c^{n+r}\) 的实现。则存在环境等距 \(T\)，使
	\[
		f_2=T\circ f_1,\qquad
		\iota_2=\mathrm dT\circ\iota_1.
	\]
	当 \(c=0\) 时，\(T\) 可取为唯一的 Euclid 刚性变换；当 \(c\ne0\) 时，\(T\) 可取为唯一一个保持上述线性环境内积、并限制到指定二次曲面分支的线性变换。
\end{proposition}
\begin{proof}
	若 \(c=0\)，每个实现给出纤维同构
	\[
		\Psi_j(Y+\xi)=\mathrm df_j(Y)+\iota_j\xi.
	\]
	Gauss 与 Weingarten 公式说明 \(X(\Psi_j\sigma)=\Psi_j(D_X\sigma)\)。因此在任一平行标架域中，\(\Psi_j\) 把平行标架送到一组常向量，故 \(Q=\Psi_2\Psi_1^{-1}\) 是常正交矩阵。这个结论是局部的，\(U\) 连通使同一个 \(Q\) 在全域成立。于是 \(\mathrm df_2=Q\,\mathrm df_1\)，故 \(f_2-Qf_1=b\) 为常向量，取 \(T(z)=Qz+b\) 即可。\(Q\) 由一点处切向和法向的完整正交分解确定，\(b\) 再由该点的位置确定，因而唯一。

	若 \(c\ne0\)，改取
	\[
		\Psi_j(Y+\xi+t s)
		=\mathrm df_j(Y)+\iota_j\xi+a t f_j.
	\]
	式\eqref{g3:eq:quadric-ambient-derivative}表明这次有 \(X(\Psi_j\sigma)=\Psi_j(\widehat D_X\sigma)\)。所以 \(Q=\Psi_2\Psi_1^{-1}\) 仍为保持环境内积的常矩阵。把 \(s\) 代入便得 \(f_2=Qf_1\)，把 \(E\) 代入得到法丛等式。它把 \(f_1(p)\) 送到 \(f_2(p)\)；在双曲情形，两点位于同一指定片，故 \(Q\) 保持这一片。位置径向、切向与法向共同张成整个线性环境空间，因此 \(Q\) 也唯一。
\end{proof}

局部唯一性使解能够沿路径继续，但“能够继续”与“绕一周后回到原解”是两个条件。先说明单连通性如何消去后者。

\begin{proposition}[单连通数据的整体实现]\label{g3:prop:submanifold-fundamental-global}
	设 \(M^n\) 连通、单连通、无边界，\(n\geq1\)，且 \(g,E,\nabla^E,B\) 是全域光滑数据，满足曲率 \(c\) 的三组相容方程。则存在定义在整个 \(M\) 上的等距浸入 \(f:M\to Q_c^{n+r}\) 及实现数据的法丛等距 \(\iota\)。它们仍满足\cref{g3:prop:submanifold-fundamental-uniqueness}的唯一性。
\end{proposition}
\begin{proof}
	先把局部平行标架延拓为整体平凡化。两个平行标架在交叠的每个连通分支上只差常矩阵。因此沿一条连续路径，可以把参数区间有限细分，使每一小段都位于一个平行标架域内，再逐段保持标架分量。增加分点不改变结果；两种选择取共同细分后，在每段连通的路径像上，常过渡矩阵也给出相同结果。这就定义了沿连续路径的纤维同构。

	若两条路径定端点同伦，把同伦参数方形细分为足够小的矩形，使每格的像都位于一个平行标架域内。这样的有限细分由参数方形的紧性取得。在每一格中，先向右再向上与先向上再向右都保持同一标架分量，因而得到相同同构。逐格用此等式替换路径，内部边的运输与反向运输抵消，便得两条边界路径的运输相同；两条端点边是常路径，不产生变化。于是运输只依赖定端点同伦类。单连通性保证同端点路径同伦，所以从固定基点的一个基出发，运输得到全域标架。在每个局部平行域内，它与已有标架只差常矩阵，因而光滑且平行。

	对 \(c\ne0\)，用这个整体平凡化代入式\eqref{g3:eq:quadric-immersion-from-frame}，立即得到全域实现。对于 \(c=0\)，整体平凡化先给出闭一形式 \(\alpha(X)=\Phi(X)\)，还需证明它有全域原函数。每个小坐标盒内，式\eqref{g3:eq:star-shaped-one-form-primitive}已经给出局部原函数；同一连通交叠上的两个原函数只差常向量。因此沿连续路径，把各小段端点的局部原函数之差相加，所得增量与细分和局部原函数选择无关；对分片光滑路径，它就是通常的积分。对定端点同伦再用前段的矩形网格，每格的四边增量和为零，内部边相消，故总增量在同伦下不变。单连通性使这个增量只依赖终点。固定基点值，遂得到全域 \(f\)，而在每个小盒上它与局部原函数只差常向量，故光滑且 \(\mathrm df=\alpha\)。其余验证与局部定理相同。
\end{proof}

这个命题不要求 \(g\) 完备，所得映射也不必单射。这里的整体结论是浸入在整个给定流形上有定义。

\ProseParagraphBreak

在非单连通流形上，同一个构造给出更具体的判据。令 \(\mathscr D_c=D\) 当 \(c=0\)，令 \(\mathscr D_c=\widehat D\) 当 \(c\ne0\)。要把局部实现连同给定法丛识别拼成全域实现，首先必须使 \(\mathscr D_c\) 沿每条闭路径的运输为恒等。这个条件等价于存在整体平行平凡化：充分性由从固定基点运输得到，必要性则由平凡化中分量不变得到。非零 \(c\) 时，这已经足够，位置由 \(s\) 的像给出；\(c=0\) 时，还必须且只须
\begin{equation}\label{g3:eq:euclidean-development-periods}
	\int_\gamma\alpha=0
	\qquad\text{对每条分片光滑闭路径 }\gamma,
	\qquad\alpha(X)=\Phi(X).
\end{equation}
必要性来自 \(\alpha=\mathrm df\)，充分性来自沿路径积分定义原函数。更换整体平行平凡化只会给 \(\alpha\) 乘一个固定线性等距，所以周期是否为零不依赖这一选择。

\ProseParagraphBreak

也可以直接观察沿环路继续一个局部实现后所得的返回等距。Euclid 情形中，它可能既旋转标架，又平移位置；即使前一部分为恒等，后一部分仍可能非零。这正是整体单值的两个不同障碍。要求法丛识别也返回，意味着整个适应标架返回；只观察浸入的像集，无法判断这一条件。

\begin{example}\label{g3:exm:flat-circle-development-period}
	在 \(M=\mathbb R/L\mathbb Z\)、\(L>0\) 上取度量 \(\mathrm ds^2\)，令 \(e=\partial_s\) 为全局单位切向场。取平凡法线丛 \(E=M\times\mathbb R\nu\)，令 \(\nabla^E\nu=0\)、\(B=0\)，并取 \(c=0\)。一维中三组曲率方程全部成立，\(D\) 也有全局平行标架 \((e,\nu)\)，因此标架沿圆周返回原值。然而相应位置一形式为 \(\alpha=(1,0)\,\mathrm ds\)，它沿正向一周的积分为 \((L,0)\ne0\)。在参数直线上得到的实现 \(f(s)=(s,0)\) 满足 \(f(s+L)=f(s)+(L,0)\)，不能下降到圆周。局部相容性、全局平行标架与位置闭合，在这个例子中分别成为可直接检验的条件。
\end{example}

% 纲要条目排为编号小节标题，并进入目录与书签；说明段仍保留于纲要注释中。

% !TeX root = ../../Book3.tex
% 纲要标题：### 4.4 子流形模型与曲率限制
% 纲要位置：几何学计划纲要.md，第 2251 行。
% 正文状态：初稿；尚未正式审核。

% 纲要内容开始
% 
% * 球面、乘积和管状模型
% * 超曲面的主曲率
% * 内在曲率相同而外在几何不同
% * 后续面积变分所需的公式
% 
% ---
% 
% 纲要内容结束

\section{子流形模型与曲率限制}
\label{g3:sec:0404}

相容方程把内在曲率、第二基本形式和法联络联系起来。具体计算时，可以从一个已知浸入求出这些量，再把它们代回方程；也可以先用曲率限制判断某种外在形状是否可能。球面、球面乘积和法向偏移提供了两种做法都能逐项核验的模型。最后，我们把平均曲率与体积密度的一阶变化联系起来。

\subsection{球面、乘积和管状模型}
\label{g3:subsec:0404:01}

在本小节的 Euclid 环境中，普通方向微分记为 \(\mathrm D\)。半径 \(a>0\) 的球面取外单位法向 \(\nu=f/a\)，则
\[
	\mathrm D_X\nu=\frac1a\mathrm df(X),\qquad
	A_\nu=-\frac1a I,\qquad
	B(X,Y)=-\frac1a g(X,Y)\nu,
	\qquad H=-\frac1a\nu.
\]
球面是全脐的，平均曲率向量指向球心。Gauss 方程却给出正截面曲率 \(a^{-2}\)，因为其中第二基本形式以两个因子的乘积出现。更换为内单位法向会反转 \(A_\nu\) 的符号，但不改变向量值的 \(B,H\) 或内在曲率。

\begin{example}\label{g3:exm:product-spheres-normal-data}
	设 \(p,q\ge1\)、\(a,b>0\)，考虑乘积浸入
	\[
		f:S^p(a)\times S^q(b)\longrightarrow
		\mathbb R^{p+1}\times\mathbb R^{q+1},\qquad f(x,y)=(x,y).
	\]
	记两个因子度量为 \(g_1,g_2\)，并在切丛另一因子上将它们补为零。法丛具有全局正交单位标架
	\[
		\nu_1=(x/a,0),\qquad \nu_2=(0,y/b).
	\]
	两者的环境导数均为切向量，故 \(\nabla^\perp\nu_1=\nabla^\perp\nu_2=0\)。直接求导得到
	\[
		B=-\frac{g_1}{a}\nu_1-\frac{g_2}{b}\nu_2,
		\qquad
		H=-\frac1{p+q}\left(\frac pa\nu_1+\frac qb\nu_2\right).
	\]
	两个形算子在乘积分解下同时对角化，彼此交换，法曲率为零。Gauss 方程对纯因子截面给出 \(a^{-2}\)、\(b^{-2}\)，对混合截面给出零；这与第三章的乘积曲率一致。当某个因子是一维时，只使用实际存在的二维截面。

	同一像还位于半径 \(L=\sqrt{a^2+b^2}\) 的球面中。相对于这个球面，单位法向可取
	\[
		\eta=\left(\frac{b}{aL}x,-\frac{a}{bL}y\right).
	\]
	它与位置向量正交，且长度为一。对两个因子的切向量求导，形算子的特征值分别是
	\[
		-\frac{b}{aL}\quad\text{（重数 \(p\)）},\qquad
		\frac{a}{bL}\quad\text{（重数 \(q\)）}.
	\]
	因此，这个乘积在球面中极小的充要条件为
	\[
		q a^2=p b^2,
		\qquad\text{即}\qquad
		a^2=\frac{p}{p+q}L^2,\quad b^2=\frac{q}{p+q}L^2.
	\]
	它在 Euclid 空间中的平均曲率向量始终非零。平均曲率依赖所选环境；切空间正交补改变后，第二基本形式所取的法向投影也随之改变。
\end{example}

沿各点的法方向移动，可以把子流形附近的空间组织成管状坐标。在 Euclid 空间中，这个构造只涉及向量加法；局部可逆性来自切向与法向的互补。

\begin{proposition}\label{g3:prop:euclidean-normal-tube}
	设 \(f:M^n\to\mathbb R^{n+r}\) 为光滑浸入，\(r\ge1\)，\(M\) 无边界。定义法丛上的映射
	\[
		\mathcal T:N_fM\longrightarrow\mathbb R^{n+r},
		\qquad \mathcal T(p,\xi)=f(p)+\xi.
	\]
	它在每个零向量 \((p,0)\) 的某个邻域上为微分同胚。若 \(M\) 紧且 \(f\) 为嵌入，则存在 \(\varepsilon>0\)，使 \(\mathcal T\) 将整个法圆盘丛 \(\{(p,\xi):|\xi|<\varepsilon\}\) 微分同胚地映到 \(f(M)\) 的一个开邻域。
\end{proposition}
\begin{proof}
	在局部法丛平凡化下，零截面处的切空间自然分成 \(\mathrm T_pM\oplus N_pf\)，且
	\[
		\mathrm d\mathcal T_{(p,0)}(X,\eta)=\mathrm df_p(X)+\eta.
	\]
	这是两个正交互补子空间到 \(\mathbb R^{n+r}\) 的线性同构。第一卷第11.3节的逆函数定理适用于同维光滑流形间的映射：一点处微分可逆便在该点附近有光滑逆。因此得到第一项结论。

	在紧情形，零截面紧，微分可逆性是开条件；有限覆盖给出共同的小半径，使 \(\mathcal T\) 在该法圆盘丛上处处为局部微分同胚。还要证明共同的小半径可以保证单射。若不能，则有两列不同点 \((p_j,\xi_j)\)、\((q_j,\eta_j)\)，满足 \(|\xi_j|,|\eta_j|\to0\) 且像相同。取子列使 \(p_j\to p\)、\(q_j\to q\)。像相同给出 \(f(p)=f(q)\)，嵌入的单射性使 \(p=q\)。在这个点的一个法丛平凡化中，两列点都收敛到 \((p,0)\)，与第一段得到的局部单射性矛盾。故足够小的半径可用。单射的局部微分同胚是到开像的微分同胚，且像包含 \(f(M)\)。
\end{proof}

上述开邻域称为一个\term[guanzhuanglinyu]{管状邻域}[tubular neighborhood]。仅有浸入时，不同分支的管状像可以相交；紧嵌入的证明同时使用了零截面的紧性与不同基点不能具有相同像。

\begin{proposition}\label{g3:prop:parallel-hypersurface}
	设 \(f:M^n\to\mathbb R^{n+1}\) 为带光滑单位法向 \(\nu\) 的光滑浸入，\(A=A_\nu\)。对常数 \(t\) 定义 \(f_t=f+t\nu\)。在 \(I-tA\) 可逆的开集上，\(f_t\) 仍为浸入，\(\nu\) 仍是其单位法向；相对于共同的参数流形，
	\begin{align}
		\mathrm df_t&=\mathrm df\circ(I-tA),
		&g_t(X,Y)&=g((I-tA)X,(I-tA)Y),\label{g3:eq:parallel-hypersurface-metric}\\
		A_t&=(I-tA)^{-1}A,
		&\mu_{g_t}&=|\det(I-tA)|\mu_g.\label{g3:eq:parallel-hypersurface-shape}
	\end{align}
\end{proposition}
\begin{proof}
	法丛是一维的，\(\langle\nu,\nu\rangle=1\) 使 \(\nabla^\perp\nu=0\)，故 \(\mathrm d\nu=-\mathrm df\circ A\)。对 \(f_t\) 求导便得到第一式。由于 \(\mathrm df\) 单射，\(\mathrm df_t\) 单射恰等价于 \(I-tA\) 可逆；其像仍是原切超平面，故 \(\nu\) 仍垂直于像。内积给出度量公式。对新浸入使用 Weingarten 公式，
	\[
		-\mathrm df(A X)=\mathrm d\nu(X)
		=-\mathrm df_t(A_tX)
		=-\mathrm df((I-tA)A_tX),
	\]
	从而得到 \(A_t\)。最后在正交基中取度量矩阵行列式及其正平方根，即得密度公式。
\end{proof}

对于外法向的半径 \(a\) 球面，\(I-tA=(1+t/a)I\)，所以外移后半径为 \(a+t\)，切向尺度乘以 \(1+t/a\)，密度乘以 \((1+t/a)^n\)。\cref{g3:fig:normal-offset}画出圆周的对应关系。

\begin{figure}[htbp]
	\centering
	\input{Figures/Ch04-NormalOffset}
	\caption[圆周的法向偏移]{圆周的外法向偏移。相同角参数的切向微分从内圆移到外圆后放大为原来的 \(1+t/a\) 倍。}
	\label{g3:fig:normal-offset}
\end{figure}

\subsection{超曲面的主曲率}
\label{g3:subsec:0404:02}

余维一时，每个法向值都由一个实数确定。选定单位法向 \(\nu\) 后，写 \(B(X,Y)=g(AX,Y)\nu\)。自伴算子 \(A\) 在每一点有正交单位特征基，其特征值 \(\kappa_1,\ldots,\kappa_n\) 称为该法向下的\term[zhuqulv]{主曲率}[principal curvatures]，相应特征方向称为\term[zhufangxiang]{主方向}[principal directions]。重特征值处未必能把单个主方向光滑选出，下面的逐点公式不需要这种选择。

\begin{proposition}\label{g3:prop:principal-curvature-contractions}
	设 \(f:(M^n,g)\to\bar M^{n+1}\) 是光滑等距浸入，\(n\ge2\)，环境截面曲率恒为 \(c\)。在一点选单位法向 \(\nu\) 及主方向正交单位基 \((e_i)\)，令 \(H_\nu=\langle H,\nu\rangle=n^{-1}\sum_i\kappa_i\)。则
	\begin{align}
		K(e_i,e_j)&=c+\kappa_i\kappa_j\quad(i\ne j),\label{g3:eq:principal-sectional-curvature}\\
		\operatorname{Ric}(e_i,e_j)
		&=\bigl((n-1)c+nH_\nu\kappa_i-\kappa_i^2\bigr)\delta_{ij},\label{g3:eq:principal-ricci-curvature}\\
		\operatorname{Scal}
		&=n(n-1)c+n^2H_\nu^2-|A|^2.\label{g3:eq:hypersurface-scalar-curvature}
	\end{align}
\end{proposition}
\begin{proof}
	Gauss 方程在这组基中成为
	\[
		R_{ijkl}
		=c(\delta_{ik}\delta_{jl}-\delta_{il}\delta_{jk})
		+\kappa_i\kappa_j(\delta_{ik}\delta_{jl}-\delta_{il}\delta_{jk}).
	\]
	取 \(k=i,l=j\) 得到截面公式。固定 \(i\)，把它对 \(j\ne i\) 相加，得到
	\[
		\operatorname{Ric}(e_i,e_i)
		=(n-1)c+\kappa_i\Bigl(\sum_j\kappa_j-\kappa_i\Bigr).
	\]
	非对角分量由同一四阶公式为零。再对 \(i\) 求和，使用 \(|A|^2=\sum_i\kappa_i^2\)，即得标量曲率式。
\end{proof}

这些公式也给出外在形状的限制。例如 Euclid 超曲面若在一点平坦，则 \(\kappa_i\kappa_j=0\) 对所有 \(i\ne j\) 成立，故该点形算子的秩至多为一。若每个截面曲率都严格为正，则任意两个主曲率的乘积均为正，所有主曲率同号。另一方面，全脐点的 \(A=H_\nu I\) 使所有截面曲率等于 \(c+H_\nu^2\)。

\ProseParagraphBreak

法向偏移的主曲率可由式\eqref{g3:eq:parallel-hypersurface-shape}逐个读取：
\[
	\kappa_i(t)=\frac{\kappa_i}{1-t\kappa_i}.
\]
当某个 \(1-t\kappa_i\) 为零时，相应切向方向被 \(\mathrm df_t\) 压为零，偏移参数化失去浸入性。这是一个局部微分条件。即使这些因子都非零，不同基点的法线仍可能在远处相交，因此共同管状半径还需要控制整体单射性。

\subsection{内在曲率相同而外在几何不同}
\label{g3:subsec:0404:03}

内在度量记录沿子流形测量的长度，第二基本形式则记录切空间在环境中的转动。同一度量可以对应不同的第二基本形式；局部基本定理的唯一性因而需要完整的相容数据。

\begin{example}\label{g3:exm:plane-cylinder-isometry}
	在一个平面开矩形 \(U\subset\mathbb R^2\) 上，考虑
	\[
		f_0(u,v)=(u,v,0),\qquad
		f_a(u,v)=\bigl(a\cos(u/a),a\sin(u/a),v\bigr),\quad a>0.
	\]
	两者的诱导度量都是 \(\mathrm du^2+\mathrm dv^2\)，所以具有相同的内在距离局部结构和零 Gauss 曲率。对圆柱取外法向
	\[
		\nu(u,v)=\bigl(\cos(u/a),\sin(u/a),0\bigr),
	\]
	则 \(A_\nu\partial_u=-a^{-1}\partial_u\)、\(A_\nu\partial_v=0\)，而平面上 \(B=0\)。故两者不可能在任何开集上由一个 Euclid 刚性运动联系：环境等距保持第二基本形式，尤其保持其范数。这里一个范数为零，另一个为 \(a^{-1}\)。

	若允许 \(u\) 遍历整个实轴，\(f_a(u+2\pi a,v)=f_a(u,v)\)。因此这个局部等距参数化在整个平面上不单射；把参数域限制到适当的小矩形后，才得到到一块圆柱面的等距微分同胚。
\end{example}

上述区别在高余维中还包括法联络。即使给定相同的 \(g\) 和对称型的分量，换法标架时也必须同时按同一个丛等距变换法联络与 \(B\)。三组相容方程正是对这套几何数据的条件；只比较一张矩阵或只比较内在曲率，不能使用基本定理的唯一性结论。

\subsection{后续面积变分所需的公式}
\label{g3:subsec:0404:04}

当浸入随参数移动时，切向速度主要改变参数点的分布，法向速度才直接与平均曲率配对。先计算诱导度量及其体积密度的导数，便能在任意维数、任意余维得到面积的一阶公式。这里“面积”指 \(n\) 维诱导体积；对非单射浸入按参数流形积分，保留重数。

\begin{theorem}[面积的一阶变分]\label{g3:thm:first-variation-submanifold-area}
	设 \(F:(-\varepsilon,\varepsilon)\times M^n\to(\bar M,\bar g)\) 光滑，每个 \(f_t=F(t,\cdot)\) 都是浸入，\(n\ge1\)。记 \(g_t=f_t^*\bar g\)、\(g=g_0\)，并将初速度分解为 \(V=\partial_tF|_0=\mathrm df_0(W)+V^\perp\)。则
	\begin{align}
		\dot g_0(X,Y)
		&=g(\nabla_XW,Y)+g(X,\nabla_YW)
		-2\langle B(X,Y),V^\perp\rangle,\label{g3:eq:variation-induced-metric}\\
		\left.\frac{\mathrm d}{\mathrm dt}\right|_0\mu_{g_t}
		&=\bigl(\operatorname{div}_gW-n\langle H,V^\perp\rangle\bigr)\mu_g.\label{g3:eq:variation-induced-density}
	\end{align}
	若 \(\Omega\subset M\) 是具有光滑边界的紧区域，\(\eta\) 是 \(\partial\Omega\) 在 \(M\) 内的外单位法向，则
	\begin{equation}\label{g3:eq:first-variation-area-integral}
		\left.\frac{\mathrm d}{\mathrm dt}\right|_0\operatorname{Vol}_{g_t}(\Omega)
		=\int_{\partial\Omega}g(W,\eta)\mu_{\partial g}
		-n\int_\Omega\langle H,V^\perp\rangle\mu_g.
	\end{equation}
\end{theorem}
\begin{proof}
	把 \(X,Y\) 延拓成与时间无关的参数向量场。环境联络度量相容且无挠，因而沿二参数坐标面有 \(\bar\nabla_t\mathrm df_t(X)=\bar\nabla_X\partial_tF\)。逐项求导得到
	\[
		\dot g_0(X,Y)
		=\langle\bar\nabla_XV,\mathrm df_0(Y)\rangle
		+\langle\mathrm df_0(X),\bar\nabla_YV\rangle.
	\]
	代入 \(V=\mathrm df_0(W)+V^\perp\)，用 Gauss 分解和 Weingarten 公式提取切向分量，便得度量变化式。

	在任一固定坐标中，密度系数为 \(\sqrt{\det g_t}\)。可逆矩阵的行列式求导给出
	\[
		\left.\frac{\mathrm d}{\mathrm dt}\right|_0\sqrt{\det g_t}
		=\frac12\operatorname{tr}_g(\dot g_0)\sqrt{\det g}.
	\]
	在一点取正交单位基求迹，前两项各贡献 \(\operatorname{div}_gW\)，最后一项贡献 \(-2n\langle H,V^\perp\rangle\)。这证明密度公式。紧区域上可在积分号下求导，再对 \(W\) 使用\cref{g3:thm:divergence-theorem}，得到边界通量。密度与该散度定理均不要求 \(M\) 可定向。
\end{proof}

\begin{corollary}\label{g3:cor:minimal-area-stationarity}
	设 \(f:M^n\to\bar M\) 为光滑浸入，\(n\ge1\)，\(M\) 无边界。则 \(H=0\) 当且仅当对于任意具有光滑边界的紧区域 \(\Omega\subset M\)，以及任意满足 \(\operatorname{supp}V\subset\operatorname{int}\Omega\) 的上述光滑变分，都有 \(\left.\frac{\mathrm d}{\mathrm dt}\right|_0\operatorname{Vol}_{g_t}(\Omega)=0\)。
\end{corollary}
\begin{proof}
	若 \(H=0\)，在包含初速度支集的紧区域上应用式\eqref{g3:eq:first-variation-area-integral}，边界项和内部项都为零。反过来，若某点 \(H\ne0\)，在其附近选一个非负光滑截断函数 \(\chi\)，支集紧含于一个坐标片，且在该点附近为正。取法向初速度 \(V=\chi H\)。它可以由变分
	\[
		f_t(p)=\exp^{\bar g}_{f(p)}(tV(p))
	\]
	实现：支集紧使共同的小时间可用，指数映射的联合光滑性见\cref{g3:prop:exponential-local}；浸入性是开条件，在这个紧支集上缩小时间便仍成立，支集外则保持原映射。此时一阶变分为 \(-n\int\chi|H|^2\mu_g<0\)，与驻定性矛盾。因此 \(H\) 必处处为零。
\end{proof}

这一结论说明“极小”的一阶意义是面积驻定。是否在所有邻近变分下取得局部最小值，还需要二阶变化的信息。球面外移的检验则直接固定本章符号：若 \(V=\nu\)、\(H=-\nu/a\)，密度导数为 \(n\mu_g/a>0\)，与半径增大一致。

\begin{proposition}\label{g3:prop:immersion-position-laplacian}
	设 \(f:(M^n,g)\to\mathbb R^{n+r}\) 为光滑等距浸入，\(n\ge1\)，对位置向量的各坐标分别使用 \(\Delta_g=-\operatorname{div}_g\operatorname{grad}_g\)。则
	\begin{equation}\label{g3:eq:position-laplacian}
		\Delta_g f=-nH,\qquad
		\Delta_g|f|^2=-2n-2n\langle f,H\rangle.
	\end{equation}
	若像位于 \(S^{n+r-1}(a)\)，以 \(H^S\) 表示相对于球面的平均曲率向量，则
	\begin{equation}\label{g3:eq:spherical-minimal-eigenmap}
		H^{\mathbb R}=H^S-\frac f{a^2},\qquad
		H^S=0\ \Longleftrightarrow\ \Delta_g f=\frac n{a^2}f.
	\end{equation}
\end{proposition}
\begin{proof}
	对任一坐标函数 \(f^\alpha\)，Hessian 的定义给出
	\[
		H_{f^\alpha}(X,Y)
		=\bigl(\mathrm D_X\mathrm df(Y)-\mathrm df(\nabla_XY)\bigr)^\alpha
		=B(X,Y)^\alpha.
	\]
	取负迹即得 \(\Delta_g f=-nH\)。由乘积法则，
	\[
		\Delta_g|f|^2=2\langle f,\Delta_g f\rangle-2|\mathrm df|^2.
	\]
	等距浸入使 \(|\mathrm df|^2=n\)，故得第二式。若像在半径 \(a\) 的球面内，先将环境导数投影到球面的切空间，再投影到 \(M\) 的切空间；球面自身的第二基本形式为 \(-g f/a^2\)。于是
	\[
		B^{\mathbb R}(X,Y)=B^S(X,Y)-\frac{g(X,Y)}{a^2}f.
	\]
	取平均迹得到两个平均曲率的关系，再代入位置向量的 Laplace 公式，便得最后的等价式。
\end{proof}

\begin{corollary}\label{g3:cor:no-compact-euclidean-minimal}
	正维紧无边界光滑流形不存在到 Euclid 空间的极小浸入。
\end{corollary}
\begin{proof}
	若有这样的浸入，以诱导度量计算，由 \(H=0\) 及式\eqref{g3:eq:position-laplacian}，有 \(\Delta_g|f|^2=-2n\)。在 \(M\) 上积分，散度定理使左边为零，右边却是 \(-2n\operatorname{Vol}_g(M)<0\)，矛盾。
\end{proof}

球面中的极小乘积没有这一矛盾：其 Euclid 平均曲率为 \(-f/a^2\)，在位置向量公式中恰好抵消常数项。这样，局部的切法分解、相容方程及整体积分在同一个模型上相互核对。


\begin{chaptersummary}
	切向与法向投影把环境导数分成内在 Levi-Civita 联络、第二基本形式、形算子及法联络。第二基本形式的对称性来自无挠性，形算子是它与指定法向配对后的自伴算子，平均曲率向量则是其平均迹。高余维法空间的变化具有自己的曲率，不能由单个形算子代替。

	\ProseParagraphBreak

	Gauss、Codazzi 和 Ricci 方程分别比较切曲率、第二基本形式的导数与法曲率。在空间形式中，它们恰好使所构造的切法直和联络或其增广联络平坦。积分平坦标架，再积分位置数据，就得到局部浸入；同一数据的两次构造相差环境等距。整体实现还要求沿闭路的标架返回以及 Euclid 情形的位置周期相容，局部方程本身不会消去这些障碍。

	\ProseParagraphBreak

	模型计算展示了内在和外在信息的不同作用：平面与圆柱可以局部等距，球面乘积可以在球面内极小而在 Euclid 空间中不极小，法向偏移的正则性则由形算子的特征值控制。面积的一阶变化由边界通量与平均曲率配对给出；无边界紧支变分下，\(H=0\) 等价于面积驻定。位置向量的 Laplace 公式进一步把这一外在条件转化为可积分的函数等式。
\end{chaptersummary}

\BookExercises

\begin{exercise}\label{g3:ex:graph-prescribed-second-fundamental-form}
	设 \(h:U\to\mathbb R^r\) 光滑，其中 \(U\subset\mathbb R^n\) 为原点邻域，\(n\ge1\)，且 \(h(0)=0\)、\(\mathrm dh_0=0\)。考虑图像浸入
	\[
		f(x)=(x,h(x))\in\mathbb R^{n+r},
	\]
	并给 \(U\) 诱导度量。
	\begin{enumerate}
		\item 求原点处的度量、联络系数、第二基本形式及平均曲率向量。将形算子写成 \(h\) 的二阶偏导矩阵。
		\item 给定任意对称双线性映射 \(C:\mathbb R^n\times\mathbb R^n\to\mathbb R^r\)，构造一个这样的图像，使其原点第二基本形式等于 \(C\)。当 \(n\ge2\)、\(r\ge1\) 时，给出 \(H(0)=0\) 而 \(B_0\ne0\) 的具体例子。
	\end{enumerate}
\end{exercise}
\begin{solution}
	1. 记 \(h_i^\alpha=\partial_i h^\alpha\)。坐标切向量及度量为
	\[
		f_i=(e_i,h_i),\qquad
		g_{ij}=\delta_{ij}+\sum_{\alpha=1}^r h_i^\alpha h_j^\alpha.
	\]
	因为 \(h_i^\alpha(0)=0\)，有 \(g_{ij}(0)=\delta_{ij}\)，且
	\[
		\partial_kg_{ij}(0)
		=\sum_\alpha
			\bigl(h_{ik}^\alpha h_j^\alpha+h_i^\alpha h_{jk}^\alpha\bigr)(0)=0.
	\]
	由 Christoffel 公式，\(\Gamma^k_{ij}(0)=0\)。原点切空间是水平的 \(\mathbb R^n\times\{0\}\)，法空间是竖直的 \(\{0\}\times\mathbb R^r\)，而
	\[
		\partial_i\partial_jf(0)
		=\left(0,\bigl(\partial_i\partial_jh^\alpha(0)\bigr)_{\alpha=1}^r\right)
	\]
	已经法向。因此
	\[
		B_0(e_i,e_j)
		=\left(0,\bigl(\partial_i\partial_jh^\alpha(0)\bigr)_\alpha\right),
		\qquad
		H(0)=\frac1n
			\left(0,\bigl(\sum_i\partial_i^2h^\alpha(0)\bigr)_\alpha\right).
	\]
	对法向量 \(\xi=(0,\xi^1,\ldots,\xi^r)\)，形算子在水平正交基中的矩阵为
	\[
		(A_\xi)_{ij}
		=\sum_{\alpha=1}^r\xi^\alpha\partial_i\partial_jh^\alpha(0).
	\]
	矩阵对称，对应 \(B\) 的对称性。这里原点处 \(\partial g=0\) 已足以计算联络与 \(B\)，不需要把整个图像坐标识别为法坐标。

	2. 写 \(C(e_i,e_j)=C_{ij}\in\mathbb R^r\)，并取
	\[
		h(x)=\frac12C(x,x)=\frac12\sum_{i,j}C_{ij}x^ix^j.
	\]
	对称性给出 \(\partial_i\partial_jh=C_{ij}\)，且 \(h(0)=0\)、\(\mathrm dh_0=0\)。图像的前 \(n\) 个分量就是 \(x\)，所以处处为浸入；第 1 问表明其原点第二基本形式恰为给定的 \(C\)。

	当 \(n\ge2\)、\(r\ge1\) 时，取
	\[
		h^1(x)=\frac12\bigl((x^1)^2-(x^2)^2\bigr),\qquad
		h^\alpha(x)=0\quad(\alpha>1).
	\]
	记 \(\varepsilon_1\) 为 \(\mathbb R^r\) 的第一个标准基向量。原点处 \(B(e_1,e_1)=(0,\varepsilon_1)\)、\(B(e_2,e_2)=-(0,\varepsilon_1)\)，其他对角分量为零，因此 \(H(0)=0\)，但 \(B_0\ne0\)。这显示一点处平均曲率为零允许不同切向方向的法向弯曲相互抵消。
\end{solution}

\begin{exercise}\label{g3:ex:normal-frame-change}
	设 \(f:(M,g)\to(\bar M,\bar g)\) 为光滑等距浸入，在开集 \(U\subset M\) 上取正交单位法标架 \(\nu=(\nu_1,\ldots,\nu_r)\)，写
	\(\nabla^\perp\nu=\nu\eta\)、\(F^\perp=\mathrm d\eta+\eta\wedge\eta\)。
	\begin{enumerate}
		\item 若 \(\nu'=\nu Q\)，其中 \(Q:U\to\operatorname{O}(r)\) 光滑，证明
		\[
			\eta'=Q^{-1}\eta Q+Q^{-1}\mathrm dQ,\qquad
			(F^\perp)'=Q^{-1}F^\perp Q.
		\]
		\item 设法丛秩为二，在定向正交单位法标架中写
		\[
			\eta=\begin{pmatrix}0&-\alpha\\\alpha&0\end{pmatrix}.
		\]
		若 \(Q\) 为角度 \(\phi\) 的转动矩阵，求 \(\alpha'\)。进一步，若 \(U\) 是二维坐标矩形，证明在 \(U\) 上存在由这种转动得到的法平行标架，当且仅当 \(\mathrm d\alpha=0\)。充分性中直接构造 \(\phi\)。
	\end{enumerate}
\end{exercise}
\begin{solution}
	1. 联络的乘积规则给出
	\[
		\nabla^\perp(\nu Q)=\nu(\eta Q+\mathrm dQ)
		=\nu Q\bigl(Q^{-1}\eta Q+Q^{-1}\mathrm dQ\bigr),
	\]
	所以括号中就是新标架的联络矩阵。曲率对法向输入为函数线性，故
	\[
		R^\perp(X,Y)(\nu Q)=\nu F^\perp(X,Y)Q
		=\nu Q\bigl(Q^{-1}F^\perp(X,Y)Q\bigr).
	\]
	这证明曲率矩阵按共轭变换，并由此再见 \(R^\perp\) 是否为零与标架无关。

	2. 记
	\[
		J=\begin{pmatrix}0&-1\\1&0\end{pmatrix},\qquad
		Q=\begin{pmatrix}\cos\phi&-\sin\phi\\
			\sin\phi&\cos\phi\end{pmatrix}.
	\]
	\(JQ=QJ\)，且 \(Q^{-1}\mathrm dQ=J\,\mathrm d\phi\)，所以
	\[
		\eta'=J(\alpha+\mathrm d\phi),\qquad
		\alpha'=\alpha+\mathrm d\phi.
	\]
	又因 \(\alpha\wedge\alpha=0\)，有 \(F^\perp=J\,\mathrm d\alpha\)。若转动后标架法平行，则 \(\alpha'=0\)，从而 \(\mathrm d\alpha=0\)。

	反过来，在坐标矩形上写 \(\alpha=a(x,y)\,\mathrm dx+b(x,y)\,\mathrm dy\)，并固定其中一点 \((x_0,y_0)\)。闭性意味着 \(\partial_xb=\partial_ya\)。令
	\[
		\phi(x,y)
		=-\int_{x_0}^x a(t,y_0)\,\mathrm dt
			-\int_{y_0}^y b(x,s)\,\mathrm ds.
	\]
	矩形保证两段积分道路都在坐标域内。直接求导，
	\[
		\partial_y\phi=-b(x,y),\qquad
		\partial_x\phi=-a(x,y_0)
			-\int_{y_0}^y\partial_xb(x,s)\,\mathrm ds=-a(x,y).
	\]
	故 \(\mathrm d\phi=-\alpha\)，从而 \(\alpha'=0\)，所得正交单位法标架处处法平行。矩形条件在这里提供了单值的原函数构造；只有一点处曲率为零，并不足以使这一构造适用于整个邻域。
\end{solution}

\begin{exercise}\label{g3:ex:helix-normal-parallel-frame}
	设 \(a>0\)、\(b\ne0\)、\(c=\sqrt{a^2+b^2}\)，考虑 Euclid 空间中的单位速曲线
	\[
		\gamma(s)=\left(a\cos\frac sc,\ a\sin\frac sc,\ \frac bc s\right),
		\qquad s\in\mathbb R.
	\]
	记 \(T=\gamma'\)，并定义
	\[
		\nu_1(s)=\left(-\cos\frac sc,-\sin\frac sc,0\right),\qquad
		\nu_2(s)=\left(\frac bc\sin\frac sc,-\frac bc\cos\frac sc,\frac ac\right).
	\]
	\begin{enumerate}
		\item 验证 \((T,\nu_1,\nu_2)\) 为正交单位标架，求第二基本形式、平均曲率向量及法联络矩阵。
		\item 构造整个参数直线上的正交单位法平行标架。说明法平行不等于作为 \(\mathbb R^3\) 中的向量保持不变。
		\item 证明此曲线的法联络曲率为零，但其平均曲率向量不法平行，虽然 \(|H|\) 为常数。
	\end{enumerate}
\end{exercise}
\begin{solution}
	1. 求导得到
	\[
		T=\left(-\frac ac\sin\frac sc,\frac ac\cos\frac sc,\frac bc\right).
	\]
	由 \(a^2+b^2=c^2\)，三个向量的长度都为一，互相内积为零。令
	\[
		\kappa=\frac a{c^2},\qquad \tau=\frac b{c^2}.
	\]
	直接求导并在此标架中分解，得到
	\[
		T'=\kappa\nu_1,\qquad
		\nu_1'=-\kappa T+\tau\nu_2,\qquad
		\nu_2'=-\tau\nu_1.
	\]
	参数直线取诱导度量 \(\mathrm ds^2\)。于是
	\[
		B(\partial_s,\partial_s)=H=\kappa\nu_1,\qquad
		\nabla^\perp_{\partial_s}\nu_1=\tau\nu_2,\qquad
		\nabla^\perp_{\partial_s}\nu_2=-\tau\nu_1.
	\]
	行标架 \(\nu=(\nu_1,\nu_2)\) 的法联络矩阵为
	\[
		\eta=\begin{pmatrix}0&-\tau\\\tau&0\end{pmatrix}\mathrm ds.
	\]

	2. 取转动角 \(\phi(s)=-\tau s\)，并令
	\[
		\begin{aligned}
			E_1&=\cos(\tau s)\nu_1-\sin(\tau s)\nu_2,\\
			E_2&=\sin(\tau s)\nu_1+\cos(\tau s)\nu_2.
		\end{aligned}
	\]
	这是处处正交的转动，所以 \((E_1,E_2)\) 仍为正交单位法标架。利用第 1 问的导数式，各个法向项逐对抵消，得到
	\[
		E_1'=-\kappa\cos(\tau s)T,\qquad
		E_2'=-\kappa\sin(\tau s)T.
	\]
	两个导数都切向，故 \(\nabla^\perp_{\partial_s}E_1=
	\nabla^\perp_{\partial_s}E_2=0\)。但 \(E_1'(0)=-\kappa T(0)\ne0\)，所以法平行向量场仍可以作为环境向量改变；改变的部分恰好全部切向。

	3. 法曲率对两个切向输入反对称，一维切空间中任意两个输入线性相关，故 \(R^\perp=0\)。在给定标架中也可直接计算：\(\mathrm d\eta=0\)，且 \(\eta\wedge\eta\) 含因子 \(\mathrm ds\wedge\mathrm ds=0\)。另一方面，
	\[
		|H|=\kappa,\qquad
		\nabla^\perp_{\partial_s}H
		=\kappa\tau\nu_2\ne0,
	\]
	因为 \(\kappa>0\)、\(\tau\ne0\)。法联络平坦、某个法向场长度恒定、该法向场法平行，是三个不同的条件。
\end{solution}

% 第三卷第四章习题草稿：一般余维收缩、极小浸入、二维法曲率。
% 合入章文件后由章内环境连续编号。

\begin{exercise}\label{g3:ex:submanifold-ricci-scalar-contractions}
	设 \(F\colon(M^n,g)\to\bar M^{n+q}\) 为到常截面曲率 \(c\) 空间的光滑等距浸入，\(n\ge2\)。令 \(B\) 为第二基本形式，\(H=\operatorname{tr}_gB/n\)，并对任一正交单位基 \((e_i)\) 定义
	\[
		\lVert B\rVert^2=\sum_{i,j}\lVert B(e_i,e_j)\rVert^2,
		\qquad B^\circ(X,Y)=B(X,Y)-g(X,Y)H.
	\]
	由 Gauss 方程求 \(\operatorname{Ric}\) 与 \(\operatorname{Scal}\) 的一般余维表达式，证明
	\[
		\operatorname{Scal}\le n(n-1)(c+\lVert H\rVert^2),
	\]
	并用 \(B^\circ\) 描述逐点取等的充要条件。最后说明最小维数 \(n=2\) 时所得不等式的几何含义。
\end{exercise}
\begin{solution}
	固定一点，并在该点取正交单位基 \((e_i)\)。按照本卷的曲率次序，\(\operatorname{Ric}(X,Y)=\sum_i\mathrm{Rm}(e_i,X,e_i,Y)\)。对 \eqref{g3:eq:gauss-space-form} 收缩，环境项为
	\[
		c\sum_i\bigl(g(X,Y)-g(e_i,Y)g(X,e_i)\bigr)
		=(n-1)c\,g(X,Y).
	\]
	第二基本形式的两个项分别给出
	\[
		\sum_i\langle B(e_i,e_i),B(X,Y)\rangle
		=n\langle H,B(X,Y)\rangle,
	\]
	以及 \(-\sum_i\langle B(X,e_i),B(Y,e_i)\rangle\)。因此
	\[
		\operatorname{Ric}(X,Y)
		=(n-1)c\,g(X,Y)+n\langle H,B(X,Y)\rangle
		-\sum_i\langle B(X,e_i),B(Y,e_i)\rangle.
	\]
	再取迹得到
	\[
		\operatorname{Scal}
		=n(n-1)c+n^2\lVert H\rVert^2-\lVert B\rVert^2.
	\]
	以上均为张量的度量收缩，所以不依赖所取正交单位基。

	展开无迹部分的平方，利用 \(\sum_iB(e_i,e_i)=nH\)，有
	\[
		\begin{aligned}
		\lVert B^\circ\rVert^2
		&=\sum_{i,j}\lVert B(e_i,e_j)-\delta_{ij}H\rVert^2\\
		&=\lVert B\rVert^2-2\sum_i\langle B(e_i,e_i),H\rangle
		+n\lVert H\rVert^2\\
		&=\lVert B\rVert^2-n\lVert H\rVert^2.
		\end{aligned}
	\]
	于是精确的差值是
	\[
		\operatorname{Scal}
		=n(n-1)(c+\lVert H\rVert^2)-\lVert B^\circ\rVert^2.
	\]
	范数平方非负给出题设不等式，且某点取等当且仅当 \(B^\circ=0\)，即该点满足 \(B(X,Y)=g(X,Y)H\)，正是脐点条件。当 \(n=2\) 时，\(\operatorname{Scal}=2K\)，所以
	\[
		K=c+\lVert H\rVert^2-\tfrac12\lVert B^\circ\rVert^2
		\le c+\lVert H\rVert^2.
	\]
	因此在环境曲率和平均曲率向量给定时，偏离脐性的部分精确降低了内在 Gauss 曲率。
\end{solution}

\begin{exercise}\label{g3:ex:minimal-euclidean-ricci-nullity}
	设 \(F\colon(M^n,g)\to\mathbb R^{n+q}\) 为光滑极小等距浸入，即 \(H=0\)。证明 \(\operatorname{Ric}\) 为半负定，并求出其作为双线性型的核
	\[
		\ker\operatorname{Ric}_p
		=\{v\in\mathrm{T}_pM:\operatorname{Ric}_p(v,w)=0
		\text{ 对每个 }w\in\mathrm{T}_pM\}.
	\]
	进一步证明 \(\operatorname{Ric}_p=0\) 当且仅当 \(B_p=0\)，并推出：若一个极小 Euclid 浸入的诱导度量平直，则它全测地。解释为何删除极小性后最后一个结论不再成立。
\end{exercise}
\begin{solution}
	由上一题的 Ricci 收缩公式，令 \(c=0\)、\(H=0\)，便有
	\[
		\operatorname{Ric}_p(v,w)
		=-\sum_i\langle B_p(v,e_i),B_p(w,e_i)\rangle.
	\]
	特别地，\(\operatorname{Ric}_p(v,v)=-\sum_i\lVert B_p(v,e_i)\rVert^2\le0\)。如果 \(v\) 属于 Ricci 型的核，取 \(w=v\) 可知每个 \(B_p(v,e_i)\) 都为零，故由线性性 \(B_p(v,w)=0\) 对任意 \(w\) 成立。反过来，若这一条件成立，所显示的双线性公式立即给出 \(v\in\ker\operatorname{Ric}_p\)。因此
	\[
		\ker\operatorname{Ric}_p
		=\{v\in\mathrm{T}_pM:B_p(v,w)=0\text{ 对所有 }w\in\mathrm{T}_pM\}.
	\]
	若 \(\operatorname{Ric}_p=0\)，上述核就是整个切空间，所以 \(B_p=0\)；反向蕴含也由同一公式得到。

	若诱导度量平直，则每点的 Ricci 张量为零，因而 \(B=0\)。由 \cref{g3:prop:totally-geodesic-characterization}，浸入全测地。这里用到了极小性：例如半径为 \(a>0\) 的圆柱
	\[
		F(s,t)=(a\cos(s/a),a\sin(s/a),t)
	\]
	诱导度量为 \(\mathrm ds^2+\mathrm dt^2\)，但采用外法向时，\(B(\partial_s,\partial_s)=-\nu/a\)、\(B(\partial_t,\partial_t)=B(\partial_s,\partial_t)=0\)，所以 \(B\ne0\) 且 \(H=-\nu/(2a)\ne0\)。平直性本身不排除这种弯曲。
\end{solution}

\begin{exercise}\label{g3:ex:minimal-surface-normal-curvature-bound}
	设 \(F\colon M^2\to\mathbb R^4\) 为光滑极小等距浸入，在一点 \(p\) 选取切空间和法空间的正交单位基 \((e_1,e_2)\)、\((\nu_1,\nu_2)\)。由于极小性，可写
	\[
		A_{\nu_1}=\begin{pmatrix}a&b\\b&-a\end{pmatrix},
		\qquad
		A_{\nu_2}=\begin{pmatrix}c&d\\d&-c\end{pmatrix}.
	\]
	记 \(K\) 为内在 Gauss 曲率，并相对于所选两组基定义
	\[
		K^\perp=\langle R^\perp(e_1,e_2)\nu_1,\nu_2\rangle.
	\]
	计算 \(K\) 与 \(K^\perp\)，证明 \(\lvert K^\perp\rvert\le-K\)，并给出取等条件。将等号条件解释为法平面内曲线
	\[
		\theta\longmapsto
		B(\cos\theta\,e_1+\sin\theta\,e_2,
		  \cos\theta\,e_1+\sin\theta\,e_2)
	\]
	的几何性质，允许该曲线退化为一点。
\end{exercise}
\begin{solution}
	把法平面用所选基识别为 \(\mathbb R^2\)，则
	\[
		u\coloneqq B(e_1,e_1)=(a,c),\qquad
		v\coloneqq B(e_1,e_2)=(b,d),\qquad
		B(e_2,e_2)=-u.
	\]
	Gauss 方程给出
	\[
		K=-\lVert u\rVert^2-\lVert v\rVert^2
		=-(a^2+b^2+c^2+d^2).
	\]
	Ricci 方程则给出 \(K^\perp=g([A_{\nu_1},A_{\nu_2}]e_1,e_2)\)。直接相乘有
	\[
		[A_{\nu_1},A_{\nu_2}]
		=\begin{pmatrix}0&2(ad-bc)\\2(bc-ad)&0\end{pmatrix},
	\]
	所以在本卷的符号约定下，
	\[
		K^\perp=2(bc-ad)=-2\det(u,v).
	\]
	反转任一组基的定向会改变 \(K^\perp\) 的符号，但不影响下面的绝对值不等式。利用平面行列式的估计和平方差，有
	\[
		\lvert K^\perp\rvert
		=2\lvert\det(u,v)\rvert
		\le2\lVert u\rVert\lVert v\rVert
		\le\lVert u\rVert^2+\lVert v\rVert^2=-K.
	\]
	两步同时取等当且仅当 \(u\perp v\) 且 \(\lVert u\rVert=\lVert v\rVert\)，其中包括 \(u=v=0\)。用矩阵分量写为
	\[
		ab+cd=0,\qquad a^2+c^2=b^2+d^2.
	\]
	最后利用双线性和 \(B(e_2,e_2)=-u\)，题设曲线为
	\[
		u\cos(2\theta)+v\sin(2\theta).
	\]
	它是单位圆在以 \(u,v\) 为列向量的线性映射下的像，通常为椭圆。若 \(u,v\) 正交且等长，则其像为以原点为圆心的圆；两者都为零时退化为原点。反过来，若此像是半径 \(\rho>0\) 的圆，则在 \(2\theta=0,\pi/2,\pi/4\) 处比较到原点的距离，先得 \(\lVert u\rVert=\lVert v\rVert=\rho\)，再得 \(\langle u,v\rangle=0\)。退化为一点时，其中心为原点且 \(u=v=0\)。故题设不等式取等恰好对应这一个圆形或点形的情形。
\end{solution}

\begin{exercise}\label{g3:ex:normal-data-gauge-equivalence}
	设 \((M,g)\) 为无边界光滑黎曼流形，\(E,E'\to M\) 为同秩光滑 Euclid 向量丛，\(J:E\to E'\) 是覆盖恒等映射的光滑纤维等距。给定 \(E\) 上的度量联络 \(\nabla^E\) 与对称型 \(B\)，在 \(E'\) 上定义
	\[
		B'(X,Y)=J B(X,Y),\qquad
		\nabla^{E'}_X\xi'=J\bigl(\nabla^E_X(J^{-1}\xi')\bigr).
	\]
	\begin{enumerate}
		\item 证明新旧数据同时满足给定 \(c\) 的三组相容方程，并求两个数据联络之间的关系。
		\item 若 \((f,\iota)\) 实现原数据，写出实现新数据的同一个浸入的法丛识别。说明为什么旋转法标架时不能只旋转 \(B\) 的分量而保持联络矩阵不变。
	\end{enumerate}
\end{exercise}
\begin{solution}
	1. 由定义及 \(J\) 保持内积，\(\nabla^{E'}\) 满足联络的乘积规则并保持新内积。对 \(\xi\in\Gamma(E)\)，形算子满足
	\[
		g(A'_{J\xi}X,Y)
		=\langle JB(X,Y),J\xi\rangle_{E'}
		=g(A_\xi X,Y),
	\]
	所以 \(A'_{J\xi}=A_\xi\)。连续使用两次联络定义，得到
	\[
		R^{E'}(X,Y)(J\xi)=J(R^E(X,Y)\xi),
		\qquad(\nabla B')(X;Y,Z)=J((\nabla B)(X;Y,Z)).
	\]
	Gauss 式中的内积不变，Codazzi 的两边同乘 \(J\)，Ricci 式的两边也由上述曲率与形算子关系保持。因此相容性等价。

	令 \(L=\operatorname{id}_{\mathrm TM}\oplus J\)。对式\eqref{g3:eq:submanifold-data-connection}逐项代入，得到
	\[
		D'_X(L\sigma)=L(D_X\sigma).
	\]
	若 \(c\ne0\)，把 \(L\) 扩充为在 \(s\) 上取恒等，同样有 \(\widehat D'_XL=L\widehat D_X\)。所以两个数据联络只是经同一个纤维等距识别后的同一联络；平坦性及沿闭路径的返回条件同时成立。

	2. 取 \(\iota'=\iota\circ J^{-1}\)。则
	\[
		\iota'B'=\iota B=B_f,
		\qquad
		\nabla^{\perp,f}_X(\iota'\xi')
		=\iota'(\nabla^{E'}_X\xi'),
	\]
	故 \((f,\iota')\) 实现新数据。

	再用矩阵说明旋转标架的影响。设 \(\nu\) 是同一个法丛中的行正交标架，\(\nabla^E\nu=\nu\omega\)，并换为 \(\nu'=\nu O\)，其中 \(O:M\to\operatorname O(r)\) 光滑。由乘积规则，
	\[
		\omega'=O^{-1}\omega O+O^{-1}\mathrm dO.
	\]
	若 \(B=\nu b\)，其新分量为 \(b'=O^{-1}b\)。额外的 \(O^{-1}\mathrm dO\) 记录法标架自身的变化，必须与 \(B\) 分量的变化同时计入。例如在区间上，原联络矩阵为零，取
	\[
		O(s)=\begin{pmatrix}\cos s&-\sin s\\ \sin s&\cos s\end{pmatrix},
	\]
	则新矩阵为 \(\left(\begin{smallmatrix}0&-1\\1&0\end{smallmatrix}\right)\mathrm ds\)，并不为零。把它强行保持为零，会换掉法联络，因而改变给定的几何数据。
\end{solution}

\begin{exercise}\label{g3:ex:fundamental-theorem-smooth-families}
	设 \(U_0\subset\mathbb R^n\) 是含原点的开集，\(\lambda\) 属于 \(0\) 的有限维参数邻域。固定 \(c\in\mathbb R\)，在 \(U_0\) 上给定一族数据 \((g_\lambda,E,\langle\ ,\ \rangle_\lambda,\nabla^{E,\lambda},B_\lambda)\)，其中 \(E=U_0\times\mathbb R^r\)，所有局部分量对 \((x,\lambda)\) 联合光滑，各纤维内积正定，每组数据都满足曲率 \(c\) 的相容方程。
	\begin{enumerate}
		\item 在原点给定光滑依赖 \(\lambda\) 的初始线性等距：\(c=0\) 时从 \((\mathrm TM\oplus E)_0\) 到 \(\mathbb R^{n+r}\)，\(c\ne0\) 时从增广纤维 \((\mathrm TM\oplus E\oplus\mathbb Rs)_0\) 到相应的线性环境空间。Euclid 情形再给定光滑的初始位置；\(c<0\) 时要求初始等距将 \(s_0\) 送到最后一坐标为正的单位向量。证明缩小坐标域与参数域后，可以选择实现 \((f_\lambda,\iota_\lambda)\) 对 \((x,\lambda)\) 联合光滑，并具有所给初值。
		\item 若 \(\lambda\) 为实参数，且 \(g_\lambda=g\) 不变，令 \(V=\left.\partial_\lambda f_\lambda\right|_{\lambda=0}\)。证明
		\[
			\langle\nabla^{Q}_X V,\mathrm df_0(Y)\rangle
			+\langle\mathrm df_0(X),\nabla^{Q}_Y V\rangle=0.
		\]
		将 \(V=\mathrm df_0(Z)+\eta\) 分成切向与法向后，把这一式子写成 \(Z,\eta,B_{f_0}\) 的关系。
	\end{enumerate}
\end{exercise}
\begin{solution}
	1. 因度量与纤维内积对参数光滑且在原点正定，可在共同小邻域中对固定基作逐次正交归一化；每一步只涉及光滑运算与正数开平方，故所得标架对 \((x,\lambda)\) 联合光滑。相应的 Levi-Civita 联络系数由度量、逆矩阵与一阶导数给出，故也联合光滑。于是式\eqref{g3:eq:submanifold-data-connection}或式\eqref{g3:eq:augmented-submanifold-connection}的所有系数均联合光滑。

	在\cref{g3:lem:flat-connection-local-frame}的逐轴构造中，将 \(\lambda\) 与尚未行进的坐标一起作为参数。常微分方程的光滑依赖给出共同小盒与参数邻域中的光滑解矩阵。构造只有有限个坐标步骤，因此有限次缩小后，所有步骤都有同一个定义域。给定的初始线性等距也光滑依赖参数，所以最终 \(\Phi_\lambda\) 联合光滑。

	若 \(c=0\)，则 \(\alpha_\lambda(X)=\Phi_\lambda(X)\) 联合光滑，在固定区间 \([0,1]\) 上使用式\eqref{g3:eq:star-shaped-one-form-primitive}积分，给出联合光滑的 \(f_\lambda\)。若 \(c\ne0\)，直接取 \(f_\lambda=a^{-1}\Phi_\lambda(s)\)；双曲情形的初始同构取向指定片，缩小后仍在该片内。两种情形都取 \(\iota_\lambda=\Phi_\lambda|_E\)，从而满足全部要求。唯一性由初始标架及位置已经固定得到。

	2. 对 \(f_\lambda^*g_Q=g\) 求参数导数。环境联络保持度量，并且无挠性给出
	\[
		\left.\nabla_{\partial_\lambda}^{Q}\bigl(\mathrm df_\lambda(X)\bigr)\right|_0
		=\nabla_X^{Q}V.
	\]
	这个交换也可在环境坐标中直接检查：二阶普通偏导交换，联络修正项中的系数关于两个下标对称。度量右侧不随参数变化，于是得到题述等式。

	由 Gauss 与 Weingarten 公式，\(\nabla_X^Q V\) 的切向部分为
	\[
		\mathrm df_0(\nabla_XZ-A_\eta X).
	\]
	代入后得到
	\[
		g(\nabla_XZ,Y)+g(X,\nabla_YZ)
		=2\langle B_{f_0}(X,Y),\eta\rangle.
	\]
	这给出保持诱导度量的一阶必要条件：切向变化的对称导数恰好补偿法向变化通过第二基本形式产生的度量变化。此处结论来自已经存在的光滑实现族，并没有断言每个满足这一线性方程的 \(V\) 都能积分成等距浸入族。
\end{solution}

\begin{exercise}\label{g3:ex:closed-curve-data-monodromy}
	设 \(M=\mathbb R/L\mathbb Z\)，\(L>0\)，取度量 \(\mathrm ds^2\) 及单位切向场 \(e=\partial_s\)。令 \(E=M\times\mathbb R^2\) 有全局正交单位标架 \(\nu_1,\nu_2\)，给定常数 \(\kappa>0\)、\(\tau\in\mathbb R\)，规定
	\[
		B(e,e)=\kappa\nu_1,\qquad
		\nabla^E_e\nu_1=\tau\nu_2,\qquad
		\nabla^E_e\nu_2=-\tau\nu_1.
	\]
	\begin{enumerate}
		\item 验证这些数据满足 \(c=0\) 的全部相容方程，并求数据联络 \(D\) 在标架 \((e,\nu_1,\nu_2)\) 中的矩阵。
		\item 令 \(q=\sqrt{\kappa^2+\tau^2}\)。验证参数直线上的曲线
		\[
			f(s)=\left(\frac{\kappa}{q^2}\cos(qs),
			\frac{\kappa}{q^2}\sin(qs),\frac{\tau}{q}s\right)
		\]
		实现提升后的数据。求这些数据在 \(M\) 上存在带所给法丛识别的整体浸入 \(M\to\mathbb R^3\) 的充要条件。
		\item 证明线性标架沿一周返回原值当且仅当 \(qL\in2\pi\mathbb Z\)。说明此条件单独成立为何还不能保证位置闭合。
	\end{enumerate}
\end{exercise}
\begin{solution}
	1. 在一维上，曲率二形式为零，Gauss 与 Codazzi 两边关于最前两个方向反对称，因而均为零。形算子在一维切空间上都是数乘，彼此交换，所以 Ricci 方程也成立。\(A_{\nu_1}e=\kappa e\)、\(A_{\nu_2}e=0\)，数据联络因此满足
	\[
		D_ee=\kappa\nu_1,\quad
		D_e\nu_1=-\kappa e+\tau\nu_2,\quad
		D_e\nu_2=-\tau\nu_1.
	\]
	按行标架、列分量约定，其矩阵为 \(\Omega\,\mathrm ds\)，其中
	\[
		\Omega=\begin{pmatrix}
			0&-\kappa&0\\
			\kappa&0&-\tau\\
			0&\tau&0
		\end{pmatrix}.
	\]

	2. 直接求导，记
	\[
		\begin{aligned}
		T&=f'=\left(-\frac\kappa q\sin(qs),\frac\kappa q\cos(qs),\frac\tau q\right),\\
		N_1&=(-\cos(qs),-\sin(qs),0),\\
		N_2&=\left(\frac\tau q\sin(qs),-\frac\tau q\cos(qs),\frac\kappa q\right).
		\end{aligned}
	\]
	这三向量是正交单位基，并且
	\[
		T'=\kappa N_1,\qquad
		N_1'=-\kappa T+\tau N_2,\qquad
		N_2'=-\tau N_1.
	\]
	取 \(\iota(\nu_i)=N_i\)，即可逐项读出所给的第二基本形式及法联络，故该曲线实现数据。

	若存在圆周上的整体实现，将它沿 \(\mathbb R\to\mathbb R/L\mathbb Z\) 提升到参数直线。\cref{g3:prop:submanifold-fundamental-uniqueness}说明它与上述 \(f\) 只差一个刚性变换，因此必须有 \(f(s+L)=f(s)\)。第三坐标之差为 \(\tau L/q\)，所以必有 \(\tau=0\)。此时 \(q=\kappa\)，前两个坐标周期要求 \(\kappa L\in2\pi\mathbb Z\)。反之，这两个条件成立时，\(f,T,N_1,N_2\) 都为 \(L\) 周期，故曲线和法丛识别同时下降到圆周。因此充要条件为
	\[
		\tau=0,\qquad \kappa L=2\pi m\quad\text{对某个整数 }m\geq1.
	\]
	当 \(m>1\) 时，所得浸入多次绕过同一个圆，其成立并不要求单射。

	3. \(\Omega\) 的核由 \((\tau,0,\kappa)\) 张成；在其正交补上，直接计算得到 \(\Omega^2=-q^2I\)。因此指数 \(\exp(-L\Omega)\) 在核上为恒等，在正交补上是转角大小为 \(qL\) 的平面旋转。这正是平行方程 \(u'=-\Omega u\) 沿一周的解矩阵，故返回为恒等恰当且仅当 \(qL\in2\pi\mathbb Z\)。

	若满足这一条件而 \(\tau\ne0\)，则三向量 \(T,N_1,N_2\) 都已经返回，但
	\[
		f(s+L)-f(s)=\left(0,0,\frac{\tau L}{q}\right)\ne0.
	\]
	标架的旋转障碍已经消失，位置仍有非零平移周期。这个差向量正是式\eqref{g3:eq:euclidean-development-periods}中的闭路径积分，说明必须同时检查两个返回条件。
\end{solution}

\begin{exercise}\label{g3:ex:prescribed-cylinder-bending}
	给定光滑函数 \(k:\mathbb R\to\mathbb R\)，构造光滑等距浸入 \(f:(\mathbb R^2,\mathrm ds^2+\mathrm dz^2)\to\mathbb R^3\)，使它在某个全局单位法向下的形算子为 \(\operatorname{diag}(k(s),0)\)。证明诱导度量完备，并说明当 \(k\) 非常数时，这个浸入为何在内在平坦的同时具有变化的外在曲率。
\end{exercise}
\begin{solution}
	令
	\[
		\theta(s)=\int_0^s k(\tau)\,\mathrm d\tau,\qquad
		\gamma(s)=\left(\int_0^s\cos\theta(\tau)\,\mathrm d\tau,
		\int_0^s\sin\theta(\tau)\,\mathrm d\tau\right).
	\]
	于是 \(\gamma'=(\cos\theta,\sin\theta)\) 为单位向量，且
	\[
		\gamma''=k(s)(-\sin\theta,\cos\theta).
	\]
	取 \(f(s,z)=(\gamma(s),z)\) 及
	\[
		\nu(s,z)=(-\sin\theta(s),\cos\theta(s),0).
	\]
	两个坐标导数长度均为一且彼此正交，所以诱导度量就是 \(\mathrm ds^2+\mathrm dz^2\)。二阶导数只有 \(f_{ss}=k\nu\) 非零，故 \(B(\partial_s,\partial_s)=k\nu\)、其余分量为零。等价地，\(\partial_s\nu=-k\partial_sf\)，由 Weingarten 公式得到所求形算子。

	Euclid 平面的 Cauchy 序列逐坐标收敛，故诱导距离完备。Gauss 曲率恒为两个主曲率之积 \(k(s)\cdot0=0\)，而平均曲率向量为 \(H=k(s)\nu/2\)。非常数的连续函数 \(k\) 不可能具有恒定的绝对值，因此 \(|B|=|k(s)|\) 也非常数。内在平坦与外在弯曲的变化同时出现。该构造给出浸入；曲线 \(\gamma\) 的可能自交不会破坏已验证的微分单射性或参数流形上的完备性。
\end{solution}

\begin{exercise}\label{g3:ex:two-sided-tube-volume}
	设 \(M^n\subset\mathbb R^{n+1}\) 为紧、无边界、带全局单位法向 \(\nu\) 的光滑嵌入超曲面，\(n\ge2\)。对充分小的 \(\varepsilon>0\)，令
	\[
		\mathcal U_\varepsilon=\{p+t\nu(p):p\in M,\ |t|<\varepsilon\}.
	\]
	证明其 Euclid 体积满足
	\[
		\operatorname{Vol}_{n+1}(\mathcal U_\varepsilon)
		=2\varepsilon\operatorname{Vol}_g(M)
		+\frac{\varepsilon^3}{3}\int_M\operatorname{Scal}_g\,\mu_g
		+O(\varepsilon^5).
	\]
	证明当 \(n=2\) 或 \(n=3\) 时，上式没有余项。以 \(\mathbb R^3\) 中的半径 \(a\) 球面作直接核验。
\end{exercise}
\begin{solution}
	由\cref{g3:prop:euclidean-normal-tube}，足够小的法向参数给出整个管状邻域的单射坐标。令 \(A=A_\nu\)。由于 \(M\) 紧，缩小 \(\varepsilon\) 后所有 \(I-tA\) 都正定，行列式为正。\(t\) 方向的速度是 \(\nu\)，与所有切向速度正交且长度为一，所以管状坐标中的体积密度为
	\[
		\det(I-tA)\,\mu_g\,\mathrm dt.
	\]
	在一点按主曲率展开，
	\[
		\det(I-tA)=1-t\sum_i\kappa_i
		+t^2\sum_{i<j}\kappa_i\kappa_j+\cdots+(-t)^n\prod_i\kappa_i.
	\]
	在对称区间上积分时，所有奇次幂都消失；二次项系数由 Euclid 超曲面的 Gauss 方程满足
	\[
		\sum_{i<j}\kappa_i\kappa_j
		=\frac12\operatorname{Scal}_g.
	\]
	因此二次项积分恰为 \(\varepsilon^3\operatorname{Scal}_g/3\)。其余非零偶次项至少四次，系数是 \(A\) 的连续对称多项式，在紧 \(M\) 上有界，积分总量为 \(O(\varepsilon^5)\)。当 \(n=2,3\) 时，多项式没有四次或更高的偶次项，所以结果精确。

	对半径 \(a\) 的二维球面，\(\operatorname{Vol}_g(M)=4\pi a^2\)、\(\operatorname{Scal}=2/a^2\)。公式给出
	\[
		8\pi a^2\varepsilon+\frac{8\pi}{3}\varepsilon^3.
	\]
	直接减去半径 \(a-\varepsilon\) 与 \(a+\varepsilon\) 的球体体积，同样得到
	\[
		\frac{4\pi}{3}\bigl((a+\varepsilon)^3-(a-\varepsilon)^3\bigr)
		=8\pi a^2\varepsilon+\frac{8\pi}{3}\varepsilon^3.
	\]
\end{solution}

\begin{exercise}\label{g3:ex:ambient-conformal-submanifold}
	设 \(f:M^n\to(\bar M,\bar g)\) 是光滑浸入，\(n\ge1\)，环境度量变为 \(\widetilde{\bar g}=e^{2u}\bar g\)。将切、法子丛按同一个向量空间识别，令 \(g=f^*\bar g\)。求新第二基本形式、平均曲率向量及法联络与原数据的关系。余维一时再求单位法向 \(\widetilde\nu=e^{-u}\nu\) 对应的形算子。最后，在 \(\mathbb R^{n+1}\setminus\{0\}\) 取 \(u=-\log|x|\)，判断所有以原点为中心的球面在新环境度量下是否全测地。
\end{exercise}
\begin{solution}
	共形变化保持正交关系，因而法子丛及切法投影不变。由第一章的共形联络公式\eqref{g3:eq:conformal-connection}，
	\[
		\widetilde{\bar\nabla}_XV
		=\bar\nabla_XV+X(u)V+V(u)\mathrm df(X)
		-\bar g(\mathrm df(X),V)\operatorname{grad}_{\bar g}u.
	\]
	这里 \(X(u)\) 表示沿 \(f\) 的切向导数。令 \(V=\mathrm df(Y)\) 并取法向分量，得到
	\[
		\widetilde B(X,Y)
		=B(X,Y)-g(X,Y)(\operatorname{grad}_{\bar g}u)^\perp.
	\]
	诱导度量为 \(\widetilde g=e^{2u\circ f}g\)，取平均迹时还须乘以 \(e^{-2u\circ f}\)，故
	\[
		\widetilde H=e^{-2u\circ f}
		\bigl(H-(\operatorname{grad}_{\bar g}u)^\perp\bigr).
	\]
	若 \(V=\xi\) 为法向量场，则正交性使最后一项为零，中间的 \(\xi(u)\mathrm df(X)\) 为切向量，因而
	\[
		\widetilde\nabla^\perp_X\xi
		=\nabla^\perp_X\xi+X(u)\xi.
	\]
	该式中的同一个 \(\xi\) 没有经过单位化。对于余维一的单位法向，形算子的定义及上式给出
	\[
		\widetilde A_{\widetilde\nu}
		=e^{-u\circ f}\bigl(A_\nu-\nu(u)I\bigr).
	\]
	例如在半径 \(a\) 的球面上，外法向满足 \(A_\nu=-I/a\)。取 \(u=-\log|x|\)，有 \(\nu(u)=-1/a\)，故新形算子为零，球面全测地。也可从环境度量直接核对：令 \(r=|x|\)、\(s=\log r\)，则
	\[
		|x|^{-2}g_{\mathrm{Euc}}
		=\frac{\mathrm dr^2}{r^2}+g_{S^n(1)}
		=\mathrm ds^2+g_{S^n(1)}.
	\]
	原来的同心球面变成乘积中的等 \(s\) 切片，正与计算相符。
\end{solution}

% !TeX root = ../../Book3.tex
% 纲要标题：## 第5章　完整群、平行结构与黎曼分解
% 纲要位置：几何学计划纲要.md，第 2157 行。
% 正文状态：未撰写；本文件目前是框架。

% 纲要内容开始
% 
% > 本章目标与证明范围：从平行移动得到完整群，研究其如何决定平行张量和局部分解。de Rham 的整体版本保留单连通及完备性；特殊完整群分类单列引用而不倒推基础结果。
% 
% 纲要内容结束

\chapter{完整群、平行结构与黎曼分解}
\label{g3:ch:05}
\BookChapterNumber{g3:ch:05}

\begin{chapterabstract}
待撰写
\end{chapterabstract}

\begin{learninggoals}
待撰写
\end{learninggoals}

待撰写

% !TeX root = ../../Book3.tex
% 纲要标题：### 5.1 完整群的李群结构与曲率
% 纲要位置：几何学计划纲要.md，第 2264 行。
% 正文状态：未撰写；本文件目前是框架。

% 纲要内容开始
% 
% * 平行运输的路径群
% * 受限完整群和整体完整群
% * 局部平行截面及浸入李群结构
% * 小矩形运输的二阶展开
% * 沿路径搬运的曲率算子
% * Ambrose–Singer 生成结论及其证明
% 
% 纲要内容结束

\section{完整群的李群结构与曲率}
\label{g3:sec:0501}

待撰写

\subsection{平行运输的路径群}
\label{g3:subsec:0501:01}
待撰写

\subsection{受限完整群和整体完整群}
\label{g3:subsec:0501:02}
待撰写

\subsection{局部平行截面及浸入李群结构}
\label{g3:subsec:0501:03}
待撰写

\subsection{小矩形运输的二阶展开}
\label{g3:subsec:0501:04}
待撰写

\subsection{沿路径搬运的曲率算子}
\label{g3:subsec:0501:05}
待撰写

\subsection{Ambrose–Singer 生成结论及其证明}
\label{g3:subsec:0501:06}
待撰写
% 纲要条目排为编号小节标题，并进入目录与书签；说明段仍保留于纲要注释中。

% !TeX root = ../../Book3.tex
% 纲要标题：### 5.2 平行张量与分布
% 纲要位置：几何学计划纲要.md，第 2273 行。
% 正文状态：未撰写；本文件目前是框架。

% 纲要内容开始
% 
% * 不变向量与平行截面
% * 不变子空间和正交分布
% * 平行分布的可积性
% * 局部乘积的构造
% 
% 纲要内容结束

\section{平行张量与分布}
\label{g3:sec:0502}

待撰写

\subsection{不变向量与平行截面}
\label{g3:subsec:0502:01}
待撰写

\subsection{不变子空间和正交分布}
\label{g3:subsec:0502:02}
待撰写

\subsection{平行分布的可积性}
\label{g3:subsec:0502:03}
待撰写

\subsection{局部乘积的构造}
\label{g3:subsec:0502:04}
待撰写
% 纲要条目排为编号小节标题，并进入目录与书签；说明段仍保留于纲要注释中。

% !TeX root = ../../Book3.tex
% 纲要标题：### 5.3 de Rham 分解
% 纲要位置：几何学计划纲要.md，第 2280 行。
% 正文状态：未撰写；本文件目前是框架。

% 纲要内容开始
% 
% * 局部分解的存在
% * 单连通完备情形的整体乘积
% * 因子的不可约性与唯一性
% * 一般商空间中的扭曲现象
% 
% 纲要内容结束

\section{de Rham 分解}
\label{g3:sec:0503}

待撰写

\subsection{局部分解的存在}
\label{g3:subsec:0503:01}
待撰写

\subsection{单连通完备情形的整体乘积}
\label{g3:subsec:0503:02}
待撰写

\subsection{因子的不可约性与唯一性}
\label{g3:subsec:0503:03}
待撰写

\subsection{一般商空间中的扭曲现象}
\label{g3:subsec:0503:04}
待撰写
% 纲要条目排为编号小节标题，并进入目录与书签；说明段仍保留于纲要注释中。

% !TeX root = ../../Book3.tex
% 纲要标题：### 5.4 齐性与对称空间
% 纲要位置：几何学计划纲要.md，第 2287 行。
% 正文状态：未撰写；本文件目前是框架。
% 规划任务：保留现有齐性度量、点反射、对称分解与曲率模型，补充最大平坦、秩及紧非紧对偶的构造。完整核验球面、双曲空间与一个经典矩阵商的相应几何；局部 Lie 代数对偶与整体覆盖、离散商分别处理，不据局部对应推断任意整体商相同。

% 纲要内容开始
% 
% > 补充写作范围与外部输入：规划任务：保留现有齐性度量、点反射、对称分解与曲率模型，补充最大平坦、秩及紧非紧对偶的构造。完整核验球面、双曲空间与一个经典矩阵商的相应几何；局部 Lie 代数对偶与整体覆盖、离散商分别处理，不据局部对应推断任意整体商相同。
% 
% * 齐性度量及不变联络
% * 对称空间的点反射
% * Lie 代数分解与曲率
% * 紧型、非紧型和欧氏模型
% 
% * 极大交换子空间与最大平坦
% * 对称空间的秩及乘积结构
% * 紧型与非紧型的对偶
% 
% 
% 纲要内容结束

\section{齐性与对称空间}
\label{g3:sec:0504}

待撰写

\subsection{齐性度量及不变联络}
\label{g3:subsec:0504:01}
待撰写

\subsection{对称空间的点反射}
\label{g3:subsec:0504:02}
待撰写

\subsection{Lie 代数分解与曲率}
\label{g3:subsec:0504:03}
待撰写

\subsection{紧型、非紧型和欧氏模型}
\label{g3:subsec:0504:04}
待撰写

\subsection{极大交换子空间与最大平坦}
\label{g3:subsec:0504:05}
待撰写

\subsection{对称空间的秩及乘积结构}
\label{g3:subsec:0504:06}
待撰写

\subsection{紧型与非紧型的对偶}
\label{g3:subsec:0504:07}
待撰写

% 纲要条目排为编号小节标题，并进入目录与书签；说明段仍保留于纲要注释中。

% !TeX root = ../../Book3.tex
% 纲要标题：### 5.5* 特殊完整群的几何内容
% 纲要位置：几何学计划纲要.md，第 2301 行。
% 正文状态：未撰写；本文件目前是框架。

% 纲要内容开始
% 
% * Kähler、超 Kähler 与平行形式
% * G₂、Spin(7) 的代数模型
% * 分类结论的条件与外部输入
% 
% ---
% 
% 纲要内容结束

\OptionalSection{特殊完整群的几何内容}{g3:sec:0505}

待撰写

\subsection{Kähler、超 Kähler 与平行形式}
\label{g3:subsec:0505:01}
待撰写

\subsection{\texorpdfstring{\(G_2\)、Spin(7) 的代数模型}{G₂、Spin(7) 的代数模型}}
\label{g3:subsec:0505:02}
待撰写

\subsection{分类结论的条件与外部输入}
\label{g3:subsec:0505:03}
待撰写
% 纲要条目排为编号小节标题，并进入目录与书签；说明段仍保留于纲要注释中。


\begin{chaptersummary}
待撰写
\end{chaptersummary}

\BookExercises
待撰写


% !TeX root = ../Book3.tex
% 纲要标题：# 第二编　变分、曲率比较与度量空间
% 纲要位置：几何学计划纲要.md，第 2199 行。
% 正文状态：未撰写；本文件目前是框架。

% 纲要内容开始
% 
% 纲要内容结束

\BookPart{变分、曲率比较与度量空间}{g3:part:02}
\BookPartGuide
待撰写

% !TeX root = ../../Book3.tex
% 纲要标题：## 第6章　测地线变分、Jacobi 场与指标定理
% 纲要位置：几何学计划纲要.md，第 2324 行。
% 正文状态：未撰写；本文件目前是框架。

% 纲要内容开始
% 
% > 本章目标与证明范围：逐项导出端点项和指标形式，证明共轭点与二次变分的关系。路径空间语言只引入证明实际需要的部分，不以无限维直觉代替边值问题。
% 
% 纲要内容结束

\chapter{测地线变分、Jacobi 场与指标定理}
\label{g3:ch:06}
\BookChapterNumber{g3:ch:06}

\begin{chapterabstract}
待撰写
\end{chapterabstract}

\begin{learninggoals}
待撰写
\end{learninggoals}

待撰写

% !TeX root = ../../Book3.tex
% 纲要标题：### 6.1 长度与能量的变分
% 纲要位置：几何学计划纲要.md，第 2205 行。
% 正文状态：未撰写；本文件目前是框架。

% 纲要内容开始
% 
% * 变分场和协变导数交换
% * 第一变分及端点项
% * 第二变分和曲率项
% * 固定端点与自由端点的区别
% 
% 纲要内容结束

\section{长度与能量的变分}
\label{g3:sec:0601}

待撰写

\subsection{变分场和协变导数交换}
\label{g3:subsec:0601:01}
待撰写

\subsection{第一变分及端点项}
\label{g3:subsec:0601:02}
待撰写

\subsection{第二变分和曲率项}
\label{g3:subsec:0601:03}
待撰写

\subsection{固定端点与自由端点的区别}
\label{g3:subsec:0601:04}
待撰写
% 纲要条目排为编号小节标题，并进入目录与书签；说明段仍保留于纲要注释中。

% !TeX root = ../../Book3.tex
% 纲要标题：### 6.2 Jacobi 场与指数映射
% 纲要位置：几何学计划纲要.md，第 2335 行。
% 正文状态：未撰写；本文件目前是框架。

% 纲要内容开始
% 
% * Jacobi 方程的推导
% * 初值、Wronskian 配对与解空间
% * 指数映射的微分
% * 共轭点及其重数
% 
% 纲要内容结束

\section{Jacobi 场与指数映射}
\label{g3:sec:0602}

待撰写

\subsection{Jacobi 方程的推导}
\label{g3:subsec:0602:01}
待撰写

\subsection{初值、Wronskian 配对与解空间}
\label{g3:subsec:0602:02}
待撰写

\subsection{指数映射的微分}
\label{g3:subsec:0602:03}
待撰写

\subsection{共轭点及其重数}
\label{g3:subsec:0602:04}
待撰写
% 纲要条目排为编号小节标题，并进入目录与书签；说明段仍保留于纲要注释中。

% !TeX root = ../../Book3.tex
% 纲要标题：### 6.3 指标形式和最短性
% 纲要位置：几何学计划纲要.md，第 2342 行。
% 正文状态：未撰写；本文件目前是框架。

% 纲要内容开始
% 
% * 分段光滑场与指标形式
% * Jacobi 场的极小性质
% * 共轭点前后的符号变化
% * 割点与共轭点的关系
% 
% 纲要内容结束

\section{指标形式和最短性}
\label{g3:sec:0603}

待撰写

\subsection{分段光滑场与指标形式}
\label{g3:subsec:0603:01}
待撰写

\subsection{Jacobi 场的极小性质}
\label{g3:subsec:0603:02}
待撰写

\subsection{共轭点前后的符号变化}
\label{g3:subsec:0603:03}
待撰写

\subsection{割点与共轭点的关系}
\label{g3:subsec:0603:04}
待撰写
% 纲要条目排为编号小节标题，并进入目录与书签；说明段仍保留于纲要注释中。

% !TeX root = ../../Book3.tex
% 纲要标题：### 6.4 Morse 指标定理
% 纲要位置：几何学计划纲要.md，第 2226 行。
% 正文状态：未撰写；本文件目前是框架。

% 纲要内容开始
% 
% * 有限维分段测地线约化
% * 共轭时刻的指标跳跃
% * 指标等于共轭点重数之和
% * 端点退化与零空间
% * 焦点和子流形边界条件
% 
% ---
% 
% 纲要内容结束

\section{Morse 指标定理}
\label{g3:sec:0604}

待撰写

\subsection{有限维分段测地线约化}
\label{g3:subsec:0604:01}
待撰写

\subsection{共轭时刻的指标跳跃}
\label{g3:subsec:0604:02}
待撰写

\subsection{指标等于共轭点重数之和}
\label{g3:subsec:0604:03}
待撰写

\subsection{端点退化与零空间}
\label{g3:subsec:0604:04}
待撰写

\subsection{焦点和子流形边界条件}
\label{g3:subsec:0604:05}
待撰写
% 纲要条目排为编号小节标题，并进入目录与书签；说明段仍保留于纲要注释中。


\begin{chaptersummary}
待撰写
\end{chaptersummary}

\BookExercises
待撰写

% !TeX root = ../../Book3.tex
% 纲要标题：## 第7章　Rauch、Riccati 与距离函数比较
% 纲要位置：几何学计划纲要.md，第 2359 行。
% 正文状态：未撰写；本文件目前是框架。

% 纲要内容开始
% 
% > 本章目标与证明范围：比较定理按共同初值、有效区间、不等式方向与等号条件组织。由 Jacobi 场逐步转为 Hessian、Laplace 和体积密度，而非并列罗列四个定理。
% 
% 纲要内容结束

\chapter{Rauch、Riccati 与距离函数比较}
\label{g3:ch:07}
\BookChapterNumber{g3:ch:07}

\begin{chapterabstract}
待撰写
\end{chapterabstract}

\begin{learninggoals}
待撰写
\end{learninggoals}

待撰写

% !TeX root = ../../Book3.tex
% 纲要标题：### 7.1 指标形式的比较方法
% 纲要位置：几何学计划纲要.md，第 2363 行。
% 正文状态：未撰写；本文件目前是框架。

% 纲要内容开始
% 
% * 模型 Jacobi 场
% * 相同端点条件与变分极小性
% * Rauch 比较的证明
% * 共轭半径和微分长度估计
% 
% 纲要内容结束

\section{指标形式的比较方法}
\label{g3:sec:0701}

待撰写

\subsection{模型 Jacobi 场}
\label{g3:subsec:0701:01}
待撰写

\subsection{相同端点条件与变分极小性}
\label{g3:subsec:0701:02}
待撰写

\subsection{Rauch 比较的证明}
\label{g3:subsec:0701:03}
待撰写

\subsection{共轭半径和微分长度估计}
\label{g3:subsec:0701:04}
待撰写
% 纲要条目排为编号小节标题，并进入目录与书签；说明段仍保留于纲要注释中。

% !TeX root = ../../Book3.tex
% 纲要标题：### 7.2 矩阵 Riccati 方程
% 纲要位置：几何学计划纲要.md，第 2370 行。
% 正文状态：未撰写；本文件目前是框架。

% 纲要内容开始
% 
% * 径向形算子及奇异初值
% * 矩阵与标量微分不等式
% * 迹估计和 Cauchy–Schwarz 的等号
% * 比较解的爆破与有效半径
% 
% 纲要内容结束

\section{矩阵 Riccati 方程}
\label{g3:sec:0702}

待撰写

\subsection{径向形算子及奇异初值}
\label{g3:subsec:0702:01}
待撰写

\subsection{矩阵与标量微分不等式}
\label{g3:subsec:0702:02}
待撰写

\subsection{迹估计和 Cauchy–Schwarz 的等号}
\label{g3:subsec:0702:03}
待撰写

\subsection{比较解的爆破与有效半径}
\label{g3:subsec:0702:04}
待撰写
% 纲要条目排为编号小节标题，并进入目录与书签；说明段仍保留于纲要注释中。

% !TeX root = ../../Book3.tex
% 纲要标题：### 7.3 Hessian 和 Laplace 比较
% 纲要位置：几何学计划纲要.md，第 2377 行。
% 正文状态：未撰写；本文件目前是框架。

% 纲要内容开始
% 
% * 距离函数的径向 Hessian
% * 截面界与 Hessian 比较
% * Ricci 界与迹比较
% * 本书非负 Laplace 符号的换算
% 
% 纲要内容结束

\section{Hessian 和 Laplace 比较}
\label{g3:sec:0703}

待撰写

\subsection{距离函数的径向 Hessian}
\label{g3:subsec:0703:01}
待撰写

\subsection{截面界与 Hessian 比较}
\label{g3:subsec:0703:02}
待撰写

\subsection{Ricci 界与迹比较}
\label{g3:subsec:0703:03}
待撰写

\subsection{本书非负 Laplace 符号的换算}
\label{g3:subsec:0703:04}
待撰写
% 纲要条目排为编号小节标题，并进入目录与书签；说明段仍保留于纲要注释中。

% !TeX root = ../../Book3.tex
% 纲要标题：### 7.4 割迹与刚性
% 纲要位置：几何学计划纲要.md，第 2384 行。
% 正文状态：未撰写；本文件目前是框架。

% 纲要内容开始
% 
% * 割迹之外的光滑结论
% * 支撑函数及屏障解释
% * 必要的分布版本所需工具
% * 等号传播及局部空间形式
% 
% ---
% 
% 纲要内容结束

\section{割迹与刚性}
\label{g3:sec:0704}

待撰写

\subsection{割迹之外的光滑结论}
\label{g3:subsec:0704:01}
待撰写

\subsection{支撑函数及屏障解释}
\label{g3:subsec:0704:02}
待撰写

\subsection{必要的分布版本所需工具}
\label{g3:subsec:0704:03}
待撰写

\subsection{等号传播及局部空间形式}
\label{g3:subsec:0704:04}
待撰写
% 纲要条目排为编号小节标题，并进入目录与书签；说明段仍保留于纲要注释中。


\begin{chaptersummary}
待撰写
\end{chaptersummary}

\BookExercises
待撰写

% !TeX root = ../../Book3.tex
% 纲要标题：## 第8章　Ricci 曲率、体积与整体限制
% 纲要位置：几何学计划纲要.md，第 2393 行。
% 正文状态：未撰写；本文件目前是框架。

% 纲要内容开始
% 
% > 本章目标与证明范围：完整建立直径和体积比较及其覆盖应用。Cheeger–Gromoll 分裂的证明移到已具备弱最大值和 Bochner 工具的位置，不在这里隐藏分析先修。
% 
% 纲要内容结束

\chapter{Ricci 曲率、体积与整体限制}
\label{g3:ch:08}
\BookChapterNumber{g3:ch:08}

\begin{chapterabstract}
待撰写
\end{chapterabstract}

\begin{learninggoals}
待撰写
\end{learninggoals}

待撰写

% !TeX root = ../../Book3.tex
% 纲要标题：### 8.1 Bonnet–Myers 定理
% 纲要位置：几何学计划纲要.md，第 2397 行。
% 正文状态：未撰写；本文件目前是框架。

% 纲要内容开始
% 
% * 测试变分场与指标估计
% * 直径界和紧性
% * 万有覆盖上的应用
% * 基本群有限性
% 
% 纲要内容结束

\section{Bonnet–Myers 定理}
\label{g3:sec:0801}

待撰写

\subsection{测试变分场与指标估计}
\label{g3:subsec:0801:01}
待撰写

\subsection{直径界和紧性}
\label{g3:subsec:0801:02}
待撰写

\subsection{万有覆盖上的应用}
\label{g3:subsec:0801:03}
待撰写

\subsection{基本群有限性}
\label{g3:subsec:0801:04}
待撰写
% 纲要条目排为编号小节标题，并进入目录与书签；说明段仍保留于纲要注释中。

% !TeX root = ../../Book3.tex
% 纲要标题：### 8.2 Bishop–Gromov 体积比较
% 纲要位置：几何学计划纲要.md，第 2404 行。
% 正文状态：未撰写；本文件目前是框架。

% 纲要内容开始
% 
% * 极坐标 Jacobian 和割点截断
% * 径向密度单调性
% * 球体积比及球壳比较
% * 等号与刚性
% 
% 纲要内容结束

\section{Bishop–Gromov 体积比较}
\label{g3:sec:0802}

待撰写

\subsection{极坐标 Jacobian 和割点截断}
\label{g3:subsec:0802:01}
待撰写

\subsection{径向密度单调性}
\label{g3:subsec:0802:02}
待撰写

\subsection{球体积比及球壳比较}
\label{g3:subsec:0802:03}
待撰写

\subsection{等号与刚性}
\label{g3:subsec:0802:04}
待撰写
% 纲要条目排为编号小节标题，并进入目录与书签；说明段仍保留于纲要注释中。

% !TeX root = ../../Book3.tex
% 纲要标题：### 8.3 增长、覆盖和基本群
% 纲要位置：几何学计划纲要.md，第 2288 行。
% 正文状态：未撰写；本文件目前是框架。

% 纲要内容开始
% 
% * 非紧体积增长
% * 覆盖空间中的体积计数
% * 有限生成子群的增长约束
% * 欧氏、双曲及乘积模型
% 
% 纲要内容结束

\section{增长、覆盖和基本群}
\label{g3:sec:0803}

待撰写

\subsection{非紧体积增长}
\label{g3:subsec:0803:01}
待撰写

\subsection{覆盖空间中的体积计数}
\label{g3:subsec:0803:02}
待撰写

\subsection{有限生成子群的增长约束}
\label{g3:subsec:0803:03}
待撰写

\subsection{欧氏、双曲及乘积模型}
\label{g3:subsec:0803:04}
待撰写
% 纲要条目排为编号小节标题，并进入目录与书签；说明段仍保留于纲要注释中。

% !TeX root = ../../Book3.tex
% 纲要标题：### 8.4 局部比较的极限
% 纲要位置：几何学计划纲要.md，第 2295 行。
% 正文状态：未撰写；本文件目前是框架。

% 纲要内容开始
% 
% * 非负 Ricci 与非负截面曲率的区别
% * 塌缩时体积信息的变化
% * 分裂定理所需的 Busemann 函数问题
% 
% ---
% 
% 纲要内容结束

\section{局部比较的极限}
\label{g3:sec:0804}

待撰写

\subsection{非负 Ricci 与非负截面曲率的区别}
\label{g3:subsec:0804:01}
待撰写

\subsection{塌缩时体积信息的变化}
\label{g3:subsec:0804:02}
待撰写

\subsection{分裂定理所需的 Busemann 函数问题}
\label{g3:subsec:0804:03}
待撰写
% 纲要条目排为编号小节标题，并进入目录与书签；说明段仍保留于纲要注释中。


\begin{chaptersummary}
待撰写
\end{chaptersummary}

\BookExercises
待撰写

% !TeX root = ../../Book3.tex
% 纲要标题：## 第9章　截面曲率、空间形式与 Toponogov
% 纲要位置：几何学计划纲要.md，第 2426 行。
% 正文状态：未撰写；本文件目前是框架。

% 纲要内容开始
% 
% > 本章目标与证明范围：把上界和下界的不同三角几何分开。Cartan–Hadamard 与空间形式分类给完整证明；Toponogov 的选定全局版本有独立的局部到整体论证。
% 
% 纲要内容结束

\chapter{截面曲率、空间形式与 Toponogov}
\label{g3:ch:09}
\BookChapterNumber{g3:ch:09}

\begin{chapterabstract}
待撰写
\end{chapterabstract}

\begin{learninggoals}
待撰写
\end{learninggoals}

待撰写

% !TeX root = ../../Book3.tex
% 纲要标题：### 9.1 非正曲率与 Cartan–Hadamard
% 纲要位置：几何学计划纲要.md，第 2430 行。
% 正文状态：未撰写；本文件目前是框架。

% 纲要内容开始
% 
% * 无共轭点和指数映射的局部微分性质
% * 覆盖论证及单连通条件
% * 全局唯一测地线
% * 距离函数凸性和拓扑后果
% 
% 纲要内容结束

\section{非正曲率与 Cartan–Hadamard}
\label{g3:sec:0901}

待撰写

\subsection{无共轭点和指数映射的局部微分性质}
\label{g3:subsec:0901:01}
待撰写

\subsection{覆盖论证及单连通条件}
\label{g3:subsec:0901:02}
待撰写

\subsection{全局唯一测地线}
\label{g3:subsec:0901:03}
待撰写

\subsection{距离函数凸性和拓扑后果}
\label{g3:subsec:0901:04}
待撰写
% 纲要条目排为编号小节标题，并进入目录与书签；说明段仍保留于纲要注释中。

% !TeX root = ../../Book3.tex
% 纲要标题：### 9.2 常曲率的局部与整体分类
% 纲要位置：几何学计划纲要.md，第 2314 行。
% 正文状态：未撰写；本文件目前是框架。

% 纲要内容开始
% 
% * Jacobi 方程确定局部度量
% * 单连通完备模型
% * 离散自由等距群商
% * 球面、平坦与双曲空间形式
% 
% 纲要内容结束

\section{常曲率的局部与整体分类}
\label{g3:sec:0902}

待撰写

\subsection{Jacobi 方程确定局部度量}
\label{g3:subsec:0902:01}
待撰写

\subsection{单连通完备模型}
\label{g3:subsec:0902:02}
待撰写

\subsection{离散自由等距群商}
\label{g3:subsec:0902:03}
待撰写

\subsection{球面、平坦与双曲空间形式}
\label{g3:subsec:0902:04}
待撰写
% 纲要条目排为编号小节标题，并进入目录与书签；说明段仍保留于纲要注释中。

% !TeX root = ../../Book3.tex
% 纲要标题：### 9.3* Toponogov 三角比较
% 纲要位置：几何学计划纲要.md，第 2444 行。
% 正文状态：未撰写；本文件目前是框架。

% 纲要内容开始
% 
% * 铰链比较与模型角
% * 局部比较及拼接引理
% * 全局三角比较的证明
% * 正曲率模型中的周长限制
% * 等号和刚性情形
% 
% 纲要内容结束

\OptionalSection{Toponogov 三角比较}{g3:sec:0903}

待撰写

\subsection{铰链比较与模型角}
\label{g3:subsec:0903:01}
待撰写

\subsection{局部比较及拼接引理}
\label{g3:subsec:0903:02}
待撰写

\subsection{全局三角比较的证明}
\label{g3:subsec:0903:03}
待撰写

\subsection{正曲率模型中的周长限制}
\label{g3:subsec:0903:04}
待撰写

\subsection{等号和刚性情形}
\label{g3:subsec:0903:05}
待撰写
% 纲要条目排为编号小节标题，并进入目录与书签；说明段仍保留于纲要注释中。

% !TeX root = ../../Book3.tex
% 纲要标题：### 9.4 三角比较的几何后果
% 纲要位置：几何学计划纲要.md，第 2329 行。
% 正文状态：未撰写；本文件目前是框架。

% 纲要内容开始
% 
% * 凸集、测地三角形和投影
% * 球面型直径约束
% * CAT 上界与 Alexandrov 下界的模型对照
% 
% ---
% 
% 纲要内容结束

\section{三角比较的几何后果}
\label{g3:sec:0904}

待撰写

\subsection{凸集、测地三角形和投影}
\label{g3:subsec:0904:01}
待撰写

\subsection{球面型直径约束}
\label{g3:subsec:0904:02}
待撰写

\subsection{CAT 上界与 Alexandrov 下界的模型对照}
\label{g3:subsec:0904:03}
待撰写
% 纲要条目排为编号小节标题，并进入目录与书签；说明段仍保留于纲要注释中。


\begin{chaptersummary}
待撰写
\end{chaptersummary}

\BookExercises
待撰写

% !TeX root = ../../Book3.tex
% 纲要标题：## 第10章　长度空间、Gromov–Hausdorff 距离与预紧性
% 纲要位置：几何学计划纲要.md，第 2460 行。
% 正文状态：未撰写；本文件目前是框架。

% 纲要内容开始
% 
% > 本章目标与证明范围：完整建立度量极限工具，先以可计算退化模型解释收敛，再证明覆盖数控制的预紧性。一般 GH 收敛不蕴含拓扑或光滑结构稳定。
% 
% 纲要内容结束

\chapter{长度空间、Gromov–Hausdorff 距离与预紧性}
\label{g3:ch:10}
\BookChapterNumber{g3:ch:10}

\begin{chapterabstract}
待撰写
\end{chapterabstract}

\begin{learninggoals}
待撰写
\end{learninggoals}

待撰写

% !TeX root = ../../Book3.tex
% 纲要标题：### 10.1 长度空间与测地空间
% 纲要位置：几何学计划纲要.md，第 2341 行。
% 正文状态：未撰写；本文件目前是框架。

% 纲要内容开始
% 
% * 曲线长度和内在度量
% * 完备、局部紧及测地存在
% * 锥、商和粘合空间
% 
% 纲要内容结束

\section{长度空间与测地空间}
\label{g3:sec:1001}

待撰写

\subsection{曲线长度和内在度量}
\label{g3:subsec:1001:01}
待撰写

\subsection{完备、局部紧及测地存在}
\label{g3:subsec:1001:02}
待撰写

\subsection{锥、商和粘合空间}
\label{g3:subsec:1001:03}
待撰写
% 纲要条目排为编号小节标题，并进入目录与书签；说明段仍保留于纲要注释中。

% !TeX root = ../../Book3.tex
% 纲要标题：### 10.2 GH 距离的构造
% 纲要位置：几何学计划纲要.md，第 2470 行。
% 正文状态：未撰写；本文件目前是框架。

% 纲要内容开始
% 
% * 公共度量空间中的 Hausdorff 距离
% * 对应与畸变
% * ε-逼近映射的等价描述
% * 零距离与等距类
% 
% 纲要内容结束

\section{GH 距离的构造}
\label{g3:sec:1002}

待撰写

\subsection{公共度量空间中的 Hausdorff 距离}
\label{g3:subsec:1002:01}
待撰写

\subsection{对应与畸变}
\label{g3:subsec:1002:02}
待撰写

\subsection{\texorpdfstring{\(\varepsilon\)-逼近映射的等价描述}{\GeometryPdfEpsilon{}-逼近映射的等价描述}}
\label{g3:subsec:1002:03}
待撰写

\subsection{零距离与等距类}
\label{g3:subsec:1002:04}
待撰写
% 纲要条目排为编号小节标题，并进入目录与书签；说明段仍保留于纲要注释中。

% !TeX root = ../../Book3.tex
% 纲要标题：### 10.3 预紧性和极限构造
% 纲要位置：几何学计划纲要.md，第 2477 行。
% 正文状态：未撰写；本文件目前是框架。

% 纲要内容开始
% 
% * 有限网及距离矩阵
% * 对角子列构造极限
% * 统一覆盖数的预紧性判据
% * 完备性与紧性的保持
% 
% 纲要内容结束

\section{预紧性和极限构造}
\label{g3:sec:1003}

待撰写

\subsection{有限网及距离矩阵}
\label{g3:subsec:1003:01}
待撰写

\subsection{对角子列构造极限}
\label{g3:subsec:1003:02}
待撰写

\subsection{统一覆盖数的预紧性判据}
\label{g3:subsec:1003:03}
待撰写

\subsection{完备性与紧性的保持}
\label{g3:subsec:1003:04}
待撰写
% 纲要条目排为编号小节标题，并进入目录与书签；说明段仍保留于纲要注释中。

% !TeX root = ../../Book3.tex
% 纲要标题：### 10.4 非紧极限与切锥
% 纲要位置：几何学计划纲要.md，第 2484 行。
% 正文状态：未撰写；本文件目前是框架。

% 纲要内容开始
% 
% * 带基点收敛和球的控制
% * 重缩放及切锥
% * 扁平环面塌缩
% * 切锥不唯一的概念与反例范围
% 
% 纲要内容结束

\section{非紧极限与切锥}
\label{g3:sec:1004}

待撰写

\subsection{带基点收敛和球的控制}
\label{g3:subsec:1004:01}
待撰写

\subsection{重缩放及切锥}
\label{g3:subsec:1004:02}
待撰写

\subsection{扁平环面塌缩}
\label{g3:subsec:1004:03}
待撰写

\subsection{切锥不唯一的概念与反例范围}
\label{g3:subsec:1004:04}
待撰写
% 纲要条目排为编号小节标题，并进入目录与书签；说明段仍保留于纲要注释中。

% !TeX root = ../../Book3.tex
% 纲要标题：### 10.5* 合成曲率的基本模型
% 纲要位置：几何学计划纲要.md，第 2491 行。
% 正文状态：未撰写；本文件目前是框架。

% 纲要内容开始
% 
% * CAT 空间的比较定义
% * Alexandrov 下界及三角条件
% * 光滑曲率与合成条件的连接
% 
% ---
% 
% 纲要内容结束

\OptionalSection{合成曲率的基本模型}{g3:sec:1005}

待撰写

\subsection{CAT 空间的比较定义}
\label{g3:subsec:1005:01}
待撰写

\subsection{Alexandrov 下界及三角条件}
\label{g3:subsec:1005:02}
待撰写

\subsection{光滑曲率与合成条件的连接}
\label{g3:subsec:1005:03}
待撰写
% 纲要条目排为编号小节标题，并进入目录与书签；说明段仍保留于纲要注释中。


\begin{chaptersummary}
待撰写
\end{chaptersummary}

\BookExercises
待撰写


% !TeX root = ../Book3.tex
% 纲要标题：# 第三编　函数空间、椭圆理论与全局分析
% 纲要位置：几何学计划纲要.md，第 2499 行。
% 正文状态：未撰写；本文件目前是框架。

% 纲要内容开始
% 
% 纲要内容结束

\BookPart{函数空间、椭圆理论与全局分析}{g3:part:03}
\BookPartGuide
待撰写

% !TeX root = ../../Book3.tex
% 纲要标题：## 第11章　Sobolev 空间、紧性与几何不等式
% 纲要位置：几何学计划纲要.md，第 2378 行。
% 正文状态：未撰写；本文件目前是框架。

% 纲要内容开始
% 
% > 本章目标与证明范围：必要的泛函分析和流形上的局部化在本章系统建立。嵌入的指数、丛值版本和几何常数分别说明，为后续估计提供真正可使用的工具。
% 
% 纲要内容结束

\chapter{Sobolev 空间、紧性与几何不等式}
\label{g3:ch:11}
\BookChapterNumber{g3:ch:11}

\begin{chapterabstract}
待撰写
\end{chapterabstract}

\begin{learninggoals}
待撰写
\end{learninggoals}

待撰写

% !TeX root = ../../Book3.tex
% 纲要标题：### 11.1 积分、分布和弱微分
% 纲要位置：几何学计划纲要.md，第 2382 行。
% 正文状态：未撰写；本文件目前是框架。

% 纲要内容开始
% 
% * 体积密度及 Lᵖ 截面
% * 测试截面和分布
% * 弱导数与分部积分
% * 局部平凡化和坐标变换
% 
% 纲要内容结束

\section{积分、分布和弱微分}
\label{g3:sec:1101}

待撰写

\subsection{\texorpdfstring{体积密度及 \(L^p\) 截面}{体积密度及 Lᵖ 截面}}
\label{g3:subsec:1101:01}
待撰写

\subsection{测试截面和分布}
\label{g3:subsec:1101:02}
待撰写

\subsection{弱导数与分部积分}
\label{g3:subsec:1101:03}
待撰写

\subsection{局部平凡化和坐标变换}
\label{g3:subsec:1101:04}
待撰写
% 纲要条目排为编号小节标题，并进入目录与书签；说明段仍保留于纲要注释中。

% !TeX root = ../../Book3.tex
% 纲要标题：### 11.2 Sobolev 截面的构造
% 纲要位置：几何学计划纲要.md，第 2512 行。
% 正文状态：未撰写；本文件目前是框架。

% 纲要内容开始
% 
% * 整数阶范数和等价性
% * 卷积、截断与光滑稠密性
% * 乘法、复合及所需链式法则
% * 弱拓扑、反射性和弱紧性
% 
% 纲要内容结束

\section{Sobolev 截面的构造}
\label{g3:sec:1102}

待撰写

\subsection{整数阶范数和等价性}
\label{g3:subsec:1102:01}
待撰写

\subsection{卷积、截断与光滑稠密性}
\label{g3:subsec:1102:02}
待撰写

\subsection{乘法、复合及所需链式法则}
\label{g3:subsec:1102:03}
待撰写

\subsection{弱拓扑、反射性和弱紧性}
\label{g3:subsec:1102:04}
待撰写
% 纲要条目排为编号小节标题，并进入目录与书签；说明段仍保留于纲要注释中。

% !TeX root = ../../Book3.tex
% 纲要标题：### 11.3 Sobolev 与 Morrey 嵌入
% 纲要位置：几何学计划纲要.md，第 2396 行。
% 正文状态：未撰写；本文件目前是框架。

% 纲要内容开始
% 
% * 欧氏基本积分不等式
% * Sobolev 指数及尺度
% * Morrey 连续性及 Hölder 指数
% * 临界指数的端点限制
% * 紧流形和丛值截面的粘合
% 
% 纲要内容结束

\section{Sobolev 与 Morrey 嵌入}
\label{g3:sec:1103}

待撰写

\subsection{欧氏基本积分不等式}
\label{g3:subsec:1103:01}
待撰写

\subsection{Sobolev 指数及尺度}
\label{g3:subsec:1103:02}
待撰写

\subsection{Morrey 连续性及 Hölder 指数}
\label{g3:subsec:1103:03}
待撰写

\subsection{临界指数的端点限制}
\label{g3:subsec:1103:04}
待撰写

\subsection{紧流形和丛值截面的粘合}
\label{g3:subsec:1103:05}
待撰写
% 纲要条目排为编号小节标题，并进入目录与书签；说明段仍保留于纲要注释中。

% !TeX root = ../../Book3.tex
% 纲要标题：### 11.4 Poincaré 不等式与几何常数
% 纲要位置：几何学计划纲要.md，第 2527 行。
% 正文状态：未撰写；本文件目前是框架。

% 纲要内容开始
% 
% * 零边界和零均值两种版本
% * 连通性及常数核
% * 从局部球到整体的覆盖论证
% * 体积、尺度和几何控制的作用
% 
% 纲要内容结束

\section{Poincaré 不等式与几何常数}
\label{g3:sec:1104}

待撰写

\subsection{零边界和零均值两种版本}
\label{g3:subsec:1104:01}
待撰写

\subsection{连通性及常数核}
\label{g3:subsec:1104:02}
待撰写

\subsection{从局部球到整体的覆盖论证}
\label{g3:subsec:1104:03}
待撰写

\subsection{体积、尺度和几何控制的作用}
\label{g3:subsec:1104:04}
待撰写
% 纲要条目排为编号小节标题，并进入目录与书签；说明段仍保留于纲要注释中。

% !TeX root = ../../Book3.tex
% 纲要标题：### 11.5 Rellich 紧性
% 纲要位置：几何学计划纲要.md，第 2411 行。
% 正文状态：未撰写；本文件目前是框架。

% 纲要内容开始
% 
% * 平移控制与紧性判据
% * 有界域上的紧嵌入
% * 紧流形上的全局紧性
% * 临界指数及非紧空间中的失效
% 
% 纲要内容结束

\section{Rellich 紧性}
\label{g3:sec:1105}

待撰写

\subsection{平移控制与紧性判据}
\label{g3:subsec:1105:01}
待撰写

\subsection{有界域上的紧嵌入}
\label{g3:subsec:1105:02}
待撰写

\subsection{紧流形上的全局紧性}
\label{g3:subsec:1105:03}
待撰写

\subsection{临界指数及非紧空间中的失效}
\label{g3:subsec:1105:04}
待撰写
% 纲要条目排为编号小节标题，并进入目录与书签；说明段仍保留于纲要注释中。

% !TeX root = ../../Book3.tex
% 纲要标题：### 11.6 边界与迹的最低工具
% 纲要位置：几何学计划纲要.md，第 2541 行。
% 正文状态：未撰写；本文件目前是框架。

% 纲要内容开始
% 
% * 边界展平和延拓算子
% * 迹映射的连续性
% * 零边界 Sobolev 空间
% * 边界正则性的作用
% 
% ---
% 
% 纲要内容结束

\section{边界与迹的最低工具}
\label{g3:sec:1106}

待撰写

\subsection{边界展平和延拓算子}
\label{g3:subsec:1106:01}
待撰写

\subsection{迹映射的连续性}
\label{g3:subsec:1106:02}
待撰写

\subsection{零边界 Sobolev 空间}
\label{g3:subsec:1106:03}
待撰写

\subsection{边界正则性的作用}
\label{g3:subsec:1106:04}
待撰写
% 纲要条目排为编号小节标题，并进入目录与书签；说明段仍保留于纲要注释中。


\begin{chaptersummary}
待撰写
\end{chaptersummary}

\BookExercises
待撰写

% !TeX root = ../../Book3.tex
% 纲要标题：## 第12章　弱解与局部椭圆估计
% 纲要位置：几何学计划纲要.md，第 2550 行。
% 正文状态：未撰写；本文件目前是框架。

% 纲要内容开始
% 
% > 本章目标与证明范围：把可解性、内部估计和正则性分开。完整建立后续 Hodge 与 Cauchy–Riemann 模块实际使用的 L² 和所选 Lᵖ 估计；标量最大值原理不得替代系统论证。
% 
% 纲要内容结束

\chapter{弱解与局部椭圆估计}
\label{g3:ch:12}
\BookChapterNumber{g3:ch:12}

\begin{chapterabstract}
待撰写
\end{chapterabstract}

\begin{learninggoals}
待撰写
\end{learninggoals}

待撰写

% !TeX root = ../../Book3.tex
% 纲要标题：### 12.1 能量极小化与弱解
% 纲要位置：几何学计划纲要.md，第 2431 行。
% 正文状态：未撰写；本文件目前是框架。

% 纲要内容开始
% 
% * Hilbert 投影和 Riesz 表示
% * Lax–Milgram 的证明
% * 约束最小化与弱下半连续
% * Poisson 方程的存在与唯一性
% 
% 纲要内容结束

\section{能量极小化与弱解}
\label{g3:sec:1201}

待撰写

\subsection{Hilbert 投影和 Riesz 表示}
\label{g3:subsec:1201:01}
待撰写

\subsection{Lax–Milgram 的证明}
\label{g3:subsec:1201:02}
待撰写

\subsection{约束最小化与弱下半连续}
\label{g3:subsec:1201:03}
待撰写

\subsection{Poisson 方程的存在与唯一性}
\label{g3:subsec:1201:04}
待撰写
% 纲要条目排为编号小节标题，并进入目录与书签；说明段仍保留于纲要注释中。

% !TeX root = ../../Book3.tex
% 纲要标题：### 12.2 主符号和常系数估计
% 纲要位置：几何学计划纲要.md，第 2561 行。
% 正文状态：未撰写；本文件目前是框架。

% 纲要内容开始
% 
% * 微分算子、形式伴随与主符号
% * 椭圆性、强椭圆性及 Laplace 型系统
% * Fourier 变换的最低工具
% * 常系数 L² 估计和微分阶数
% 
% 纲要内容结束

\section{主符号和常系数估计}
\label{g3:sec:1202}

待撰写

\subsection{微分算子、形式伴随与主符号}
\label{g3:subsec:1202:01}
待撰写

\subsection{椭圆性、强椭圆性及 Laplace 型系统}
\label{g3:subsec:1202:02}
待撰写

\subsection{Fourier 变换的最低工具}
\label{g3:subsec:1202:03}
待撰写

\subsection{\texorpdfstring{常系数 \(L^2\) 估计和微分阶数}{常系数 L² 估计和微分阶数}}
\label{g3:subsec:1202:04}
待撰写
% 纲要条目排为编号小节标题，并进入目录与书签；说明段仍保留于纲要注释中。

% !TeX root = ../../Book3.tex
% 纲要标题：### 12.3 变系数估计的局部化
% 纲要位置：几何学计划纲要.md，第 2568 行。
% 正文状态：未撰写；本文件目前是框架。

% 纲要内容开始
% 
% * 冻结系数与小球选择
% * 截断交换子和低阶项吸收
% * 差商方法与弱解的二阶正则性
% * 椭圆系统的同阶估计
% * 自举和分布解的光滑性
% 
% 纲要内容结束

\section{变系数估计的局部化}
\label{g3:sec:1203}

待撰写

\subsection{冻结系数与小球选择}
\label{g3:subsec:1203:01}
待撰写

\subsection{截断交换子和低阶项吸收}
\label{g3:subsec:1203:02}
待撰写

\subsection{差商方法与弱解的二阶正则性}
\label{g3:subsec:1203:03}
待撰写

\subsection{椭圆系统的同阶估计}
\label{g3:subsec:1203:04}
待撰写

\subsection{自举和分布解的光滑性}
\label{g3:subsec:1203:05}
待撰写
% 纲要条目排为编号小节标题，并进入目录与书签；说明段仍保留于纲要注释中。

% !TeX root = ../../Book3.tex
% 纲要标题：### 12.4 Lᵖ 椭圆估计
% 纲要位置：几何学计划纲要.md，第 2576 行。
% 正文状态：未撰写；本文件目前是框架。

% 纲要内容开始
% 
% * 基本解和奇异积分核
% * Calderón–Zygmund 分解及有界性模块
% * 常系数估计与 Fourier 乘子
% * 变系数冻结、截断和吸收
% * 分布解到 W^{m,p} 的正则性
% * 系统和系数光滑性的明确范围
% 
% 纲要内容结束

\section{\texorpdfstring{\(L^p\)}{Lᵖ} 椭圆估计}
\label{g3:sec:1204}

待撰写

\subsection{基本解和奇异积分核}
\label{g3:subsec:1204:01}
待撰写

\subsection{Calderón–Zygmund 分解及有界性模块}
\label{g3:subsec:1204:02}
待撰写

\subsection{常系数估计与 Fourier 乘子}
\label{g3:subsec:1204:03}
待撰写

\subsection{变系数冻结、截断和吸收}
\label{g3:subsec:1204:04}
待撰写

\subsection{\texorpdfstring{分布解到 \(W^{m,p}\) 的正则性}{分布解到 W\textasciicircum{}\{m,p\} 的正则性}}
\label{g3:subsec:1204:05}
待撰写

\subsection{系统和系数光滑性的明确范围}
\label{g3:subsec:1204:06}
待撰写
% 纲要条目排为编号小节标题，并进入目录与书签；说明段仍保留于纲要注释中。

% !TeX root = ../../Book3.tex
% 纲要标题：### 12.5 Schauder 方法
% 纲要位置：几何学计划纲要.md，第 2462 行。
% 正文状态：未撰写；本文件目前是框架。

% 纲要内容开始
% 
% * Hölder 半范数及尺度
% * 常系数的振荡控制
% * 局部插值不等式
% * 变系数估计与高阶自举
% 
% 纲要内容结束

\section{Schauder 方法}
\label{g3:sec:1205}

待撰写

\subsection{Hölder 半范数及尺度}
\label{g3:subsec:1205:01}
待撰写

\subsection{常系数的振荡控制}
\label{g3:subsec:1205:02}
待撰写

\subsection{局部插值不等式}
\label{g3:subsec:1205:03}
待撰写

\subsection{变系数估计与高阶自举}
\label{g3:subsec:1205:04}
待撰写
% 纲要条目排为编号小节标题，并进入目录与书签；说明段仍保留于纲要注释中。

% !TeX root = ../../Book3.tex
% 纲要标题：### 12.6 一阶 Cauchy–Riemann 型算子
% 纲要位置：几何学计划纲要.md，第 2592 行。
% 正文状态：未撰写；本文件目前是框架。

% 纲要内容开始
% 
% * 通过椭圆主符号获得局部 W¹ᵖ 估计
% * p>2 时的连续性及非线性复合
% * 后续闭曲面与条带的不同全局条件
% 
% 纲要内容结束

\section{一阶 Cauchy–Riemann 型算子}
\label{g3:sec:1206}

待撰写

\subsection{\texorpdfstring{通过椭圆主符号获得局部 \(W^{1,p}\) 估计}{通过椭圆主符号获得局部 W¹ᵖ 估计}}
\label{g3:subsec:1206:01}
待撰写

\subsection{\texorpdfstring{\(p>2\) 时的连续性及非线性复合}{p>2 时的连续性及非线性复合}}
\label{g3:subsec:1206:02}
待撰写

\subsection{后续闭曲面与条带的不同全局条件}
\label{g3:subsec:1206:03}
待撰写
% 纲要条目排为编号小节标题，并进入目录与书签；说明段仍保留于纲要注释中。

% !TeX root = ../../Book3.tex
% 纲要标题：### 12.7 标量最大值与 Harnack 方法
% 纲要位置：几何学计划纲要.md，第 2598 行。
% 正文状态：未撰写；本文件目前是框架。

% 纲要内容开始
% 
% * 弱最大值原理及零阶项符号
% * 强最大值原理和屏障
% * Hopf 边界引理的条件
% * Moser 迭代与所选 Harnack 不等式
% * 系统不具备同一结论的原因
% 
% ---
% 
% 纲要内容结束

\section{标量最大值与 Harnack 方法}
\label{g3:sec:1207}

待撰写

\subsection{弱最大值原理及零阶项符号}
\label{g3:subsec:1207:01}
待撰写

\subsection{强最大值原理和屏障}
\label{g3:subsec:1207:02}
待撰写

\subsection{Hopf 边界引理的条件}
\label{g3:subsec:1207:03}
待撰写

\subsection{Moser 迭代与所选 Harnack 不等式}
\label{g3:subsec:1207:04}
待撰写

\subsection{系统不具备同一结论的原因}
\label{g3:subsec:1207:05}
待撰写
% 纲要条目排为编号小节标题，并进入目录与书签；说明段仍保留于纲要注释中。


\begin{chaptersummary}
待撰写
\end{chaptersummary}

\BookExercises
待撰写

% !TeX root = ../../Book3.tex
% 纲要标题：## 第13章　全局椭圆理论与非线性共同工具
% 纲要位置：几何学计划纲要.md，第 2608 行。
% 正文状态：未撰写；本文件目前是框架。

% 纲要内容开始
% 
% > 本章目标与证明范围：从局部估计和紧嵌入推出闭值域、有限核余核，再构造 Green 算子及椭圆复形的调和模型。第13.5—13.6节独立建立非线性共同工具，供 Yamabe、Einstein、几何流、伪全纯方程与形变理论调用；只学线性椭圆与 Hodge 的读者可以先跳过这两节。
% 
% 纲要内容结束

\chapter{全局椭圆理论与非线性共同工具}
\label{g3:ch:13}
\BookChapterNumber{g3:ch:13}

\begin{chapterabstract}
待撰写
\end{chapterabstract}

\begin{learninggoals}
待撰写
\end{learninggoals}

待撰写

% !TeX root = ../../Book3.tex
% 纲要标题：### 13.1 紧流形上的全局估计
% 纲要位置：几何学计划纲要.md，第 2489 行。
% 正文状态：未撰写；本文件目前是框架。

% 纲要内容开始
% 
% * 有限覆盖与局部估计的粘合
% * 含低阶范数的先验不等式
% * 核的有限维性
% * 正交核后加强的估计
% 
% 纲要内容结束

\section{紧流形上的全局估计}
\label{g3:sec:1301}

待撰写

\subsection{有限覆盖与局部估计的粘合}
\label{g3:subsec:1301:01}
待撰写

\subsection{含低阶范数的先验不等式}
\label{g3:subsec:1301:02}
待撰写

\subsection{核的有限维性}
\label{g3:subsec:1301:03}
待撰写

\subsection{正交核后加强的估计}
\label{g3:subsec:1301:04}
待撰写
% 纲要条目排为编号小节标题，并进入目录与书签；说明段仍保留于纲要注释中。

% !TeX root = ../../Book3.tex
% 纲要标题：### 13.2 闭值域与 Fredholm 选择
% 纲要位置：几何学计划纲要.md，第 2496 行。
% 正文状态：未撰写；本文件目前是框架。

% 纲要内容开始
% 
% * 闭值域的矛盾紧性证明
% * 形式伴随的正则性
% * 余核与伴随核
% * 可解条件和 Fredholm 指数
% * 指数对连续扰动的稳定性
% 
% 纲要内容结束

\section{闭值域与 Fredholm 选择}
\label{g3:sec:1302}

待撰写

\subsection{闭值域的矛盾紧性证明}
\label{g3:subsec:1302:01}
待撰写

\subsection{形式伴随的正则性}
\label{g3:subsec:1302:02}
待撰写

\subsection{余核与伴随核}
\label{g3:subsec:1302:03}
待撰写

\subsection{可解条件和 Fredholm 指数}
\label{g3:subsec:1302:04}
待撰写

\subsection{指数对连续扰动的稳定性}
\label{g3:subsec:1302:05}
待撰写
% 纲要条目排为编号小节标题，并进入目录与书签；说明段仍保留于纲要注释中。

% !TeX root = ../../Book3.tex
% 纲要标题：### 13.3 Green 算子和谱
% 纲要位置：几何学计划纲要.md，第 2504 行。
% 正文状态：未撰写；本文件目前是框架。

% 纲要内容开始
% 
% * 广义逆及有限秩投影
% * 自伴非负算子的紧预解式
% * 特征值离散性和光滑特征截面
% * 谱展开的范数与正则性
% 
% 纲要内容结束

\section{Green 算子和谱}
\label{g3:sec:1303}

待撰写

\subsection{广义逆及有限秩投影}
\label{g3:subsec:1303:01}
待撰写

\subsection{自伴非负算子的紧预解式}
\label{g3:subsec:1303:02}
待撰写

\subsection{特征值离散性和光滑特征截面}
\label{g3:subsec:1303:03}
待撰写

\subsection{谱展开的范数与正则性}
\label{g3:subsec:1303:04}
待撰写
% 纲要条目排为编号小节标题，并进入目录与书签；说明段仍保留于纲要注释中。

% !TeX root = ../../Book3.tex
% 纲要标题：### 13.4 椭圆复形
% 纲要位置：几何学计划纲要.md，第 2634 行。
% 正文状态：未撰写；本文件目前是框架。

% 纲要内容开始
% 
% * 符号序列的正合性
% * 相应 Laplace 算子的椭圆性
% * 正合、余正合与调和分解
% * 有限维上同调及代表元
% * de Rham 和 Dolbeault 的共同机制
% 
% 纲要内容结束

\section{椭圆复形}
\label{g3:sec:1304}

待撰写

\subsection{符号序列的正合性}
\label{g3:subsec:1304:01}
待撰写

\subsection{相应 Laplace 算子的椭圆性}
\label{g3:subsec:1304:02}
待撰写

\subsection{正合、余正合与调和分解}
\label{g3:subsec:1304:03}
待撰写

\subsection{有限维上同调及代表元}
\label{g3:subsec:1304:04}
待撰写

\subsection{de Rham 和 Dolbeault 的共同机制}
\label{g3:subsec:1304:05}
待撰写
% 纲要条目排为编号小节标题，并进入目录与书签；说明段仍保留于纲要注释中。

% !TeX root = ../../Book3.tex
% 纲要标题：### 13.5 Banach 微分与非线性隐函数方法
% 纲要位置：几何学计划纲要.md，第 2642 行。
% 正文状态：未撰写；本文件目前是框架。

% 纲要内容开始
% 
% 本节是非线性几何的共同工具，独立于第24章的 Yamabe 专题；逆函数、分裂满射和 Fredholm 约化三种情形分别列明假设。
% 
% * Fréchet 微分、连续可微性与链式法则
% * Banach 空间逆函数与隐函数定理的定量证明
% * 分裂满射的零点集及有限维核的局部参数
% * Fredholm 线性化的核、像和余核分解
% * Lyapunov–Schmidt 有限维约化与障碍映射
% * 压缩映射、二次余项和参数一致性
% 
% 纲要内容结束

\section{Banach 微分与非线性隐函数方法}
\label{g3:sec:1305}

待撰写

\subsection{Fréchet 微分、连续可微性与链式法则}
\label{g3:subsec:1305:01}
待撰写

\subsection{Banach 空间逆函数与隐函数定理的定量证明}
\label{g3:subsec:1305:02}
待撰写

\subsection{分裂满射的零点集及有限维核的局部参数}
\label{g3:subsec:1305:03}
待撰写

\subsection{Fredholm 线性化的核、像和余核分解}
\label{g3:subsec:1305:04}
待撰写

\subsection{Lyapunov–Schmidt 有限维约化与障碍映射}
\label{g3:subsec:1305:05}
待撰写

\subsection{压缩映射、二次余项和参数一致性}
\label{g3:subsec:1305:06}
待撰写
% 纲要条目排为编号小节标题，并进入目录与书签；说明段仍保留于纲要注释中。

% !TeX root = ../../Book3.tex
% 纲要标题：### 13.6 Sobolev 映射空间与非线性算子
% 纲要位置：几何学计划纲要.md，第 2653 行。
% 正文状态：未撰写；本文件目前是框架。

% 纲要内容开始
% 
% 函数空间、微分阶数、系数正则性及参数范围同时固定；不对任意 Sobolev 指数无条件宣称复合算子光滑。
% 
% * 紧 $m$ 维源上的 $W^{k,p}$ 映射空间与 $kp>m$ 的连续嵌入
% * 指数图册、丛值截面及局部坐标转换
% * Sobolev 乘法、复合及微分损失的估计
% * 非线性微分算子的可微性和目标空间选择
% 
%   曲面上的 Cauchy–Riemann 模块固定 $p>2$，使用 $W^{1,p}\to L^p$ 的映射；高阶族与规范群作用另行核验所需指数，不将微分损失隐藏在“光滑作用”中。
% 
% * 椭圆自举、先验界与弱解的正则性恢复
% * 直接方法、约束临界点与规范对称的区别
% 
% ---
% 
% 纲要内容结束

\section{Sobolev 映射空间与非线性算子}
\label{g3:sec:1306}

待撰写

\subsection{\texorpdfstring{紧 \(m\) 维源上的 \(W^{k,p}\) 映射空间与 \(kp>m\) 的连续嵌入}{紧 m 维源上的 W\textasciicircum{}\{k,p\} 映射空间与 kp>m 的连续嵌入}}
\label{g3:subsec:1306:01}
待撰写

\subsection{指数图册、丛值截面及局部坐标转换}
\label{g3:subsec:1306:02}
待撰写

\subsection{Sobolev 乘法、复合及微分损失的估计}
\label{g3:subsec:1306:03}
待撰写

\subsection{非线性微分算子的可微性和目标空间选择}
\label{g3:subsec:1306:04}
待撰写

\subsection{椭圆自举、先验界与弱解的正则性恢复}
\label{g3:subsec:1306:05}
待撰写

\subsection{直接方法、约束临界点与规范对称的区别}
\label{g3:subsec:1306:06}
待撰写
% 纲要条目排为编号小节标题，并进入目录与书签；说明段仍保留于纲要注释中。


\begin{chaptersummary}
待撰写
\end{chaptersummary}

\BookExercises
待撰写

% !TeX root = ../../Book3.tex
% 纲要标题：## 第14章　椭圆边值问题与参数族
% 纲要位置：几何学计划纲要.md，第 2546 行。
% 正文状态：未撰写；本文件目前是框架。

% 纲要内容开始
% 
% > 本章目标与证明范围：把原来一节里的边界、参数和系统拆为有明确对象的整章。主证明限定光滑紧区域上的 Dirichlet、Neumann 及所选丛值问题；一般伪微分边界理论只作准确出口。
% 
% 纲要内容结束

\chapter{椭圆边值问题与参数族}
\label{g3:ch:14}
\BookChapterNumber{g3:ch:14}

\begin{chapterabstract}
待撰写
\end{chapterabstract}

\begin{learninggoals}
待撰写
\end{learninggoals}

待撰写

% !TeX root = ../../Book3.tex
% 纲要标题：### 14.1 边界上的弱问题
% 纲要位置：几何学计划纲要.md，第 2673 行。
% 正文状态：未撰写；本文件目前是框架。

% 纲要内容开始
% 
% * Green 公式和边界配对
% * Dirichlet 变分问题
% * Neumann 兼容条件及常数核
% * 混合条件和角点的额外问题
% 
% 纲要内容结束

\section{边界上的弱问题}
\label{g3:sec:1401}

待撰写

\subsection{Green 公式和边界配对}
\label{g3:subsec:1401:01}
待撰写

\subsection{Dirichlet 变分问题}
\label{g3:subsec:1401:02}
待撰写

\subsection{Neumann 兼容条件及常数核}
\label{g3:subsec:1401:03}
待撰写

\subsection{混合条件和角点的额外问题}
\label{g3:subsec:1401:04}
待撰写
% 纲要条目排为编号小节标题，并进入目录与书签；说明段仍保留于纲要注释中。

% !TeX root = ../../Book3.tex
% 纲要标题：### 14.2 边界正则性
% 纲要位置：几何学计划纲要.md，第 2557 行。
% 正文状态：未撰写；本文件目前是框架。

% 纲要内容开始
% 
% * 边界展平与切向差商
% * 法向导数的方程恢复
% * 齐次及非齐次边界数据
% * 高阶边界估计
% * 不满足边界条件的内部估计不能替代整体结论
% 
% 纲要内容结束

\section{边界正则性}
\label{g3:sec:1402}

待撰写

\subsection{边界展平与切向差商}
\label{g3:subsec:1402:01}
待撰写

\subsection{法向导数的方程恢复}
\label{g3:subsec:1402:02}
待撰写

\subsection{齐次及非齐次边界数据}
\label{g3:subsec:1402:03}
待撰写

\subsection{高阶边界估计}
\label{g3:subsec:1402:04}
待撰写

\subsection{不满足边界条件的内部估计不能替代整体结论}
\label{g3:subsec:1402:05}
待撰写
% 纲要条目排为编号小节标题，并进入目录与书签；说明段仍保留于纲要注释中。

% !TeX root = ../../Book3.tex
% 纲要标题：### 14.3 系统与补条件
% 纲要位置：几何学计划纲要.md，第 2688 行。
% 正文状态：未撰写；本文件目前是框架。

% 纲要内容开始
% 
% * 半空间常系数边界模型
% * 椭圆补条件的含义
% * Hodge 的绝对与相对边界条件
% * Fredholm 边值问题的所选版本
% 
% 纲要内容结束

\section{系统与补条件}
\label{g3:sec:1403}

待撰写

\subsection{半空间常系数边界模型}
\label{g3:subsec:1403:01}
待撰写

\subsection{椭圆补条件的含义}
\label{g3:subsec:1403:02}
待撰写

\subsection{Hodge 的绝对与相对边界条件}
\label{g3:subsec:1403:03}
待撰写

\subsection{Fredholm 边值问题的所选版本}
\label{g3:subsec:1403:04}
待撰写
% 纲要条目排为编号小节标题，并进入目录与书签；说明段仍保留于纲要注释中。

% !TeX root = ../../Book3.tex
% 纲要标题：### 14.4 参数依赖与跳跃
% 纲要位置：几何学计划纲要.md，第 2572 行。
% 正文状态：未撰写；本文件目前是框架。

% 纲要内容开始
% 
% * 固定核维数下的投影和 Green 算子
% * 参数椭圆估计和连续性
% * 核维数跳跃及谱流动机
% * 参数化求解与几何族
% 
% ---
% 
% 纲要内容结束

\section{参数依赖与跳跃}
\label{g3:sec:1404}

待撰写

\subsection{固定核维数下的投影和 Green 算子}
\label{g3:subsec:1404:01}
待撰写

\subsection{参数椭圆估计和连续性}
\label{g3:subsec:1404:02}
待撰写

\subsection{核维数跳跃及谱流动机}
\label{g3:subsec:1404:03}
待撰写

\subsection{参数化求解与几何族}
\label{g3:subsec:1404:04}
待撰写
% 纲要条目排为编号小节标题，并进入目录与书签；说明段仍保留于纲要注释中。


\begin{chaptersummary}
待撰写
\end{chaptersummary}

\BookExercises
待撰写

% !TeX root = ../../Book3.tex
% 纲要标题：## 第15章　Hodge 理论、Bochner 与非紧分裂
% 纲要位置：几何学计划纲要.md，第 2581 行。
% 正文状态：未撰写；本文件目前是框架。

% 纲要内容开始
% 
% > 本章目标与证明范围：闭流形 Hodge 定理在全局椭圆工具之后完整证明。随后将 Weitzenböck 用于曲率消失；非紧分裂独立补齐截断、弱调和与 Busemann 的工具，不照搬紧情形积分。
% 
% 纲要内容结束

\chapter{Hodge 理论、Bochner 与非紧分裂}
\label{g3:ch:15}
\BookChapterNumber{g3:ch:15}

\begin{chapterabstract}
待撰写
\end{chapterabstract}

\begin{learninggoals}
待撰写
\end{learninggoals}

待撰写

% !TeX root = ../../Book3.tex
% 纲要标题：### 15.1 形式上的椭圆结构
% 纲要位置：几何学计划纲要.md，第 2585 行。
% 正文状态：未撰写；本文件目前是框架。

% 纲要内容开始
% 
% * Hodge 星和形式伴随
% * Hodge Laplace 的主符号
% * 能量恒等式及调和条件
% * 定向与密度约定
% 
% 纲要内容结束

\section{形式上的椭圆结构}
\label{g3:sec:1501}

待撰写

\subsection{Hodge 星和形式伴随}
\label{g3:subsec:1501:01}
待撰写

\subsection{Hodge Laplace 的主符号}
\label{g3:subsec:1501:02}
待撰写

\subsection{能量恒等式及调和条件}
\label{g3:subsec:1501:03}
待撰写

\subsection{定向与密度约定}
\label{g3:subsec:1501:04}
待撰写
% 纲要条目排为编号小节标题，并进入目录与书签；说明段仍保留于纲要注释中。

% !TeX root = ../../Book3.tex
% 纲要标题：### 15.2 Hodge 定理的证明
% 纲要位置：几何学计划纲要.md，第 2715 行。
% 正文状态：未撰写；本文件目前是框架。

% 纲要内容开始
% 
% * 正合、余正合及调和分解
% * 每个 de Rham 类的调和代表
% * 唯一性、有限维性和自然配对
% * Poincaré 对偶的解析再解释
% 
% 纲要内容结束

\section{Hodge 定理的证明}
\label{g3:sec:1502}

待撰写

\subsection{正合、余正合及调和分解}
\label{g3:subsec:1502:01}
待撰写

\subsection{每个 de Rham 类的调和代表}
\label{g3:subsec:1502:02}
待撰写

\subsection{唯一性、有限维性和自然配对}
\label{g3:subsec:1502:03}
待撰写

\subsection{Poincaré 对偶的解析再解释}
\label{g3:subsec:1502:04}
待撰写
% 纲要条目排为编号小节标题，并进入目录与书签；说明段仍保留于纲要注释中。

% !TeX root = ../../Book3.tex
% 纲要标题：### 15.3 Weitzenböck 与 Bochner
% 纲要位置：几何学计划纲要.md，第 2599 行。
% 正文状态：未撰写；本文件目前是框架。

% 纲要内容开始
% 
% * 联络 Laplace 和曲率作用
% * 一形式上的 Ricci 项
% * 正 Ricci 下的消失
% * 非负曲率下的平行性
% * 高阶形式中不同曲率算子的要求
% 
% 纲要内容结束

\section{Weitzenböck 与 Bochner}
\label{g3:sec:1503}

待撰写

\subsection{联络 Laplace 和曲率作用}
\label{g3:subsec:1503:01}
待撰写

\subsection{一形式上的 Ricci 项}
\label{g3:subsec:1503:02}
待撰写

\subsection{正 Ricci 下的消失}
\label{g3:subsec:1503:03}
待撰写

\subsection{非负曲率下的平行性}
\label{g3:subsec:1503:04}
待撰写

\subsection{高阶形式中不同曲率算子的要求}
\label{g3:subsec:1503:05}
待撰写
% 纲要条目排为编号小节标题，并进入目录与书签；说明段仍保留于纲要注释中。

% !TeX root = ../../Book3.tex
% 纲要标题：### 15.4 非紧分析中的截断
% 纲要位置：几何学计划纲要.md，第 2730 行。
% 正文状态：未撰写；本文件目前是框架。

% 纲要内容开始
% 
% * 完备度量上的截断函数
% * 积分边界项的控制
% * L² 调和形式与额外增长条件
% * 弱次调和函数的正则化
% 
% 纲要内容结束

\section{非紧分析中的截断}
\label{g3:sec:1504}

待撰写

\subsection{完备度量上的截断函数}
\label{g3:subsec:1504:01}
待撰写

\subsection{积分边界项的控制}
\label{g3:subsec:1504:02}
待撰写

\subsection{\texorpdfstring{\(L^2\) 调和形式与额外增长条件}{L² 调和形式与额外增长条件}}
\label{g3:subsec:1504:03}
待撰写

\subsection{弱次调和函数的正则化}
\label{g3:subsec:1504:04}
待撰写
% 纲要条目排为编号小节标题，并进入目录与书签；说明段仍保留于纲要注释中。

% !TeX root = ../../Book3.tex
% 纲要标题：### 15.5* Cheeger–Gromoll 分裂
% 纲要位置：几何学计划纲要.md，第 2614 行。
% 正文状态：未撰写；本文件目前是框架。

% 纲要内容开始
% 
% * 直线与两侧 Busemann 函数
% * 弱 Laplace 比较与最大值原理
% * Busemann 的调和性和光滑性
% * Bochner 公式得到平行单位梯度
% * 整体等距乘积的构造
% 
% ---
% 
% 纲要内容结束

\OptionalSection{Cheeger–Gromoll 分裂}{g3:sec:1505}

待撰写

\subsection{直线与两侧 Busemann 函数}
\label{g3:subsec:1505:01}
待撰写

\subsection{弱 Laplace 比较与最大值原理}
\label{g3:subsec:1505:02}
待撰写

\subsection{Busemann 的调和性和光滑性}
\label{g3:subsec:1505:03}
待撰写

\subsection{Bochner 公式得到平行单位梯度}
\label{g3:subsec:1505:04}
待撰写

\subsection{整体等距乘积的构造}
\label{g3:subsec:1505:05}
待撰写
% 纲要条目排为编号小节标题，并进入目录与书签；说明段仍保留于纲要注释中。


\begin{chaptersummary}
待撰写
\end{chaptersummary}

\BookExercises
待撰写

% !TeX root = ../../Book3.tex
% 纲要标题：## 第16章　热核、谱渐近与几何系数
% 纲要位置：几何学计划纲要.md，第 2624 行。
% 正文状态：未撰写；本文件目前是框架。

% 纲要内容开始
% 
% > 本章目标与证明范围：谱展开构造和短时局部参数构造分别证明，再比较二者。Weyl 渐近和逆谱问题在明确的余项与 Tauber 型输入之后展开。
% 
% 纲要内容结束

\chapter{热核、谱渐近与几何系数}
\label{g3:ch:16}
\BookChapterNumber{g3:ch:16}

\begin{chapterabstract}
待撰写
\end{chapterabstract}

\begin{learninggoals}
待撰写
\end{learninggoals}

待撰写

% !TeX root = ../../Book3.tex
% 纲要标题：### 16.1 热半群和谱构造
% 纲要位置：几何学计划纲要.md，第 2628 行。
% 正文状态：未撰写；本文件目前是框架。

% 纲要内容开始
% 
% * 非负自伴 Laplace 的指数
% * 热方程解与平滑作用
% * 核的存在、对称性和半群律
% * 最大值与质量保持
% 
% 纲要内容结束

\section{热半群和谱构造}
\label{g3:sec:1601}

待撰写

\subsection{非负自伴 Laplace 的指数}
\label{g3:subsec:1601:01}
待撰写

\subsection{热方程解与平滑作用}
\label{g3:subsec:1601:02}
待撰写

\subsection{核的存在、对称性和半群律}
\label{g3:subsec:1601:03}
待撰写

\subsection{最大值与质量保持}
\label{g3:subsec:1601:04}
待撰写
% 纲要条目排为编号小节标题，并进入目录与书签；说明段仍保留于纲要注释中。

% !TeX root = ../../Book3.tex
% 纲要标题：### 16.2 短时热核的局部构造
% 纲要位置：几何学计划纲要.md，第 2635 行。
% 正文状态：未撰写；本文件目前是框架。

% 纲要内容开始
% 
% * 测地距离和 Gaussian 主项
% * 运输方程与系数递推
% * 截断参数核的误差
% * Volterra 修正及余项控制
% * 局部构造与谱核的一致性
% 
% 纲要内容结束

\section{短时热核的局部构造}
\label{g3:sec:1602}

待撰写

\subsection{测地距离和 Gaussian 主项}
\label{g3:subsec:1602:01}
待撰写

\subsection{运输方程与系数递推}
\label{g3:subsec:1602:02}
待撰写

\subsection{截断参数核的误差}
\label{g3:subsec:1602:03}
待撰写

\subsection{Volterra 修正及余项控制}
\label{g3:subsec:1602:04}
待撰写

\subsection{局部构造与谱核的一致性}
\label{g3:subsec:1602:05}
待撰写
% 纲要条目排为编号小节标题，并进入目录与书签；说明段仍保留于纲要注释中。

% !TeX root = ../../Book3.tex
% 纲要标题：### 16.3 热迹和谱渐近
% 纲要位置：几何学计划纲要.md，第 2766 行。
% 正文状态：未撰写；本文件目前是框架。

% 纲要内容开始
% 
% * 对角系数与体积、标量曲率
% * 积分余项与热迹展开
% * 从热迹到 Weyl 律的所需 Tauber 引理
% * 特征值增长的归一化
% 
% 纲要内容结束

\section{热迹和谱渐近}
\label{g3:sec:1603}

待撰写

\subsection{对角系数与体积、标量曲率}
\label{g3:subsec:1603:01}
待撰写

\subsection{积分余项与热迹展开}
\label{g3:subsec:1603:02}
待撰写

\subsection{从热迹到 Weyl 律的所需 Tauber 引理}
\label{g3:subsec:1603:03}
待撰写

\subsection{特征值增长的归一化}
\label{g3:subsec:1603:04}
待撰写
% 纲要条目排为编号小节标题，并进入目录与书签；说明段仍保留于纲要注释中。

% !TeX root = ../../Book3.tex
% 纲要标题：### 16.4 谱估计及逆谱限制
% 纲要位置：几何学计划纲要.md，第 2650 行。
% 正文状态：未撰写；本文件目前是框架。

% 纲要内容开始
% 
% * Rayleigh 商和极小极大原则
% * 曲率、直径及特征值的所选估计
% * 球面和环面的谱
% * 等谱不等距及可恢复信息
% 
% ---
% 
% 纲要内容结束

\section{谱估计及逆谱限制}
\label{g3:sec:1604}

待撰写

\subsection{Rayleigh 商和极小极大原则}
\label{g3:subsec:1604:01}
待撰写

\subsection{曲率、直径及特征值的所选估计}
\label{g3:subsec:1604:02}
待撰写

\subsection{球面和环面的谱}
\label{g3:subsec:1604:03}
待撰写

\subsection{等谱不等距及可恢复信息}
\label{g3:subsec:1604:04}
待撰写
% 纲要条目排为编号小节标题，并进入目录与书签；说明段仍保留于纲要注释中。


\begin{chaptersummary}
待撰写
\end{chaptersummary}

\BookExercises
待撰写

% !TeX root = ../../Book3.tex
% 纲要标题：## 第17章　Gauss–Bonnet–Chern、Dirac 与指标观点
% 纲要位置：几何学计划纲要.md，第 2659 行。
% 正文状态：未撰写；本文件目前是框架。

% 纲要内容开始
% 
% > 本章目标与证明范围：高维 Euler 公式保留完整证明。自旋与一般指标理论区分“构造可以证明”“局部公式需要更多分析”“一般指标定理外引”三个层次，不能把 Euler 特例扩大成一般证明。
% 
% 纲要内容结束

\chapter{Gauss–Bonnet–Chern、Dirac 与指标观点}
\label{g3:ch:17}
\BookChapterNumber{g3:ch:17}

\begin{chapterabstract}
待撰写
\end{chapterabstract}

\begin{learninggoals}
待撰写
\end{learninggoals}

待撰写

% !TeX root = ../../Book3.tex
% 纲要标题：### 17.1 Euler 形式的构造
% 纲要位置：几何学计划纲要.md，第 2663 行。
% 正文状态：未撰写；本文件目前是框架。

% 纲要内容开始
% 
% * Pfaffian 的线性代数
% * 曲率形式和归一化
% * Chern–Weil 闭性及联络无关性
% 
% 纲要内容结束

\section{Euler 形式的构造}
\label{g3:sec:1701}

待撰写

\subsection{Pfaffian 的线性代数}
\label{g3:subsec:1701:01}
待撰写

\subsection{曲率形式和归一化}
\label{g3:subsec:1701:02}
待撰写

\subsection{Chern–Weil 闭性及联络无关性}
\label{g3:subsec:1701:03}
待撰写
% 纲要条目排为编号小节标题，并进入目录与书签；说明段仍保留于纲要注释中。

% !TeX root = ../../Book3.tex
% 纲要标题：### 17.2 球丛上的转移证明
% 纲要位置：几何学计划纲要.md，第 2669 行。
% 正文状态：未撰写；本文件目前是框架。

% 纲要内容开始
% 
% * 拉回丛的典范截面
% * 转移形式及其外微分
% * 纤维积分的归一化
% * 孤立零点附近的边界积分
% * Poincaré–Hopf 与整体 Euler 公式
% 
% 纲要内容结束

\section{球丛上的转移证明}
\label{g3:sec:1702}

待撰写

\subsection{拉回丛的典范截面}
\label{g3:subsec:1702:01}
待撰写

\subsection{转移形式及其外微分}
\label{g3:subsec:1702:02}
待撰写

\subsection{纤维积分的归一化}
\label{g3:subsec:1702:03}
待撰写

\subsection{孤立零点附近的边界积分}
\label{g3:subsec:1702:04}
待撰写

\subsection{Poincaré–Hopf 与整体 Euler 公式}
\label{g3:subsec:1702:05}
待撰写
% 纲要条目排为编号小节标题，并进入目录与书签；说明段仍保留于纲要注释中。

% !TeX root = ../../Book3.tex
% 纲要标题：### 17.3* Dirac 型算子与自旋
% 纲要位置：几何学计划纲要.md，第 2800 行。
% 正文状态：未撰写；本文件目前是框架。

% 纲要内容开始
% 
% * Clifford 代数和 Clifford 丛
% * spin 提升与自旋子丛
% * Dirac 算子及符号
% * Lichnerowicz 公式与正标量曲率障碍
% 
% 纲要内容结束

\OptionalSection{Dirac 型算子与自旋}{g3:sec:1703}

待撰写

\subsection{Clifford 代数和 Clifford 丛}
\label{g3:subsec:1703:01}
待撰写

\subsection{spin 提升与自旋子丛}
\label{g3:subsec:1703:02}
待撰写

\subsection{Dirac 算子及符号}
\label{g3:subsec:1703:03}
待撰写

\subsection{Lichnerowicz 公式与正标量曲率障碍}
\label{g3:subsec:1703:04}
待撰写
% 纲要条目排为编号小节标题，并进入目录与书签；说明段仍保留于纲要注释中。

% !TeX root = ../../Book3.tex
% 纲要标题：### 17.4* 一般指标定理的准确接口
% 纲要位置：几何学计划纲要.md，第 2807 行。
% 正文状态：未撰写；本文件目前是框架。

% 纲要内容开始
% 
% * 椭圆符号与解析指标
% * 拓扑 K 类和推前
% * Atiyah–Singer 的所选陈述
% * de Rham、signature 与 spin 模型的比较
% 
% ---
% 
% 纲要内容结束

\OptionalSection{一般指标定理的准确接口}{g3:sec:1704}

待撰写

\subsection{椭圆符号与解析指标}
\label{g3:subsec:1704:01}
待撰写

\subsection{拓扑 K 类和推前}
\label{g3:subsec:1704:02}
待撰写

\subsection{Atiyah–Singer 的所选陈述}
\label{g3:subsec:1704:03}
待撰写

\subsection{de Rham、signature 与 spin 模型的比较}
\label{g3:subsec:1704:04}
待撰写
% 纲要条目排为编号小节标题，并进入目录与书签；说明段仍保留于纲要注释中。


\begin{chaptersummary}
待撰写
\end{chaptersummary}

\BookExercises
待撰写


% !TeX root = ../Book3.tex
% 纲要标题：# 第四编　非线性映射、极小对象与校准
% 纲要位置：几何学计划纲要.md，第 2693 行。
% 正文状态：未撰写；本文件目前是框架。

% 纲要内容开始
% 
% 纲要内容结束

\BookPart{非线性映射、极小对象与校准}{g3:part:04}
\BookPartGuide
待撰写

% !TeX root = ../../Book3.tex
% 纲要标题：## 第18章　调和映射的变分、正则性与集中
% 纲要位置：几何学计划纲要.md，第 2818 行。
% 正文状态：未撰写；本文件目前是框架。

% 纲要内容开始
% 
% > 本章目标与证明范围：把光滑临界方程、极小映射、平稳映射和一般弱解分开。完整完成所选小能量正则性模块；热流在抛物工具之后处理，避免局部椭圆正则性与长期存在混同。
% 
% 纲要内容结束

\chapter{调和映射的变分、正则性与集中}
\label{g3:ch:18}
\BookChapterNumber{g3:ch:18}

\begin{chapterabstract}
待撰写
\end{chapterabstract}

\begin{learninggoals}
待撰写
\end{learninggoals}

待撰写

% !TeX root = ../../Book3.tex
% 纲要标题：### 18.1 能量与几何变分
% 纲要位置：几何学计划纲要.md，第 2699 行。
% 正文状态：未撰写；本文件目前是框架。

% 纲要内容开始
% 
% * 拉回丛上的协变微分
% * 张力场和第一变分
% * 第二变分及目标曲率
% * 测地线作为一维原型
% 
% 纲要内容结束

\section{能量与几何变分}
\label{g3:sec:1801}

待撰写

\subsection{拉回丛上的协变微分}
\label{g3:subsec:1801:01}
待撰写

\subsection{张力场和第一变分}
\label{g3:subsec:1801:02}
待撰写

\subsection{第二变分及目标曲率}
\label{g3:subsec:1801:03}
待撰写

\subsection{测地线作为一维原型}
\label{g3:subsec:1801:04}
待撰写
% 纲要条目排为编号小节标题，并进入目录与书签；说明段仍保留于纲要注释中。

% !TeX root = ../../Book3.tex
% 纲要标题：### 18.2 弱解与可接受的极限
% 纲要位置：几何学计划纲要.md，第 2829 行。
% 正文状态：未撰写；本文件目前是框架。

% 纲要内容开始
% 
% * Sobolev 映射的定义
% * 目标嵌入及内在独立性
% * 极小、平稳和弱调和的不同条件
% * 弱极限不保持所有约束的现象
% 
% 纲要内容结束

\section{弱解与可接受的极限}
\label{g3:sec:1802}

待撰写

\subsection{Sobolev 映射的定义}
\label{g3:subsec:1802:01}
待撰写

\subsection{目标嵌入及内在独立性}
\label{g3:subsec:1802:02}
待撰写

\subsection{极小、平稳和弱调和的不同条件}
\label{g3:subsec:1802:03}
待撰写

\subsection{弱极限不保持所有约束的现象}
\label{g3:subsec:1802:04}
待撰写
% 纲要条目排为编号小节标题，并进入目录与书签；说明段仍保留于纲要注释中。

% !TeX root = ../../Book3.tex
% 纲要标题：### 18.3 曲面源上的小能量正则性
% 纲要位置：几何学计划纲要.md，第 2836 行。
% 正文状态：未撰写；本文件目前是框架。

% 纲要内容开始
% 
% 完整证明的主版本固定为二维圆盘内的光滑调和映射、光滑紧黎曼目标及内域估计。先对光滑解族建立一致界，不把所有高维弱解都归入同一小能量定理。
% 
% * 二维能量缩放与临界性
% * 光滑调和映射的 Bochner 微分不等式
% * 小能量阈值、均值估计和梯度界
% 
%   证明存在由目标几何决定的 $\varepsilon_0,C>0$，使 $E(u;B_r)<\varepsilon_0$ 时 $\sup_{B_{r/2}}|du|^2\le Cr^{-2}E(u;B_r)$；在曲面坐标中追踪背景度量的控制常数。
% 
% * 内域高阶估计及光滑解族的紧性
% * 能量集中点、重缩放及奇点形成
% * 一般弱解与高维部分正则性的独立范围
% 
%   二维一般弱调和映射的正则性及高维极小、平稳映射的部分正则性分别列出外部定理的准确对象和假设；单调性公式仅用于已经验证相应驻值条件的映射。
% 
% 纲要内容结束

\section{曲面源上的小能量正则性}
\label{g3:sec:1803}

待撰写

\subsection{二维能量缩放与临界性}
\label{g3:subsec:1803:01}
待撰写

\subsection{光滑调和映射的 Bochner 微分不等式}
\label{g3:subsec:1803:02}
待撰写

\subsection{小能量阈值、均值估计和梯度界}
\label{g3:subsec:1803:03}
待撰写

\subsection{内域高阶估计及光滑解族的紧性}
\label{g3:subsec:1803:04}
待撰写

\subsection{能量集中点、重缩放及奇点形成}
\label{g3:subsec:1803:05}
待撰写

\subsection{一般弱解与高维部分正则性的独立范围}
\label{g3:subsec:1803:06}
待撰写
% 纲要条目排为编号小节标题，并进入目录与书签；说明段仍保留于纲要注释中。

% !TeX root = ../../Book3.tex
% 纲要标题：### 18.4 曲面映射与复结构
% 纲要位置：几何学计划纲要.md，第 2852 行。
% 正文状态：未撰写；本文件目前是框架。

% 纲要内容开始
% 
% * 能量的共形不变性
% * Hopf 微分及全纯性
% * 共形调和映射与极小曲面
% * 气泡和同伦类丢失的模型
% 
% ---
% 
% 纲要内容结束

\section{曲面映射与复结构}
\label{g3:sec:1804}

待撰写

\subsection{能量的共形不变性}
\label{g3:subsec:1804:01}
待撰写

\subsection{Hopf 微分及全纯性}
\label{g3:subsec:1804:02}
待撰写

\subsection{共形调和映射与极小曲面}
\label{g3:subsec:1804:03}
待撰写

\subsection{气泡和同伦类丢失的模型}
\label{g3:subsec:1804:04}
待撰写
% 纲要条目排为编号小节标题，并进入目录与书签；说明段仍保留于纲要注释中。


\begin{chaptersummary}
待撰写
\end{chaptersummary}

\BookExercises
待撰写

% !TeX root = ../../Book3.tex
% 纲要标题：## 第19章　极小子流形、稳定性与校准
% 纲要位置：几何学计划纲要.md，第 2738 行。
% 正文状态：未撰写；本文件目前是框架。

% 纲要内容开始
% 
% > 本章目标与证明范围：面积的驻值、稳定和全局极小分别建立。主线从变分公式进入局部方程及单调性，再通过校准给出真正的全局极小证书。
% 
% 纲要内容结束

\chapter{极小子流形、稳定性与校准}
\label{g3:ch:19}
\BookChapterNumber{g3:ch:19}

\begin{chapterabstract}
待撰写
\end{chapterabstract}

\begin{learninggoals}
待撰写
\end{learninggoals}

待撰写

% !TeX root = ../../Book3.tex
% 纲要标题：### 19.1 面积的一阶与二阶变分
% 纲要位置：几何学计划纲要.md，第 2865 行。
% 正文状态：未撰写；本文件目前是框架。

% 纲要内容开始
% 
% * 法向变分与平均曲率
% * 边界项和固定边界
% * Jacobi 稳定算子
% * 高余维法联络和曲率项
% 
% 纲要内容结束

\section{面积的一阶与二阶变分}
\label{g3:sec:1901}

待撰写

\subsection{法向变分与平均曲率}
\label{g3:subsec:1901:01}
待撰写

\subsection{边界项和固定边界}
\label{g3:subsec:1901:02}
待撰写

\subsection{Jacobi 稳定算子}
\label{g3:subsec:1901:03}
待撰写

\subsection{高余维法联络和曲率项}
\label{g3:subsec:1901:04}
待撰写
% 纲要条目排为编号小节标题，并进入目录与书签；说明段仍保留于纲要注释中。

% !TeX root = ../../Book3.tex
% 纲要标题：### 19.2 极小图与局部模型
% 纲要位置：几何学计划纲要.md，第 2872 行。
% 正文状态：未撰写；本文件目前是框架。

% 纲要内容开始
% 
% * 面积积分及 Euler–Lagrange 方程
% * 一致椭圆性的梯度条件
% * 平面、悬链面及螺旋面
% * 共形参数和 Weierstrass 表示的范围
% 
% 纲要内容结束

\section{极小图与局部模型}
\label{g3:sec:1902}

待撰写

\subsection{面积积分及 Euler–Lagrange 方程}
\label{g3:subsec:1902:01}
待撰写

\subsection{一致椭圆性的梯度条件}
\label{g3:subsec:1902:02}
待撰写

\subsection{平面、悬链面及螺旋面}
\label{g3:subsec:1902:03}
待撰写

\subsection{共形参数和 Weierstrass 表示的范围}
\label{g3:subsec:1902:04}
待撰写
% 纲要条目排为编号小节标题，并进入目录与书签；说明段仍保留于纲要注释中。

% !TeX root = ../../Book3.tex
% 纲要标题：### 19.3 单调性和稳定不等式
% 纲要位置：几何学计划纲要.md，第 2756 行。
% 正文状态：未撰写；本文件目前是框架。

% 纲要内容开始
% 
% * 径向变分及面积比
% * 环境曲率修正
% * 稳定性与测试函数
% * 切锥和密度的几何动机
% 
% 纲要内容结束

\section{单调性和稳定不等式}
\label{g3:sec:1903}

待撰写

\subsection{径向变分及面积比}
\label{g3:subsec:1903:01}
待撰写

\subsection{环境曲率修正}
\label{g3:subsec:1903:02}
待撰写

\subsection{稳定性与测试函数}
\label{g3:subsec:1903:03}
待撰写

\subsection{切锥和密度的几何动机}
\label{g3:subsec:1903:04}
待撰写
% 纲要条目排为编号小节标题，并进入目录与书签；说明段仍保留于纲要注释中。

% !TeX root = ../../Book3.tex
% 纲要标题：### 19.4 校准及全局极小性
% 纲要位置：几何学计划纲要.md，第 2763 行。
% 正文状态：未撰写；本文件目前是框架。

% 纲要内容开始
% 
% * 余范数和校准平面
% * 闭性与同调竞争类
% * 校准对象的体积极小证明
% * 复子流形与 Kähler 校准
% * 特殊 Lagrange 模型及非唯一性
% 
% ---
% 
% 纲要内容结束

\section{校准及全局极小性}
\label{g3:sec:1904}

待撰写

\subsection{余范数和校准平面}
\label{g3:subsec:1904:01}
待撰写

\subsection{闭性与同调竞争类}
\label{g3:subsec:1904:02}
待撰写

\subsection{校准对象的体积极小证明}
\label{g3:subsec:1904:03}
待撰写

\subsection{复子流形与 Kähler 校准}
\label{g3:subsec:1904:04}
待撰写

\subsection{特殊 Lagrange 模型及非唯一性}
\label{g3:subsec:1904:05}
待撰写
% 纲要条目排为编号小节标题，并进入目录与书签；说明段仍保留于纲要注释中。


\begin{chaptersummary}
待撰写
\end{chaptersummary}

\BookExercises
待撰写


% !TeX root = ../Book3.tex
% 纲要标题：# 第五编　几何测度、广义边界与 Plateau 问题
% 纲要位置：几何学计划纲要.md，第 2773 行。
% 正文状态：未撰写；本文件目前是框架。

% 纲要内容开始
% 
% 纲要内容结束

\BookPart{几何测度、广义边界与 Plateau 问题}{g3:part:05}
\BookPartGuide
待撰写

% !TeX root = ../../Book3.tex
% 纲要标题：## 第20章　Hausdorff 测度、可整流性与面积公式
% 纲要位置：几何学计划纲要.md，第 2922 行。
% 正文状态：未撰写；本文件目前是框架。

% 纲要内容开始
% 
% > 本章目标与证明范围：先建立可测几何的真正局部工具，再引入弱曲面。Rademacher、面积和余面积各有独立证明步骤，近似切平面不是光滑切平面的名义推广。
% 
% 纲要内容结束

\chapter{Hausdorff 测度、可整流性与面积公式}
\label{g3:ch:20}
\BookChapterNumber{g3:ch:20}

\begin{chapterabstract}
待撰写
\end{chapterabstract}

\begin{learninggoals}
待撰写
\end{learninggoals}

待撰写

% !TeX root = ../../Book3.tex
% 纲要标题：### 20.1 测度及覆盖方法
% 纲要位置：几何学计划纲要.md，第 2779 行。
% 正文状态：未撰写；本文件目前是框架。

% 纲要内容开始
% 
% * Hausdorff 外测度与归一化
% * Vitali 型覆盖引理
% * 密度、微分测度与零测集
% * 光滑子流形面积的一致性
% 
% 纲要内容结束

\section{测度及覆盖方法}
\label{g3:sec:2001}

待撰写

\subsection{Hausdorff 外测度与归一化}
\label{g3:subsec:2001:01}
待撰写

\subsection{Vitali 型覆盖引理}
\label{g3:subsec:2001:02}
待撰写

\subsection{密度、微分测度与零测集}
\label{g3:subsec:2001:03}
待撰写

\subsection{光滑子流形面积的一致性}
\label{g3:subsec:2001:04}
待撰写
% 纲要条目排为编号小节标题，并进入目录与书签；说明段仍保留于纲要注释中。

% !TeX root = ../../Book3.tex
% 纲要标题：### 20.2 Lipschitz 可微与线性化
% 纲要位置：几何学计划纲要.md，第 2786 行。
% 正文状态：未撰写；本文件目前是框架。

% 纲要内容开始
% 
% * 一维绝对连续函数
% * 逐方向导数和 Rademacher 定理
% * 近似微分与密度点
% * Lipschitz 映射的局部分解
% 
% 纲要内容结束

\section{Lipschitz 可微与线性化}
\label{g3:sec:2002}

待撰写

\subsection{一维绝对连续函数}
\label{g3:subsec:2002:01}
待撰写

\subsection{逐方向导数和 Rademacher 定理}
\label{g3:subsec:2002:02}
待撰写

\subsection{近似微分与密度点}
\label{g3:subsec:2002:03}
待撰写

\subsection{Lipschitz 映射的局部分解}
\label{g3:subsec:2002:04}
待撰写
% 纲要条目排为编号小节标题，并进入目录与书签；说明段仍保留于纲要注释中。

% !TeX root = ../../Book3.tex
% 纲要标题：### 20.3 面积公式
% 纲要位置：几何学计划纲要.md，第 2793 行。
% 正文状态：未撰写；本文件目前是框架。

% 纲要内容开始
% 
% * 单射 Lipschitz 片的 Jacobian
% * 重数函数和可测性
% * 由线性模型到一般映射
% * 非单射情形的分片求和
% * 几何面积与参数面积的区别
% 
% 纲要内容结束

\section{面积公式}
\label{g3:sec:2003}

待撰写

\subsection{单射 Lipschitz 片的 Jacobian}
\label{g3:subsec:2003:01}
待撰写

\subsection{重数函数和可测性}
\label{g3:subsec:2003:02}
待撰写

\subsection{由线性模型到一般映射}
\label{g3:subsec:2003:03}
待撰写

\subsection{非单射情形的分片求和}
\label{g3:subsec:2003:04}
待撰写

\subsection{几何面积与参数面积的区别}
\label{g3:subsec:2003:05}
待撰写
% 纲要条目排为编号小节标题，并进入目录与书签；说明段仍保留于纲要注释中。

% !TeX root = ../../Book3.tex
% 纲要标题：### 20.4 余面积公式
% 纲要位置：几何学计划纲要.md，第 2948 行。
% 正文状态：未撰写；本文件目前是框架。

% 纲要内容开始
% 
% * 线性投影模型
% * 切片的测度与 Jacobian
% * 局部分解和积分交换
% * 光滑水平集及非光滑切片
% 
% 纲要内容结束

\section{余面积公式}
\label{g3:sec:2004}

待撰写

\subsection{线性投影模型}
\label{g3:subsec:2004:01}
待撰写

\subsection{切片的测度与 Jacobian}
\label{g3:subsec:2004:02}
待撰写

\subsection{局部分解和积分交换}
\label{g3:subsec:2004:03}
待撰写

\subsection{光滑水平集及非光滑切片}
\label{g3:subsec:2004:04}
待撰写
% 纲要条目排为编号小节标题，并进入目录与书签；说明段仍保留于纲要注释中。

% !TeX root = ../../Book3.tex
% 纲要标题：### 20.5 可整流集的切向结构
% 纲要位置：几何学计划纲要.md，第 2955 行。
% 正文状态：未撰写；本文件目前是框架。

% 纲要内容开始
% 
% * 可数 Lipschitz 参数片
% * 近似切平面的存在唯一性
% * 密度和重数
% * 纯不可整流部分的意义
% 
% ---
% 
% 纲要内容结束

\section{可整流集的切向结构}
\label{g3:sec:2005}

待撰写

\subsection{可数 Lipschitz 参数片}
\label{g3:subsec:2005:01}
待撰写

\subsection{近似切平面的存在唯一性}
\label{g3:subsec:2005:02}
待撰写

\subsection{密度和重数}
\label{g3:subsec:2005:03}
待撰写

\subsection{纯不可整流部分的意义}
\label{g3:subsec:2005:04}
待撰写
% 纲要条目排为编号小节标题，并进入目录与书签；说明段仍保留于纲要注释中。


\begin{chaptersummary}
待撰写
\end{chaptersummary}

\BookExercises
待撰写

% !TeX root = ../../Book3.tex
% 纲要标题：## 第21章　BV、有限周长与弱 Gauss–Green
% 纲要位置：几何学计划纲要.md，第 2964 行。
% 正文状态：未撰写；本文件目前是框架。

% 纲要内容开始
% 
% > 本章目标与证明范围：把有限周长从原先一节扩为完整的边界理论。分布导数、紧性、约化边界及 Gauss–Green 的联系有连续证明；不把拓扑边界当作测度论边界。
% 
% 纲要内容结束

\chapter{BV、有限周长与弱 Gauss–Green}
\label{g3:ch:21}
\BookChapterNumber{g3:ch:21}

\begin{chapterabstract}
待撰写
\end{chapterabstract}

\begin{learninggoals}
待撰写
\end{learninggoals}

待撰写

% !TeX root = ../../Book3.tex
% 纲要标题：### 21.1 BV 函数与全变差
% 纲要位置：几何学计划纲要.md，第 2968 行。
% 正文状态：未撰写；本文件目前是框架。

% 纲要内容开始
% 
% * 分布导数的 Radon 测度表示
% * 变差的对偶定义
% * 光滑逼近与变差收敛
% * 截断、复合及链式公式的范围
% 
% 纲要内容结束

\section{BV 函数与全变差}
\label{g3:sec:2101}

待撰写

\subsection{分布导数的 Radon 测度表示}
\label{g3:subsec:2101:01}
待撰写

\subsection{变差的对偶定义}
\label{g3:subsec:2101:02}
待撰写

\subsection{光滑逼近与变差收敛}
\label{g3:subsec:2101:03}
待撰写

\subsection{截断、复合及链式公式的范围}
\label{g3:subsec:2101:04}
待撰写
% 纲要条目排为编号小节标题，并进入目录与书签；说明段仍保留于纲要注释中。

% !TeX root = ../../Book3.tex
% 纲要标题：### 21.2 紧性和下半连续
% 纲要位置：几何学计划纲要.md，第 2828 行。
% 正文状态：未撰写；本文件目前是框架。

% 纲要内容开始
% 
% * 局部 BV 紧性
% * L¹ 收敛下的变差下半连续
% * 支集与整体紧性
% * 有限周长集合的直接方法
% 
% 纲要内容结束

\section{紧性和下半连续}
\label{g3:sec:2102}

待撰写

\subsection{局部 BV 紧性}
\label{g3:subsec:2102:01}
待撰写

\subsection{\texorpdfstring{\(L^1\) 收敛下的变差下半连续}{L¹ 收敛下的变差下半连续}}
\label{g3:subsec:2102:02}
待撰写

\subsection{支集与整体紧性}
\label{g3:subsec:2102:03}
待撰写

\subsection{有限周长集合的直接方法}
\label{g3:subsec:2102:04}
待撰写
% 纲要条目排为编号小节标题，并进入目录与书签；说明段仍保留于纲要注释中。

% !TeX root = ../../Book3.tex
% 纲要标题：### 21.3 测度论边界的定义
% 纲要位置：几何学计划纲要.md，第 2982 行。
% 正文状态：未撰写；本文件目前是框架。

% 纲要内容开始
% 
% * 密度为零、一和中间值的点
% * 约化边界与广义法向
% * 拓扑边界与测度论边界的区别
% 
% 纲要内容结束

\section{测度论边界的定义}
\label{g3:sec:2103}

待撰写

\subsection{密度为零、一和中间值的点}
\label{g3:subsec:2103:01}
待撰写

\subsection{约化边界与广义法向}
\label{g3:subsec:2103:02}
待撰写

\subsection{拓扑边界与测度论边界的区别}
\label{g3:subsec:2103:03}
待撰写
% 纲要条目排为编号小节标题，并进入目录与书签；说明段仍保留于纲要注释中。

% !TeX root = ../../Book3.tex
% 纲要标题：### 21.4 约化边界结构定理的证明
% 纲要位置：几何学计划纲要.md，第 2841 行。
% 正文状态：未撰写；本文件目前是框架。

% 纲要内容开始
% 
% * 相对等周估计和密度控制
% * 放大族的 BV 紧性
% * 约化点的半空间极限
% * 法向控制和近似切平面
% * Lipschitz 片覆盖的构造
% * 周长测度的表示
% * 可整流性与零测例外集
% 
% 纲要内容结束

\section{约化边界结构定理的证明}
\label{g3:sec:2104}

待撰写

\subsection{相对等周估计和密度控制}
\label{g3:subsec:2104:01}
待撰写

\subsection{放大族的 BV 紧性}
\label{g3:subsec:2104:02}
待撰写

\subsection{约化点的半空间极限}
\label{g3:subsec:2104:03}
待撰写

\subsection{法向控制和近似切平面}
\label{g3:subsec:2104:04}
待撰写

\subsection{Lipschitz 片覆盖的构造}
\label{g3:subsec:2104:05}
待撰写

\subsection{周长测度的表示}
\label{g3:subsec:2104:06}
待撰写

\subsection{可整流性与零测例外集}
\label{g3:subsec:2104:07}
待撰写
% 纲要条目排为编号小节标题，并进入目录与书签；说明段仍保留于纲要注释中。

% !TeX root = ../../Book3.tex
% 纲要标题：### 21.5 弱 Gauss–Green 与等周问题
% 纲要位置：几何学计划纲要.md，第 2998 行。
% 正文状态：未撰写；本文件目前是框架。

% 纲要内容开始
% 
% * 积分分部公式
% * 切片和 coarea 的 BV 版本
% * 等周不等式的所选证明
% * 固定体积极小周长的模型
% 
% ---
% 
% 纲要内容结束

\section{弱 Gauss–Green 与等周问题}
\label{g3:sec:2105}

待撰写

\subsection{积分分部公式}
\label{g3:subsec:2105:01}
待撰写

\subsection{切片和 coarea 的 BV 版本}
\label{g3:subsec:2105:02}
待撰写

\subsection{等周不等式的所选证明}
\label{g3:subsec:2105:03}
待撰写

\subsection{固定体积极小周长的模型}
\label{g3:subsec:2105:04}
待撰写
% 纲要条目排为编号小节标题，并进入目录与书签；说明段仍保留于纲要注释中。


\begin{chaptersummary}
待撰写
\end{chaptersummary}

\BookExercises
待撰写

% !TeX root = ../../Book3.tex
% 纲要标题：## 第22章　流、切片与 Federer–Fleming 紧性
% 纲要位置：几何学计划纲要.md，第 3007 行。
% 正文状态：未撰写；本文件目前是框架。

% 纲要内容开始
% 
% > 本章目标与证明范围：从测试形式的对偶构造流，区分一般、正规、可整流和整数流。变形定理、整数性闭包和紧性均为正文证明任务；仅用弱紧性不能完成整数 Plateau。
% 
% 纲要内容结束

\chapter{流、切片与 Federer–Fleming 紧性}
\label{g3:ch:22}
\BookChapterNumber{g3:ch:22}

\begin{chapterabstract}
待撰写
\end{chapterabstract}

\begin{learninggoals}
待撰写
\end{learninggoals}

待撰写

% !TeX root = ../../Book3.tex
% 纲要标题：### 22.1 流及其几何运算
% 纲要位置：几何学计划纲要.md，第 2864 行。
% 正文状态：未撰写；本文件目前是框架。

% 纲要内容开始
% 
% * 测试形式、连续性和边界
% * 质量、局部质量与支集
% * 推前、乘积和同伦公式
% * 光滑定向子流形及抵消现象
% 
% 纲要内容结束

\section{流及其几何运算}
\label{g3:sec:2201}

待撰写

\subsection{测试形式、连续性和边界}
\label{g3:subsec:2201:01}
待撰写

\subsection{质量、局部质量与支集}
\label{g3:subsec:2201:02}
待撰写

\subsection{推前、乘积和同伦公式}
\label{g3:subsec:2201:03}
待撰写

\subsection{光滑定向子流形及抵消现象}
\label{g3:subsec:2201:04}
待撰写
% 纲要条目排为编号小节标题，并进入目录与书签；说明段仍保留于纲要注释中。

% !TeX root = ../../Book3.tex
% 纲要标题：### 22.2 可整流流与整数流
% 纲要位置：几何学计划纲要.md，第 3018 行。
% 正文状态：未撰写；本文件目前是框架。

% 纲要内容开始
% 
% * 切平面、取向与整数重数
% * 边界质量和正规流
% * 整数流的定义
% * 取向与非定向模型的区别
% 
% 纲要内容结束

\section{可整流流与整数流}
\label{g3:sec:2202}

待撰写

\subsection{切平面、取向与整数重数}
\label{g3:subsec:2202:01}
待撰写

\subsection{边界质量和正规流}
\label{g3:subsec:2202:02}
待撰写

\subsection{整数流的定义}
\label{g3:subsec:2202:03}
待撰写

\subsection{取向与非定向模型的区别}
\label{g3:subsec:2202:04}
待撰写
% 纲要条目排为编号小节标题，并进入目录与书签；说明段仍保留于纲要注释中。

% !TeX root = ../../Book3.tex
% 纲要标题：### 22.3 切片、锥与平坦范数
% 纲要位置：几何学计划纲要.md，第 3025 行。
% 正文状态：未撰写；本文件目前是框架。

% 纲要内容开始
% 
% * Lipschitz 切片及边界公式
% * 锥构造和填充估计
% * 平坦范数及其弱收敛含义
% * 质量收敛、平坦收敛与弱收敛
% 
% 纲要内容结束

\section{切片、锥与平坦范数}
\label{g3:sec:2203}

待撰写

\subsection{Lipschitz 切片及边界公式}
\label{g3:subsec:2203:01}
待撰写

\subsection{锥构造和填充估计}
\label{g3:subsec:2203:02}
待撰写

\subsection{平坦范数及其弱收敛含义}
\label{g3:subsec:2203:03}
待撰写

\subsection{质量收敛、平坦收敛与弱收敛}
\label{g3:subsec:2203:04}
待撰写
% 纲要条目排为编号小节标题，并进入目录与书签；说明段仍保留于纲要注释中。

% !TeX root = ../../Book3.tex
% 纲要标题：### 22.4 Federer–Fleming 变形定理
% 纲要位置：几何学计划纲要.md，第 2885 行。
% 正文状态：未撰写；本文件目前是框架。

% 纲要内容开始
% 
% * 立方网格与逐维径向投影
% * 质量控制和边界相容
% * 多面体逼近及误差填充
% * 网格尺度和统一常数
% 
% 纲要内容结束

\section{Federer–Fleming 变形定理}
\label{g3:sec:2204}

待撰写

\subsection{立方网格与逐维径向投影}
\label{g3:subsec:2204:01}
待撰写

\subsection{质量控制和边界相容}
\label{g3:subsec:2204:02}
待撰写

\subsection{多面体逼近及误差填充}
\label{g3:subsec:2204:03}
待撰写

\subsection{网格尺度和统一常数}
\label{g3:subsec:2204:04}
待撰写
% 纲要条目排为编号小节标题，并进入目录与书签；说明段仍保留于纲要注释中。

% !TeX root = ../../Book3.tex
% 纲要标题：### 22.5 整数流闭包与紧性的证明
% 纲要位置：几何学计划纲要.md，第 2892 行。
% 正文状态：未撰写；本文件目前是框架。

% 纲要内容开始
% 
% * 正规流的弱紧性
% 
%   在共同紧支和质量、边界质量界下，先获得对测试形式收敛的子列。
% 
% * 变形估计与统一平坦控制
% 
%   利用多面体逼近把弱序列转化为可控制的几何误差。
% 
% * 切片及维数归纳
% 
%   证明几乎所有切片保留相容的整数信息，逐维恢复整性。
% 
% * 极限的可整流性
% 
%   此处单列闭包定理真正需要的切向结构论证，不从弱收敛直接断言。
% 
% * 整数重数的保持
% 
%   区分可整流性与重数属于整数这两项结论。
% 
% * 边界相容及下半连续
% 
%   证明极限边界正确，并核验质量下半连续。
% 
% * 紧性定理及非紧推广的边界
% 
%   先完成固定紧支版本，再说明局部质量界与质量逃逸的区别。
% 
% ---
% 
% 纲要内容结束

\section{整数流闭包与紧性的证明}
\label{g3:sec:2205}

待撰写

\subsection{正规流的弱紧性}
\label{g3:subsec:2205:01}
待撰写

\subsection{变形估计与统一平坦控制}
\label{g3:subsec:2205:02}
待撰写

\subsection{切片及维数归纳}
\label{g3:subsec:2205:03}
待撰写

\subsection{极限的可整流性}
\label{g3:subsec:2205:04}
待撰写

\subsection{整数重数的保持}
\label{g3:subsec:2205:05}
待撰写

\subsection{边界相容及下半连续}
\label{g3:subsec:2205:06}
待撰写

\subsection{紧性定理及非紧推广的边界}
\label{g3:subsec:2205:07}
待撰写
% 纲要条目排为编号小节标题，并进入目录与书签；说明段仍保留于纲要注释中。


\begin{chaptersummary}
待撰写
\end{chaptersummary}

\BookExercises
待撰写

% !TeX root = ../../Book3.tex
% 纲要标题：## 第23章　Plateau 存在性、正则性与变分流形
% 纲要位置：几何学计划纲要.md，第 3071 行。
% 正文状态：未撰写；本文件目前是框架。

% 纲要内容开始
% 
% > 本章目标与证明范围：在给定可填充的紧支整数边界和质量有限的竞争类中，完整证明欧氏 Plateau 存在性。存在性、内部正则性、边界正则性及弱模型选择分开；一般奇异集界逐版本外引。
% 
% 纲要内容结束

\chapter{Plateau 存在性、正则性与变分流形}
\label{g3:ch:23}
\BookChapterNumber{g3:ch:23}

\begin{chapterabstract}
待撰写
\end{chapterabstract}

\begin{learninggoals}
待撰写
\end{learninggoals}

待撰写

% !TeX root = ../../Book3.tex
% 纲要标题：### 23.1 极小化问题的准确设置
% 纲要位置：几何学计划纲要.md，第 2928 行。
% 正文状态：未撰写；本文件目前是框架。

% 纲要内容开始
% 
% * 边界流、竞争类和非空性
% * 锥或显式填充提供初始竞争者
% * 凸投影控制支集
% * 极小化序列及质量界
% 
% 纲要内容结束

\section{极小化问题的准确设置}
\label{g3:sec:2301}

待撰写

\subsection{边界流、竞争类和非空性}
\label{g3:subsec:2301:01}
待撰写

\subsection{锥或显式填充提供初始竞争者}
\label{g3:subsec:2301:02}
待撰写

\subsection{凸投影控制支集}
\label{g3:subsec:2301:03}
待撰写

\subsection{极小化序列及质量界}
\label{g3:subsec:2301:04}
待撰写
% 纲要条目排为编号小节标题，并进入目录与书签；说明段仍保留于纲要注释中。

% !TeX root = ../../Book3.tex
% 纲要标题：### 23.2 存在性的完整证明
% 纲要位置：几何学计划纲要.md，第 3082 行。
% 正文状态：未撰写；本文件目前是框架。

% 纲要内容开始
% 
% * 整数流紧性获得极限
% * 边界算子的连续性
% * 质量下半连续
% * 极小性和已定竞争类
% * 拓扑类型为何可能变化
% 
% 纲要内容结束

\section{存在性的完整证明}
\label{g3:sec:2302}

待撰写

\subsection{整数流紧性获得极限}
\label{g3:subsec:2302:01}
待撰写

\subsection{边界算子的连续性}
\label{g3:subsec:2302:02}
待撰写

\subsection{质量下半连续}
\label{g3:subsec:2302:03}
待撰写

\subsection{极小性和已定竞争类}
\label{g3:subsec:2302:04}
待撰写

\subsection{拓扑类型为何可能变化}
\label{g3:subsec:2302:05}
待撰写
% 纲要条目排为编号小节标题，并进入目录与书签；说明段仍保留于纲要注释中。

% !TeX root = ../../Book3.tex
% 纲要标题：### 23.3 单调性、切锥与正则点
% 纲要位置：几何学计划纲要.md，第 2943 行。
% 正文状态：未撰写；本文件目前是框架。

% 纲要内容开始
% 
% * 第一变分和驻值
% * 密度及放大极限
% * 平面切锥与正则性问题
% * 整数流内部及边界正则性定理的具体范围
% 
% 纲要内容结束

\section{单调性、切锥与正则点}
\label{g3:sec:2303}

待撰写

\subsection{第一变分和驻值}
\label{g3:subsec:2303:01}
待撰写

\subsection{密度及放大极限}
\label{g3:subsec:2303:02}
待撰写

\subsection{平面切锥与正则性问题}
\label{g3:subsec:2303:03}
待撰写

\subsection{整数流内部及边界正则性定理的具体范围}
\label{g3:subsec:2303:04}
待撰写
% 纲要条目排为编号小节标题，并进入目录与书签；说明段仍保留于纲要注释中。

% !TeX root = ../../Book3.tex
% 纲要标题：### 23.4* 变分流形的另一种弱化
% 纲要位置：几何学计划纲要.md，第 2950 行。
% 正文状态：未撰写；本文件目前是框架。

% 纲要内容开始
% 
% * Grassmann 丛上的测度
% * 一阶变分和广义平均曲率
% * 无取向极限及重数
% * 与流的边界、抵消信息比较
% * Allard 正则性：精确输入与代表应用
% 
% ---
% 
% 纲要内容结束

\OptionalSection{变分流形的另一种弱化}{g3:sec:2304}

待撰写

\subsection{Grassmann 丛上的测度}
\label{g3:subsec:2304:01}
待撰写

\subsection{一阶变分和广义平均曲率}
\label{g3:subsec:2304:02}
待撰写

\subsection{无取向极限及重数}
\label{g3:subsec:2304:03}
待撰写

\subsection{与流的边界、抵消信息比较}
\label{g3:subsec:2304:04}
待撰写

\subsection{Allard 正则性：精确输入与代表应用}
\label{g3:subsec:2304:05}
待撰写
% 纲要条目排为编号小节标题，并进入目录与书签；说明段仍保留于纲要注释中。

% !TeX root = ../../Book3.tex
% 纲要标题：### 23.5*　扫掠、宽度与非极小临界点
% 纲要位置：几何学计划纲要.md，第 3107 行。

% 纲要内容开始
% 
% > 写作范围与外部输入：规划任务：在第20—22章的弱极限工具与第23.4节变分流形之后引入几何 min-max。完整证明黎曼二球面上 Birkhoff 非平凡闭测地线的代表模块，交代离散缩短、宽度为正和极限不塌缩为常值；这不是第28章的平均曲率流。再在闭环境维数3—7及指定系数、扫掠空间和收敛拓扑下，准确外引所选 Almgren–Pitts 极小超曲面存在性，分别说明重数、正则性和指标。有限维山路模型只作动机，不替代几何存在证明；第16.4节谱的极小极大原则与本节另行辨析。
% 
% * 扫掠族、同伦约束与宽度
% * 分段测地线的缩短与 Birkhoff 模型
% * 二球面闭测地线存在性的完整证明
% * 超曲面扫掠与驻值变分流形
% * Almgren–Pitts 存在性的准确版本
% * 维数、重数、正则性与指标的区别
% 
% ---
% 
% 纲要内容结束
% 正文状态：未撰写；本文件目前是框架。
% 规划任务：在第20—22章的弱极限工具与第23.4节变分流形之后引入几何 min-max。完整证明黎曼二球面上 Birkhoff 非平凡闭测地线的代表模块，交代离散缩短、宽度为正和极限不塌缩为常值；这不是第28章的平均曲率流。再在闭环境维数3—7及指定系数、扫掠空间和收敛拓扑下，准确外引所选 Almgren–Pitts 极小超曲面存在性，分别说明重数、正则性和指标。有限维山路模型只作动机，不替代几何存在证明；第16.4节谱的极小极大原则与本节另行辨析。

\OptionalSection{扫掠、宽度与非极小临界点}{g3:sec:2305}

待撰写

\subsection{扫掠族、同伦约束与宽度}
\label{g3:subsec:2305:01}
待撰写

\subsection{分段测地线的缩短与 Birkhoff 模型}
\label{g3:subsec:2305:02}
待撰写

\subsection{二球面闭测地线存在性的完整证明}
\label{g3:subsec:2305:03}
待撰写

\subsection{超曲面扫掠与驻值变分流形}
\label{g3:subsec:2305:04}
待撰写

\subsection{Almgren–Pitts 存在性的准确版本}
\label{g3:subsec:2305:05}
待撰写

\subsection{维数、重数、正则性与指标的区别}
\label{g3:subsec:2305:06}
待撰写


\begin{chaptersummary}
待撰写
\end{chaptersummary}

\BookExercises
待撰写


% !TeX root = ../Book3.tex
% 纲要标题：# 第六编　非线性方法、几何流与奇异极限
% 纲要位置：几何学计划纲要.md，第 3120 行。
% 正文状态：未撰写；本文件目前是框架。

% 纲要内容开始
% 
% 纲要内容结束

\BookPart{非线性方法、几何流与奇异极限}{g3:part:06}
\BookPartGuide
待撰写

% !TeX root = ../../Book3.tex
% 纲要标题：## 第24章*　Yamabe 问题与临界变分
% 纲要位置：几何学计划纲要.md，第 2962 行。
% 正文状态：未撰写；本文件目前是框架。

% 纲要内容开始
% 
% > 本章目标与证明范围：以共形标量曲率为贯穿问题，完整处理非正 Yamabe 情形及严格低于球面阈值时的紧性论证。正质量等完成一般正情形所需的深层输入明确外引，不把全部存在性挤进“典型结果”。
% 
% 纲要内容结束

\OptionalChapter{Yamabe 问题与临界变分}{g3:ch:24}
\BookChapterNumber{g3:ch:24}

\begin{chapterabstract}
待撰写
\end{chapterabstract}

\begin{learninggoals}
待撰写
\end{learninggoals}

待撰写

% !TeX root = ../../Book3.tex
% 纲要标题：### 24.1 Yamabe 问题的分析设置与规范化
% 纲要位置：几何学计划纲要.md，第 2966 行。
% 正文状态：未撰写；本文件目前是框架。

% 纲要内容开始
% 
% 本章不再首次建立通用非线性工具；它们的证明位置统一前移至第三编。
% 
% * 调用第13.5—13.6节的 Banach 微分、Sobolev 复合与约束变分
% * 固定共形类、正因子及体积约束
% * 弱极小化问题与光滑正解之间的转换
% * 零阶项符号、核与唯一性的不同规范
% 
% 纲要内容结束

\section{Yamabe 问题的分析设置与规范化}
\label{g3:sec:2401}

待撰写

\subsection{调用第13.5—13.6节的 Banach 微分、Sobolev 复合与约束变分}
\label{g3:subsec:2401:01}
待撰写

\subsection{固定共形类、正因子及体积约束}
\label{g3:subsec:2401:02}
待撰写

\subsection{弱极小化问题与光滑正解之间的转换}
\label{g3:subsec:2401:03}
待撰写

\subsection{零阶项符号、核与唯一性的不同规范}
\label{g3:subsec:2401:04}
待撰写
% 纲要条目排为编号小节标题，并进入目录与书签；说明段仍保留于纲要注释中。

% !TeX root = ../../Book3.tex
% 纲要标题：### 24.2 共形变化与变分方程
% 纲要位置：几何学计划纲要.md，第 2975 行。
% 正文状态：未撰写；本文件目前是框架。

% 纲要内容开始
% 
% * n≥3 时标量曲率变换
% * 共形 Laplace 的归一化
% * Yamabe 泛函和临界指数
% * Euler–Lagrange 方程及正性
% 
% 纲要内容结束

\section{共形变化与变分方程}
\label{g3:sec:2402}

待撰写

\subsection{\texorpdfstring{\(n\ge3\) 时标量曲率变换}{n≥3 时标量曲率变换}}
\label{g3:subsec:2402:01}
待撰写

\subsection{共形 Laplace 的归一化}
\label{g3:subsec:2402:02}
待撰写

\subsection{Yamabe 泛函和临界指数}
\label{g3:subsec:2402:03}
待撰写

\subsection{Euler–Lagrange 方程及正性}
\label{g3:subsec:2402:04}
待撰写
% 纲要条目排为编号小节标题，并进入目录与书签；说明段仍保留于纲要注释中。

% !TeX root = ../../Book3.tex
% 纲要标题：### 24.3 非正情形的存在与唯一性
% 纲要位置：几何学计划纲要.md，第 2982 行。
% 正文状态：未撰写；本文件目前是框架。

% 纲要内容开始
% 
% * 次临界逼近或直接极小化
% * 非正阈值下的紧性
% * 最大值原理与正解
% * 规范体积后的唯一性范围
% 
% 纲要内容结束

\section{非正情形的存在与唯一性}
\label{g3:sec:2403}

待撰写

\subsection{次临界逼近或直接极小化}
\label{g3:subsec:2403:01}
待撰写

\subsection{非正阈值下的紧性}
\label{g3:subsec:2403:02}
待撰写

\subsection{最大值原理与正解}
\label{g3:subsec:2403:03}
待撰写

\subsection{规范体积后的唯一性范围}
\label{g3:subsec:2403:04}
待撰写
% 纲要条目排为编号小节标题，并进入目录与书签；说明段仍保留于纲要注释中。

% !TeX root = ../../Book3.tex
% 纲要标题：### 24.4 临界紧性的失效
% 纲要位置：几何学计划纲要.md，第 2989 行。
% 正文状态：未撰写；本文件目前是框架。

% 纲要内容开始
% 
% * Euclid 气泡与尺度不变性
% * 集中紧性的所选引理
% * 低于球面常数时排除集中
% * 严格阈值下的极小值实现
% 
% 纲要内容结束

\section{临界紧性的失效}
\label{g3:sec:2404}

待撰写

\subsection{Euclid 气泡与尺度不变性}
\label{g3:subsec:2404:01}
待撰写

\subsection{集中紧性的所选引理}
\label{g3:subsec:2404:02}
待撰写

\subsection{低于球面常数时排除集中}
\label{g3:subsec:2404:03}
待撰写

\subsection{严格阈值下的极小值实现}
\label{g3:subsec:2404:04}
待撰写
% 纲要条目排为编号小节标题，并进入目录与书签；说明段仍保留于纲要注释中。

% !TeX root = ../../Book3.tex
% 纲要标题：### 24.5 一般存在性的剩余输入
% 纲要位置：几何学计划纲要.md，第 2996 行。
% 正文状态：未撰写；本文件目前是框架。

% 纲要内容开始
% 
% * Aubin 型测试函数的几何意义
% * 低维或局部共形平坦情形的额外困难
% * 正质量定理与 Schoen 方法的精确引用
% * 局部极小、全局极小和非唯一解
% 
% ---
% 
% 纲要内容结束

\section{一般存在性的剩余输入}
\label{g3:sec:2405}

待撰写

\subsection{Aubin 型测试函数的几何意义}
\label{g3:subsec:2405:01}
待撰写

\subsection{低维或局部共形平坦情形的额外困难}
\label{g3:subsec:2405:02}
待撰写

\subsection{正质量定理与 Schoen 方法的精确引用}
\label{g3:subsec:2405:03}
待撰写

\subsection{局部极小、全局极小和非唯一解}
\label{g3:subsec:2405:04}
待撰写
% 纲要条目排为编号小节标题，并进入目录与书签；说明段仍保留于纲要注释中。


\begin{chaptersummary}
待撰写
\end{chaptersummary}

\BookExercises
待撰写

% !TeX root = ../../Book3.tex
% 纲要标题：## 第25章*　Einstein 方程、规范与局部模问题
% 纲要位置：几何学计划纲要.md，第 3005 行。
% 正文状态：未撰写；本文件目前是框架。

% 纲要内容开始
% 
% > 本章目标与证明范围：与 Yamabe 分章：Einstein 方程是模微分同胚的系统，而非标量共形方程。完整建立线性化、规范固定和有限维局部障碍模型；一般存在与分类不在本章无条件承诺。 非线性共同输入定位于 III.13.5—13.6，不以完整阅读 Yamabe 章为先修。
% 
% 纲要内容结束

\OptionalChapter{Einstein 方程、规范与局部模问题}{g3:ch:25}
\BookChapterNumber{g3:ch:25}

\begin{chapterabstract}
待撰写
\end{chapterabstract}

\begin{learninggoals}
待撰写
\end{learninggoals}

待撰写

% !TeX root = ../../Book3.tex
% 纲要标题：### 25.1 度量空间和曲率变分
% 纲要位置：几何学计划纲要.md，第 3169 行。
% 正文状态：未撰写；本文件目前是框架。

% 纲要内容开始
% 
% * 对称二张量的切空间
% * Ricci 与标量曲率的线性化
% * 微分同胚方向的退化
% * Einstein–Hilbert 泛函和体积约束
% 
% 纲要内容结束

\section{度量空间和曲率变分}
\label{g3:sec:2501}

待撰写

\subsection{对称二张量的切空间}
\label{g3:subsec:2501:01}
待撰写

\subsection{Ricci 与标量曲率的线性化}
\label{g3:subsec:2501:02}
待撰写

\subsection{微分同胚方向的退化}
\label{g3:subsec:2501:03}
待撰写

\subsection{Einstein–Hilbert 泛函和体积约束}
\label{g3:subsec:2501:04}
待撰写
% 纲要条目排为编号小节标题，并进入目录与书签；说明段仍保留于纲要注释中。

% !TeX root = ../../Book3.tex
% 纲要标题：### 25.2 规范固定与椭圆化
% 纲要位置：几何学计划纲要.md，第 3016 行。
% 正文状态：未撰写；本文件目前是框架。

% 纲要内容开始
% 
% * 散度、Bianchi 算子和规范条件
% * 规范化 Einstein 算子的主符号
% * 切片所需的椭圆方程
% * 无穷小等距与剩余对称
% 
% 纲要内容结束

\section{规范固定与椭圆化}
\label{g3:sec:2502}

待撰写

\subsection{散度、Bianchi 算子和规范条件}
\label{g3:subsec:2502:01}
待撰写

\subsection{规范化 Einstein 算子的主符号}
\label{g3:subsec:2502:02}
待撰写

\subsection{切片所需的椭圆方程}
\label{g3:subsec:2502:03}
待撰写

\subsection{无穷小等距与剩余对称}
\label{g3:subsec:2502:04}
待撰写
% 纲要条目排为编号小节标题，并进入目录与书签；说明段仍保留于纲要注释中。

% !TeX root = ../../Book3.tex
% 纲要标题：### 25.3 无穷小形变及障碍
% 纲要位置：几何学计划纲要.md，第 3023 行。
% 正文状态：未撰写；本文件目前是框架。

% 纲要内容开始
% 
% * Lichnerowicz 型算子
% * 核、余核和规范等价
% * 隐函数定理的可解条件
% * 有限维约化及二阶障碍
% 
% 纲要内容结束

\section{无穷小形变及障碍}
\label{g3:sec:2503}

待撰写

\subsection{Lichnerowicz 型算子}
\label{g3:subsec:2503:01}
待撰写

\subsection{核、余核和规范等价}
\label{g3:subsec:2503:02}
待撰写

\subsection{隐函数定理的可解条件}
\label{g3:subsec:2503:03}
待撰写

\subsection{有限维约化及二阶障碍}
\label{g3:subsec:2503:04}
待撰写
% 纲要条目排为编号小节标题，并进入目录与书签；说明段仍保留于纲要注释中。

% !TeX root = ../../Book3.tex
% 纲要标题：### 25.4 模型和几何限制
% 纲要位置：几何学计划纲要.md，第 3030 行。
% 正文状态：未撰写；本文件目前是框架。

% 纲要内容开始
% 
% * 空间形式、乘积及 Ricci 平坦模型
% * 自同构导致的模空间奇点
% * 紧性、退化及更深存在定理
% 
% ---
% 
% 纲要内容结束

\section{模型和几何限制}
\label{g3:sec:2504}

待撰写

\subsection{空间形式、乘积及 Ricci 平坦模型}
\label{g3:subsec:2504:01}
待撰写

\subsection{自同构导致的模空间奇点}
\label{g3:subsec:2504:02}
待撰写

\subsection{紧性、退化及更深存在定理}
\label{g3:subsec:2504:03}
待撰写
% 纲要条目排为编号小节标题，并进入目录与书签；说明段仍保留于纲要注释中。


\begin{chaptersummary}
待撰写
\end{chaptersummary}

\BookExercises
待撰写

% !TeX root = ../../Book3.tex
% 纲要标题：## 第26章*　抛物方法与调和映射热流
% 纲要位置：几何学计划纲要.md，第 3038 行。
% 正文状态：未撰写；本文件目前是框架。

% 纲要内容开始
% 
% > 本章目标与证明范围：作为几何流路线的工具入口，可不经过 Yamabe 和 Einstein 章。先证明紧光滑情形所需线性及拟线性短时理论，再在非正目标曲率下完成 Eells–Sampson 型代表模块。 非线性共同输入定位于 III.13.5—13.6，不以完整阅读 Yamabe 章为先修。
% 
% 纲要内容结束

\OptionalChapter{抛物方法与调和映射热流}{g3:ch:26}
\BookChapterNumber{g3:ch:26}

\begin{chapterabstract}
待撰写
\end{chapterabstract}

\begin{learninggoals}
待撰写
\end{learninggoals}

待撰写

% !TeX root = ../../Book3.tex
% 纲要标题：### 26.1 线性抛物方程
% 纲要位置：几何学计划纲要.md，第 3202 行。
% 正文状态：未撰写；本文件目前是框架。

% 纲要内容开始
% 
% * 热核和 Duhamel 公式
% * 抛物 Hölder 或 Sobolev 空间
% * 线性估计与初值相容
% * 抛物最大值原理
% 
% 纲要内容结束

\section{线性抛物方程}
\label{g3:sec:2601}

待撰写

\subsection{热核和 Duhamel 公式}
\label{g3:subsec:2601:01}
待撰写

\subsection{抛物 Hölder 或 Sobolev 空间}
\label{g3:subsec:2601:02}
待撰写

\subsection{线性估计与初值相容}
\label{g3:subsec:2601:03}
待撰写

\subsection{抛物最大值原理}
\label{g3:subsec:2601:04}
待撰写
% 纲要条目排为编号小节标题，并进入目录与书签；说明段仍保留于纲要注释中。

% !TeX root = ../../Book3.tex
% 纲要标题：### 26.2 拟线性短时存在
% 纲要位置：几何学计划纲要.md，第 3049 行。
% 正文状态：未撰写；本文件目前是框架。

% 纲要内容开始
% 
% * 强抛物性及线性化
% * 迭代或逆函数方法
% * 短时唯一性和光滑依赖
% * 延拓判据与几何先验界
% 
% 纲要内容结束

\section{拟线性短时存在}
\label{g3:sec:2602}

待撰写

\subsection{强抛物性及线性化}
\label{g3:subsec:2602:01}
待撰写

\subsection{迭代或逆函数方法}
\label{g3:subsec:2602:02}
待撰写

\subsection{短时唯一性和光滑依赖}
\label{g3:subsec:2602:03}
待撰写

\subsection{延拓判据与几何先验界}
\label{g3:subsec:2602:04}
待撰写
% 纲要条目排为编号小节标题，并进入目录与书签；说明段仍保留于纲要注释中。

% !TeX root = ../../Book3.tex
% 纲要标题：### 26.3 调和映射热流的基本演化
% 纲要位置：几何学计划纲要.md，第 3216 行。
% 正文状态：未撰写；本文件目前是框架。

% 纲要内容开始
% 
% * 张力场热方程及几何不变性
% * 能量下降和时间导数的积分控制
% * 梯度 Bochner 演化
% * 非正目标曲率的符号作用
% 
% 纲要内容结束

\section{调和映射热流的基本演化}
\label{g3:sec:2603}

待撰写

\subsection{张力场热方程及几何不变性}
\label{g3:subsec:2603:01}
待撰写

\subsection{能量下降和时间导数的积分控制}
\label{g3:subsec:2603:02}
待撰写

\subsection{梯度 Bochner 演化}
\label{g3:subsec:2603:03}
待撰写

\subsection{非正目标曲率的符号作用}
\label{g3:subsec:2603:04}
待撰写
% 纲要条目排为编号小节标题，并进入目录与书签；说明段仍保留于纲要注释中。

% !TeX root = ../../Book3.tex
% 纲要标题：### 26.4 长期存在的先验控制
% 纲要位置：几何学计划纲要.md，第 3223 行。
% 正文状态：未撰写；本文件目前是框架。

% 纲要内容开始
% 
% * 闭源上的最大值估计
% * 紧目标下的局部坐标控制
% * 高阶抛物估计
% * 由延拓判据排除有限时奇点
% 
% 纲要内容结束

\section{长期存在的先验控制}
\label{g3:sec:2604}

待撰写

\subsection{闭源上的最大值估计}
\label{g3:subsec:2604:01}
待撰写

\subsection{紧目标下的局部坐标控制}
\label{g3:subsec:2604:02}
待撰写

\subsection{高阶抛物估计}
\label{g3:subsec:2604:03}
待撰写

\subsection{由延拓判据排除有限时奇点}
\label{g3:subsec:2604:04}
待撰写
% 纲要条目排为编号小节标题，并进入目录与书签；说明段仍保留于纲要注释中。

% !TeX root = ../../Book3.tex
% 纲要标题：### 26.5 调和代表的极限构造
% 纲要位置：几何学计划纲要.md，第 3230 行。
% 正文状态：未撰写；本文件目前是框架。

% 纲要内容开始
% 
% * 时间子列上的速度衰减
% * 统一高阶界及光滑子列
% * 极限满足调和方程
% * 同伦类的保持
% * 全时收敛和唯一性所需的附加假设
% 
% ---
% 
% 纲要内容结束

\section{调和代表的极限构造}
\label{g3:sec:2605}

待撰写

\subsection{时间子列上的速度衰减}
\label{g3:subsec:2605:01}
待撰写

\subsection{统一高阶界及光滑子列}
\label{g3:subsec:2605:02}
待撰写

\subsection{极限满足调和方程}
\label{g3:subsec:2605:03}
待撰写

\subsection{同伦类的保持}
\label{g3:subsec:2605:04}
待撰写

\subsection{全时收敛和唯一性所需的附加假设}
\label{g3:subsec:2605:05}
待撰写
% 纲要条目排为编号小节标题，并进入目录与书签；说明段仍保留于纲要注释中。


\begin{chaptersummary}
待撰写
\end{chaptersummary}

\BookExercises
待撰写

% !TeX root = ../../Book3.tex
% 纲要标题：## 第27章*　Ricci 流、规范消除与奇点模型
% 纲要位置：几何学计划纲要.md，第 3240 行。
% 正文状态：未撰写；本文件目前是框架。

% 纲要内容开始
% 
% > 本章目标与证明范围：紧光滑初始度量的短时存在、唯一性和基本演化为完整模块。熵、非塌缩和手术按实际选取的深层结果区分证明与引用，不以二维算例替代高维理论。
% 
% 纲要内容结束

\OptionalChapter{Ricci 流、规范消除与奇点模型}{g3:ch:27}
\BookChapterNumber{g3:ch:27}

\begin{chapterabstract}
待撰写
\end{chapterabstract}

\begin{learninggoals}
待撰写
\end{learninggoals}

待撰写

% !TeX root = ../../Book3.tex
% 纲要标题：### 27.1 Ricci 流的几何退化
% 纲要位置：几何学计划纲要.md，第 3244 行。
% 正文状态：未撰写；本文件目前是框架。

% 纲要内容开始
% 
% * 方程、缩放和微分同胚不变性
% * 线性化中的非抛物方向
% * DeTurck 向量场与严格抛物方程
% 
% 纲要内容结束

\section{Ricci 流的几何退化}
\label{g3:sec:2701}

待撰写

\subsection{方程、缩放和微分同胚不变性}
\label{g3:subsec:2701:01}
待撰写

\subsection{线性化中的非抛物方向}
\label{g3:subsec:2701:02}
待撰写

\subsection{DeTurck 向量场与严格抛物方程}
\label{g3:subsec:2701:03}
待撰写
% 纲要条目排为编号小节标题，并进入目录与书签；说明段仍保留于纲要注释中。

% !TeX root = ../../Book3.tex
% 纲要标题：### 27.2 短时理论的证明
% 纲要位置：几何学计划纲要.md，第 3250 行。
% 正文状态：未撰写；本文件目前是框架。

% 纲要内容开始
% 
% * 调用第26章的拟线性存在
% * 时间依赖微分同胚的构造
% * 拉回恢复 Ricci 流
% * 唯一性的规范比较
% * 曲率界下的延拓判据
% 
% 纲要内容结束

\section{短时理论的证明}
\label{g3:sec:2702}

待撰写

\subsection{调用第26章的拟线性存在}
\label{g3:subsec:2702:01}
待撰写

\subsection{时间依赖微分同胚的构造}
\label{g3:subsec:2702:02}
待撰写

\subsection{拉回恢复 Ricci 流}
\label{g3:subsec:2702:03}
待撰写

\subsection{唯一性的规范比较}
\label{g3:subsec:2702:04}
待撰写

\subsection{曲率界下的延拓判据}
\label{g3:subsec:2702:05}
待撰写
% 纲要条目排为编号小节标题，并进入目录与书签；说明段仍保留于纲要注释中。

% !TeX root = ../../Book3.tex
% 纲要标题：### 27.3 曲率演化和保持性质
% 纲要位置：几何学计划纲要.md，第 3098 行。
% 正文状态：未撰写；本文件目前是框架。

% 纲要内容开始
% 
% * 体积、标量及 Ricci 演化
% * 完整曲率张量的反应项
% * 标量最大值与张量最大值的区别
% * 所选曲率锥的保持
% 
% 纲要内容结束

\section{曲率演化和保持性质}
\label{g3:sec:2703}

待撰写

\subsection{体积、标量及 Ricci 演化}
\label{g3:subsec:2703:01}
待撰写

\subsection{完整曲率张量的反应项}
\label{g3:subsec:2703:02}
待撰写

\subsection{标量最大值与张量最大值的区别}
\label{g3:subsec:2703:03}
待撰写

\subsection{所选曲率锥的保持}
\label{g3:subsec:2703:04}
待撰写
% 纲要条目排为编号小节标题，并进入目录与书签；说明段仍保留于纲要注释中。

% !TeX root = ../../Book3.tex
% 纲要标题：### 27.4 放大、孤立子和熵
% 纲要位置：几何学计划纲要.md，第 3105 行。
% 正文状态：未撰写；本文件目前是框架。

% 纲要内容开始
% 
% * 奇点尺度及抛物重缩放
% * 收缩、稳态和膨胀孤立子
% * 单调泛函的构造或精确引用
% * 非塌缩定理的输入与用途
% 
% 纲要内容结束

\section{放大、孤立子和熵}
\label{g3:sec:2704}

待撰写

\subsection{奇点尺度及抛物重缩放}
\label{g3:subsec:2704:01}
待撰写

\subsection{收缩、稳态和膨胀孤立子}
\label{g3:subsec:2704:02}
待撰写

\subsection{单调泛函的构造或精确引用}
\label{g3:subsec:2704:03}
待撰写

\subsection{非塌缩定理的输入与用途}
\label{g3:subsec:2704:04}
待撰写
% 纲要条目排为编号小节标题，并进入目录与书签；说明段仍保留于纲要注释中。

% !TeX root = ../../Book3.tex
% 纲要标题：### 27.5 曲面与齐性算例
% 纲要位置：几何学计划纲要.md，第 3112 行。
% 正文状态：未撰写；本文件目前是框架。

% 纲要内容开始
% 
% * 常曲率曲面的显式演化
% * 归一化和非归一化流
% * 齐性度量的常微分方程
% * 几何化和手术理论的明确边界
% 
% ---
% 
% 纲要内容结束

\section{曲面与齐性算例}
\label{g3:sec:2705}

待撰写

\subsection{常曲率曲面的显式演化}
\label{g3:subsec:2705:01}
待撰写

\subsection{归一化和非归一化流}
\label{g3:subsec:2705:02}
待撰写

\subsection{齐性度量的常微分方程}
\label{g3:subsec:2705:03}
待撰写

\subsection{几何化和手术理论的明确边界}
\label{g3:subsec:2705:04}
待撰写
% 纲要条目排为编号小节标题，并进入目录与书签；说明段仍保留于纲要注释中。


\begin{chaptersummary}
待撰写
\end{chaptersummary}

\BookExercises
待撰写

% !TeX root = ../../Book3.tex
% 纲要标题：## 第28章*　平均曲率流、单调性与弱演化
% 纲要位置：几何学计划纲要.md，第 3121 行。
% 正文状态：未撰写；本文件目前是框架。

% 纲要内容开始
% 
% > 本章目标与证明范围：完成光滑紧初始子流形的短时理论、面积下降及欧氏 Huisken 单调公式。水平集和 Brakke 弱演化有独立定义；奇点后唯一性不由单调量自动得到。
% 
% 纲要内容结束

\OptionalChapter{平均曲率流、单调性与弱演化}{g3:ch:28}
\BookChapterNumber{g3:ch:28}

\begin{chapterabstract}
待撰写
\end{chapterabstract}

\begin{learninggoals}
待撰写
\end{learninggoals}

待撰写

% !TeX root = ../../Book3.tex
% 纲要标题：### 28.1 方程与光滑短时理论
% 纲要位置：几何学计划纲要.md，第 3125 行。
% 正文状态：未撰写；本文件目前是框架。

% 纲要内容开始
% 
% * 参数退化及法向速度
% * 图像方程的抛物性
% * 规范与短时存在
% * 面积下降及边界条件
% 
% 纲要内容结束

\section{方程与光滑短时理论}
\label{g3:sec:2801}

待撰写

\subsection{参数退化及法向速度}
\label{g3:subsec:2801:01}
待撰写

\subsection{图像方程的抛物性}
\label{g3:subsec:2801:02}
待撰写

\subsection{规范与短时存在}
\label{g3:subsec:2801:03}
待撰写

\subsection{面积下降及边界条件}
\label{g3:subsec:2801:04}
待撰写
% 纲要条目排为编号小节标题，并进入目录与书签；说明段仍保留于纲要注释中。

% !TeX root = ../../Book3.tex
% 纲要标题：### 28.2 演化公式与比较
% 纲要位置：几何学计划纲要.md，第 3292 行。
% 正文状态：未撰写；本文件目前是框架。

% 纲要内容开始
% 
% * 第二基本形式和平均曲率的演化
% * 超曲面避免相交原理
% * 凸性或平均凸性的保持条件
% * 曲线缩短流和球面模型
% 
% 纲要内容结束

\section{演化公式与比较}
\label{g3:sec:2802}

待撰写

\subsection{第二基本形式和平均曲率的演化}
\label{g3:subsec:2802:01}
待撰写

\subsection{超曲面避免相交原理}
\label{g3:subsec:2802:02}
待撰写

\subsection{凸性或平均凸性的保持条件}
\label{g3:subsec:2802:03}
待撰写

\subsection{曲线缩短流和球面模型}
\label{g3:subsec:2802:04}
待撰写
% 纲要条目排为编号小节标题，并进入目录与书签；说明段仍保留于纲要注释中。

% !TeX root = ../../Book3.tex
% 纲要标题：### 28.3 Huisken 单调公式
% 纲要位置：几何学计划纲要.md，第 3139 行。
% 正文状态：未撰写；本文件目前是框架。

% 纲要内容开始
% 
% * 后向热核及加权面积
% * 微分和分部积分
% * 平方项与等号条件
% * 自收缩子和奇点放大
% 
% 纲要内容结束

\section{Huisken 单调公式}
\label{g3:sec:2803}

待撰写

\subsection{后向热核及加权面积}
\label{g3:subsec:2803:01}
待撰写

\subsection{微分和分部积分}
\label{g3:subsec:2803:02}
待撰写

\subsection{平方项与等号条件}
\label{g3:subsec:2803:03}
待撰写

\subsection{自收缩子和奇点放大}
\label{g3:subsec:2803:04}
待撰写
% 纲要条目排为编号小节标题，并进入目录与书签；说明段仍保留于纲要注释中。

% !TeX root = ../../Book3.tex
% 纲要标题：### 28.4 弱演化的不同对象
% 纲要位置：几何学计划纲要.md，第 3306 行。
% 正文状态：未撰写；本文件目前是框架。

% 纲要内容开始
% 
% * 水平集方程及粘性解
% * Brakke 不等式及测度族
% * 质量丢失、重数和肥化
% * 存在、唯一性及正则性的不同结论范围
% 
% ---
% 
% 纲要内容结束

\section{弱演化的不同对象}
\label{g3:sec:2804}

待撰写

\subsection{水平集方程及粘性解}
\label{g3:subsec:2804:01}
待撰写

\subsection{Brakke 不等式及测度族}
\label{g3:subsec:2804:02}
待撰写

\subsection{质量丢失、重数和肥化}
\label{g3:subsec:2804:03}
待撰写

\subsection{存在、唯一性及正则性的不同结论范围}
\label{g3:subsec:2804:04}
待撰写
% 纲要条目排为编号小节标题，并进入目录与书签；说明段仍保留于纲要注释中。


\begin{chaptersummary}
待撰写
\end{chaptersummary}

\BookExercises
待撰写

% !TeX root = ../../Book3.tex
% 纲要标题：## 第29章*　几何收敛、塌缩与奇异结构
% 纲要位置：几何学计划纲要.md，第 3155 行。
% 正文状态：未撰写；本文件目前是框架。

% 纲要内容开始
% 
% > 本章目标与证明范围：从坐标控制建立光滑预紧性，再比较 GH 和度量测度极限。曲率导数、注入半径及体积假设各有用途；所选更深正则性结论精确外引。
% 
% 纲要内容结束

\OptionalChapter{几何收敛、塌缩与奇异结构}{g3:ch:29}
\BookChapterNumber{g3:ch:29}

\begin{chapterabstract}
待撰写
\end{chapterabstract}

\begin{learninggoals}
待撰写
\end{learninggoals}

待撰写

% !TeX root = ../../Book3.tex
% 纲要标题：### 29.1 光滑紧性的方法
% 纲要位置：几何学计划纲要.md，第 3159 行。
% 正文状态：未撰写；本文件目前是框架。

% 纲要内容开始
% 
% * 谐和坐标或法坐标的统一控制
% * 过渡映射与对角子列
% * 曲率导数及注入半径条件下的紧性
% * Cheeger–Gromov 收敛与微分同胚选择
% 
% 纲要内容结束

\section{光滑紧性的方法}
\label{g3:sec:2901}

待撰写

\subsection{谐和坐标或法坐标的统一控制}
\label{g3:subsec:2901:01}
待撰写

\subsection{过渡映射与对角子列}
\label{g3:subsec:2901:02}
待撰写

\subsection{曲率导数及注入半径条件下的紧性}
\label{g3:subsec:2901:03}
待撰写

\subsection{Cheeger–Gromov 收敛与微分同胚选择}
\label{g3:subsec:2901:04}
待撰写
% 纲要条目排为编号小节标题，并进入目录与书签；说明段仍保留于纲要注释中。

% !TeX root = ../../Book3.tex
% 纲要标题：### 29.2 塌缩与维数变化
% 纲要位置：几何学计划纲要.md，第 3326 行。
% 正文状态：未撰写；本文件目前是框架。

% 纲要内容开始
% 
% * 扁平环面和圆丛塌缩
% * 体积、直径和曲率的不同控制
% * 局部纤维结构的深层定理接口
% * 极限维数与拓扑维数
% 
% 纲要内容结束

\section{塌缩与维数变化}
\label{g3:sec:2902}

待撰写

\subsection{扁平环面和圆丛塌缩}
\label{g3:subsec:2902:01}
待撰写

\subsection{体积、直径和曲率的不同控制}
\label{g3:subsec:2902:02}
待撰写

\subsection{局部纤维结构的深层定理接口}
\label{g3:subsec:2902:03}
待撰写

\subsection{极限维数与拓扑维数}
\label{g3:subsec:2902:04}
待撰写
% 纲要条目排为编号小节标题，并进入目录与书签；说明段仍保留于纲要注释中。

% !TeX root = ../../Book3.tex
% 纲要标题：### 29.3 Ricci 下界的极限
% 纲要位置：几何学计划纲要.md，第 3333 行。
% 正文状态：未撰写；本文件目前是框架。

% 纲要内容开始
% 
% * 体积比较及非塌缩条件
% * 测度的归一化和收敛
% * 正则点、切锥及奇异集
% * 所选结构定理的准确范围
% 
% 纲要内容结束

\section{Ricci 下界的极限}
\label{g3:sec:2903}

待撰写

\subsection{体积比较及非塌缩条件}
\label{g3:subsec:2903:01}
待撰写

\subsection{测度的归一化和收敛}
\label{g3:subsec:2903:02}
待撰写

\subsection{正则点、切锥及奇异集}
\label{g3:subsec:2903:03}
待撰写

\subsection{所选结构定理的准确范围}
\label{g3:subsec:2903:04}
待撰写
% 纲要条目排为编号小节标题，并进入目录与书签；说明段仍保留于纲要注释中。

% !TeX root = ../../Book3.tex
% 纲要标题：### 29.4 合成曲率与弱微分
% 纲要位置：几何学计划纲要.md，第 3180 行。
% 正文状态：未撰写；本文件目前是框架。

% 纲要内容开始
% 
% * Cheeger 能量的动机
% * 最优输运和曲率维数条件
% * 光滑流形与合成框架的比较
% * 稳定性结果和进一步正则性问题
% 
% ---
% 
% 纲要内容结束

\section{合成曲率与弱微分}
\label{g3:sec:2904}

待撰写

\subsection{Cheeger 能量的动机}
\label{g3:subsec:2904:01}
待撰写

\subsection{最优输运和曲率维数条件}
\label{g3:subsec:2904:02}
待撰写

\subsection{光滑流形与合成框架的比较}
\label{g3:subsec:2904:03}
待撰写

\subsection{稳定性结果和进一步正则性问题}
\label{g3:subsec:2904:04}
待撰写
% 纲要条目排为编号小节标题，并进入目录与书签；说明段仍保留于纲要注释中。


\begin{chaptersummary}
待撰写
\end{chaptersummary}

\BookExercises
待撰写


% !TeX root = ../Book3.tex
% 纲要标题：# 第三卷附录
% 纲要位置：几何学计划纲要.md，第 3189 行。
% 正文状态：未撰写；本文件目前是框架。

% 纲要内容开始
% 
% 纲要内容结束

\appendix
\renewcommand{\thechapter}{\Alph{chapter}}
\ctexset{chapter={number=\Alph{chapter}}}
\BookAppendixPart{g3:part:appendices}

% !TeX root = ../../Book3.tex
% 纲要标题：## 附录 A　伪黎曼与 Lorentz 几何
% 纲要位置：几何学计划纲要.md，第 3351 行。
% 正文状态：未撰写；本文件目前是框架。

% 纲要内容开始
% 
% > 从不定度量出发说明哪些黎曼证明失效；不将完备性、距离和因果性混用。
% 
% * **局部结构与测地线**
%   * 不定配对及 Levi-Civita 联络
%   * 时间、空间和零向量
%   * 测地线和光锥
%   * 度量完备性的不同含义
% 
% * **因果结构与整体问题**
%   * 时间定向和因果曲线
%   * Cauchy 超曲面及全局双曲性
%   * Minkowski 与典型时空模型
%   * 奇点定理的条件和外部输入
% 
% ---
% 
% 纲要内容结束

\chapter{伪黎曼与 Lorentz 几何}
\label{g3:app:A}
\BookChapterNumber{g3:app:A}

待撰写

\section{局部结构与测地线}
\label{g3:sec:A01}
待撰写

\subsection{不定配对及 Levi-Civita 联络}
\label{g3:subsec:A01:01}
待撰写

\subsection{时间、空间和零向量}
\label{g3:subsec:A01:02}
待撰写

\subsection{测地线和光锥}
\label{g3:subsec:A01:03}
待撰写

\subsection{度量完备性的不同含义}
\label{g3:subsec:A01:04}
待撰写

\section{因果结构与整体问题}
\label{g3:sec:A02}
待撰写

\subsection{时间定向和因果曲线}
\label{g3:subsec:A02:01}
待撰写

\subsection{Cauchy 超曲面及全局双曲性}
\label{g3:subsec:A02:02}
待撰写

\subsection{Minkowski 与典型时空模型}
\label{g3:subsec:A02:03}
待撰写

\subsection{奇点定理的条件和外部输入}
\label{g3:subsec:A02:04}
待撰写

% !TeX root = ../../Book3.tex
% 纲要标题：## 附录 B　规范场、指标与特殊完整群
% 纲要位置：几何学计划纲要.md，第 3209 行。
% 正文状态：未撰写；本文件目前是框架。

% 纲要内容开始
% 
% > 利用正文丛联络和椭圆工具建立可计算的规范场模型；深层紧性和模空间理论精确引用。
% 
% * **Yang–Mills 变分**
%   * 联络空间及规范群
%   * 曲率能量和 Euler–Lagrange 方程
%   * 四维自对偶分解
%   * 拓扑下界与瞬子模型
% 
% * **椭圆化和模空间**
%   * Coulomb 规范的局部问题
%   * 形变复形与指标
%   * Uhlenbeck 紧性的准确输入
%   * 可去奇点及气泡能量
% 
% * **自旋与平行结构**
%   * Dirac 核和指标约束
%   * 平行自旋子与特殊完整群
%   * Kähler、超 Kähler 和例外模型
% 
% ---
% 
% 纲要内容结束

\chapter{规范场、指标与特殊完整群}
\label{g3:app:B}
\BookChapterNumber{g3:app:B}

待撰写

\section{Yang–Mills 变分}
\label{g3:sec:B01}
待撰写

\subsection{联络空间及规范群}
\label{g3:subsec:B01:01}
待撰写

\subsection{曲率能量和 Euler–Lagrange 方程}
\label{g3:subsec:B01:02}
待撰写

\subsection{四维自对偶分解}
\label{g3:subsec:B01:03}
待撰写

\subsection{拓扑下界与瞬子模型}
\label{g3:subsec:B01:04}
待撰写

\section{椭圆化和模空间}
\label{g3:sec:B02}
待撰写

\subsection{Coulomb 规范的局部问题}
\label{g3:subsec:B02:01}
待撰写

\subsection{形变复形与指标}
\label{g3:subsec:B02:02}
待撰写

\subsection{Uhlenbeck 紧性的准确输入}
\label{g3:subsec:B02:03}
待撰写

\subsection{可去奇点及气泡能量}
\label{g3:subsec:B02:04}
待撰写

\section{自旋与平行结构}
\label{g3:sec:B03}
待撰写

\subsection{Dirac 核和指标约束}
\label{g3:subsec:B03:01}
待撰写

\subsection{平行自旋子与特殊完整群}
\label{g3:subsec:B03:02}
待撰写

\subsection{Kähler、超 Kähler 和例外模型}
\label{g3:subsec:B03:03}
待撰写

% !TeX root = ../../Book3.tex
% 纲要标题：## 附录 C　非光滑度量测度空间
% 纲要位置：几何学计划纲要.md，第 3232 行。
% 正文状态：未撰写；本文件目前是框架。

% 纲要内容开始
% 
% > 与正文第29章互补；从一个实际能量和运输问题出发，不堆列合成几何缩写。
% 
% * **运输与测地结构**
%   * 耦合、Kantorovich 问题和 Wasserstein 距离
%   * 存在性的紧性方法
%   * 测度空间中的测地线
% 
% * **弱梯度与能量**
%   * 上梯度及松弛能量
%   * Sobolev 空间和 Hilbert 性
%   * 热流作为能量梯度流
% 
% * **曲率与极限**
%   * 熵凸性和曲率维数条件
%   * RCD 中附加的二次能量要求
%   * 模型和收敛稳定性
%   * 深层奇异集理论的边界
% 
% ---
% 
% 纲要内容结束

\chapter{非光滑度量测度空间}
\label{g3:app:C}
\BookChapterNumber{g3:app:C}

待撰写

\section{运输与测地结构}
\label{g3:sec:C01}
待撰写

\subsection{耦合、Kantorovich 问题和 Wasserstein 距离}
\label{g3:subsec:C01:01}
待撰写

\subsection{存在性的紧性方法}
\label{g3:subsec:C01:02}
待撰写

\subsection{测度空间中的测地线}
\label{g3:subsec:C01:03}
待撰写

\section{弱梯度与能量}
\label{g3:sec:C02}
待撰写

\subsection{上梯度及松弛能量}
\label{g3:subsec:C02:01}
待撰写

\subsection{Sobolev 空间和 Hilbert 性}
\label{g3:subsec:C02:02}
待撰写

\subsection{热流作为能量梯度流}
\label{g3:subsec:C02:03}
待撰写

\section{曲率与极限}
\label{g3:sec:C03}
待撰写

\subsection{熵凸性和曲率维数条件}
\label{g3:subsec:C03:01}
待撰写

\subsection{RCD 中附加的二次能量要求}
\label{g3:subsec:C03:02}
待撰写

\subsection{模型和收敛稳定性}
\label{g3:subsec:C03:03}
待撰写

\subsection{深层奇异集理论的边界}
\label{g3:subsec:C03:04}
待撰写



\backmatter
\cleardoublepage
\pagenumbering{roman}
\bookmarksetup{startatroot}
% !TeX root = ../Book3.tex
\cleardoublepage
% !TeX root = ../../Book3.tex
\chapter*{后记}
\addcontentsline{toc}{chapter}{后记}
待撰写

\cleardoublepage
% !TeX root = ../../Book3.tex
\PrintBookBibliography
% 仅在核读并引用真实资料后补入书目。

\cleardoublepage
% !TeX root = ../../Book3.tex
\PrintBookIndex{chinese}{中文索引}

\cleardoublepage
% !TeX root = ../../Book3.tex
\PrintBookIndex{foreign}{外文索引}

\cleardoublepage
% !TeX root = ../../Book3.tex
\PrintBookIndex{symbols}{符号索引}


\end{document}
